\documentclass[11pt]{article}
\usepackage[T1]{fontenc}
\usepackage[utf8]{inputenc}
\usepackage{lmodern,amsmath,amssymb,amsthm,mathtools}
\usepackage[a4paper,margin=29mm]{geometry}
\usepackage{microtype,booktabs,array,longtable,enumitem}
\usepackage[hidelinks]{hyperref}
\usepackage{xurl}
\setlength{\emergencystretch}{2em}
\setlist[enumerate]{leftmargin=*,itemsep=.35em}
\setlist[itemize]{leftmargin=*,itemsep=.25em}
\newtheorem{theorem}{Theorem}[section]
\newtheorem{proposition}[theorem]{Proposition}
\newtheorem{lemma}[theorem]{Lemma}
\newtheorem{corollary}[theorem]{Corollary}
\theoremstyle{definition}
\newtheorem{definition}[theorem]{Definition}
\newtheorem{example}[theorem]{Example}
\newtheorem{question}[theorem]{Question}
\theoremstyle{remark}
\newtheorem{remark}[theorem]{Remark}
\newcommand{\Q}{\mathbf Q}
\newcommand{\Z}{\mathbf Z}
\newcommand{\R}{\mathbf R}
\newcommand{\C}{\mathbf C}
\newcommand{\Qbar}{\overline{\mathbf Q}}
\newcommand{\A}{\mathcal A}
\newcommand{\cl}{\operatorname{cl}}
\newcommand{\Nm}{\operatorname{Nm}}
\newcommand{\Spec}{\operatorname{Spec}}
\newcommand{\Hilb}{\operatorname{Hilb}}
\newcommand{\Gal}{\operatorname{Gal}}
\newcommand{\workbenchurl}{https://github.com/KokunoYumeto/algebraic-cycles-and-motives-workbench}
\newcommand{\recordurl}{https://github.com/KokunoYumeto/mathematics-workbenches}
\title{Hodge's conjecture for abelian varieties as a uniform bound over finite fields, after Br\"unjes and Serp\'e}
\author{}
\date{}
\begin{document}
\maketitle
\begin{abstract}
Br\"unjes and Serp\'e established a criterion for lifting algebraic cycle classes from finite fields to characteristic zero under a uniform bound on their complexity. We express the rational Hodge conjecture for complex abelian varieties as one uniform degree bound in each special family: it is enough to test the bound at infinitely many good reductions of every CM fibre. The degree bound is uniform in the CM fibre; the places, the auxiliary prime and the integer clearing denominators may vary. Proper parameter spaces give the lifting criterion, and a positive pairing built from the Poincar\'e bundle bounds the denominator. We also study quantitative detection of Weil classes and the eventual behaviour of minimum cycle degrees under reduction. These results do not prove the Hodge conjecture. The mathematics and text were produced by GPT-6 Astra, with separate re-proofs by the same model and earlier checks by Claude Opus 5.5 at high effort; no human has checked them.
\end{abstract}

\section{The conjecture and the uniformity it requires}
\label{sec:introduction}

The rational Hodge conjecture says that every rational cohomology class of type $(p,p)$ on a smooth projective complex variety is a rational linear combination of codimension-$p$ algebraic cycles. Even for abelian varieties, this statement asks for algebraic cycles in situations where no construction is known. We show that, for abelian varieties, the existence question is equivalent to uniform boundedness in special families, and that this boundedness can be tested over finite fields.

The arithmetic starting point is Br\"unjes--Serp\'e \cite[Definition 2.3(ii), p.~4; Theorem 3.4, p.~8]{BS}. They bound both the effective numerator cycles and the integer denominator, use a fixed auxiliary prime, and work with cycles over finite residue fields. The form developed here permits cycles over algebraic closures of residue fields, permits the auxiliary prime to vary, and makes the denominator bound automatic. For abelian varieties an algebraic pairing gives an explicit linear bound for that denominator. The passage from one variety to special families relates this uniformity principle to Hodge's conjecture for all complex abelian varieties.

The complex argument also has a precise antecedent. In the opening page of Cattani--Deligne--Kaplan \cite[p.~483]{CDK}, the implication from the Hodge conjecture to algebraicity of Hodge loci is explained by bounded numerator degrees, a bounded denominator, Baire's theorem and Chow parameter spaces. Their Theorem~1.1 proves algebraicity of the Hodge locus unconditionally. We use the bounded-cycle argument for a global flat class and add arithmetic specialization. Bhattacharjee--Bhattacharya \cite[Lemmas~4.2--4.3 and Theorem~4.5]{BBDense} and \cite[Theorems~16.9 and 16.17]{BBAlgebraic} also formulate closed bounded-cycle loci and dense-set criteria; those recent preprints concern the qualitative complex step. Voisin \cite[Propositions~0.2 and 0.5, Theorem~0.6]{Voisin07} studies fields of definition of Hodge loci, which is a distinct issue from the cycle parameter spaces used below.

\subsection{Conventions and the main theorem}

If $(A,\theta)$ is a polarized abelian $g$-fold and $Z$ is an effective cycle of pure dimension $d$, set
\[
 \deg_\theta Z=\int_Z\theta^d.
\]
An effective integral cycle means a finite sum of integral subvarieties with nonnegative integer coefficients; the zero cycle is allowed. The word ``integral'' here refers to coefficients and does not require the cycle to be irreducible. When an irreducible reduced subvariety is intended, we say so explicitly.

\begin{definition}\label{def:special}
A \emph{special family with a Hodge class} consists of a connected component $S$ of a Shimura variety at neat level, a map
\[
 S\longrightarrow\mathcal A_{g,N},\qquad N\geq3,
\]
induced by a morphism of Shimura data, the pulled-back principally polarized abelian scheme $f:\A\to S$ with polarization $\theta$, and a nonzero global section
\[
 t\in H^0(S,R^{2p}f_*\Q(p))
\]
of type $(0,0)$ after the Tate twist, or equivalently type $(p,p)$ before the twist, at every point. For $s\in S(\C)$ define
\[
 c_t(s)=\min\left\{\max\{M,\deg_\theta Z^+,\deg_\theta Z^-\}:
 M t_s=\cl(Z^+)-\cl(Z^-),\ M\in\Z_{>0}\right\}.
\]
The cycles in this minimum are effective integral cycles of codimension $p$. The value is $+\infty$ if the set is empty. In this article a CM point of $S$ means a special point of its Shimura datum. Its image in the Siegel variety is then a CM abelian variety.
\end{definition}

A number-field model of a fibre means a polarized abelian variety over a number field $L\subset\C$ together with a specified identification of its complex fibre. It does not mean that $t_s$ is Galois invariant. At a good place $v$ of $L$, a geometric place above $v$ and a prime $\ell\ne\operatorname{char}k(v)$ determine the usual comparison and smooth proper specialization of $t_s$ to $\ell$-adic cohomology. The hypothesis on the map of Shimura data ensures that CM points have number-field models.

\begin{theorem}[Family criterion]\label{thm:family}
For a special family $(S,\theta,t)$, the following conditions are equivalent.
\begin{enumerate}[label=\textnormal{(\roman*)}]
\item Every $t_s$, for $s\in S(\C)$, is algebraic.
\item The function $c_t$ is bounded on $S(\C)$.
\item The function $c_t$ is bounded on one Zariski-dense subset of $S(\C)$.
\item There is an integer $D\geq1$ such that, for every CM point $x$ and every number-field model of its fibre, infinitely many good places $v$ admit a geometric place above $v$, a prime $\ell\ne\operatorname{char}k(v)$, an integer $M\geq1$, and effective integral cycles $Z_v^\pm$ over $\overline{k(v)}$ satisfying
\[
 \operatorname{sp}_{v,\ell}(t_x)
 =\frac{\cl_\ell(Z_v^+)-\cl_\ell(Z_v^-)}{M},
 \qquad \deg_\theta Z_v^\pm\leq D.
\]
There is no bound on $M$ in this condition. The choices other than $D$ may depend on $x$ and $v$.
\end{enumerate}
In \textnormal{(iii)} any chosen Zariski-dense subset can be used, including a rational Hecke orbit in the connected component. In \textnormal{(iv)} the good places may be required to belong to any prescribed infinite set. Equivalently, the condition holds at all but finitely many good places of every model, with one uniform bound on the numerator degrees and on $M$, for every allowable $\ell$ and geometric place. ``Every model'' can be replaced by ``some model.''
\end{theorem}

\begin{theorem}[Hodge's conjecture as boundedness]\label{thm:global}
The rational Hodge conjecture holds for every complex abelian variety in every codimension if and only if, in every special family of Definition~\ref{def:special}, $c_t$ is bounded on the CM points. It is also equivalent to condition \textnormal{(iv)} of Theorem~\ref{thm:family} for every such family.
\end{theorem}

The arithmetic bound therefore has the order of quantifiers
\[
 \text{for every family}\quad\exists D\quad\forall x\text{ CM}
 \quad\text{for infinitely many }v\quad\exists(M,Z_v^+,Z_v^-).
\]
Replacing $\exists D\,\forall x$ by $\forall x\,\exists D_x$ gives a different statement: it tests algebraicity separately at the CM points and does not give the conjecture on the other fibres. OpenAI \cite[Theorem~1.1]{OpenAICM} and van Hoften--Sempliner \cite[Theorem~I]{vHS}, following that work, give recent unrefereed proofs of the Hodge conjecture for CM abelian varieties. If those proofs are accepted, pointwise existence at CM points is available; the uniformity in $x$ remains exactly the issue in Theorem~\ref{thm:global}. Neither CM theorem is an input to its proof.

\subsection{Why the argument works}

For a fixed degree bound, integral subvarieties have only finitely many Hilbert polynomials. Products of projective Hilbert schemes therefore parametrize all pairs of cycles of those degrees. The equation specifying their cohomology class cuts out unions of connected components. Their images in the base are closed. A countable union of proper closed subvarieties cannot cover a complex variety, while one closed set containing a dense set is the whole variety. These two facts give the complex part of the theorem.

The arithmetic part uses the same finite parameter space over the ring of integers of one fibre. Its components contained in special fibres lie above finitely many places. Outside those places every point specializes from characteristic zero. The cycle class follows this specialization, so one suitable prime outside the exceptional set already suffices. Finally, an algebraic quadratic pairing is positive on all real Hodge classes of the required degree. It bounds $M$ in terms of the numerator degrees and is constant along the family. Primitivity of $t$ is unnecessary.

The proof uses the boundedness and representability theorems of Grothendieck and Kleiman, the trace, comparison and specialization theorems of SGA~4 and SGA~$4\frac12$, and the standard geometry of Shimura varieties. The particular loci below are images of spaces of algebraic cycles and are constructed directly. The argument does not use the Tate conjecture, independence of $\ell$, or a theorem on absolute Hodge classes.

\subsection{Further quantitative statements}

The rest of the paper has three parts.  The first proves the uniform-bound
criterion and its specialization dichotomy.  The second makes a single
Weil detector effective, transports detectors between lattices, and derives
degree bounds from specified geometric inputs.  The third studies the
remaining existence questions: it gives constructions with precisely stated
deformation gaps, proves obstructions for particular constructions, and
develops a separate detector criterion for Mumford fourfolds.

For a polarized abelian variety of Weil type, an imaginary quadratic field $K$ acts with signature $(n,n)$ on a variety of dimension $2n$. Its rational plane of Weil classes is
$W_K=\bigwedge_K^{2n}H^1(A,\Q)$. Let $\delta_W$ be the minimum degree of an irreducible subvariety with nonzero projection to $W_K$. Theorem~\ref{thm:extraction} proves, for every integral polarization compatible with the action, every nonzero $t\in W_K$, every $\delta\in\Z_{>0}$ with $\delta t$ integral, and every endomorphism $a\in K$ with $a/\bar a$ not a root of unity,
\[
 \delta_W\leq c_t\leq
 \left\lceil\frac{2\,16^n}{n!}\right\rceil
 \delta\lceil\sqrt T\rceil\Nm(a)^{n(2n+1)}\delta_W^2,
 \qquad T=(-1)^n\int t^2.
\]
Compatibility means that the Rosati involution restricts to conjugation on $K$. Infinite values have their usual extended-real interpretation. Thus one detector gives the whole Weil plane at quadratic cost.

Theorem~\ref{thm:small-detectors}, using the cycles of Schoen and van Geemen, gives detector degrees $48$ for even principal Gaussian fourfolds, $192$ for odd principal Gaussian fourfolds, and $36$ for principal Eisenstein fourfolds, all with full integral maximal-order action. In the Eisenstein case $3(W_K\cap H^4(A,\Z))$ is algebraic and
\[
 c_t\leq\max\{3\delta,36\lfloor\delta\sqrt{2T}/3\rfloor\}
 \leq17\delta\lceil\sqrt T\rceil.
\]
The even and odd Gaussian lattices, the integral normalization and the effective cycles are specified in Section~\ref{sec:small}.

Theorem~\ref{thm:stabilization}, after Schoen, says that bounded detectors in every connected split Weil family of dimension $2n+2$ over a fixed imaginary quadratic field imply bounded detectors in every Weil family of dimension $2n$ over that field, for every determinant. No class-number restriction is needed. If $\mathcal O_K=\Z[\tau]$ has class number one, acts integrally, and $\Delta=\chi(L)$, Theorem~\ref{thm:explicit-stabilization} gives
\[
 \delta_W(A,L)\leq\frac{\max\{1,\Nm(\tau)/\Delta\}}{n+1}D_{n+1}(\Delta^2),
\]
where $D_{n+1}(x)$ bounds detectors for all split Weil $(2n+2)$-folds of polarization product $x$. The premise includes every polarization type of that product. The numerator of the displayed factor is one for the Gaussian and Eisenstein fields.

Schoen's cyclic cycles and the dominance theorems of Schoen and Koike give the explicit sixfold estimates in Theorem~\ref{thm:sixfold-suppliers}: $746496$ for the even principal Gaussian lattice, $5971968$ for the odd one, and $30375$ for the principal Eisenstein lattice. Theorem~\ref{thm:principal-quotient} gives a principal hyperbolic quotient over every imaginary quadratic field. Pullback of detectors costs $t_K^nf^n\Delta$, where $\operatorname{Tr}(\mathcal O_K)=t_K\Z$, $f$ is the conductor of the acting order and $\Delta$ is the polarization volume. The ideal classes can also be removed at the additional cost $m_K^n$, where $m_K\le2\sqrt{|\operatorname{disc}K|}/\pi$. Thus one principal family with lattice $U^n$ suffices for each field and dimension; for class number one the additional factor is one. Consequently, for a compatible polarized Weil fourfold of any determinant with integral action by $\mathcal O_f=\Z+f\mathcal O_K$ and $\Delta=\chi(L)$, Theorem~\ref{thm:all-determinant-fourfolds} gives
\[
 \delta_W\le1990656 f^6\Delta^2\quad(K=\Q(i)),
 \qquad\delta_W\le10125 f^6\Delta^2\quad(K=\Q(\omega)).
\]
The proof keeps the source cycles effective and uses published geometric inputs.

Markman \cite[Theorem~1.2]{MarkmanICM} proves algebraicity for split Weil sixfolds over every imaginary quadratic field, and for all Weil fourfolds; the published account gives a proof sketch and points to the detailed preprint. The field restriction to nine class-number-one fields belongs to particular numerical estimates, not to this qualitative theorem. Assuming Perry's unrefereed equivariant semiregularity theorem, Theorem~\ref{thm:all-field-sixfolds} gives an explicit bound for every imaginary quadratic field: if $K=\Q(\sqrt{-d})$, $d$ is positive squarefree, $\mathcal O_f$ acts integrally, and the compatible polarization has volume $\Delta$, then every split Weil sixfold satisfies
\[
 \delta_W<10^{179}d^{53}f^3\Delta.
\]
For class number one the coefficient improves to $10^{178}d^{51}$. The sharper constant is an explicitly evaluated Chern-class degree. Split eightfolds and higher dimensions require further arguments.

The $G_{34}$ family in Severn's unrefereed preprint gives a different input in dimension eight. Theorem~\ref{thm:eisenstein-eightfolds} writes out an explicit integer $B_0$ and proves, from that construction, the bounds
\[
 \begin{aligned}
 \delta_W&\le B_0f^4\Delta
 &&\text{for split Eisenstein eightfolds},\\
 \delta_W&\le\frac{B_0}{4}\frac{f^9\Delta^2}{b}
 &&\text{for all Eisenstein sixfolds},
 \end{aligned}
\]
where $b$ is the index of maximal-order saturation. Schoen's cyclic-cover cycles, Severn's exceptional-divisor correction, Brauer induction and Milne's invariant theory carry this proof. The degree calculation keeps integral orbit multiplicities and then uses the principal quotient to reduce the transport cost.

Theorem~\ref{thm:ghosts} applies to any fixed smooth projective variety over a number field with an ample line bundle and any rational Betti class. Its minimum numerator degree at good reductions is eventually the characteristic-zero minimum if that is finite; otherwise it exceeds every fixed bound outside finitely many places. The multiplier and auxiliary prime may vary. The same dichotomy holds for the cost that also bounds the multiplier.

These results concern algebraic representatives, not merely Hodge-theoretic tensors or motivated classes. Andr\'e's deformation theorem and his theorem that abelian Hodge classes are motivated \cite[Theorems~0.5 and 0.6.2, pp.~8--9]{Andre96} establish powerful propagation results in the motivated category. His CM reduction \cite[Lemma~6.3.2, p.~32]{Andre96} uses split Weil classes over general CM fields. Neither assertion supplies the uniform algebraic cycles required here, and the present criterion does not reduce all abelian Hodge classes to imaginary-quadratic Weil classes.

Two restricted eightfold constructions admit precise tests. Theorem~\ref{thm:theta-rank} shows that, for $V$ with basis indexed by $\mathbf F_2^4$, every nonzero alternating four-form invariant under translations and character signs induces a pairing on $\bigwedge^2V$ of rank at least $24$, and that $24$ occurs. Corollary~\ref{cor:all-eightfold-quadrics} identifies this $35$-dimensional space with the entire quadratic ideal of the minus projection on a nonempty open subset of the even Gaussian eightfold family. Thus no quadratic relation there can provide the rank-$10$ quadric needed by the direct ruling-class extension of van Geemen's argument. Theorem~\ref{thm:localized-support} proves a different obstruction: a perfect complex supported on translates of abelian subvarieties meeting in codimension at least two cannot deform in a direction with nonzero normal class on a component of nonzero generic virtual multiplicity, even with an infinitesimal gerbe twist. On norm-one scalar graphs in $E_i^8$ this excludes all six skew normal directions, including for arbitrary extensions and thickenings on those supports. The general eightfold existence problem remains open.


Two other results give exact tests for possible eightfold constructions. Theorem~\ref{thm:pure-weil-ext} applies to a perfect complex $E$ on a balanced Weil $2n$-fold, $n\ge2$, with a rational algebraic $B$-field and an exactly pure nonzero normalized Weil character. For a coherent finite translation--tensor action and an equivariant structure, its ordinary semiregularity trace on invariant obstructions is injective if and only if
\[
 \dim\operatorname{Ext}^2(E,E)^G=2n(2n-1).
\]
On an eightfold this number is $56$. Proposition~\ref{prop:sixteen-line-character} constructs an integral sixteen-term character on the principal Gaussian family with precisely that normalized Weil character, but Proposition~\ref{prop:sixteen-line-gluing} proves that every indicated two-level gluing has invariant second self-Ext dimension $448$. Theorem~\ref{thm:all-point-products} excludes every external product of point-ideal characters, for every partition of dimension eight, and all their rational Fourier--Mukai orbits as sources of a nonzero full-family Weil component.

For finite covers of curves, Theorem~\ref{thm:finite-cover-rank} says that a faithful irreducible coefficient representation with nonrational character, eight-dimensional parabolic multiplicity of type $(4,4)$, and branch moving ranks at least two needs coefficient rank at least twenty-six to have a sixteen-dimensional period image. Theorem~\ref{thm:gaussian-reflections} excludes dominance when a Gaussian representation with exact character field contains a reflection. Hence the rank bound applies to every dominant source of that specified Gaussian kind. Theorem~\ref{thm:scalar-standard-all} further excludes scalar-standard coefficients in every degree: their period variations factor through rank-one local systems on curves with moduli dimension at most fifteen. These statements restrict construction methods and prove no additional case of algebraicity.


Theorem~\ref{thm:simple-hodge-bundle} gives simple vector bundles on
every CM elliptic product $E^{2n}$, for every imaginary quadratic
field, whose Chern character is a nonzero prescribed primitive
middle class plus an explicit polynomial in the polarization.
For a Weil class and $n\ge3$, the Hochschild action has rank
$7n^2-2n$, equal to $104$ on an eightfold; the general rank formula
is \cite[Theorem~A6]{BBK3Sixfolds}.
The construction uses positive multiplicities and ordinary extensions.
Theorem~\ref{thm:tree-extension-obstruction} proves that its specified
tree of permutation orbits cannot deform in every Weil direction,
even after its filtration is forgotten.  Additional gluing data and
the sufficient second-Ext equality remain to be found.
Theorem~\ref{thm:sign-cube-obstruction} extends the obstruction to a
specified family of coordinate-sign networks with many cycles.
Theorem~\ref{thm:gaussian-positive-trace-network} gives a different,
explicit Gaussian reflection network whose positive multiplicities
have a nonzero Weil character and solve every matrix-valued diagonal
trace equation.  The remaining off-diagonal extension equations and
strong semiregularity are not established, so this construction does
not yet give a deforming detector.

For a polarized Hodge-generic Mumford fourfold $A$, the exceptional
codimension-two classes on $A^2$ form a plane identified with the
trace-zero part of its defining totally real cubic field. Theorem
\ref{thm:mumford-cubic} proves
$\delta_{\mathrm{exc}}\leq c_t\leq C_t\delta_{\mathrm{exc}}^3$
for each fixed nonzero class $t$ in that plane, where
$\delta_{\mathrm{exc}}$ is the least degree of an irreducible
codimension-two detector for the product polarization. The constants
are uniform in a fixed polarized family with global $t$.
Corollary~\ref{cor:mumford-family} therefore reduces algebraicity
throughout that family to bounded detectors on infinitely many CM
fibres. These exceptional classes are not Weil classes; the theorem
does not construct their first detector.

When $A$ is the Jacobian of one of the genus-four curves of Bouchet--Hanselman--Pieper--Schiavone, Theorem~\ref{thm:mumford-surface} gives a further sufficient condition. Twist the polarization on its rank-nine K3-type Hodge structure by a totally positive $\lambda\in F\setminus\Q^\times F^{\times2}$, and realize it by a projective K3 surface $S_\lambda$. One algebraic surface correspondence from $S_\lambda$ to $C\times C$ with nonzero transcendental action makes the entire exceptional plane algebraic. The extraction uses a nonconstant weighted Casimir tensor. A dominant rational map from $C\times C$ to this chosen $S_\lambda$ cannot provide it. The theorem refines the conditional K3 approach in their Section~5.2; it does not construct the first correspondence.

A different conditional reduction is Theorem~\ref{thm:mumford-coniveau}. Suppose the quaternion algebra over $F$ is split by $F(\sqrt{-3})$ and has split corestriction to $\Q$. On a common finite cover of the Mumford curve there is a split Weil sixfold $B$ and an explicitly algebraic projector onto a rank-sixteen substructure $P\subset H^3(B,\Q)$ with $P(1)\simeq H^1(A,\Q)^{\oplus2}$. At a Hodge-generic point, one nonzero geometrically coniveau-one part of $P$ implies the rational Hodge conjecture for every power of $A$, including the exceptional plane on $A^2$. The known Weil cycles on $B$ give $P$ as a summand of the middle cohomology of supporting threefolds; the additional divisor support remains an open generalized Hodge question.


\section{Proper parameter spaces for bounded cycles}
\label{sec:parameters}

This section isolates the boundedness and cycle-class facts needed in both characteristics. We keep the use of Hilbert schemes explicit because arbitrary nonreduced subschemes of bounded degree need not have bounded Hilbert polynomials.

\begin{lemma}[Cycle parameter spaces]\label{lem:parameters}
Let $T$ be noetherian and let $f:X\to T$ be smooth and projective of relative dimension $g$, with a relatively ample line bundle of class $\theta$. Fix $0\leq d\leq g$ and $D\geq1$. There is a projective $T$-scheme $H_D$ whose geometric points give all pairs of effective integral $d$-cycles of $\theta$-degree at most $D$. Each universal pair has cycle classes which, for $\ell$ invertible on the base, are sections of the pullback of $R^{2g-2d}f_*\Q_\ell(g-d)$. These sections commute with base change. The corresponding assertion holds in Betti cohomology over $\C$. The assertion also holds for a polarized abelian scheme without a chosen line bundle representing its polarization.
\end{lemma}

\begin{proof}
A sufficiently high fixed power of the relatively ample line bundle is relatively very ample. In the abelian case the polarization morphism $\lambda:\A\to\A^\vee$ and the normalized Poincar\'e bundle give, by pullback under $(1,\lambda)$, a relatively ample line bundle of class $2\theta$; one uses a fixed very ample power of this bundle instead. The choice changes the degree bound by one fixed factor and has no effect on the argument.

By SGA~6, XIII, Definition~6.9 and Corollary~6.11 \cite{SGA6}, the Hilbert polynomials of integral $d$-dimensional subvarieties of bounded degree belong to a finite set. An effective cycle of $\theta$-degree at most $D$ contains at most $D$ integral components counted with their multiplicities, since every nonzero integral subvariety has positive integer degree. Enumerate all tuples of these finitely many polynomials for which the sum of the leading degrees is at most $D$ in the $\theta$ normalization. Include the empty tuple for the zero cycle.

For each pair of tuples take the corresponding fibre product of relative Hilbert schemes. Their finite disjoint union is projective over $T$ by FGA, Expos\'e~221, Theorem~3.2 and its specialization to the Hilbert functor \cite{FGA}. At any point take the top-dimensional fundamental cycles of the universal subschemes and add them. All such cycles have the specified degree, including at nonreduced limits, and every effective cycle occurs by using its integral components with repetition. The construction does not claim to parametrize every nonreduced subscheme of bounded degree.

For a universal subscheme $Z$ flat over its parameter scheme $H$, its fibres have dimension $d$. Write $f_H:X_H\to H$ and $h:Z\to H$. Restriction followed by trace gives
\[
 R^{2d}(f_H)_*\Q_\ell(d)\longrightarrow
 R^{2d}h_*\Q_\ell(d)\xrightarrow{\operatorname{Tr}_h}\Q_\ell.
\]
The trace exists and commutes with base change by SGA~4, XVIII, Theorem~2.9 \cite{SGA4}. Poincar\'e duality for the smooth proper map $f_H$, in XVIII, Theorem~3.2.5, identifies this functional with a section of $R^{2g-2d}(f_H)_*\Q_\ell(g-d)$. On every geometric fibre it is the class of the fundamental cycle, with scheme-theoretic multiplicities. Add these sections over the tuple. Their compatibility with algebraic cycle classes is SGA~$4\frac12$, Cycle, construction~2.3.1, Definition~2.3.2, Lemma~2.3.7 and Theorem~2.3.8 \cite{SGA45}. Comparison on the smooth complex fibres is SGA~4, XI, Theorem~4.4, together with the cycle-class compatibility. This also proves the last assertion.
\end{proof}

\begin{lemma}[Closed class loci]\label{lem:closed}
Let $f:X\to S$ be a smooth projective family with a relatively ample class over a connected smooth complex algebraic variety, and let $t$ be a nonzero global rational cohomology section. For each integer $C\geq1$, the locus
\[
 S_C=\{s\in S(\C):c_t(s)\leq C\}
\]
is Zariski closed.
\end{lemma}

\begin{proof}
On the parameter scheme $H_C$ of Lemma~\ref{lem:parameters}, and for a fixed $1\leq M\leq C$, the difference of the two cycle-class sections minus $Mt$ is a section of a local system. Its zero locus is a union of connected components. There are finitely many such components, and each is algebraic and both open and closed. Thus the zero locus is proper over $S$, and its image is closed. The union of these images for the finitely many possible $M$ is precisely $S_C$.
\end{proof}

\begin{lemma}[An unrestricted multiplier is automatically bounded]\label{lem:components-multiplier}
In the setting of Lemma~\ref{lem:closed}, fix $D$. Among all equations
\[
 Mt_s=\cl(Z^+)-\cl(Z^-),\qquad
 M\in\Z_{>0},\quad \deg Z^\pm\leq D,
\]
only finitely many values of $M$ occur, as $s$ and the cycles vary. In particular the locus defined by a bound on the two numerator degrees alone is Zariski closed, even without a Hodge-type assumption on $t$.
\end{lemma}

\begin{proof}
On $H_D$, let $\gamma$ be the difference of the universal cycle classes. Both $\gamma$ and the pullback of $t$ are flat sections. If $\gamma=Mt$ at one point of a connected component, the equality holds throughout that component. Two distinct values $M,M'$ on the same component would give $(M-M')t=0$, contrary to the nonvanishing of the global flat section $t$. There are finitely many connected components of $H_D$, so there are finitely many possible multipliers. The desired locus is the image of the union of the components on which one of these equalities holds; this union is proper over $S$.
\end{proof}

The preceding lemma is qualitative and applies to arbitrary smooth projective families. The pairing in the next section supplies an explicit multiplier bound in the abelian case, without needing to count components of a parameter space.

\begin{proposition}[Uniformity over a complex base]\label{prop:complex}
In Lemma~\ref{lem:closed}, assume that $t$ is a Hodge class on every fibre. Algebraicity on every fibre, boundedness of $c_t$ on all fibres, and boundedness of $c_t$ on one Zariski-dense set are equivalent.
\end{proposition}

\begin{proof}
Algebraicity says that $S(\C)=\bigcup_{C\geq1}S_C$. Since $S$ is smooth and connected, it is irreducible. A proper algebraic closed subset has empty interior in its complex manifold. A nonempty relatively compact analytic coordinate ball is a Baire space after passing to a smaller closed ball; it cannot be a countable union of such subsets. Hence some $S_C$ equals $S$. Conversely, a closed subset $S_C$ containing a Zariski-dense subset equals $S$, and finite $c_t(s)$ means algebraicity. Boundedness on all fibres implies boundedness on any subset.
\end{proof}

\section{A positive algebraic pairing bounds the denominator}
\label{sec:positive}

We now bound the multiplier using a pairing which remains available in positive characteristic. This is the part that turns a degree bound into a complexity bound.

On $A\times A$ let $\theta_1,\theta_2$ denote the pullbacks of $\theta$. For addition $m$ and difference $a$, put
\[
 H=\theta_1+\theta_2,\qquad P=m^*\theta-H.
\]
The theorem of the square gives $a^*\theta=H-P$ and $m^*\theta=H+P$. In particular $H+P$ and $H-P$ are nef.

\begin{lemma}[The Poincar\'e pairing]\label{lem:positive}
Let $t\ne0$ be a real Hodge class of codimension $p$ on a polarized complex abelian $g$-fold and put $d=g-p$. Then
\[
 q_t=\left|\int_{A\times A}(t\boxtimes t)P^{2d}\right|>0.
\]
If $t$ and $\theta$ vary as flat sections in a connected family of Hodge classes, $q_t$ is constant. For effective $d$-cycles $Z,Z'$ over any algebraically closed field,
\begin{equation}\label{eq:nef}
 \left|\int_{Z\times Z'}P^{2d}\right|
 \leq \binom{2d}{d}\deg_\theta Z\deg_\theta Z'.
\end{equation}
Consequently, if $Mt=\cl(Z^+)-\cl(Z^-)$, then
\begin{equation}\label{eq:denominator}
 M^2q_t\leq \binom{2d}{d}
       (\deg_\theta Z^++\deg_\theta Z^-)^2.
\end{equation}
The same inequality holds for a smooth proper specialization of the equation.
\end{lemma}

\begin{proof}
Let $V=H_1(A,\R)$, let $J$ be its complex structure, and let $E$ be the Riemann form of $\theta$. Choose the convention in which $G(u,v)=E(u,Jv)$ is positive definite. Let $x\in\bigwedge^{2d}V$ be Poincar\'e dual to $t$. In symplectic coordinates, the mixed part of $m^*\theta$ is the bilinear form $E$ between the two copies of $V$. Expanding its $2d$-th power gives
\[
 \left|\int_{A\times A}(t\boxtimes t)P^{2d}\right|
 =(2d)!\,\left|(\bigwedge^{2d}E)(x,x)\right|.
\]
Here the exterior-power form is defined on decomposable vectors by
\[
 (\bigwedge^r E)(u_1\wedge\cdots\wedge u_r,
 v_1\wedge\cdots\wedge v_r)=\det(E(u_i,v_j)).
\]
The factorial counts the permutations of the mixed factors; the sign disappears under the absolute value. Since $x$ has homological type $(-d,-d)$, $\bigwedge^{2d}J$ fixes $x$. Therefore
\[
 (\bigwedge^{2d}E)(x,x)
 =(\bigwedge^{2d}G)(x,x)>0.
\]
The last inequality is positive definiteness of the exterior-power metric. Constancy follows because the original integral contains only flat classes.

Write $U=H+P$ and $V'=H-P$. All mixed intersections of $U$ and $V'$ with the effective cycle $Z\times Z'$ are nonnegative. Expanding $(U-V')^{2d}$ and applying the triangle inequality, then using $U+V'=2H$, gives
\[
 \left|\int_{Z\times Z'}P^{2d}\right|
 \leq\int_{Z\times Z'}H^{2d}
 =\binom{2d}{d}\deg_\theta Z\deg_\theta Z'.
\]
Only the term of degree $d$ on each factor survives in the final expansion. This argument is intersection theory with nef divisor classes and works in every characteristic. Apply it to the four products arising from $(Z^+-Z^-)\times(Z^+-Z^-)$ to obtain \eqref{eq:denominator}. Smooth proper specialization preserves the pairing and the intersection numbers, so it preserves the same positive rational number $q_t$ and the same inequality.
\end{proof}

\begin{corollary}\label{cor:denominator}
If both numerator degrees are at most $D$ in a connected family, then
\[
 M\leq 2D\sqrt{\binom{2d}{d}/q_t}.
\]
Thus bounded numerator degrees alone are equivalent to bounded $c_t$ on any subset of that family.
\end{corollary}

\begin{example}
On a principally polarized elliptic curve, take $t=r\theta$, with $r\in\Q_{>0}$, and $p=g=1$. Then $d=0$ and $q_t=r^2$. An equation $Mt=[Z^+]-[Z^-]$ says $Mr=\deg Z^+-\deg Z^-$, so the denominator inequality becomes the elementary estimate $Mr\leq\deg Z^++\deg Z^-$. The same argument in higher dimension uses the Poincar\'e pairing instead of the degree map.
\end{example}

\section{Lifting from one nonexceptional reduction}
\label{sec:lifting}

Properness turns infinitely many reductions into a characteristic-zero cycle. The exceptional set depends on the variety and the degree bound; no bound for its size or largest prime is asserted.

\begin{proposition}[Lifting bounded numerator degrees]\label{prop:lift}
Let $(X,\theta)$ be a smooth projective variety with an ample line bundle over a number field $L\subset\C$, let $t\in H^{2p}(X(\C),\Q(p))$, and fix $D\geq1$. There is a finite set $E_D$ of finite places of $L$ such that the following holds. If $v\notin E_D$ is a good place, and for one geometric place over $v$ and one allowable prime $\ell$ the specialization of $t$ equals
\[
 (\cl_\ell Z^+-\cl_\ell Z^-)/M
\]
with $M\in\Z_{>0}$ and effective cycles of degrees at most $D$, then there are effective cycles $Y^\pm$ on $X_{\Qbar}$ of degrees at most $D$ with
\[
 Mt=\cl(Y^+)-\cl(Y^-)
\]
in Betti cohomology. Neither Galois invariance nor absolute Hodgeness of $t$ is required.
\end{proposition}

\begin{proof}
Spread the smooth projective variety and its ample line bundle over $R=\mathcal O_L[1/b]$ for a nonzero integer $b$. Construct $H_D\to\Spec R$ by Lemma~\ref{lem:parameters}. It has finitely many irreducible components. Those which do not dominate $\Spec R$ are supported over finitely many closed points; put these places, and the excluded places of the model, in $E_D$.

Let $z\in H_D(\overline{k(v)})$ represent the two given cycles. It lies on the closure of a component of the generic fibre. To see the specialization with its geometric choices, pass to the strict henselization $R^{\mathrm{sh}}$ at the chosen geometric place. The reduced structure on a horizontal irreducible component of $H_D$ is torsion-free over the Dedekind base, hence flat. After this base change, every point on it is a specialization of a point in its generic fibre: flat morphisms of finite presentation are open, so every open neighbourhood meets the generic fibre. Choose a geometric generization $\eta$ of $z$ there.

Smooth proper base change with constant coefficients, SGA~4 XVI, Corollary~2.2, identifies the cohomology of the generic and special fibres over the strictly henselian trait. These coefficients satisfy the local-constancy hypothesis added in the SGA~$4\frac12$ erratum. The class $t_\ell$ extends to a constant section on this trait; it need not extend to the original arithmetic base. The universal cycle-class sections commute with specialization. Their equality with $Mt_\ell$ at $z$ therefore holds at $\eta$.

The generic parameter component containing $\eta$, after passage to an algebraically closed characteristic-zero field, descends to a component over $\Qbar$. That component has a $\Qbar$-point, and the classes of its universal cycles are constant on each connected component because the ambient variety is fixed. Choose such a point on the component of $\eta$. Its cycles have the same degrees and satisfy the same $\ell$-adic equation on $X_{\Qbar}$. Comparison with Betti cohomology, and injectivity of extension of scalars from $\Q$ to $\Q_\ell$, give the stated rational equation. All choices can be made compatibly with the original embedding $L\subset\C$ and geometric place. Inertia acts trivially on the cohomology of a good fibre, so the choice of a compatible generization does not alter the class being lifted.
\end{proof}

\begin{proposition}[Spreading an algebraic class]\label{prop:spread}
In Proposition~\ref{prop:lift}, if an equation $Mt=\cl(Y^+)-\cl(Y^-)$ with numerator degrees at most $D$ holds over $\C$, the cycles can be chosen over $\Qbar$ with the same degrees. At all but finitely many good places of any model the specialized equation holds over the residue algebraic closure, for every allowable $\ell$ and geometric place, with the same $M$ and the same degree bound.
\end{proposition}

\begin{proof}
The complex point of $H_D$ giving the cycles belongs to a connected component of $H_{D,\Qbar}$ after extension of scalars. That component has a $\Qbar$-point and has constant cycle classes, giving the descent. Its point is defined over a finite extension $L'/L$. Properness of $H_D$ extends it over every valuation ring of $L'$ above the chosen arithmetic base. Taking the fundamental cycles of the special fibres preserves the degrees. Compatibility of cycle classes with specialization gives the equation. The argument works with the fixed geometric place by choosing a place of $L'$ above it. Only finitely many places were excluded in making the model and its parameter scheme.
\end{proof}

\section{Special families and the proof of the equivalence}
\label{sec:shimura}

Special families connect the complex boundedness statement to all abelian varieties, because every rational Hodge class occurs in one and its CM points are dense.

\begin{lemma}[Every class occurs in a special family]\label{lem:shimura}
Every nonzero rational Hodge class on a complex abelian variety is, after pullback by an isogeny from a principally polarized abelian variety, the value of a section in Definition~\ref{def:special}. CM points are dense in each connected Shimura variety, and a rational Hecke orbit in a fixed connected component is dense.
\end{lemma}

\begin{proof}
A polarized complex abelian variety admits an isogeny from a principally polarized abelian variety. One elementary description is to take a symplectic basis $e_i,f_i$ of its homology lattice with $E(e_i,f_j)=d_i\delta_{ij}$, choose an integer $m$ divisible by all $d_i$, and pass to the sublattice generated by $e_i$ and $(m/d_i)f_i$. The form $E/m$ is integral and unimodular on that sublattice and remains a positive Riemann form. The corresponding torus $B$ is principally polarized and maps by an isogeny $\phi$ to the original variety. If $\phi^*t$ is algebraic, then so is $t$, because $\phi_*\phi^*t=(\deg\phi)t$.

Let $h$ describe the polarized rational Hodge structure on $H_1(B,\Q)$ and let $G$ be its Mumford--Tate group inside the corresponding group of symplectic similitudes. The pair $(G,X)$, where $X$ is the $G(\R)$-conjugacy class of $h$, is a Shimura datum of Hodge type. Indeed the Hodge types in its adjoint representation are $(-1,1),(0,0),(1,-1)$; polarization gives the Cartan involution; and minimality of the Mumford--Tate group excludes an adjoint rational factor on which $h$ is trivial. The inclusion in the symplectic similitude group is a morphism of Shimura data. At neat level it gives the required algebraic map to $\mathcal A_{g,N}$ \cite[\S\S5--7]{MilneSV}.

The tensor $\phi^*t$ with its Tate twist is fixed by $G$. It therefore gives a global section on the component containing $B$, and conjugation of $h$ by $G(\R)$ preserves its Hodge type. This proves the first assertion.

A Shimura datum has special points, namely homomorphisms factoring through rational tori. Their existence follows by approximating a regular element in a suitable real Cartan algebra by a rational regular element and taking its rational centralizer; the resulting real torus is conjugate to the original one. Real approximation for connected rational algebraic groups says that $G(\Q)$ is dense in $G(\R)$ \cite[Theorem~5.4]{MilneSV}. Conjugating one special point by these rational points therefore gives a dense set in each relevant connected component. More generally the same approximation applied to any point gives density of its rational Hecke orbit. A morphism of Shimura data sends special points to special points. In the Siegel variety these are CM abelian varieties and are defined over number fields. This is the only arithmetic property of special points used here.
\end{proof}

\begin{proof}[Proof of Theorems~\ref{thm:family} and \ref{thm:global}]
Proposition~\ref{prop:complex} proves the first three equivalences in the family theorem. If $c_t\leq C$ everywhere, Proposition~\ref{prop:spread} applied to each CM fibre gives the arithmetic condition with $D=C$, even at all but finitely many places and with $M\leq C$. This also gives the assertion for any prescribed infinite set of good places.

Conversely, assume the arithmetic condition with bound $D$. At each CM point, infinitely many places guarantee one outside the finite exceptional set of Proposition~\ref{prop:lift}. That proposition gives a characteristic-zero equation with numerator degrees at most $D$. Corollary~\ref{cor:denominator} gives
\[
 c_t(x)\leq
 \max\left\{D,\left\lceil 2D\sqrt{\binom{2g-2p}{g-p}/q_t}\right\rceil\right\}.
\]
This bound is independent of $x$ because $q_t$ is constant in the family. CM density and the closed-locus argument prove algebraicity at every point. The argument used one auxiliary prime at one suitable place for each $x$ and made no comparison between distinct auxiliary primes. It also shows that one model is enough.

If the rational Hodge conjecture holds for all abelian varieties, condition (i) holds for every special family, so all the stated boundedness conditions follow. In the other direction, Lemma~\ref{lem:shimura} places an arbitrary Hodge class, after an isogeny, in such a family. The family theorem makes that pullback algebraic, and isogeny pushforward makes the original class algebraic.
\end{proof}

\section{Minimum degrees and reductions}
\label{sec:ghosts}

The same properness argument describes exactly what bounded-degree reduction can and cannot hide for one fixed variety and one fixed class.

Let $(X,\theta)$ be a smooth projective variety with an ample line bundle over a number field and fix a rational Betti class $t$. Define $b_0$ to be the least $D$ such that $Mt=\cl(Z^+)-\cl(Z^-)$ for some $M\geq1$ and effective cycles of degrees at most $D$ over $\C$. Set $b_0=+\infty$ if there is no such equation. Define $b_v$ similarly on the geometric good reduction, allowing a choice of geometric place and auxiliary prime as in Theorem~\ref{thm:family}.

\begin{theorem}[Eventual minimum-degree dichotomy]\label{thm:ghosts}
If $b_0<\infty$, then $b_v=b_0$ for all but finitely many good places $v$. If $b_0=\infty$, then for each integer $D$ only finitely many good places satisfy $b_v\leq D$. The same statements hold when the cost also bounds $M$, as in $c_t$.
\end{theorem}

\begin{proof}
For $t=0$ all the numerator-degree minima are zero, and all the minima for the cost bounding $M$ are one, so both assertions are immediate. Suppose $t\ne0$. When $b_0<\infty$, choose a characteristic-zero equation realizing that minimum. Proposition~\ref{prop:spread} gives $b_v\leq b_0$ outside finitely many places. If $b_0\geq2$, Proposition~\ref{prop:lift} with $D=b_0-1$ gives $b_v\geq b_0$ outside another finite set. If $b_0=1$, smooth proper specialization makes $b_v=0$ impossible. If $b_0=\infty$, infinitely many places with $b_v\leq D$ would give a characteristic-zero equation by Proposition~\ref{prop:lift}, a contradiction. For the cost bounding $M$, the same lifting preserves the specific multiplier and the same spreading preserves it, so the argument is identical.
\end{proof}

Thus a class which is algebraic at every good reduction but not in characteristic zero would have reduction degrees escaping every fixed bound. This is a conditional description of such a class, not an existence assertion. Moreover, for fixed $A$ and $D$, one good place outside the finite set $E_D$ suffices for lifting. Making that finite set effective and uniform in a CM family is a separate problem.

\section{Weil classes and their detecting cycles}
\label{sec:weil}
This section studies the quantitative form of algebraicity for the two-dimensional space of Weil classes. Let $K$ be an imaginary quadratic field acting on a polarized abelian variety $A$ of dimension $2n$, with signature $(n,n)$ and with the Rosati involution restricting to complex conjugation on $K$. Its Weil space is the rational plane
\[
 W_K(A)=\bigwedge_K^{2n}H^1(A,\Q)\subset H^{2n}(A,\Q).
\]
It consists of Hodge classes. Write $e_W$ for its projector in the eigenspace decomposition for the $K$-action, and write $\delta_W(A,\theta)$ for the least degree of an irreducible reduced $n$-dimensional subvariety $Y$ with $e_W\cl(Y)\ne0$, or $+\infty$ if no such subvariety exists. Such a subvariety will be called a detector.

The principal open question for the Weil-type part of this article is whether $\delta_W$ is bounded on the CM points in every connected split Weil family, in every dimension and over every imaginary quadratic field. Here ``split'' means that the rational Hermitian space of dimension $2n$ and signature $(n,n)$ is a sum of hyperbolic planes. Its determinant is $(-1)^n$ modulo norms from $K$; the positive signed discriminant is therefore $1$ modulo norms. We use ordinary intersection degrees throughout.

\subsection{One detector gives the entire plane}

\begin{theorem}[Quadratic extraction]\label{thm:extraction}
Let $n\geq1$ and let $A$, $\theta$ and $K$ be as above, with $\theta$ integral and ample. Let $a\in K$ act by an endomorphism of $A$, assume $a/\bar a$ is not a root of unity, and put $q=\Nm_{K/\Q}(a)$. For a nonzero $t\in W_K(A)$, choose $\delta\in\Z_{>0}$ with $\delta t$ integral and set
\[
 T=(-1)^n\int_A t^2>0.
\]
Then
\begin{equation}\label{eq:sharper-extraction}
 \delta_W\leq c_t\leq
 \left\lceil\frac{2\,16^n}{n!}\right\rceil
 \delta\lceil\sqrt T\rceil\,
 q^{n(2n+1)}\delta_W^2.
\end{equation}
The assertion is valid for every polarization type. In particular it implies
\begin{equation}\label{eq:original-extraction}
 c_t\leq
 \left\lceil\frac{8\,16^n}{n!}\right\rceil
 \delta\lceil\sqrt T\rceil\,
 q^{2n(n+1)}\delta_W^2.
\end{equation}
If no detector exists, both $\delta_W$ and $c_t$ are infinite. Thus the existence of one detector is equivalent to algebraicity of the entire rational Weil plane.
\end{theorem}

The proof has two parts. A polynomial in endomorphism pullbacks removes every middle eigenspace except the Weil plane. Two integral vectors in that positive plane then span a sublattice of finite index. Clearing that index, rather than estimating an arbitrary inverse matrix, keeps the cost quadratic.

\begin{lemma}[Norm and degree]\label{lem:weil-norm}
The Weil plane is primitive for $\theta$, and
\[
 B(u,v)=(-1)^n\int_Au\smile v
\]
is a positive definite integral form on its integral lattice. With
$v_\theta=\int_A\theta^{2n}/(2n)!\in\Z_{>0}$ and
$\|u\|^2=B(u,u)$, every effective $n$-cycle $Y$ satisfies
\[
 \|e_W\cl(Y)\|\leq\frac{\deg_\theta Y}{n!\sqrt{v_\theta}}.
\]
For any chosen signed presentation $C=C^+-C^-$ whose class lies in $W$, its norm is at most
$(\deg C^++\deg C^-)/(n!\sqrt{v_\theta})$.
\end{lemma}

\begin{proof}
Over $\C$ write $V=H^1(A,\C)=V_\sigma\oplus V_{\bar\sigma}$ for the two eigenspaces of $K$, each of dimension $2n$. The signature hypothesis makes the extreme determinant lines
\[
 W_\C=\bigwedge^{2n}V_\sigma\oplus\bigwedge^{2n}V_{\bar\sigma}
\]
of Hodge type $(n,n)$. Rosati compatibility gives $a^*\theta=q\theta$, so $\theta$ lies in $V_\sigma\otimes V_{\bar\sigma}$. Its product with either extreme line vanishes for dimension reasons. This proves primitivity.

The primitive Hodge-star formula gives $*w=(-1)^nw$ for $w\in W_\R$, and consequently $B(w,w)=\int_A w\wedge*w>0$. This is the middle-degree Hodge--Riemann form, with the stated sign \cite[Theorem~5.42 and (5.43)]{Ballmann}. The cup product is integral on integral cohomology.

Use the translation-invariant K\"ahler metric determined by $\theta$. Harmonic forms are translation-invariant, so a vector $w$ of cohomological norm one has pointwise norm $v_\theta^{-1/2}$. Wirtinger's formula gives $\operatorname{vol}(Y)=\deg_\theta(Y)/n!$, including for singular $Y$ by integration on its regular locus. Therefore
\[
 \left|\int_Yw\right|\leq\frac{\deg_\theta Y}{n!\sqrt{v_\theta}}.
\]
The cup-product projection onto $W$ is the metric projection: for $w\in W$, its metric pairing with any middle-degree class $u$ is $\int_Au\wedge*w=B(u,w)$. Taking the supremum over unit $w$ proves the assertion for an effective cycle, and the triangle inequality proves it for a signed cycle.
\end{proof}

\begin{proof}[Proof of Theorem~\ref{thm:extraction}]
The trace of $a$ is integral, and hence $\bar a=\operatorname{Tr}_{K/\Q}(a)-a$ is also an endomorphism. On $H^{2n}$ set
\[
 S=a^*+\bar a^*,\qquad
 \beta_j=q^j\bigl(a^{2n-2j}+\bar a^{2n-2j}\bigr)
 \quad(0\leq j\leq n).
\]
Here $S$ denotes a sum of pullback operators. On the summand with $j$ factors from $V_{\bar\sigma}$ and $2n-j$ from $V_\sigma$, its eigenvalue is $\beta_j=\beta_{2n-j}$. These numbers are integers, and $|\beta_j|\leq2q^n$. If $z=a/\bar a$, equality of two distinct $\beta_j$, for $0\leq j\leq n$, would give $z^{j-k}=1$ or $z^{2n-j-k}=1$. This is excluded by the hypothesis. Thus
\[
 P(X)=\prod_{j=1}^n(X-\beta_j)\in\Z[X],\qquad
 P(S)=P(\beta_0)e_W,
\]
and $P(\beta_0)\ne0$. The middle eigenspaces are orthogonal to the extreme lines because a nonzero top wedge must have exactly $2n$ factors from each $K$-eigenspace. This verifies that this operator is the projector specified above, multiplied by a nonzero integer.

For a chosen signed cycle presentation put
\[
 \mathcal D(C)=\deg_\theta C^++\deg_\theta C^-.
\]
Pullback by $a$ preserves effectivity and multiplies the degree of an $n$-cycle by $q^n$: indeed $\deg(a)=q^{2n}$ and $a^*\theta=q\theta$. It follows that
\[
 \mathcal D((S-\beta_j)C)
 \leq(2q^n+|\beta_j|)\mathcal D(C)
 \leq4q^n\mathcal D(C).
\]
Choose an irreducible detector $Y$ of degree $d=\delta_W$. The integral signed cycle $C=P(S)Y$ has nonzero class $u\in W$ and has a presentation with
\[
 \mathcal D(C)\leq L:=4^nq^{n^2}d,
 \qquad \|u\|\leq L/n!.
\]
Its pullback $a^*C$ has mass at most $q^nL$ and class $v=a^*u$ with $\|v\|=q^n\|u\|$. The vectors $u,v$ are independent over $\Q$: the eigenvalues $a^{2n},\bar a^{2n}$ of $a^*|_W$ are nonreal conjugates, so this rational operator has no rational eigenline.

Let $\mathcal W=W\cap H^{2n}(A,\Z)$. It is a saturated rank-two lattice with integral positive Gram matrix. Its Euclidean covolume $c$ satisfies $c^2\in\Z_{>0}$, so $c\geq1$. Give $W_\R$ the orientation for which $\det(u,v)>0$, and let
\[
 m=[\mathcal W:\Z u+\Z v].
\]
Then $mc=\det(u,v)\leq\|u\|\|v\|$, and $m\mathcal W\subset\Z u+\Z v$. Therefore there exist integers $r,s$ such that
\[
 m\delta t=ru+sv.
\]
Taking determinants with $v$ and with $u$ gives
\[
 r=\frac{\delta\det(t,v)}c,\qquad
 s=\frac{\delta\det(u,t)}c,
\]
and hence
\[
 |r|\leq\delta\sqrt T\,\|v\|,
 \qquad |s|\leq\delta\sqrt T\,\|u\|.
\]
The signed cycle $rC+s a^*C$ represents $M t$, where $M=\delta m$ is a positive integer. Grouping positive and negative coefficients gives effective integral cycles satisfying
\[
 M\leq\frac{\delta q^nL^2}{(n!)^2},\qquad
 \deg Z^\pm\leq|r|L+|s|q^nL
 \leq\frac{2\delta\sqrt T\,q^nL^2}{n!}.
\]
Substitution of $L$ proves \eqref{eq:sharper-extraction}. Since $q\geq1$, it also proves \eqref{eq:original-extraction}.

Finally, if $Mt=\cl(Z^+)-\cl(Z^-)$, at least one irreducible component of one of these cycles has nonzero Weil projection. Its degree is no greater than that of the corresponding effective cycle. Thus $\delta_W\leq c_t$. If no detector exists, this observation excludes every finite representative.
\end{proof}

\begin{remark}[An endomorphism available in every order]
If $K=\Q(\sqrt{-m})$, with $m>0$ squarefree, and the endomorphism order contains $\Z+f\mathcal O_K$, one may take $a=f(2+\sqrt{-m})$. Then $q=f^2(m+4)$. Outside $\Q(i)$ and $\Q(\sqrt{-3})$, the only roots of unity in $K$ are $\pm1$, whereas $a/\bar a$ is neither. In the two exceptional fields its real part is respectively $3/5$ and $1/7$, so it is again not a root of unity.
\end{remark}

\subsection{Bounded detector loci}

\begin{proposition}\label{prop:detector-loci}
In a connected polarized family with a fixed integral quadratic action of signature $(n,n)$, the locus $\{s:\delta_W(s)\leq D\}$ is Zariski closed. The following are equivalent:
\begin{enumerate}[label=\textnormal{(\roman*)}]
\item every Weil class on every fibre is algebraic;
\item $\delta_W$ is bounded on every fibre;
\item $\delta_W$ is bounded on one Zariski-dense set.
\end{enumerate}
For a special Weil family these conditions are also equivalent to one uniform detector-degree bound at infinitely many good reductions of every CM fibre, allowing the auxiliary prime and geometric place to vary.
\end{proposition}

\begin{proof}
On the projective parameter space of effective cycles of degree at most $D$, the projected universal class $e_W\gamma$ is a section of the pulled-back Weil local system. Its nonvanishing locus is a union of connected components, hence is also closed and proper over the base. Its image is exactly the stated locus: if an effective sum detects, some irreducible component detects. This proves closedness.

If every Weil class is algebraic, every fibre has a detector. The closed loci therefore cover the complex base; Baire's theorem gives one bound. Conversely, a bound on a dense set gives a bound everywhere by closedness, and Theorem~\ref{thm:extraction} makes every Weil class algebraic. Alternatively, the integral Weil local system has finite monodromy, because it preserves the positive integral form $B$; after a finite \'etale cover it has a basis of global rational sections to which Theorem~\ref{thm:family} applies.

For the arithmetic assertion, extend the quadratic endomorphisms over an arithmetic open of each CM fibre. The projector is a rational polynomial in their pullbacks and commutes with specialization. The proof of Proposition~\ref{prop:lift} therefore lifts a bounded-degree special cycle with nonzero projected class to a bounded-degree characteristic-zero detector, at one place outside the finite exceptional set. Nonvanishing is preserved by the specialization isomorphism. In the reverse direction, descend and spread a characteristic-zero detector as in Proposition~\ref{prop:spread}. This gives detectors at all but finitely many good places.
\end{proof}

\subsection{An explicit Gaussian extraction constant}

\begin{proposition}\label{prop:gaussian-extraction}
On a Gaussian Weil fourfold with integral $\Z[i]$-action, a detector of degree $D$ gives
\[
 c_t\leq61\,440\,\delta\lceil\sqrt T\rceil D^2.
\]
\end{proposition}

\begin{proof}
Put $I=[i]^*$ and $J=[1+i]^*$. On the five summands of $H^4$ indexed by the number of factors from one $K$-eigenspace, direct substitution gives
\[
 (1+I)(4-J)=16e_W.
\]
For a detector $Y$, set $y=16e_W[Y]$. It is represented by
\[
 4Y+4[i]^*Y-[1+i]^*Y-[i(1+i)]^*Y,
\]
whose two effective parts each have degree $8D$. The norm lemma gives $d:=B(y,y)\leq64D^2$. The operator $A=[2+i]^*$ has eigenvalues $-7\pm24i$ on $W$. Hence $z=(A+7)y$ satisfies
\[
 B(y,z)=0,\qquad B(z,z)=576d,
\]
and has signed effective parts of degree at most $256D$. For the integers $b_1=B(\delta t,y)$ and $b_2=B(\delta t,z)$, orthogonal expansion gives
\[
 (576\delta d)t=576b_1y+b_2z.
\]
Write $t=x y/\sqrt d+x' z/(24\sqrt d)$, so $x^2+(x')^2=T$. The degree of each signed part is at most
\[
 \delta D\sqrt d\,(4608|x|+6144|x'|)
 \leq7680\delta\sqrt T D\sqrt d
 \leq61\,440\delta\sqrt T D^2.
\]
The multiplier is at most $36\,864\delta D^2$, which is bounded by the asserted expression.
\end{proof}

\section{Stabilization after Schoen}
\label{sec:stabilization}

This section explains how bounded detectors in the next split dimension control every determinant in the preceding dimension. Schoen's product argument is stated for fourfolds and sixfolds in \cite[Proposition~10, pp.~332--333]{Schoen98}; the general product and contraction statement is also recorded by Voisin \cite[Lemma~2.9]{VoisinB}. We include the degree estimate, for which the polarization and the chosen effective divisors matter.

\begin{theorem}[Qualitative stabilization]\label{thm:stabilization}
Fix an imaginary quadratic field $K$ and $n\geq1$. Suppose that the detector degree is bounded in every connected split Weil family of dimension $2n+2$. Then it is bounded in every connected Weil family of dimension $2n$, of every determinant. The bound may depend on the family and its polarization. No class-number hypothesis is required.
\end{theorem}

\begin{proof}
The determinant class of a fixed connected Weil family is constant. Write it as $(-1)^n d$ modulo norms, with $d=r/s>0$ and $r,s$ positive integers. On $B=E_K\times E_K$, where $E_K=\C/\mathcal O_K$, use the action $a\cdot(x,y)=(ax,\bar a y)$ and the polarization $\lambda=r\vartheta_1+s\vartheta_2$. Its Hermitian determinant is $-rs$ modulo norms, hence $-d$ because $rs=d s^2$. The product $A\times B$, with polarization $\Theta=\theta+\lambda$, has determinant $(-1)^{n+1}d^2$, so its Hermitian space is split by the rational classification \cite[Proposition~4.1 and Corollary~4.2]{Deligne82}. All this data is fixed in the family.

The extreme $K$-eigenlines give an inclusion
\[
 W_K(A\times B)\subset W_K(A)\otimes W_K(B).
\]
Choose an integral basis $\eta_1,\eta_2$ of $W_K(B)\cap H^2(B,\Z)$. The Lefschetz $(1,1)$-theorem makes them divisor classes. For a sufficiently large fixed $N$, the classes
\[
 d_0=N\lambda,\qquad d_1=N\lambda+\eta_1,
 \qquad d_2=N\lambda+\eta_2
\]
are very ample. Choose $b_B>0$ such that $b_B\lambda-d_j$ is nef for each $j$.

Let $Y$ be an effective detector on the split product, of degree at most $D$. Choose effective divisors $D_j$ in the classes $d_j$, translating them to meet $Y$ properly after pullback, and set
\[
 C_j=p_{A*}(Y\cdot p_B^*D_j).
\]
These are effective $n$-cycles. If all three had zero Weil projection, the projection formula and $d_i-d_0=\eta_i$ would make $[Y]$ orthogonal to every $w\otimes\eta_i$, with $w\in W_K(A)$. It would then be orthogonal to $W_K(A\times B)$, contradicting detection. Thus some $C_j$ detects. Mixed nef intersections give
\[
 \deg_\theta C_j
 \leq b_B\int_Yp_A^*\theta^n p_B^*\lambda
 \leq\frac{b_B}{n+1}\int_Y\Theta^{n+1}
 \leq\frac{b_BD}{n+1}.
\]
The factor $n+1$ is the coefficient of $\theta^n\lambda$ in the binomial expansion. Taking an irreducible detecting component proves a uniform bound on $A$.

If the hypothesis is formulated only for principal polarizations, pass to a fixed isogenous principally polarized split family. The choice of a maximal isotropic subgroup of the polarization kernel becomes constant after a finite \'etale cover. The rational $K$-action transports across the isogeny, and an order in $K$ acts integrally. For an isogeny $f:C\to C'$ with $f^*\theta'=k^2\Theta$, a middle-dimensional cycle satisfies
\[
 \deg_\Theta f^*Y'
 =\frac{\deg f}{k^{2n+2}}\deg_{\theta'}Y'.
\]
This is a constant factor in the family. Pullback preserves a nonzero Weil projection, so the preceding contraction applies. The finite cover is surjective, and therefore the bound descends as a fibrewise existence assertion.
\end{proof}

\begin{theorem}[An explicit contraction bound]\label{thm:explicit-stabilization}
Assume $\mathcal O_K=\Z[\tau]$ has class number one and acts integrally on a polarized Weil $2n$-fold $(A,L)$. Put $\Delta=\chi(L)$ and
$\kappa_K(\Delta)=\max\{1,\Nm(\tau)/\Delta\}$. If every split Weil $(2n+2)$-fold of polarization product $x$ has a detector of degree at most $D_{n+1}(x)$, then
\[
 \delta_W(A,L)\leq\frac{\kappa_K(\Delta)}{n+1}D_{n+1}(\Delta^2).
\]
For $K=\Q(i)$ or $\Q(\omega)$ choose $\tau=i$ or $\omega$, respectively; then $\kappa_K(\Delta)=1$. The premise ranges over all polarization types of the specified product, not only principal polarizations.
\end{theorem}

\begin{proof}
Class number one makes $H_1(A,\Z)$ free over $\mathcal O_K$. Normalize the Hermitian form associated to the alternating Riemann form $E$ by
\[
 h(x,y)=E(x,\tau y)-\bar\tau E(x,y).
\]
In an integral $\mathcal O_K$-basis it has determinant $(-1)^n\Delta$. Indeed for one diagonal block $h=(a)$, the alternating form on $(v,\tau v)$ has matrix $\left(\begin{smallmatrix}0&a\\-a&0\end{smallmatrix}\right)$. Both sides of $\det_\Z E=(\det_Kh)^2$ transform by the same norm factor under a change of $K$-basis. The sign is fixed by signature and $|\det_Kh|=\chi(L)$.

Take $B=E_K^2$ with balanced action and $\lambda=\vartheta_1+\Delta\vartheta_2$, where $\vartheta_1$ is the class of $\{0\}\times E_K$ and $\vartheta_2$ that of $E_K\times\{0\}$. Its determinant is $-\Delta$, so $A\times B$ is split and has polarization product $\Delta^2$. Let $Y$ be a detector there, of degree at most $D=D_{n+1}(\Delta^2)$ for $\Theta=\theta+\lambda$, and put
\[
 I=\int_Y\theta^n\lambda.
\]
All terms in the expansion of $\Theta^{n+1}$ have nonnegative intersection with $Y$, since both factors are nef. Hence $(n+1)I\le D$.

Intersect with general translates of the coordinate elliptic divisors and push to $A$, as in the preceding proof. Each contraction has degree at most $I$, since $\vartheta_j\le\lambda$ in the nef order. If either has nonzero Weil projection, an irreducible component gives the result.

Suppose both projected contractions vanish. For $u=1,\tau$ use the two graphs $\Gamma_u=\{x=u y\}$ and $\Gamma_{-u}$. Their classes satisfy
\begin{equation}\label{eq:opposite-graphs}
 [\Gamma_u]+[\Gamma_{-u}]
 =2(\vartheta_1+\Nm(u)\vartheta_2)
 \le2\kappa_K(\Delta)\lambda
\end{equation}
in the nef order. The equality follows either by adding the Hermitian matrices
$\left(\begin{smallmatrix}1&\mp u\\\mp\bar u&\Nm(u)\end{smallmatrix}\right)$ or from the equation $x\mp uy=0$. The axes and the graphs of $1,\tau$ span $\operatorname{NS}(B)_\Q$: the axes give the diagonal entries and $1,\tau$ give a rational basis for the off-diagonal entry. Their off-diagonal parts span $W_K(B)$.

Because $Y$ detects $W_K(A\times B)$, contraction against some element of $W_K(B)$ has nonzero projection to $W_K(A)$. Thus one of the two graph pairs has a detecting member. Its two projected contractions are opposite by \eqref{eq:opposite-graphs} and the vanishing of the axis contractions, so both detect. Choose the member of smaller degree. Its degree is at most the average, which by \eqref{eq:opposite-graphs} is at most $\kappa_K(\Delta)I$. One of its irreducible components detects and has no larger degree. This proves the bound.
\end{proof}

\begin{remark}[The polarization range]
Full integral maximal-order action and principal polarization over a class-number-one field force $\Delta=1$ and determinant $(-1)^n$, hence the split determinant class. Statements covering every determinant must therefore allow nonprincipal polarizations or rational actions transported through isogenies. The preceding theorem does so explicitly.
\end{remark}

\section{Small fourfold detectors after Schoen and van Geemen}
\label{sec:small}
This section computes degrees and integral coefficients for the geometric cycles of Schoen and van Geemen. The construction of the cycles is their input; the estimates keep track of the polarization, the isogeny and the saturated Weil lattice.

\begin{theorem}\label{thm:small-detectors}
Let $(A,\theta)$ be a principally polarized Weil fourfold with full integral maximal-order action. For $K=\Q(i)$ the detector degree is at most $48$ on the even Hermitian lattice $U\oplus U$, and at most $192$ on the odd lattice $I_{2,2}$. For $K=\Q(\sqrt{-3})$ it is at most $36$. In the last case, if $\mathcal W=W_K(A)\cap H^4(A,\Z)$, then $3\mathcal W$ consists of integral algebraic cycle classes. For $0\ne t\in W_K(A)$, $\delta\in\Z_{>0}$ with $\delta t$ integral, and $T=\int_A t^2$, one has
\begin{equation}\label{eq:eisenstein-cost}
 c_t\leq\max\left\{3\delta,
 36\left\lfloor\frac{\delta\sqrt{2T}}3\right\rfloor\right\}
 \leq17\delta\lceil\sqrt T\rceil.
\end{equation}
Here $U=\left(\begin{smallmatrix}0&1\\1&0\end{smallmatrix}\right)$ and $I_{r,s}$ is the diagonal Hermitian form with $r$ entries $1$ and $s$ entries $-1$.
\end{theorem}

\subsection{The Gaussian bound in two lattice types}

Van Geemen \cite[Proposition~3.6 and Theorem~3.7, pp.~623--625]{vG96} constructs a dominant rational map from a general member of the even principal Gaussian family to a smooth quadric fourfold in $\mathbf P^5$. The six defining sections belong to $|2\theta|$ and have finite base locus; the inverse image of a suitable ruling plane has nonzero Weil projection. This yields the degree bound with a short residual-intersection argument.

Choose such a plane $\Pi$ generally in its ruling, and choose a linear $\mathbf P^3$ through it. Its intersection with the quadric is the union of $\Pi$ and a residual ruling plane. The two pulled-back hyperplanes have no common divisor for general choices: away from the finite base locus this follows by avoiding the finitely many divisorial components contracted by the rational map. Their proper intersection is an effective surface cycle of class $4\theta^2$ and degree
\[
 \int_A4\theta^4=96.
\]
It contains a detecting component $Y$ from the inverse image of $\Pi$. If $\deg Y\leq48$, this is the desired detector. Otherwise subtract one copy of $Y$ from the complete-intersection cycle. The residual effective cycle has degree less than $48$ and has nonzero Weil projection, since the full class $4\theta^2$ has zero Weil projection. One of its components detects. Proposition~\ref{prop:detector-loci} extends the bound from the general member to every member of the even component.

For the odd lattice put
\[
 P=\begin{pmatrix}1&1\\1&-1\end{pmatrix},\qquad F=P\oplus P.
\]
Direct multiplication gives $P^*UP=2I_{1,1}$ and $\Nm(\det F)=16$. Transporting the complex structure by $F$ gives a $K$-linear isogeny
\[
 f:A_{\rm odd}\longrightarrow A_{\rm even},\qquad
 \deg f=16,\quad f^*\theta_{\rm even}=2\theta_{\rm odd}.
\]
For an effective surface cycle $Z$ on the target,
\[
 \deg_{\theta_{\rm odd}} f^*Z
 =\frac{16}{2^2}\deg_{\theta_{\rm even}}Z.
\]
Pullback preserves nonzero Weil projection, because it is a $K$-linear cohomology isomorphism. Pulling back an even detector and taking a detecting component proves the bound $192$.

Combining these bounds with Proposition~\ref{prop:gaussian-extraction} gives, respectively,
\[
 c_t\leq141\,557\,760\,\delta\lceil\sqrt T\rceil,
 \qquad
 c_t\leq2\,264\,924\,160\,\delta\lceil\sqrt T\rceil.
\]

\subsection{The Schoen surface and its principal quotient}

The published geometric input is the deformation of
$\Delta_J\cup(C\times C)\subset J\times J$, for a genus-two Jacobian $J$, to an embedded Schoen surface in a four-dimensional Eisenstein Weil family. We use the form of Schoen's construction \cite{Schoen07} stated by van Geemen \cite[Introduction, p.~1, and \S6.2, pp.~14--15]{vG22}, including the ambient family and its polarization. The embedded flat degeneration is also described in \cite[\S3, pp.~5--6]{Beauville}, and the four-dimensional smooth surface moduli in \cite[Theorem~1]{RRS}.

For clarity, the dimension of the surface moduli really gives dominance onto the Weil component. A smooth Schoen surface $S\subset A$ is Lagrangian for a nondegenerate holomorphic two-form and has $q(S)=4$ \cite[Introduction]{RRS}. Thus $N_{S/A}\simeq\Omega_S^1$ and $h^0(N_{S/A})=4$. Translation supplies four independent infinitesimal embedded deformations: the infinitesimal stabilizer is zero because a smooth surface of general type has no vector fields. Therefore a fixed abelian variety admits no positive-dimensional family of these embedded surfaces modulo translation. The map from the four-dimensional surface moduli to Albanese moduli has zero-dimensional fibres. Its image in the four-dimensional Weil component is dominant. This supplies an analytic open set of the component carrying the required effective surface class.

We compute its degree and Weil part at the degeneration. Let $\vartheta$ be the principal class on $J$, and on $J^2$ set
\[
 x=p_1^*\vartheta,\qquad y=p_2^*\vartheta,\qquad
 z=m^*\vartheta-x-y.
\]
The polarization $a x+b z+c y$ corresponds to the matrix
$\left(\begin{smallmatrix}a&b\\b&c\end{smallmatrix}\right)$ on the two factors. The determinant formula for its Euler characteristic, or expansion in a symplectic basis, gives
\begin{equation}\label{eq:intersection-polynomial}
 \int_{J^2}(a x+b z+c y)^4=24(ac-b^2)^2.
\end{equation}
In particular
\[
 \int x^2y^2=4,\qquad \int xyz^2=-4,\qquad
 \int z^4=24,
\]
and the other degree-four monomial intersections vanish. The central surface has class
\[
 s=[\Delta_J]+[C\times C]
   =\tfrac12(x+y-z)^2+xy.
\]
Take $\omega=(-1+\sqrt{-3})/2$ and the endomorphism and polarization
\[
 R=\begin{pmatrix}0&-1\\1&-1\end{pmatrix},\qquad
 H=2x-z+2y.
\]
The matrix of $H$ is $M=\left(\begin{smallmatrix}2&-1\\-1&2\end{smallmatrix}\right)$ and $R^tMR=M$. Its polarization type is $(1,1,3,3)$ and $H^4=216$. Equation~\eqref{eq:intersection-polynomial} gives
\[
 \int H^2s=36,\qquad \int s^2=8.
\]
For example, the two surface components contribute respectively $8$ and $28$ to the first number. The diagonal has self-intersection zero, the product has self-intersection $4$, and their mutual intersection is $2$, giving the second number.

Put $\ell_+=\omega^2x+\omega z+y$ and
$\ell_-=\omega x+\omega^2z+y$. Their squares span the two extreme $K$-eigenlines. Direct expansion gives
\[
 w_0=s-\tfrac16H^2
   =\tfrac16\bigl(\omega\ell_+^2+\omega^2\ell_-^2\bigr)\in W,
 \qquad \int w_0^2=2.
\]
The same flat class and norm hold on the deformed surface.

The polarization kernel is $\{(u,-u):u\in J[3]\}$, and $R$ acts trivially on it. A maximal isotropic subgroup of order $9$ therefore gives an Eisenstein-linear quotient
\[
 q:A_H\longrightarrow A,\qquad \deg q=9,\qquad q^*\theta=H,
\]
with principal polarization on $A$. The quotient can be chosen over a finite cover of the period family. To identify its principal lattice, specialize $J$ to $E_\omega^2$. The explicit map
\[
 F=\begin{pmatrix}1&\omega\\1&\omega^2\end{pmatrix}
       \otimes I_{E_\omega^2}
\]
has degree $9$, satisfies $F^*\theta_{E_\omega^4}=H$, and intertwines $R$ with $\operatorname{diag}(\omega,\omega^2)$. Its target is the balanced product $E_\omega^4$ with product polarization. Transporting this finite-index lattice along the period domain identifies the whole principal component.

This component covers every principal Eisenstein lattice of signature $(2,2)$. Indeed the maximal order is a principal ideal domain, and the associated integral Hermitian lattice is unimodular with determinant $1$. It is rationally split by the classification in \cite[Proposition~4.1 and Corollary~4.2]{Deligne82}. Choose a primitive integral isotropic vector $v$. Unimodularity gives $w$ with $h(v,w)=1$. If $h(w,w)=b\in\Z$, replacing $w$ by $w+b\omega v$ makes it isotropic because $\operatorname{Tr}(\omega)=-1$. Thus $v,w$ span a unimodular hyperbolic plane, which splits off integrally. Repeat on its signature-$(1,1)$ complement. The lattice is $U\oplus U$. Finally, in $I_{1,1}$ the vectors $(1,1)$ and $(1+\omega,\omega)$ have Gram matrix $U$ and determinant $-1$, so $U\oplus U\simeq I_{2,2}$.

Push the entire effective surface cycle forward, including its multiplicities, and write $Z=q_*S$. Then
\begin{equation}\label{eq:degree36-class}
 \deg_\theta Z=36,\qquad
 [Z]=\tfrac32\theta^2+w,\qquad
 w\in\mathcal W,\qquad \int_Aw^2=18.
\end{equation}
The degree follows by projection formula. For the norm use $q^*q_*s=9s$ in cohomology, so $(q_*s)^2=9s^2=72$, and subtract $(3\theta^2/2)^2=54$. The integrality of $w$ follows from that of $\theta^2/2$, which is immediate by expanding the square of an integral alternating two-form. Proper cycle parameter spaces extend an effective cycle with these same flat relations from the nonempty open set to every fibre: the components on which these relations hold are open and closed, and at least one dominates the base. Their proper image is all the base. This proves the degree-$36$ detector bound.

\subsection{The saturated Eisenstein lattice}

It remains to turn the effective cycle into the integral coefficient bound. Compute at the balanced product $E_\omega^4$. Choose integral one-form bases $x_j,y_j$, adapted to the maximal-order action, and write
\[
 \prod_{j=1}^4(x_j+\omega y_j)=u+\omega v.
\]
The all-$x$ coefficient of $(u,v)$ is $(1,0)$, and any coefficient with exactly one $y$ is $(0,1)$. Consequently a rational combination of $u,v$ is integral only when both coefficients are integers. They are a basis of the saturated Weil lattice. The identities
\[
 (u+\omega v)^2=0,\qquad
 \int(u+\omega v)(u+\omega^2v)=9
\]
give
\begin{equation}\label{eq:eisenstein-gram}
 \operatorname{Gram}(\mathcal W)=
 \begin{pmatrix}6&3\\3&6\end{pmatrix}.
\end{equation}
These are integral local systems, so the result holds throughout the family. If $\rho$ is the automorphism induced by $\omega$, then
\[
 \rho^*u=-v,\qquad \rho^*v=u-v.
\]
The norm-$18$ vectors in this lattice have coordinates
$\pm(1,1),\pm(2,-1),\pm(-1,2)$: this follows by solving
$m^2+mn+n^2=3$. The three unit translates of any such vector have differences giving, up to signs and a lattice isometry, the three directions $3u,3v,3(u-v)$. Applied to the Weil part in \eqref{eq:degree36-class}, the differences of the three degree-$36$ cycles
$Z,\rho^*Z,\rho^{2*}Z$ therefore generate exactly $3\mathcal W$.

Write $\delta t=mu+nv$. It is a sum of
\[
 N=\max\{|m|,|n|,|m+n|\}
\]
signed steps in the directions $u,v,u-v$: for equal signs use $u$ and $v$, and for opposite signs first use $u-v$ and then the remaining direction. Put $Q=m^2+mn+n^2=\delta^2T/6$. The three identities
\[
 4Q-3m^2=(m+2n)^2,\quad
 4Q-3n^2=(2m+n)^2,\quad
 4Q-3(m+n)^2=(m-n)^2
\]
show that $N\leq\delta\sqrt{2T}/3$. Each step, multiplied by $3$, is represented by a difference of two degree-$36$ effective cycles. Hence $3\delta t$ has effective numerator degrees at most $36N$. This proves the first inequality in \eqref{eq:eisenstein-cost}; the second follows from $12\sqrt2<17$ and $\lceil\sqrt T\rceil\geq1$.

\begin{example}
For the integral class $t=u$ in \eqref{eq:eisenstein-gram}, one has $T=6$ and may take $\delta=1$. A single unit-orbit difference represents $3u$, with both effective degrees $36$. The floor in \eqref{eq:eisenstein-cost} is then $\lfloor\sqrt{12}/3\rfloor=1$, exactly reproducing this presentation.
\end{example}

\section{Principal hyperbolic quotients over every imaginary quadratic field}
\label{sec:principal-quotients}
This section reduces all split integral polarization types in a fixed dimension to one principal hyperbolic family, with a degree cost linear in the polarization volume.

Let $K$ be an imaginary quadratic field, let $\mathcal O=\mathcal O_K$, and define $t_K\in\{1,2\}$ by
\[
 \operatorname{Tr}_{K/\Q}(\mathcal O)=t_K\Z.
\]
Thus $t_K=1$ when the discriminant is odd and $t_K=2$ when it is even. For a projective $\mathcal O$-module $Q$, write
\[
 Q^\vee=\operatorname{Hom}_{\bar{\mathcal O}}(Q,\mathcal O).
\]
Its elements are additive functions satisfying $p(aq)=\bar a p(q)$. The hyperbolic Hermitian module $H(Q)$ is $Q^\vee\oplus Q$, with pairing
\[
 ((p,q),(p',q'))\longmapsto p(q')+\overline{p'(q)}.
\]
We use Hermitian forms linear in the first variable.

\begin{theorem}\label{thm:principal-quotient}
Let $(A,L)$ be a split Weil $2n$-fold over $K$, with compatible polarization and integral action by $\mathcal O_f=\Z+f\mathcal O$. Put $\Delta=\chi(L)$. There is a $K$-linear isogeny to a principally polarized split Weil variety,
\[
 q:A\longrightarrow(P,\Theta),\qquad
 q^*\Theta=t_K fL,\qquad
 \deg q=(t_K f)^{2n}\Delta,
\]
where equality of polarizations means equality of their first Chern classes. The integral Hermitian lattice of $P$ is $H(Q)$ for a projective rank-$n$ module $Q$.

At most $h_K$ isomorphism types of such hyperbolic modules occur, where $h_K$ is the ideal class number. If $C_I$ bounds detectors throughout the principal family associated with $H(\mathcal O^{n-1}\oplus I)$, then
\[
 \delta_W(A,L)\le t_K^n f^n\Delta\,
       \max_{[I]\in\operatorname{Cl}(\mathcal O)} C_I.
\]
The same conclusion on CM fibres requires the bounds only on the CM points of those principal families. If $\mathcal O$ has class number one, only the lattice $U^n$ is needed.

More generally, put
\[
 m_K=\max_{[I]\in\operatorname{Cl}(\mathcal O)}
       \min_{\substack{J\subset\mathcal O\\{[J]=[I]}}}
                    \operatorname{Nm}(J).
\]
If $C_U$ bounds detectors on the single principal family with lattice $U^n$, then
\[
 \delta_W(A,L)\le(t_Kfm_K)^n\Delta C_U,
 \qquad m_K\le\frac2\pi\sqrt{|\operatorname{disc}K|}.
\]
Again the CM version holds with a supplier only on CM points.
\end{theorem}

\begin{proof}
Write $A=V_{\R}/\Lambda$, with Riemann form $E$, and choose $\tau$ such that $\mathcal O=\Z[\tau]$. Set $S=\mathcal O\Lambda=\Lambda+\tau\Lambda$ and $b=[S:\Lambda]$. The form $E'=fE$ is integral on $S$. Indeed the cross terms use the integral endomorphism $f\tau$, and Rosati compatibility gives
\[
 E(\tau x,\tau y)=\operatorname{Nm}(\tau)E(x,y).
\]
The quotient with lattice $S$ therefore has polarization volume
\[
 \Delta'=\frac{f^{2n}\Delta}{b}.
\]
On $S$ use the integral Hermitian form
\[
 H(x,y)=E'(\tau x,y)-\bar\tau E'(x,y),
 \qquad
 E'(x,y)=\operatorname{Tr}_{K/\Q}
          \frac{H(x,y)}{\tau-\bar\tau}.
\]
The first formula is linear in $x$ and conjugate-linear in $y$. The codifferent of $\mathcal O$ is $(\tau-\bar\tau)^{-1}\mathcal O$. Consequently a perfect integral Hermitian pairing gives a perfect alternating $\Z$-pairing by the second formula. This assertion holds for projective modules and requires no choice of a free basis.

Choose a maximal isotropic $K$-subspace $W$ of dimension $n$. The module $P_0=S\cap W$ is saturated: if a nonzero $a\in\mathcal O$ has $ax\in W$, then $x\in W$. The quotient $S/P_0$ is torsion-free over a Dedekind domain, hence projective. Choose a splitting
\[
 S=P_0\oplus Q.
\]
Nondegeneracy identifies $W$ with the conjugate dual of $Q_K$ by $w\mapsto H(w,\cdot)|_Q$. Let $P'$ be the inverse image of $Q^\vee$. Integrality gives $P_0\subset P'$. On $\Gamma=P'\oplus Q$, the form becomes
\[
 H(p+q,p'+q')=p(q')+\overline{p'(q)}+C(q,q'),
 \qquad C=H|_{Q\times Q}.
\]
The off-diagonal evaluation pairing is perfect, so this Hermitian form is unimodular. Its alternating form is therefore principal. Comparing the polarization volumes on $S$ and $\Gamma$ gives
\[
 [\Gamma:S]=\Delta'.
\]

Choose $\eta\in\mathcal O$ with $\eta+\bar\eta=t_K$, and let $c:Q\to Q^\vee$ be the $\mathcal O$-linear map $c(q)(q')=C(q,q')$. Enlarge the lattice once more, to
\[
 \Gamma_t=\frac1{t_K}P'\oplus Q,
\]
and give it the form $t_KH$. The map
\[
 Q^\vee\oplus Q\longrightarrow\Gamma_t,
 \qquad (p,q)\longmapsto\frac p{t_K}+q-
                                  \frac{\eta c(q)}{t_K}
\]
is an integral module isomorphism. Its off-diagonal pairing is evaluation; its pairing on the second summand is
$t_KC-\eta C-\bar\eta C=0$. It is thus an isometry from $H(Q)$.
In particular this is a principal hyperbolic lattice and
\[
 [\Gamma_t:S]=t_K^{2n}\Delta'.
\]
All these lattices lie in the same real vector space with the same complex structure. The Riemann form has only been multiplied by a positive integer, so positivity is preserved. The inclusion $\Lambda\subset\Gamma_t$ defines a holomorphic $K$-linear isogeny with
\[
 q^*\Theta=t_K fL,
 \qquad \deg q=b\,t_K^{2n}\Delta'=(t_K f)^{2n}\Delta.
\]

For an effective detecting $n$-cycle $Z$ on $P$, finite flat pullback and the projection formula give
\[
 \deg_L(q^*Z)=\frac{\deg q}{(t_K f)^n}\deg_\Theta Z
             =t_K^n f^n\Delta\deg_\Theta Z.
\]
The pullback has nonzero Weil projection because a $K$-linear isogeny identifies the rational Weil spaces. Some irreducible component is a detector, and its degree is at most the total degree.

Finally every projective rank-$n$ module over $\mathcal O$ has Steinitz form $\mathcal O^{n-1}\oplus I$. Isomorphic modules give isometric hyperbolic modules, so at most $h_K$ types are required. For each type the compatible period domain is connected; a neat arithmetic quotient gives its principal family. Passing to a finite level cover does not affect the bound.

Choose the representative $I$ integral, of norm $m$. In coordinates,
\[
 Q=\mathcal O^{n-1}\oplus I,
 \qquad Q^\vee=\mathcal O^{n-1}\oplus\bar I^{-1},
 \qquad m\bar I^{-1}=I.
\]
Thus $(p,q)\mapsto(mp,q)$ embeds $H(Q)$ into the standard lattice $\mathcal O^{2n}$ and multiplies its Hermitian form by $m$. Transporting the complex structure through this map gives an isogeny $P_I\to P_U$ with pullback polarization $m\Theta_I$. Since both polarizations are principal, its degree is $m^{2n}$. Pullback of a middle-dimensional detector therefore costs $m^n$. Taking the smallest norm in each ideal class proves the bound involving $m_K$. All isogenies here preserve CM.

For the numerical estimate, an integral ideal $J$ embedded in $\C$ has covolume $\operatorname{Nm}(J)\sqrt{|\operatorname{disc}K|}/2$. The planar convex-body theorem applied to a disk gives $0\ne\alpha\in J$ with
\[
 |\operatorname{Nm}(\alpha)|\le
 \frac2\pi\sqrt{|\operatorname{disc}K|}\operatorname{Nm}(J).
\]
Indeed a centrally symmetric convex set of area exceeding four times the covolume contains a nonzero lattice point: two points of its half-sized copy have the same image in a lattice fundamental domain, and their difference lies in the original set. Apply this to disks with radii decreasing to the stated bound. The ideal $(\alpha)J^{-1}$ is integral, lies in the inverse class of $J$, and has norm at most $2\sqrt{|\operatorname{disc}K|}/\pi$. This proves the estimate for every class.
\end{proof}

For a free maximal-order lattice with $t_K=2$, suppose in addition that its Hermitian form is even. The restriction $C$ in the proof has even diagonal. In a basis of $Q$ there is then an integral map $R:Q\to Q^\vee$ with $C(q,q')=R(q)(q')+\overline{R(q')(q)}$: choose the diagonal entries by the trace equation and choose one of each pair of off-diagonal entries to be zero. The shear $q\mapsto q-R(q)$ already makes the lattice $\Gamma$ hyperbolic, without the final scaling. The same proof gives the improved cost $f^n\Delta$ whenever the saturated maximal-order lattice with polarization $fE$ is free and even.

The theorem removes dependence on arbitrary integral polarization types and ideal classes from the existence question. It leaves one principal hyperbolic supplier for each field and dimension. It does not construct that supplier.

\section{Sixfold cycles after Schoen and Koike, and bounds for every fourfold determinant}
\label{sec:sixfolds}
This section keeps track of the degrees in Schoen's cyclic-cover construction and improves the passage between integral polarization lattices.

Write $U=\left(\begin{smallmatrix}0&1\\1&0\end{smallmatrix}\right)$ for a Hermitian hyperbolic plane. Degrees are unnormalized intersection numbers. All polarizations below are compatible with the indicated quadratic action.

\begin{theorem}\label{thm:sixfold-suppliers}
For full maximal-order action, the following detector bounds hold on the indicated principal split sixfold families:
\[
\begin{array}{lll}
K&\text{Hermitian lattice}&\delta_W\text{ at most}\\
\hline
\Q(i)&U^3&746496\\
\Q(i)&I_{3,3}&5971968\\
\Q(\omega)&U^3\cong I_{3,3}&30375.
\end{array}
\]
For an arbitrary full maximal-order split sixfold of polarization product $\Delta$, one has
\[
 \delta_W\le5971968\Delta\quad(K=\Q(i)),
 \qquad \delta_W\le30375\Delta\quad(K=\Q(\omega)).
\]
For an even Gaussian Hermitian lattice the first bound improves to $746496\Delta$.
\end{theorem}

\begin{theorem}\label{thm:all-determinant-fourfolds}
Let $(A,L)$ be a Weil fourfold over $K$, of any determinant, and put $\Delta=\chi(L)$. Suppose the order $\mathcal O_f=\Z+f\mathcal O_K$ acts integrally. Then
\[
\begin{aligned}
 \delta_W(A,L)&\le1990656\,f^6\Delta^2
 &&\text{if }K=\Q(i),\\
 \delta_W(A,L)&\le10125\,f^6\Delta^2
 &&\text{if }K=\Q(\omega).
\end{aligned}
\]
Each right-hand side may be divided by $d$, where
$d=[\mathcal O_KH_1(A,\Z):H_1(A,\Z)]$.
\end{theorem}

The geometric inputs are published: Schoen constructs the cycles \cite[Theorem~2.0 and \S3]{Schoen88}, and Schoen and Koike prove the required dominance \cite[Lemma~13, pp.~334--335]{Schoen98}, \cite[Theorem~2.1, p.~568]{Koike}. The estimates below follow from their constructions, integral polarization calculations and closedness of bounded detector loci. They do not use the recent CM theorems or equivariant semiregularity.

\subsection{The cyclic sources and their principal polarizations}

Let $C\to X$ be a connected cyclic \'etale cover of degree $m=3$ or $4$, with $g(X)=4$ and generator $\sigma$. Then $g(C)=3m+1$. Choose a topological symplectic basis on $X$ for which all monodromy lies on the first handle. The other three handles lift in $m$ disjoint copies. Their cycles satisfy
\[
 E(a_{j,r},b_{k,s})=\delta_{jk}\delta_{rs},
 \qquad E(a_{j,r},a_{k,s})=E(b_{j,r},b_{k,s})=0,
\]
and $\sigma$ shifts the sheet index. The cover of the first punctured handle is a genus-one surface with $m$ boundary circles. After the other handles are glued, its boundary classes vanish and it contributes an invariant rank-two symplectic summand. This gives an integral orthogonal decomposition into that summand and three pairs of regular representations.

For $m=4$, the saturated primitive lattice on each handle is generated by
\[
 a_0-a_2,\quad a_1-a_3,\quad b_0-b_2,\quad b_1-b_3.
\]
The two nonzero alternating pairings are both two. Thus on the primitive sixfold $B=\ker(1+\sigma^2)^0$ the restriction of the Jacobian polarization is $2\Xi$, where $\Xi$ is principal. The Gaussian Hermitian matrix for $\Xi$ on a handle has off-diagonal entries $i,-i$; changing a basis vector by a unit gives $U$. The principal lattice is $U^3$.

For $m=3$, the norm Prym $P$ has, on each handle, augmentation generators $A=a_0-a_1$ and $B=b_0-b_1$ over $\Z[\omega]$. The form
$h(x,y)=E(x,\omega y)-\bar\omega E(x,y)$ has
\[
 E(A,B)=2,\quad E(A,\omega B)=-1,
 \qquad h(A,B)=1+2\omega=\sqrt{-3}.
\]
Its matrix is $\left(\begin{smallmatrix}0&\sqrt{-3}\\-\sqrt{-3}&0\end{smallmatrix}\right)$ and its alternating type is $(1,3)$. Enlarge the lattice on each handle from $\langle A,B\rangle$ to $\langle A,B/\sqrt{-3}\rangle$. This has index three and gives matrix $U$. On three handles it defines a degree-$27$ quotient $P\to P'$ with a principal polarization $\Xi'$ satisfying
\[
 \varphi^*\Xi'=\Theta_C|_P.
\]
The positive Riemann form on the real vector space has not changed. A finite level cover chooses the kernels in families and identifies the period domains. The Eisenstein basis with columns $(1,-\omega)$ and $(1,\bar\omega)$ takes $U$ to $\operatorname{diag}(1,-1)$ and has unit determinant. Thus the principal lattice is also $I_{3,3}$.

Koike's Theorem~2.1 states that the quartic Prym map is dominant onto the nine-dimensional unitary domain; its codifferential is multiplication of inverse primitive eigendifferentials. Schoen's Lemma~13 proves the cubic dominance by exhibiting nine independent products. These statements apply to the principalized sources just described, since a fixed isogeny change of lattice does not change the period domain.

\subsection{An effective detecting component and its degree}

Put
\[
 N=\ker\bigl((\Z/m)^6\xrightarrow{\sum}\Z/m\bigr),
 \qquad G_0=N\rtimes S_6,
 \qquad W_0=C^6/G_0.
\]
The finite maps
\[
 C^6\xrightarrow{r}C^{(6)}\xrightarrow{p}W_0
 \xrightarrow{\beta}X^{(6)}
\]
have generic degrees $6!,m^5,m$. Schoen's local equation \cite[Lemma~2.2, p.~13]{Schoen88} makes $\beta$ \'etale in the unramified case. The inverse image of $|K_X|\cong\mathbf P^3$ is a disjoint union of $m$ copies of $\mathbf P^3$. Choose one, $Q$.

The nonvanishing is especially simple here. Embed $X^{(6)}$ in $X^{(7)}$ by adding a point. Since $X^{(7)}\to\operatorname{Pic}^7(X)$ is a projective bundle, the canonical $\mathbf P^3$, after this embedding, is a fibre. Its normal bundle there is trivial of rank four; the normal line of $X^{(6)}$ restricts to $\mathcal O_{\mathbf P^3}(1)$. The normal sequence gives
\[
 c(N_{\mathbf P^3/X^{(6)}})=(1+H)^{-1},
 \qquad c_3(N_{\mathbf P^3/X^{(6)}})=-H^3\ne0.
\]
The same holds for $Q\subset W_0$. For a primitive character $\chi$ of the residual deck group, the cycle
$\sum_t\chi(-t)\,\sigma^tQ$ restricts on $Q$ to this nonzero top Chern class, since the other sheets are disjoint. Thus the primitive-character projection is nonzero. This is the unramified calculation of \cite[Lemma~2.6, p.~16]{Schoen88}.

Let $Z$ be an irreducible component of the inverse image of $Q$ in $C^6$ and let $W=r(Z)$ with reduced structure. A general canonical divisor of $X$ has six distinct zeros, so the relevant action is free at the generic point above it. If $H_Z\subset G_0$ stabilizes $Z$, then
\[
 \deg(Z/Q)=|H_Z|,
 \quad \deg(Z/W)=|H_Z\cap S_6|,
 \quad e:=\deg(W/Q)\le m^5.
\]
The primitive determinant plane $U_C\subset H^1(C)^{\otimes6}$ is $G_0$-invariant: the sum-zero sheet translations act trivially, and the alternating determinant and the cohomological permutation signs cancel. Its primitive-character components are those just detected on $W_0$. For $q:C^6\to W_0$, the generic cycle identity is
\[
 \sum_{g\in G_0}g_*[Z]=\deg(Z/Q)q^*[Q].
\]
Pair with an element of $U_C$ for which the right side is nonzero. All summands on the left have the same pairing, so $Z$ itself detects. This produces an effective integral component, without interpreting cyclotomic coefficients as effective cycles.

Let $\xi$ be the incidence class on $C^{(6)}$, and let $\alpha:C^{(6)}\to J(C)$ be an Abel map. Pullback of an incidence hyperplane of $|K_X|$ has class $m\xi$. Consequently
\begin{equation}\label{eq:cyclic-incidence}
 m^3\int_W\xi^3=e,\qquad \int_W\xi^3\le m^2.
\end{equation}
Write $\vartheta=\alpha^*\Theta_C$ and let $\mathcal D$ be the reduced big diagonal of $C^{(6)}$. Then
\begin{equation}\label{eq:symmetric-theta}
 \vartheta=(g(C)+5)\xi-\tfrac12[\mathcal D].
\end{equation}
To verify this, pull back to $C^6$. The incidence class becomes the sum of the point classes. The pair diagonal is the sum of two point classes minus the degree-one cross pairing; the Abel sum pulls theta back to $g(C)$ times the sum of point classes plus these cross pairings. Sum over the fifteen pairs and use $r^*\mathcal D=2\sum_{i<j}\Delta_{ij}$. Rational pullback is injective for the finite quotient.

The class $\xi$ is ample because its finite pullback is ample, and $\vartheta$ is nef. Since $W$ is not contained in $\mathcal D$, the restricted diagonal is effective. With $b=g(C)+5$,
\[
 (b\xi)^3-\vartheta^3
 =\tfrac12[\mathcal D](b^2\xi^2+b\xi\vartheta+\vartheta^2)
\]
has nonnegative intersection with $W$, computed on a resolution if necessary. This uses effectivity of the restricted diagonal, not its nefness. Hence
\begin{equation}\label{eq:cyclic-theta-degree}
 \int_W\vartheta^3\le b^3m^2.
\end{equation}

For $m=4$ use $\phi=(1-\sigma^2)\alpha$ with image in $B$. Its pullback of the Weil plane to $C^6$ is $U_C$, up to nonzero scalars. Indeed expand the exterior product of six primitive one-forms under the Abel sum. Two forms of the same primitive character on one factor multiply to zero, since their product has nontrivial character whereas $H^2(C)$ has trivial character. The surviving alternating tensor is nonzero; $1-\sigma^2$ is invertible on the primitive eigenspaces. Projection formula shows that $\phi(W)$ has dimension three and detects.

The theorem of the square and $\Theta_C|_B=2\Xi$ give
\[
 (1-\sigma^2)^*\Xi
 =2\Theta_C-\tfrac12(1+\sigma^2)^*\Theta_C
 \le2\Theta_C.
\]
Thus its effective pushforward has degree at most
\[
 2^3\cdot18^3\cdot4^2=746496.
\]
For $m=3$, $\alpha(W)$ lies in a translate of the norm Prym. The same argument using the invertible endomorphism $1-\sigma$ proves detection, and its restricted-theta degree is at most
\[
 15^3\cdot3^2=30375.
\]
The principal quotient has pullback polarization equal to that restricted theta, so adds no degree factor. Selecting a reduced image can only decrease either bound.

The bounds are uniform on the dense cyclic-source loci. Bounded-detector closedness extends them to the whole principal components. For the odd Gaussian bound, the three-block matrix with blocks $M=\left(\begin{smallmatrix}1&1\\1&-1\end{smallmatrix}\right)$ satisfies $M^*I_{1,1}M=2U$. It gives an isogeny from the even principal family to the odd one with polarization multiplier two. Threefold degrees increase by at most eight, giving $5971968$.

These are all Gaussian principal parity types. Split off a primitive isotropic vector and a vector pairing to one in a unimodular lattice. The partner norm reduces modulo two. An even norm gives $U$ after subtracting an integral multiple of the isotropic vector; an odd norm gives $I_{1,1}$. Repeat. A mixture containing an odd block becomes $I_{3,3}$: in $\langle1\rangle\oplus U$, with basis $e,a,b$, the vectors $e+a,e-b,e+a-b$ are a unimodular orthogonal basis of norms $1,1,-1$. This converts each remaining hyperbolic block. The rational isotropic spaces exist because the Hermitian space is split.

Theorem~\ref{thm:principal-quotient}, with $n=3$, transfers the principal bounds to every split polarization: $t_K^3\Delta$ is $8\Delta$ for the Gaussian ring and $\Delta$ for the Eisenstein ring. In the even Gaussian case the final paragraph of Section~\ref{sec:principal-quotients} gives the factor $\Delta$. This finishes Theorem~\ref{thm:sixfold-suppliers}.

\subsection{All determinants and nonmaximal orders}

\begin{proof}[Proof of Theorem~\ref{thm:all-determinant-fourfolds}]
First suppose the full maximal order acts. The balanced surface of Theorem~\ref{thm:explicit-stabilization} makes $A\times E_K^2$ a split sixfold of polarization product $\Delta^2$. In both fields $\kappa_K(\Delta)=1$, so the sixfold bound and the factor $1/3$ give
\[
 \delta_W(A,L)\le\frac{5971968\Delta^2}{3}
 =1990656\Delta^2
\]
in the Gaussian case, and $30375\Delta^2/3=10125\Delta^2$ in the Eisenstein case. The Gaussian product contains a norm-one vector from the surface, so no even-lattice improvement is used.

For conductor $f$, put $\Lambda'=\mathcal O_K\Lambda=\Lambda+\tau\Lambda$ and $d=[\Lambda':\Lambda]$. The form $fE$ is integral on $\Lambda'$: cross terms use the integral endomorphism $f\tau$, while
$E(\tau x,\tau y)=\Nm(\tau)E(x,y)$ follows from Rosati compatibility. The quotient isogeny $q:A\to A'$ has a compatible polarization with
\[
 q^*\theta'=f\theta,
 \qquad \Delta'=f^4\Delta/d.
\]
Apply the maximal-order bound $C(\Delta')^2$, where $C$ is the appropriate constant. Finite flat pullback preserves integral effectivity and a nonzero Weil projection, and multiplies the surface degree by $d/f^2$. Selecting a detecting irreducible component gives
\[
 \delta_W(A,L)\le\frac d{f^2}C\left(\frac{f^4\Delta}{d}\right)^2
 =C\frac{f^6}{d}\Delta^2.
\]
\end{proof}

The principal quotient improves transfer once a principal detector is known; it supplies no principal detector in a new dimension. Cyclic \'etale covers of genus-five curves have twelve moduli, whereas the signature-$(4,4)$ unitary domain has dimension sixteen. The genus-four dominance used here therefore does not extend to eightfolds by the same unramified cyclic construction.

\section{An explicit sixfold bound over every imaginary quadratic field, after Markman and Perry}
\label{sec:effective-sixfolds}
This section makes the ideal-curve and Fourier-transform construction quantitative over every imaginary quadratic field. The geometric design is that of Markman's detailed preprint \cite[\S\S8--9]{MarkmanSecant}; its qualitative consequence is also stated in his ICM account \cite[Theorem~1.2]{MarkmanICM}. The deformation input is Perry's Theorem~1.2 \cite{Perry}, an unrefereed preprint: a perfect complex with no negative self-Ext, a coherent finite translation--tensor action, injective semiregularity on invariant obstructions, and an algebraic rational $B$-field whose entire twisted character stays Hodge deforms locally as a twisted perfect complex. We use that theorem with its hypotheses as stated. The estimates below do not use either recent CM Hodge theorem.

\begin{theorem}[With Perry's semiregularity theorem]\label{thm:all-field-sixfolds}
Let $d$ be a positive squarefree integer, let $K=\Q(\sqrt{-d})$, and let $(A,L)$ be a split Weil sixfold with compatible polarization, integral $\mathcal O_f$-action and $\Delta=\chi(L)$. With $m_K$ as in Theorem~\ref{thm:principal-quotient},
\[
 \delta_W(A,L)\le8m_K^3 C_{4d} f^3\Delta
             <10^{179}d^{53}f^3\Delta.
\]
The positive integer $C_D$ is defined explicitly in \eqref{eq:effective-degree} below. When $K$ has class number one, the sharper bound is
\[
 \delta_W(A,L)\le8C_{4d}f^3\Delta
                  <10^{178}d^{51}f^3\Delta.
\]
The detecting cycle may be chosen reduced and irreducible.
\end{theorem}

We prove the theorem in four steps. First a genus-three curve supplies quotient ideals whose obstruction trace is injective. Two distinct elliptic directions make their Fourier transform, after duality, a coherent sheaf of projective dimension at most one. An explicit regularity estimate then turns its local deformation into a globally generated vector bundle, whose third Chern class is an effective detector. Finally a fixed isogeny places this local construction in the principal hyperbolic family, and Theorem~\ref{thm:principal-quotient} transports it to every target lattice.

\subsection{The Fermat curve and its quotient ideals}
Let $C\subset\mathbf P^2$ be $x^4+y^4+z^4=0$. It is a smooth nonhyperelliptic genus-three curve. The three sign involutions give maps $f_j:C\to E_i$ of degree two. Each has four fixed points, so its quotient has genus one. One quotient is $u^2=-(t^4+1)$; the automorphism $(t,u)\mapsto(it,u)$ fixes a point and acts by $i$ on its differential, identifying its Jacobian with $E_i$. The Klein-four quotient of $C$ is the conic $X^2+Y^2+Z^2=0$, so the three nontrivial characters exhaust $H^1(C,\Q)$.

Put $X=E_i^3$ and let $H$ be its product principal polarization. The map
\[
 q:X\longrightarrow J(C),\qquad q(z_1,z_2,z_3)=\sum_j f_j^*z_j
\]
satisfies $q^*\Theta=2H$, $\deg q=8$ and $\ker q\subset X[2]$. Indeed the cross norms $f_{j,*}f_k^*$ vanish for distinct Klein characters and the diagonal norms are $[2]$. Thus the polarization matrix is diagonal $[2]$; Riemann--Roch gives the degree, and the norm identity gives the kernel assertion.

Each $f_j^*$ is injective. A nonzero kernel point would be two-torsion and would factor $C\to E_i$ through a connected \'etale double cover, forcing the remaining map to have degree one and contradicting $g(C)=3$. For an Abel image of $C$, one has the stronger identity
\begin{equation}\label{eq:bielliptic-disjointness}
 f_j^*\operatorname{Pic}^0(E_i)\cap(C-C)=\{0\}.
\end{equation}
To prove it, write $p-q=f_j^*M$ and let $\sigma$ be the involution. Then $p+\sigma q\sim q+\sigma p$. On a nonhyperelliptic curve equal degree-two divisor classes have equal effective representatives. If $p\ne q$, equality of these unordered pairs forces both points to be fixed by $\sigma$. The natural linearization on $\mathcal O_C(p-q)$ has fibre sign $-1$ at these two fixed points and $+1$ at the other two: a local parameter and its inverse both change sign. Any other linearization differs by one global character. Neither choice has trivial inertia at all four points, whereas a pullback from the quotient does. This proves \eqref{eq:bielliptic-disjointness}.

Set $D=4d$ and $k=D+1$. For $j=1,2$, choose a point
$\tau_j\in E_i[k]$ of exact order $k$, and let
$u_j:J(C)\to Y_j=J(C)/\langle f_j^*\tau_j\rangle$ be the quotient.
It embeds $C$ as a smooth curve $\bar C_j\subset Y_j$, because its
translates are disjoint by \eqref{eq:bielliptic-disjointness}.
Let $h_j\in\operatorname{NS}(Y_j)_{\Q}$ be the rational polarization
class with $u_j^*h_j=\Theta$. Then
\[
 \operatorname{ch}(I_{\bar C_j})=1-\frac{k}{2}h_j^2+\frac{k}{3}h_j^3.
\]
This follows by pulling back to the disjoint union of $k$ Abel curves, using $[C]=\Theta^2/2$ and $\chi(\mathcal O_C)=-2$.

For completeness, we check the whole obstruction trace. The restriction $H^1(\mathcal O_{Y_j})\to H^1(\mathcal O_C)$ is an isomorphism. In the normal sequence, the dual of $H^1(T_C)\to H^1(T_{Y_j}|_C)$ contains
$\operatorname{Sym}^2H^0(K_C)\to H^0(K_C^2)$. This map is an isomorphism: both dimensions are six and a plane quartic has no quadratic equation. Hence $h^0(N_{C/Y_j})=3$.
The ideal sequence and Serre duality bound $\operatorname{ext}^1(I,I)$ above by $3+3$. Translations give three independent determinant-preserving directions, while tensoring by $\operatorname{Pic}^0$ gives three independent determinant directions. Thus
\[
 \operatorname{ext}^1(I,I)=\operatorname{ext}^2(I,I)=6.
\]

The categorical trace followed by Hochschild--Kostant--Rosenberg identifies evaluation on $I$ with contraction against $\operatorname{ch}(I)$; the precise compatibility is \cite[Proposition~6.2.1, Corollary~6.3.2 and Theorem~6.4.2, author pp.~29--31]{BFHH}. The tangent bundle is trivial, so the Todd correction is one. The degree-two target is
\[
 HH_{-2}(Y_j)=H^2(\mathcal O_{Y_j})\oplus H^3(\Omega^1_{Y_j}),
\]
of dimension six. The first block is reached isomorphically through the rank term. Contraction from $H^1(T_{Y_j})$ with the nonzero multiple $-kh_j^2/2$ surjects onto the second block: contraction with $h_j$ reaches $H^2(\mathcal O_{Y_j})$, and wedging with $h_j$ identifies that space with $H^3(\Omega^1_{Y_j})$. In coordinates $h_j=\sum x_iy_i$, the latter sends $y_1y_2$ to a nonzero multiple of $x_3y_1y_2y_3$, and likewise for the other basis vectors. The trace therefore has rank six and is injective. In degree one the same calculation separates the Picard directions through the rank term and the translation directions through contraction with $h_j^2$. The degree-zero trace is the nonzero Chern character.

External product preserves this injectivity in degree two: the three K\"unneth terms of self-Ext map to the distinct Hochschild bidegrees $(2,0)$, $(1,1)$ and $(0,2)$, by tensor products of the three injective trace maps just checked. Pullback from a finite quotient identifies invariant self-Ext with the self-Ext downstairs. Trace is natural under that pullback. Thus the pullbacks of these two quotient ideals have a coherent finite descent action and injective invariant obstruction trace.

\subsection{A Fourier transform that is a sheaf after duality}
Put $A=X\times\widehat X$, and write $H_x,H_\alpha$ for the product principal classes on its factors. Write $r_j=u_jq$, $K_j=\ker r_j$, and $F_j=r_j^*I_{\bar C_j}=I_{Z_j}$. Then $|K_j|=8k$; concretely $K_1=\ker q+\langle(\tau_1,0,0)\rangle$ and $K_2=\ker q+\langle(0,\tau_2,0)\rangle$. These summands intersect trivially by injectivity of $f_j^*$. Fourier transform with the normalized Poincar\'e bundle is an equivalence by \cite[Theorem~2.2, p.~156]{Mukai}; its translation, tensor and isogeny identities are in \cite[\S3, especially (3.1) and Proposition~3.11]{Mukai}.
Let $L=\mathcal O_X(2H)$, $\mu(x,y)=(x+y,y)$, and use the inverse
Fourier kernel $\mathcal P^{-1}[3]$ in the second variable. Define
\begin{equation}\label{eq:effective-transform}
 E_{\mathrm{ref}}=
 (\operatorname{id}\times\mathrm{FM}_{\mathcal P^{-1}[3]})
 \mu^*\bigl(([-1]^*F_2\otimes L)\boxtimes(F_1\otimes L)\bigr).
\end{equation}
The order of the factors gives the following incidence description. On $X_y\times A$, put
\[
 a(y,x)=y+x,\quad
 \mathcal L=L_y\otimes a^*L\otimes\mathcal P_{y,\alpha}^{-1},
 \quad D_1=\{y\in Z_1\},\quad D_2=\{y+x\in-Z_2\}.
\]
For $\pi:X_y\times A\to A$ and $T=D_1\cap D_2$, set
\begin{equation}\label{eq:effective-duality}
 G=R\pi_*(\mathcal L I_{D_1}I_{D_2}),
 \qquad E_{\mathrm{ref}}=G[3],\qquad F=G^\vee[-1]
               =E_{\mathrm{ref}}^\vee[2].                 
\end{equation}
Every equality here is a derived-functor identification.

Here is the additional geometric fact that makes $F$ a sheaf.
The curve $C'=q^{-1}C$ is connected: the Abel map induces the
surjection $\pi_1(C)\to\pi_1(J(C))$, hence has transitive monodromy
on the fibre of the connected etale cover $q$. It is an etale
degree-eight cover of $C$, and therefore has genus $17$. Projection
of $C'$ along the first elliptic factor is injective. Indeed, two
points in one such fibre differ by a first-factor point; their
images in $C$ differ by an element of $f_1^*\operatorname{Pic}^0(E_i)$, so \eqref{eq:bielliptic-disjointness} makes those
images equal. The difference then lies in
$\ker q\cap(E_i\times0\times0)=0$, since $f_1^*$ is injective.
The same argument applies to the second factor.

Write $\eta_j$ for the pullback of the point class from the $j$th
elliptic factor of $X$, so $H=\eta_1+\eta_2+\eta_3$.
The two projected images are irreducible divisors $B_1,B_2$ on
the corresponding elliptic surfaces. Their projection maps have
degree one, so the inverse-image divisors $S_1,S_2$ on $X$
have classes
\begin{equation}\label{eq:effective-divisor-classes}
 [S_1]=4(\eta_2+\eta_3),\qquad
 [S_2]=4(\eta_1+\eta_3).                                  
\end{equation}
This follows by pushing forward
$[C']=q^*(\Theta^2/2)=2H^2$ along each elliptic projection.
There is no unaccounted projection degree. One may also check
smoothness: the projected divisor has self-intersection $32$ and
arithmetic genus $17$, equal to the genus of its degree-one
normalization $C'$, so the normalization has no singularity defect.

The $k$ translates defining $Z_1$ lie on $S_1$. On every generic
first-factor fibre of $S_1$ they supply $k$ distinct points.
The $k$ reflected translates in $-Z_2-x$ have the analogous
property on the second-factor fibres of $-S_2-x$. A section of
$\mathcal O_X(4H)\otimes P$ vanishing on both unions, where
$P\in\operatorname{Pic}^0(X)$ is arbitrary, must vanish on each
of these entire elliptic fibres: its restriction has degree four,
whereas $k=4d+1\ge5$. Hence its divisor contains both $S_1$ and
$-S_2-x$. They are distinct for every $x$, because their classes
in \eqref{eq:effective-divisor-classes} differ. The residual effective divisor would have class
\[
 4H-[S_1]-[S_2]=-4\eta_3.
\]
Its intersection with the ample class $H^2$ would be $-8$, which
is impossible for an effective divisor. Thus that space of
sections is zero, for every $x$ and every flat factor $P$.

\paragraph{An explicit derived-to-coherent calculation.}
The coordinates $(y,z=y+x)$ identify the two curve-union ideals
with pullbacks from independent factors. They are Tor independent,
and their product equals their intersection. Therefore
\begin{equation}\label{eq:effective-ideal-sequence}
 0\to\mathcal L I_{D_1}I_{D_2}\to\mathcal L
 \to\mathcal L|_{D_1}\oplus\mathcal L|_{D_2}
 \to\mathcal L|_T\to0                                    
\end{equation}
is exact on the total space.
On a fibre of $\pi$, $\mathcal L$ has class $4H$, with
$h^0=64$ and no higher cohomology. Its restriction to a component
$C'$ has degree $48$. Since $g(C')=17$, its first cohomology is zero
and its space of sections has dimension $32$. Consequently
\[
 W=\pi_*\mathcal L,\qquad
 U=\pi_*(\mathcal L|_{D_1}\oplus\mathcal L|_{D_2})
\]
are vector bundles of ranks $64$ and $64k$, respectively.

The map $f:T\to A$ is finite flat over each of its smooth divisor
images. To see the precise degree, ignore translations and write
one component as $C'\times C'\times\widehat X$, with its map to
$X$ given, up to sign, by addition. The diagonal action
$(p',q')\mapsto(p'+h,q'-h)$ of $\ker q$ is free. Its quotient is
the base change of $C\times C\to C^{(2)}\simeq\Theta$ by $q$.
The quotient map is etale of degree eight; the latter map is
finite flat of degree two onto the smooth divisor $q^{-1}\Theta$.
Thus the composite has degree $16$. Smoothness of $\Theta$ follows
from the absence of a degree-two pencil on the nonhyperelliptic
genus-three curve. Each summand of
\[
 \mathcal T=f_*(\mathcal L|_T)
\]
is a vector bundle on a smooth divisor, pushed forward to $A$.
It has projective dimension one.

Applying direct image to \eqref{eq:effective-ideal-sequence}, including the preceding higher
cohomology vanishings, represents $G$ by
\begin{equation}\label{eq:effective-three-terms}
 [W\longrightarrow U\longrightarrow\mathcal T]
 \quad\hbox{in degrees }0,1,2.                            
\end{equation}
The section argument proves that $W\to U$ is fibrewise injective.
Hence $Q=U/W$ is a vector bundle of rank $64(k-1)=r$, and
$G=[Q\to\mathcal T]$ in degrees $1,2$. No assertion that fibrewise
$H^2(G)$ vanishes has been made or used.

Finite duality gives
\[
 \mathcal T^D=\mathcal E xt^1_A(\mathcal T,\mathcal O_A)
 =\bigoplus f_{ab,*}
       ((\mathcal L|_{T_{ab}})^{-1}\otimes\omega_{T_{ab}}),
\]
with all other sheaf Ext groups zero. Dualizing the last two-term
complex and shifting by $-1$ gives an ordinary coherent sheaf and
the exact extension
\begin{equation}\label{eq:effective-extension}
 0\longrightarrow Q^*\longrightarrow F
   \longrightarrow\mathcal T^D\longrightarrow0.          
\end{equation}
Both the local extension and its finite-duality description prove
$\operatorname{pd}F\le1$; local freeness outside a prescribed
codimension is unnecessary.

\paragraph{The finite action and its trace.}
In \eqref{eq:effective-transform}, the first translation parameter belongs to $K_2$ and
the second to $K_1$. With the convention $t_x^*$ for translation,
the transported action on $E_{\mathrm{ref}}$ has point map
\begin{equation}\label{eq:effective-action}
 (x,y)\longmapsto
 \bigl((x-y,2(x+y)),\,(-2x,y)\bigr)
 \quad\hbox{in }A\times\widehat A.                       
\end{equation}
Tensoring by $L$ sends a translation parameter $x$ to translation $x$ with tensor parameter $-2x$. Pullback by $\mu$ sends the two translation coordinates to $(x-y,y)$ and the tensor coordinates to $(-2x,-2(x+y))$. The inverse Fourier kernel exchanges translation $y$ with tensor $P_y$, and tensor $P_b$ with translation $-b$. This proves the displayed map. Its last coordinate recovers $y$ and its first then recovers $x$, so it is injective, with group order $64k^2$. Its coherence is transported from finite descent. Duality replaces a translation/tensor
point $(a,b)$ by $(a,-b)$; invert the dual linearizations to obtain
the coherent action on $F$. The trace is preserved up to invertible
signs under this operation. Explicitly, derived transpose identifies
the degree-two self-Ext spaces, the Atiyah class of the dual is
the negative transpose, and graded cyclicity of the trace gives
the sign $(-1)^p$ on
$\operatorname{tr}(\xi\operatorname{At}^p/p!)$.
The reversal signs for the cohomological and differential-form
factors cancel. Thus the complete trace kernels identify, including
the invariant subspaces. The even shift $[2]$ contributes sign $+1$.
This proves the claimed strong equivariant semiregularity of $F$.
Negative self-Ext groups vanish because $F$ is a sheaf; they also
vanish through the displayed equivalence with the original ideals.

\subsection{The entire Chern character and the split Weil action}
Put $s^2=-D$, $\theta_D=4DH_x+H_\alpha$, $B=-P/2$ and $\lambda_D=(D+1)/(2D)$, where $P$ is the first Chern class of the normalized Poincar\'e bundle on $X\times\widehat X$. The factor used before Fourier transform has character
\[
 \operatorname{ch}(F_j\otimes L)=1+2H-2DH^2-\frac{4D}{3}H^3
     =\alpha_+e^{2sH}+\alpha_-e^{-2sH},
 \quad\alpha_\pm=\frac{1\pm1/s}{2}.
\]
Reflection preserves this even character.

Here is the Fourier calculation, including the signs. On one elliptic pair let $a,b$ be the first coordinates, $c,e$ the integrated coordinates, and $u,v$ the dual coordinates; normalize $\int ce=1$. For $s_1,s_2\in\{s,-s\}$ the kernel exponent is
\[
 2s_1(a+c)(b+e)+2s_2ce-(cv-eu).
\]
Its integral is
\[
 e^{2s_1ab}\{2(s_1+s_2)-(2s_1b-v)(u-2s_1a)\}.
\]
This follows by collecting the coefficient of $ce$ in the exponential. Multiplying by $e^{-(av-bu)/2}$ and expanding the remaining nilpotent terms gives
\[
 \begin{cases}
 4s_1\exp(s_1(4Dab+uv)/(4D)),&s_1=s_2,\\
 \Omega_{s_1}:=4Dab+2s_1(av-bu)-uv,&s_1=-s_2.
 \end{cases}
\]
Take the product over the three elliptic pairs. The shift $[3]$ in the inverse Fourier transform contributes a minus sign. The equal-sign terms give the polynomial
\begin{equation}\label{eq:effective-polynomial}
 \begin{split}
 P_D(z)={}&64D-8(D-1)z-2z^2+\frac{D-1}{12D}z^3\\
 &+\frac{z^4}{96D}-\frac{D-1}{3840D^2}z^5
       -\frac{z^6}{46080D^2}.
 \end{split}
\end{equation}
The opposite-sign terms give $\lambda_Dt_D$, where
\[
 t_D=-\frac12\left(\prod_{j=1}^3\Omega_{s,j}
                      +\prod_{j=1}^3\Omega_{-s,j}\right).
\]
Duality changes the sign of the codimension-$j$ character by $(-1)^j$; the shift $[2]$ has sign $+1$. Consequently the sheaf $F$ satisfies the full identity
\begin{equation}\label{eq:effective-character}
 e^{-B}\operatorname{ch}(F)=-\lambda_Dt_D+P_D(-\theta_D).
\end{equation}

The rational action of $s$ on each elliptic pair is
$J_D=\left(\begin{smallmatrix}0&1/2\\-2D&0\end{smallmatrix}\right)$.
It squares to $-D$, is skew-adjoint for $\theta_D$, and has signature $(3,3)$. Since
$\Omega_s=-(u-2sa)(v-2sb)$, its conjugate products are the two Weil determinant lines. The class $t_D$ is primitive integral: the odd powers of $s$ cancel and the coefficient of $\prod_j u_jv_j$ is one. Direct multiplication gives
\[
 \theta_Dt_D=0,\qquad \int t_D^2=-2048D^3,
 \qquad \chi(\theta_D)=64D^3.
\]
The Hermitian space is split. For example, the two real cycles in the first elliptic factor form a $K$-basis of each pair, with zero diagonal Hermitian entries and nonzero conjugate off-diagonal entries; rescaling one vector gives $U$.

Set $r=64D$ and $\mathcal H=3(H_x+H_\alpha)$. We prove the following stronger form of $560$-regularity: for every $m\ge1$ and every line bundle $M$ such that $c_1(M)-560m^2\mathcal H$ is nef,
\begin{equation}\label{eq:effective-vanishing}
 H^j(A,[m]^*F\otimes M)=0\qquad(j>0).
\end{equation}
\subsection{A regularity estimate valid after every multiplication isogeny}
All the following vanishings allow arbitrary flat factors.
Relative duality and Serre duality show, for every integer $q\ge1$,
\begin{equation}\label{eq:effective-duality-vanishing}
 H^j(A,W^*(q\mathcal H))=H^j(A,U^*(q\mathcal H))=0
 \qquad(j>0).                                             
\end{equation}
Here is the numerical positivity test, so that \eqref{eq:effective-duality-vanishing} does not
depend on a qualitative regularity assertion. Project the duality
calculation onto the fixed $y$ coordinate, or onto the fixed curve
coordinate for a summand of $U$. The line bundle on the remaining
six-dimensional fibre has matrix
\[
 \begin{pmatrix}2-3q&0\\0&-3q\end{pmatrix}
 \quad\hbox{or}\quad
 \begin{pmatrix}2-3q&1\\1&-3q\end{pmatrix}.
\]
Both are negative definite: the largest eigenvalue before
subtracting $3qI$ is at most $1+\sqrt2<3$. Only the top direct
image of this antiample line bundle can be nonzero, which gives
the required cohomology vanishing below degree six on the dual
side. This proves \eqref{eq:effective-duality-vanishing}. In particular $W^*$ and $U^*$ are
$7$-regular.

Take the evaluation sequence
\[
 0\to K_U\to\mathcal O_A(-7\mathcal H)^N\to U^*\to0.
\]
Its kernel is locally free and $14$-regular. For degrees at least
two this follows directly from \eqref{eq:effective-duality-vanishing} and positivity of
$(7-j)\mathcal H$. For degree one, the multiplication map from
the degree-seven sections, followed by six hyperplanes in the
embedding defined by $\mathcal H$, is surjective by regularity.
Compose the evaluation map with $U^*\to W^*$ and call the kernel
$J$. Then
\begin{equation}\label{eq:effective-kernel-sequence}
 0\to K_U\to J\to Q^*\to0.                              
\end{equation}
The Buchsbaum--Rim resolution of $J$ has terms
\begin{equation}\label{eq:effective-br}
 C_j=\bigwedge^{64+j}\mathcal O_A(-7\mathcal H)^N
          \otimes\operatorname{Sym}^{j-1}W\otimes\det W,
 \qquad j\ge1.                                           
\end{equation}
The first term surjects onto $J$. Local exactness follows by
splitting the surjection as $K\oplus V\to V$ and decomposing
according to the number of factors from $K$. The resulting
strands are the polynomial de Rham complex in positive degree,
contracted by the Euler vector field in characteristic zero;
the degree-one strand maps isomorphically to $K$.

For an explicit bound on \eqref{eq:effective-br}, the theorem of the square gives
\[
 W=\mathcal O_A(2H_x)\otimes b^*\mathrm{FM}(\mathcal O_X(4H)),
 \quad b=2x-\alpha,
 \quad [4]^*W=Q_0^{\oplus64},
\]
with
\begin{equation}\label{eq:effective-four-pullback}
 c_1(Q_0)=16H_x-4H_\alpha+8P,\qquad
 \operatorname{Mat}(Q_0)=\begin{pmatrix}16&8\\8&-4\end{pmatrix}.
\end{equation}
Indeed the pullback of the Fourier bundle by the polarization
isogeny of $\mathcal O(4H)$ is
$H^0(\mathcal O(4H))\otimes\mathcal O(-4H)$; the displayed
formula is its pullback by $b$ with the first line-bundle factor.
The eigenvalues of the matrix in \eqref{eq:effective-four-pullback} are $6\pm2\sqrt{41}$,
whose minimum is greater than $-7$, while
$[4]^*\mathcal H$ has matrix $48I$. The symmetric and determinant
factor in \eqref{eq:effective-br} pulls back to copies of $Q_0^{j+63}$.
Kodaira vanishing and the trace splitting of the isogeny therefore
give
\begin{equation}\label{eq:effective-term-regularity}
 \operatorname{reg}_{\mathcal H}
   (\operatorname{Sym}^{j-1}W\otimes\det W)\le j+69,
 \qquad \operatorname{reg}C_j\le517+8j.                  
\end{equation}
Only six terms are needed to bound $J$ on a sixfold: successive
cohomology sequences leave a final kernel in cohomological
degree seven. Thus
\begin{equation}\label{eq:effective-reg560}
 \operatorname{reg}J\le
 \max_{1\le j\le6}(\operatorname{reg}C_j-j+1)
 \le560,\qquad\operatorname{reg}Q^*\le560.               
\end{equation}

For the torsion dual, use $\rho:C'\times C'\to C\times C$.
On $C\times C$ write $F_1,F_2$ for the point fibres,
$S=F_1+F_2$, and $\Delta_C$ for the diagonal. Addition pulls
$\Theta$ back to $A_C=4S-\Delta_C$. The divisor
$R_C=3S-\Delta_C$ is nef. A degree-three map $C\to\mathbb P^1$
from plane projection supplies its effective residual fibre
correspondence. If its monodromy is $S_3$, that residual divisor
is irreducible and has square $2$, hence is nef. If its monodromy
is cyclic of order three, it is the sum of the two automorphism
graphs: each has square $-4$, and they intersect transversely
in the five fixed points given by Riemann--Hurwitz. Their sum
therefore has intersection $1$ with each component and nonnegative
intersection with every other curve, again proving nefness.

Put $A'=\operatorname{add}_X^*H=\rho^*A_C/2$ and
$H'_i=\rho^*(3F_i/2)$. Up to flat factors on each translated
component,
$c_1(\mathcal L_T)=2H'_1+2H'_2-P_{1,\alpha}$ and
$f^*\mathcal H=3A'+3H_\alpha$. For $q\ge4$ decompose
\begin{equation}\label{eq:effective-torsion-positive}
 \begin{split}
 qf^*\mathcal H-c_1(\mathcal L_T)
 ={}&[3qA'-3H'_1-2H'_2]+(3q-1)H_\alpha\\
    &+[H'_1+H_\alpha+P_{1,\alpha}].
 \end{split}                                               
\end{equation}
The last bracket is a nef addition pullback. The first equals
\[
 \rho^*\left[\frac{3q}{2}R_C+
       \left(\frac{3q}{2}-\frac92\right)F_1+
       \left(\frac{3q}{2}-3\right)F_2\right]
\]
and is ample. Consequently \eqref{eq:effective-torsion-positive} is ample on the smooth
fivefold $C'\times C'\times\widehat X$. Kodaira vanishing and
finite duality give $9$-regularity of $\mathcal T^D$.
Together with \eqref{eq:effective-extension} and \eqref{eq:effective-reg560}, this proves the stated
$560$-regularity of $F$.

For the stronger assertion \eqref{eq:effective-vanishing}, pull the exact constructions
back by $[m]$. The matrices proving \eqref{eq:effective-duality-vanishing} are multiplied by
$m^2$, so the assumed nef difference makes the duality
vanishings hold after tensoring by $M$. The pulled evaluation
sequence has ample middle terms and gives vanishing of
$H^j([m]^*K_U\otimes M)$ for $j\ge2$.
After one further pullback by $[4]$, the first six terms of
\eqref{eq:effective-br} are sums of line bundles whose matrices, after removal
of a nef class, are
\[
 m^2\{48(112-7j)I+(j+63)\operatorname{Mat}(Q_0)\},
 \qquad1\le j\le6.
\]
Their least eigenvalue is at least
$m^2(3360-7\cdot69)=2877m^2>0$. The same six-term truncation
and \eqref{eq:effective-kernel-sequence} give the desired vanishing for $[m]^*Q^*\otimes M$.
For the torsion term, finite flat base change gives a smooth
etale cover $T'$ of $T$; the required line bundle is the sum
of the pullback of \eqref{eq:effective-torsion-positive} with $q=4$ and a nef class.
Kodaira vanishing on $T'$ proves the last vanishing.
The extension \eqref{eq:effective-extension} proves \eqref{eq:effective-vanishing}.

\paragraph{Normalization with an explicit regularity integer.}
Keep $D=4d$, $r=64D$ and the sheaf $F$ just constructed. Set
\begin{equation}\label{eq:effective-normalization}
\begin{gathered}
 h=3\theta_D,\quad \mathcal D=\det F,\quad
 a=\frac{D-1}{8D},\\
 R=560r^2+8Dr+6=2294272D^2+6,\qquad
 B_0=[r]^*F\otimes\mathcal D^{-r}.
\end{gathered}                                             
\end{equation}
The notation $B_0$ here denotes the ordinary normalized sheaf;
the rational B-field used below is written $\mathcal B_0$.
The class $h-\mathcal H$ is nef. Formula \eqref{eq:effective-character} gives
$c_1(F)=-32DP+8(D-1)\theta_D$, and the class
$8Dh-c_1(F)$ has matrix
\[
 \begin{pmatrix}32D(2D+1)&32D\\32D&8(2D+1)\end{pmatrix},
 \qquad\det=256D(4D^2+1)>0.
\]
It is ample. For $1\le i\le6$, the line bundle
$\mathcal D^{-r}\otimes h^{R-i}$ therefore satisfies
\[
 c_1(\mathcal D^{-r}\otimes h^{R-i})-560r^2\mathcal H
 =560r^2(h-\mathcal H)+r(8Dh-c_1(F))+(6-i)h,
\]
which is ample. Assertion \eqref{eq:effective-vanishing} proves that $B_0$ is
$R$-regular. It still has projective dimension at most one.
Its complete character is
\begin{equation}\label{eq:effective-normal-character}
 \operatorname{ch}(B_0)=e^{-r^2a\theta_D}
       \{-\lambda_Dr^6t_D+P_D(-r^2\theta_D)\}.            
\end{equation}
In particular $c_1(B_0)=0$, and, writing its first two relevant
scalar coefficients relative to $h$,
\begin{equation}\label{eq:effective-a2a3}
 \begin{aligned}
 \operatorname{ch}_2(B_0)&=a_2h^2,
 &a_2&=-\frac{r^4(D+1)^2}{18D},\\
 \operatorname{ch}_3(B_0)&=-\lambda_Dr^6t_D+a_3h^3,
 &a_3&=\frac{r^6(D-1)(D+1)^2}{648D^2}.
 \end{aligned}                                             
\end{equation}

\paragraph{Deformation and the ordinary sheaf family.}
Apply Perry's Theorem~1.2 \cite{Perry} to the sheaf $F$, the
coherent finite action proved above, and the algebraic
rational B-field $\mathcal B_0=-c_1(F)/r$. The whole twisted
character is the normalized expression in \eqref{eq:effective-normal-character} before the
$[r]$ scaling: it is a polynomial in the fixed polarization
plus a nonzero Weil class. It is flat and remains of Hodge type
through the balanced Weil family. The determinant monodromy is an integral unit of $K$ of norm one and hence a root of unity; a finite cover makes the Weil class flat.

For precision about untwisting, Perry's \S6.2, proof of Theorem~1.2, pp.~26--28, begins with the $\mu_n$-gerbe representing
the B-field denominator. Here take $n=r$. Its subsequent pushout
to $\mu_m$ does not change the associated multiplicative-group gerbe. Pullback along the map from the original $\mu_r$-gerbe to this pushout is an equivalence on the weight-one categories: locally both are modules on which the band acts through its standard character, and the descent cocycle is unchanged. We therefore pull the deformed weight-one object back to the original finite-band gerbe before untwisting. On a
contractible analytic base the original class remains in the
local system $H^2(A,\mu_r)$. Multiplication by $r$ acts as
$r^2$ on this group and therefore kills it. Pullback by $[r]$
and an analytic Cech trivialization give an ordinary perfect
family. Here is its exact seed comparison, including the sign.
The seed gerbe is the $r$th-root gerbe of $\mathcal D^{-1}$;
its tautological weight-one line $\mathcal T_0$ satisfies
$\mathcal T_0^r=\mathcal D^{-1}$, and the seed object is
$F\otimes\mathcal T_0$. Write its class without an exponential
convention as the Kummer boundary
$\vartheta_0=\delta(\mathcal D^{-1})\in H^2(A,\mu_r)$.
Its rational B-field is
$c_1(\mathcal T_0)=c_1(\mathcal T_0^r)/r
=-c_1(\mathcal D)/r=\mathcal B_0$.
After pullback by $[r]$ and neutralization, the line becomes
an ordinary line $N_0$ with
$N_0^r=[r]^*\mathcal D^{-1}$. Hence
$c_1(N_0)=-r c_1(\mathcal D)$ and
$N_0=\mathcal D^{-r}\otimes P_0$ for some
$P_0\in\operatorname{Pic}^0(A)$. The ordinary seed is therefore
exactly $B_0\otimes P_0$, as asserted.
For these conventions see \cite[Definition~4.5, Lemma~4.6, Remark~4.13, Definition~4.16 and Theorem~4.17, author pp.~15--18]{Hotchkiss}. The displayed root and tensor identities give the required
comparison directly. All preceding estimates allow the flat
factor $P_0$.

The locus where the derived fibres of a relative perfect
complex have cohomology only in degree zero is open. At the
seed this property holds by \eqref{eq:effective-extension}; after shrinking the base
the ordinary family is therefore a coherent sheaf flat over
the base. This can equally be seen by the relative Tor-amplitude
criterion applied to a local finite vector-bundle resolution.
Properness permits a common neighbourhood for the whole fibre.
Analytic upper semicontinuity preserves the six regularity
vanishings of $B_0$ on a smaller neighbourhood. No stability
of all sheaves with this character is used.

\paragraph{A vector-bundle kernel supplies the effective cycle.}
At the seed and throughout that neighbourhood, evaluation gives
\begin{equation}\label{eq:effective-evaluation}
 0\longrightarrow V\longrightarrow
    \mathcal O_A(-Rh)^N\longrightarrow B_0\longrightarrow0,
 \qquad N=h^0(B_0(Rh))=\chi(B_0(Rh)).                     
\end{equation}
At the seed $V$ is locally free, because $\operatorname{pd}B_0\le1$.
This follows locally from the exact Ext sequence, or by comparing
the evaluation presentation with a projective resolution of
length one. In the family the direct image of $B_0(Rh)$ is a
rank-$N$ vector bundle, by the regularity vanishings and base
change; evaluation is surjective. Its kernel is flat, and its
locally free locus is open. The closed complement misses the
proper seed fibre, so it misses every fibre after one more
shrinking. Thus \eqref{eq:effective-evaluation} is a relative vector-bundle
presentation on the neighbourhood.

The kernel $V$ is $(R+7)$-regular. For cohomological degrees
at least two, use the sequence \eqref{eq:effective-evaluation}, the regularity of $B_0$,
and the positive line bundles $(7-i)h$ for $2\le i\le6$.
For degree one, $R$-regularity gives the surjection
\[
 H^0(B_0(Rh))\otimes\operatorname{Sym}^6H^0(\mathbb P^M,\mathcal O(1))
 \longrightarrow H^0(B_0((R+6)h)),
\]
which factors through the degree-six sections in \eqref{eq:effective-evaluation}.
It gives the remaining $H^1$ vanishing. Therefore
$V'=V((R+7)h)$ is a globally generated vector bundle.

Let $q_0=R+7$ and define the explicit rational numbers
\begin{equation}\label{eq:effective-uvw}
 \begin{split}
 u&=7N-rq_0,\\
 v&=\frac{49N}{2}-\frac{rq_0^2}{2}-a_2,\\
 w&=\frac{343N}{6}-\frac{rq_0^3}{6}-q_0a_2-a_3.
 \end{split}                                               
\end{equation}
The identity in $K_0(A)$ is
\[
 [V']=N[\mathcal O_A(7h)]-[B_0(q_0h)].
\]
Consequently $c_1(V')=uh$, $\operatorname{ch}_2(V')=vh^2$,
and $\operatorname{ch}_3(V')=wh^3+\lambda_Dr^6t_D$.
The exact Newton identity
$c_3=2\operatorname{ch}_3-c_1\operatorname{ch}_2+c_1^3/6$
therefore gives
\begin{equation}\label{eq:effective-c3}
 c_3(V')=2\lambda_Dr^6t_D+
                (2w-uv+u^3/6)h^3.                       
\end{equation}
Its Weil projection is nonzero. In particular the rank of $V'$
is at least three, since lower-rank bundles have $c_3=0$.

General sections of a globally generated vector bundle give
an effective maximal-rank-drop locus of expected codimension
three, or the empty locus, representing $c_3(V')$. This is the
elementary Thom--Porteous construction: take
$\operatorname{rank}(V')-2$ sections, use the universal surjective
evaluation to get transversality to the rank strata, and then
take a general choice. Equation \eqref{eq:effective-c3} rules out the empty case.
It is an effective cycle with integral coefficients with nonzero Weil projection,
so at least one reduced irreducible component detects the Weil
plane. No condition on the codimension of the nonfree locus of
$B_0$ is required.

Write $P_D(z)=\sum_{j=0}^6p_jz^j$, with coefficients given in
\eqref{eq:effective-polynomial}. Since the Todd class of an abelian variety is $1$ and
$\theta_Dt_D=0$, the integer $N$ is explicitly
\[
 N=(6!\,64D^3)\sum_{j=0}^6
       p_j(-r^2)^j\frac{(3R-r^2a)^{6-j}}{(6-j)!}.
\]
The exact degree of the effective cycle above is the nonnegative
integer
\begin{equation}\label{eq:effective-degree}
 C_D=(3^6\,6!\,64D^3)(2w-uv+u^3/6).                     
\end{equation}
The formula defines every quantity by rational operations on
the specified positive integer $D$. Integrality and nonnegativity
also follow from the constructed vector bundle and effective
cycle. These same Chern classes and degrees persist on the
analytic neighbourhood. GAGA makes each fibre and cycle algebraic.

\subsection{The degree and its transport to every field and lattice}
The formulas above specify $C_D$ by rational arithmetic in $D$. We first give a convenient upper bound. Put $s_0=3R-r^2a$. From \eqref{eq:effective-polynomial}, $j!|p_j|\le64D$ for $0\le j\le6$, while $0<s_0$ and $s_0+r^2<7\cdot10^6D^2$. The Euler-characteristic formula therefore yields
\[
 0<N\le4096D^4(s_0+r^2)^6<10^{45}D^{16}.
\]
Also $q_0<3\cdot10^6D^2$, $|a_2|<4\cdot10^6D^5$ and $|a_3|<5\cdot10^8D^7$. Substitution in \eqref{eq:effective-uvw} gives
\[
 |u|\le8\cdot10^{45}D^{16},\quad
 |v|\le26\cdot10^{45}D^{16},\quad
 |w|\le59\cdot10^{45}D^{16}.
\]
As $10^{45}D^{16}\ge1$, the absolute value of $2w-uv+u^3/6$ is at most $512\cdot10^{135}D^{48}$. Hence
\[
 0<C_D<10^{146}D^{51},\qquad C_{4d}<10^{177}d^{51}.
\]
Positivity follows also from the nonzero effective detecting cycle, and integrality follows from its degree.

We next give a fixed map from the source to a principal hyperbolic lattice. Put $s=\sqrt{-d}$ and let $e_d=1$ for $d\equiv1,2\pmod4$, and $e_d=2$ for $d\equiv3\pmod4$. Thus $\tau=s$ or $(1+s)/2$ generates $\mathcal O_K$, and $\tau-\bar\tau=2s/e_d$. In one elliptic pair write the real homology basis as $u=e_x$, $v=i e_x$, $w=e_y$, $z=i e_y$. The $s$-action is
\[
 su=-4dw,\quad sv=-4dz,\quad sw=u/4,\quad sz=v/4.
\]
The source polarization $h=3(16dH_x+H_y)$ has $E_h(u,v)=48d$ and $E_h(w,z)=3$. In the $K$-basis $u,v$ its normalized Hermitian matrix, linear in the first variable, is
\[
 \begin{pmatrix}0&48ds/e_d\\-48ds/e_d&0\end{pmatrix}.
\]
In the principal hyperbolic lattice $\mathcal O_K e\oplus\mathcal O_K f$, define
\[
 T(u)=4s e,\quad T(v)=12d f,\quad
 T(w)=e,\quad T(z)=-3s f.
\]
These values are integral and compatible with the $K$-action, and $H_U(Tu,Tv)=48ds$. Thus $T^*E_U=e_dE_h$. Three copies give an isogeny $\varphi:A_0\to P_U$ with $\varphi^*\Theta_U=e_dh$. In particular the pushforward of a three-cycle costs exactly $e_d^3$ in total cycle degree. Its Weil projection stays nonzero. This calculation uses no class-number assumption.

The fixed rational similitude identifies the period domains. After choosing a common finite level it sends the analytic neighbourhood where the sheaf construction works onto an analytic open subset of the principal $U^3$ family. Bounded-detector closedness therefore gives a detector of degree at most $e_d^3C_{4d}$ on every member of that family. Theorem~\ref{thm:principal-quotient}, and $e_dt_K=2$, now give
\[
 \delta_W(A,L)\le(t_Kfm_K)^3\Delta\,e_d^3C_{4d}
                 =8m_K^3C_{4d}f^3\Delta.
\]
Since $m_K\le4\sqrt d/\pi$, one has $8m_K^3<17d^{3/2}$. The asserted coarse bound follows from $C_{4d}<10^{177}d^{51}$ and $d^{105/2}\le d^{53}$. For class number one, $m_K=1$. A detecting irreducible component never has greater degree than the effective cycle containing it. This proves Theorem~\ref{thm:all-field-sixfolds}.

The exterior-algebra and integer-arithmetic calculations are independently reproducible with the programs \texttt{sixfold\_character\_check} and \texttt{sixfold\_degree\_check} in the \href{\workbenchurl}{workbench}. The proofs above work for every positive squarefree $d$. The programs check the formulas and signs, and do not replace the geometric deformation argument.

\section{An explicit Eisenstein eightfold bound after Schoen and Severn}
\label{sec:eisenstein-eightfolds}
This section bounds the degrees in the cyclic-cover transfer used for the
split Eisenstein eightfold construction. The geometric input is the
$G_{34}$ family of Severn's unrefereed preprint
\cite[Proposition~8.5, \S9 and Appendix~A]{Severn}. The calculation below
is a quantitative consequence of that construction. It does not use
Perry's theorem or either CM Hodge-conjecture manuscript.

Put
\begin{align*}
 G&=39191040,& g&=4G+1,\\
 R&=(2G+1)G!G^G,& L_0&=2gR,\\
 S_0&=G^{2g}(g!)^2,& S_B&=(L_0/2)!S_0^R,\\
 k&=L_0/2+4,&
 B_0&=(40G)^k\binom{k}{4}S_B\,8!G^8.
\end{align*}
These formulas specify integers; they do not require decimal expansion.
The constant is deliberately coarse, so that no estimate for coefficients
in an invariant-theory decomposition enters the proof.

\begin{theorem}\label{thm:eisenstein-eightfolds}
Let $K=\Q(\sqrt{-3})$, let $\mathcal O_f=\Z+f\mathcal O_K$ act
integrally on a polarized abelian variety of Weil type $(A,L_A)$, and
assume that the polarization is compatible with the action. Write
$\Delta=\chi(L_A)$. Using the $G_{34}$ geometric input just specified,
every split Weil eightfold satisfies
\[
 \delta_W(A,L_A)\le B_0 f^4\Delta.
\]
Every Weil sixfold over $K$, of arbitrary determinant, satisfies
\[
 \delta_W(A,L_A)\le \frac{B_0}{4}
          \frac{f^9\Delta^2}{b},
 \qquad b=[\mathcal O_K H_1(A,\Z):H_1(A,\Z)].
\]
The input from the preprint is a specific dominant split family of
group-algebra factors of Jacobians. The proof gives effective integral
detecting cycles and all the degree losses from that family to $A$.
\end{theorem}

\subsection{The geometric input and the factor lattice}

Here is the precise form of the geometric input. The complex reflection
group $\mathcal G=G_{34}$ has order $G$, has its faithful six-dimensional
reflection representation over $K$, and contains the scalar $\omega I$.
Ten involutory reflections $s_1,\ldots,s_{10}$ generate it. The tuple
\[
 (s_1,s_1,\ldots,s_{10},s_{10})
\]
defines a Hurwitz family of connected $\mathcal G$-covers
$C\to\mathbf P^1$. Its reflection multiplicity in $H^1(C,K)$ has
$K$-dimension eight and Hodge numbers $(4,4)$. The period map of the
resulting abelian eightfold factor has an open image in the split Weil
domain. These are the statements proved in \cite[\S9]{Severn}; the
matrices and the seventeen boundary twists are in its Appendix~A.
Riemann--Hurwitz gives $2g(C)-2=8G$, explaining the integer $g$ above.

We recall why the period certificate has the stated consequence. On the
rank-eight parabolic cohomology, the commuting logarithms $N_j$ have
$N_jN_l=0$, their images span a four-dimensional isotropic $K$-space
$I$, and they annihilate $I$. Their associated Hermitian forms on the
quotient are semidefinite of a common sign, with definite sum. They span
a rational space of dimension sixteen. Geometric monodromy logarithms
are common-sign positive multiples of these logarithms after a finite
base change. A positive sum therefore has image and kernel $I$. The
limiting mixed Hodge structure has weight-zero part $I$ of type $(0,0)$,
so its limiting $F^1$ is complementary to $I$. The rational Lie algebra
of unitary maps from the quotient to $I$ has dimension sixteen; its
complexification maps isomorphically to the tangent space at that
limiting filtration. Thus sixteen of the $N_j$ give independent tangent
vectors. Holomorphic untwisting of the period map shows that the same
minor is nonzero at nearby smooth fibres. The isotropic space $I$ also
proves that the Hermitian space is split. This is the implication in
\cite[Lemmas~8.2--8.4 and Proposition~8.5]{Severn}; its finite matrix
input, including the nonzero minor $1728$, is distinct from the degree
argument that follows.

Let $\chi$ be the character of the reflection representation, and let
$s$ be a reflection. In the rational group algebra put
\[
 \varepsilon=\frac6G\sum_{a\in\mathcal G}
       \operatorname{Tr}_{K/\Q}(\chi(a^{-1}))a,
 \qquad e=\varepsilon(1-s)/2.
\]
The first operator is the central projector for the conjugate pair of
reflection representations. The second has rank one in its matrix
algebra $M_6(K)$. Its image on $J(C)$ is an abelian eightfold, denoted
$A_C$, with the restricted polarization $\theta_C$. Character values
are algebraic integers, so their rational traces are integers and
$Ge\in\Z[\mathcal G]$. The Rosati involution sends $a$ to $a^{-1}$;
both commuting factors of $e$ are self-adjoint. The scalar $\omega I$
preserves $A_C$, giving integral $\mathcal O_K$-action.

Let $\Lambda$ be the unimodular lattice of $J(C)$ and
$\Lambda_C=\Lambda\cap e\Lambda_\Q$. This sublattice is primitive.
Every integral functional on it extends to $\Lambda$, and
self-adjointness identifies its dual lattice with $e\Lambda$.
Consequently
\[
 \Lambda_C^{\#}=e\Lambda,\qquad
 G\Lambda_C^{\#}\subset\Lambda_C.
\]
Every elementary divisor of the restricted polarization divides $G$.
In particular
\begin{equation}\label{eq:eis-source-volume}
 \chi(\theta_C)\mid G^8,
 \qquad \int_{A_C}\theta_C^8\le8!G^8.
\end{equation}

\subsection{A bounded integral relation between determinant characters}

Let $U$ be the full centralizer of $\mathcal G$ in the symplectic group
of $H^1(C,\Q)$. Over $\C$, a nonreal conjugate pair of irreducible
representations gives a general linear multiplicity group, an orthogonal
irreducible representation gives a symplectic multiplicity group, and
a quaternionic one gives an orthogonal multiplicity group. Hence the
character group of $U$ has a free determinant coordinate for each
nonreal pair and a coordinate of order two for each quaternionic
multiplicity group. The full orthogonal groups must be used here.

For a subgroup $H'\le\mathcal G$ and a linear character $\chi'$ of
$H'$, let $a(\chi')$ be the determinant character of its multiplicity
space. Frobenius reciprocity identifies this with the image of
$\operatorname{Ind}_{H'}^{\mathcal G}\chi'$ under the homomorphism
which takes differences of conjugate multiplicities and quaternionic
multiplicities modulo two. Brauer induction, in the form used in
\cite[Theorem~5.1]{Severn}, gives an integral relation
\begin{equation}\label{eq:eis-supplier-relation}
 e_\sigma=\sum_i n_i a(\chi_i),
\end{equation}
where $e_\sigma$ is the character of one target Weil line.

Characters of order one or two contribute zero. Their multiplicities
in conjugate nonreal representations agree; their multiplicities in
a quaternionic representation are even, since their isotypic subspace
inherits a nondegenerate alternating form. A character with zero
cohomological multiplicity also contributes zero. All remaining terms
come from cyclic intermediate covers
\[
 C_i'=C/\ker\chi_i\longrightarrow X_i=C/H_i'
\]
of order $m_i\ge3$. Since the original inertia groups have order two,
every nonzero branch exponent in such a cyclic cover is $m_i/2$.
The sum of the exponents is zero modulo $m_i$, so their number $s_i$ is
even and they can be paired with their negatives. If $q_i=g(X_i)$,
the primitive eigenspaces in $H^1(C_i')$ have common dimension
$h_i=2q_i-2+s_i$. A nonzero retained multiplicity has $h_i>0$, and
$h_i$ is even. Unramified intermediate covers are included.

There is a uniform bound
\begin{equation}\label{eq:eis-relation-mass}
 \sum_i|n_i|\le R.
\end{equation}
To prove it, append columns $2e_j$ for the order-two coordinates and
lift the relation to an integer column lattice. There are at most $G$
rows, and the column entries have absolute value at most $G$. Restrict
to independent rows; this does not change relations between the
columns. If their rank is $r$, choose an invertible integer matrix $B$
of $r$ columns and put $d_0=|\det B|\le r!G^r$. In the subgroup of
$\Z^r/B\Z^r$ generated by all columns, a shortest expression for the
target uses at most $d_0-1$ signed columns: a shortest path has no
repeated vertex. The residual vector has coordinates at most $d_0G$.
Cramer's rule, with cofactors at most $(r-1)!G^{r-1}$, expresses it in
the columns of $B$ with each coefficient at most $r!G^r$. The total
coefficient sum is therefore at most $(r+1)r!G^r\le R$. Removing the
auxiliary parity columns leaves \eqref{eq:eis-supplier-relation} and
cannot increase that sum.

\subsection{Integral cycles from the cyclic intermediate covers}

Fix one term, and write $C'$, $X$, $m$, $h$ for its data. We use
Schoen's construction \cite[Theorem~2.0 and Lemmas~2.3--2.6,
pp.~11--17]{Schoen88}. The primitive rational part has dimension
$\varphi(m)h$, so
\[
 h\le g(C')\le g,
 \qquad m\le G.
\]
Let $G_0$ be the kernel of the sum map on $(\Z/m)^h$, extended by
the permutation group on the $h$ factors. It has order
$m^{h-1}h!$. There are finite quotient maps
\[
 q:C'^h\longrightarrow W_0=C'^h/G_0
       \longrightarrow\operatorname{Sym}^hX,
\]
the second with residual cyclic group $\Z/m$. Schoen chooses a linear
$P_0\simeq\mathbf P^{h/2}$ in a complete linear system on $X$. The
strict transform of $P_0$ is unchanged by a sequence of blowups and
avoids the branch divisors. The cyclic cover therefore splits over it.
A component descends to $Q_0\subset W_0$ birationally over $P_0$,
and its primitive-character components in cohomology are nonzero.

We specify the exceptional-divisor issue in this last assertion.
When $4\mid m$ and all exponents are $m/2$, the prime chosen in
\cite[Sublemma~2.9, p.~19]{Schoen88} does not exist. The replacement
is the faithful-character argument of \cite[Lemma~4.4, pp.~9--10]{Severn}.
At a blowup of the intersection of a paired pair of branch divisors,
their total exponent is zero modulo $m$. The exceptional divisor is a
$\mathbf P^1$-bundle and is not itself a branch divisor. The two
strict transforms meet each fibre in two points. For a faithful
character, the eigensheaf is extension by zero of a rank-one local
system on their complement $\C^\times$, with nontrivial monodromy.
Its compactly supported cohomology vanishes: degree zero is zero,
degree two is dual to zero invariants, and the Euler characteristic is
zero. If a whole fibre lies in another branch divisor, its eigensheaf
is zero because its inertia is nontrivial. Proper base change thus
kills the faithful-character cohomology of the entire exceptional
divisor. The original centre has nontrivial inertia as well and has
zero faithful-character eigensheaf. The relative-cohomology sequence
then identifies the primitive parts before and after the blowup.
Schoen's nonzero normal-bundle intersection on the final component
therefore remains nonzero on $Q_0$. This proves the needed descent for
these exponents, including those omitted by the prime-choice argument.

Let $H$ denote the sum of point classes on $C'^h$. An incidence divisor
$D_x=x+\operatorname{Sym}^{h-1}X$ restricts to the hyperplane class on
$P_0$, and pulls back to $mH$ on $C'^h$. If $Z_0$ is an irreducible
component above $Q_0$ and has degree $e_0$ over $Q_0$, then
\[
 e_0\le |G_0|,
 \qquad \deg_H Z_0=e_0/m^{h/2}.
\]
The integral orbit sum
$\mathcal R_0=\sum_{a\in G_0}a_*Z_0$, with repetitions, has effective
degree at most $|G_0|^2/m^{h/2}$. Rational transfer for a finite
quotient gives $[\mathcal R_0]=e_0q^*[Q_0]$. This cohomological
identity does not require the quotient to be flat along $Q_0$.

For the residual action $\delta$ of $\Z/m$, form the integral signed
cycle
\[
 Y=\sum_{t=0}^{m-1}c_m(t)\delta(t)_*\mathcal R_0,
\]
where $c_m(t)$ is the Ramanujan sum. This operator is $m$ times the
rational primitive projector, and $|c_m(t)|\le\varphi(m)$. It keeps
the nonzero primitive components. In $G_0$-invariants a primitive
residual character forces all coordinate characters to be the same
primitive character. Consequently only
$H^1(C')^{\otimes h}$ occurs, and permutation invariance identifies
its component with the corresponding determinant line; the signs in
the permutation action are the cohomological signs in degree one.
Writing the mass of a signed cycle for the sum of the degrees of its
positive and negative parts, we obtain
\begin{equation}\label{eq:eis-cyclic-mass}
 \operatorname{mass}_H(Y)
 \le m\varphi(m)\frac{(m^{h-1}h!)^2}{m^{h/2}}
 \le m^{3h/2}(h!)^2\le S_0.
\end{equation}
The orbit sum in this estimate includes its integral multiplicities.

\subsection{Transfer by a finite collection of divisors}

Take $|n_i|$ copies of the cycle for each term of
\eqref{eq:eis-supplier-relation}. Their exterior product $Y_B$ lies on
a product $B$ of $\ell=\sum_i|n_i|h_i$ curves, with
$\ell\le gR\le L_0$. For the sum $H_B$ of point classes,
\begin{equation}\label{eq:eis-product-mass}
 \operatorname{mass}_{H_B}(Y_B)
 \le \frac{(\ell/2)!}{\prod_i((h_i/2)!)^{|n_i|}}S_0^R
 \le S_B.
\end{equation}
Every selected determinant component is nonzero. The perfect pairing
between inverse primitive eigenspaces on each curve gives its dual
determinant tensor. Choose the component according to the sign of
$n_i$. The product of the dual tensors with one target Weil eigenline
has degree $\ell+8$ on $B\times A_C$, is $U$-invariant by
\eqref{eq:eis-supplier-relation}, and sends $[Y_B]$ to a nonzero vector
of that line. Denote this tensor by $\Gamma$. Both the source cycles
and their dual tensors have middle degree on the curve products.

We need control of the divisors representing $\Gamma$. Let
$D$ be the image of $\Q[\mathcal G]$ in $\operatorname{End}H^1(C)$.
The semisimple-algebra lemma and the double-centralizer argument in
\cite[Lemma~3.5]{MilneLefschetz} give
$\operatorname{End}_U H^1(C)=D$. The invariant theorem for full
symplectic and orthogonal groups and paired general linear groups
\cite[Proposition~3.6 and Remark~3.7]{MilneLefschetz} says that their
exterior invariants on powers are generated in degree two.
The degree-two tensors correspond to self-adjoint matrices over $D$.
Off-diagonal entries and their adjoints are spanned by group elements;
the diagonal self-adjoint entries are spanned by $a+a^{-1}$.
It follows that these tensors are spanned by the divisor classes of
the globally generated line bundles
\begin{equation}\label{eq:eis-divisor-list}
 \operatorname{pr}_i^*(2\Theta),\qquad
 (\operatorname{pr}_i+a\operatorname{pr}_j)^*(2\Theta),
 \qquad a\in\mathcal G,
\end{equation}
on powers of $J(C)$. Repeated indices are allowed: the identity
$(1+a)^\dagger(1+a)=2+a+a^{-1}$ supplies the diagonal entries.

For a quotient $\pi:C\to C'$, the map
$C'\to J(C')\xrightarrow{\pi^*}J(C)$ identifies its first cohomology
with a rational direct summand, equivariantly for $U$. The idempotent
$e$ gives the corresponding summand for $A_C$. Thus $\Gamma$ lifts
to an exterior invariant on a power of $J(C)$. Its degree is $\ell+8$,
so it is a complex linear combination of monomials of
$c=\ell/2+4$ divisors from \eqref{eq:eis-divisor-list}. Since its
action on $Y_B$ is nonzero, some monomial has nonzero projection to
the target Weil line. No bound for the coefficients of the linear
combination is required.

We now bound every such monomial. If $a_0=\deg\pi$ and
$\eta:C'\to J(C)$ is the preceding map, the norm identity and
Riemann--Hurwitz give
\[
 \deg\eta^*\Theta=a_0g(C')
 \le g(C)-1+a_0\le5G.
\]
On $B\times A_C$, the restricted product polarization is therefore
at most $5G(H_B+\theta_C)$ in the nef order. The theorem of the
square bounds each line bundle in \eqref{eq:eis-divisor-list} by eight
times that product polarization, also for repeated indices. Hence
\[
 D_j|_{B\times A_C}\le40G(H_B+\theta_C)
\]
in the nef order. Successive general members of these globally
generated systems meet every component of $|Y_B|\times A_C$
properly. They give effective integral four-cycles. Their signed
pushforward to $A_C$ has the prescribed correspondence class, and
its degree mass is at most
\begin{align*}
 &(40G)^c\int_{|Y_B|\times A_C}
                    (H_B+\theta_C)^c\theta_C^4\\
 &\qquad=(40G)^c\binom c4
              \operatorname{mass}_{H_B}(Y_B)
              \int_{A_C}\theta_C^8
 \le B_0.
\end{align*}
Only the term with exponents $(\ell/2,4)$ contributes. Components
whose images have smaller dimension push to zero. Since the total
Weil projection is nonzero, one irreducible component of the remaining
effective cycles detects. This proves $\delta_W(A_C,\theta_C)\le B_0$.

\subsection{From the source to every polarization and determinant}

The source has maximal-order Eisenstein action and split Hermitian
space. Since $t_K=1$ and the class number is one,
Theorem~\ref{thm:principal-quotient} gives an isogeny to the principal
$U^4$ family pulling the target polarization back to $\theta_C$.
Pushing forward a detecting cycle costs at most its original degree.
The induced map between rational period domains is an isomorphism,
so the principal family still has an open image with detector bound
$B_0$. The closed bounded-detector locus of
Proposition~\ref{prop:detector-loci} extends this bound to its entire
connected component. Pullback through the principal quotient of an
arbitrary target costs $f^4\Delta$, proving the eightfold estimate.

For a maximal-order sixfold of arbitrary determinant, the balanced CM
surface in Theorem~\ref{thm:explicit-stabilization} has polarization
product $\Delta$. Its product with the sixfold is a split eightfold
of polarization product $\Delta^2$. The opposite-graph contraction
in that theorem has cost $1/4$ in the Eisenstein case. This gives
$B_0\Delta^2/4$. For an order of conductor $f$, saturation has degree
$b$, polarization multiplier $f$, and polarization product
$\Delta'=f^6\Delta/b$. Pulling back a three-cycle costs $b/f^3$.
Thus the bound is
\[
 \frac b{f^3}\frac{B_0}{4}
       \left(\frac{f^6\Delta}{b}\right)^2
 =\frac{B_0}{4}\frac{f^9\Delta^2}{b},
\]
as asserted. The factor $8!$ in $B_0$ makes $B_0/4$ an integer.

\section{Detecting components of complete intersections}
\label{sec:complete-intersections}
This section makes the geometric question concrete: one may ask for a detecting component in one fixed multiple of the polarization.

Let $(A,L)$ be of Weil type of dimension $2n$, with compatible polarization $\theta=c_1(L)$, and put $\Delta=\chi(L)$. Define $m_W(A,L)$ to be the least positive integer $m$ such that $n$ effective divisors in $|mL|$ meet properly and their intersection cycle has an irreducible detecting component. Set $m_W=+\infty$ when no such intersection exists.

\begin{proposition}\label{prop:complete-intersections}
The numbers $m_W$ and $\delta_W$ are finite simultaneously. When finite,
\[
 \delta_W\leq\frac{(2n)!\Delta}{2}\,m_W^n,
 \qquad m_W\leq3^{n+1}\delta_W.
\]
Consequently a uniform detector-degree bound in a fixed polarized family is equivalent to a uniform bound on $m_W$.
\end{proposition}

\begin{proof}
Compatibility of the polarization gives $e_W(\theta^n)=0$: with respect to the two $K$-eigenspaces, $\theta^n$ has exterior bidegree $(n,n)$, whereas the Weil lines have bidegrees $(2n,0)$ and $(0,2n)$. If a proper intersection of $n$ members of $|mL|$ has one detecting component, it must have at least two distinct detecting components. Indeed its total projected class is zero, and one nonzero projected component with positive multiplicity cannot sum to zero by itself. The total degree is $m^n(2n)!\Delta$. The smaller detecting component has degree at most half this number, proving the first inequality.

For the converse we give an elementary bound on equations. If an integral $r$-fold $Y\subset\mathbf P^N$ has degree $d$, then degree-$d$ hypersurfaces through $Y$ have set-theoretic base locus exactly $Y$. Fix $x\notin Y$. There is a linear $\mathbf P^{N-r-1}$ through $x$ disjoint from $Y$: project $Y$ from $x$ to $\mathbf P^{N-1}$, take a general $\mathbf P^{N-r-2}$ avoiding its image, and take the cone with vertex $x$. Choose a hyperplane $K$ in this linear space not containing $x$. Projection from $K$ to $\mathbf P^{r+1}$ is a morphism on $Y$ with finite fibres, since a positive-dimensional projective fibre would meet $K$. Its image is an integral hypersurface $Y'$ not containing the image of $x$. If its degree is $d'$ and the finite map has degree $e$, then $d=ed'$, so $d'\leq d$. Pull back its equation and multiply by a power of a linear form nonzero at $x$ to obtain a degree-$d$ equation vanishing on $Y$ but not at $x$. When $N=r+1$ use the equation of $Y$ itself. This proves the base-locus assertion; it does not assert generation of the ideal sheaf.

The bundle $3L$ is very ample \cite[Remark~7.3, p.~12]{MilneAV}. Embed $A$ by it and choose a detector $Y$ of degree $\delta_W$. Its embedded degree is $d=3^n\delta_W$. Degree-$d$ hypersurfaces through $Y$ restrict to a linear system on $A$ with base locus exactly $Y$. Choose $n$ members successively so that the next contains no irreducible component of the previous intersection. This is possible until the $n$th cut: the preceding components have dimension greater than $n$ and cannot be contained in $Y$. There are finitely many components, and avoiding each is the complement of a proper linear subspace of the space of sections.

More precisely, the successive intersections are local complete intersections in the smooth variety $A$, hence Cohen--Macaulay and without embedded associated primes. Avoiding their components makes the next equation a nonzerodivisor and preserves the expected pure dimension. The final intersection has pure dimension $n$ and contains $Y$, so $Y$ occurs in its fundamental cycle with positive multiplicity. Each divisor belongs to $|3dL|$. Thus $m_W\leq3d=3^{n+1}\delta_W$.
\end{proof}

The proposition replaces a free search for subvarieties by a search for components of fixed complete intersections. It supplies no fixed value of $m$ in an unproved family. The total complete intersection always has zero Weil projection; the required phenomenon is a splitting into components with nonzero projections that cancel.

\section{Degree costs in the CM constructions}
\label{sec:cm-cost}
This section identifies the part of the recent CM arguments that still needs geometric control to imply uniformity.

For fixed Galois CM field and fixed codimension, the finitely many CM types and four-factor relations can be treated at common auxiliary levels. In OpenAI's construction this is explicit in Corollary~6.7 \cite{OpenAICM}. Its Lemma~2.2 and Proposition~3.1, and van Hoften--Sempliner's Lemma~2.3.2 and Theorem~2.3.4 \cite{vHS}, then turn rational Hodge maps into morphisms by clearing lattice denominators and transport surface classes. These are pointwise constructions. The following calculation isolates a cost which remains even when the auxiliary data have been fixed.

\begin{proposition}[The degree of a transported support]\label{prop:transport}
Let $f:(A_0,L_0)\to(A,L)$ be an isogeny of complex polarized abelian varieties with $f^*c_1(L)=q\,c_1(L_0)$, where $q$ is a positive integer. For an integral $r$-dimensional subvariety $Z\subset A_0$,
\[
 \deg_L f(Z)=
 \frac{q^r\deg_{L_0}Z}
 {|\ker f\cap\operatorname{Stab}(Z)|}.
\]
If $\operatorname{Stab}(Z)$ is finite, this is at least
$q^r\deg_{L_0}Z/|\operatorname{Stab}(Z)|$.
\end{proposition}

\begin{proof}
The restriction $Z\to f(Z)$ has generic degree
$e=|\ker f\cap\operatorname{Stab}(Z)|$. Indeed a translate by a kernel element not stabilizing $Z$ meets it in a proper closed subset; remove the finitely many such intersections. Over what remains, the points identified by $f$ are exactly the stabilizer translates. Projection formula gives
$e\deg_L f(Z)=\int_Zf^*c_1(L)^r=q^r\deg_{L_0}Z$.
\end{proof}

\begin{corollary}[A limitation of a fixed collection of cycles]\label{cor:fixed-catalog}
Fix finitely many positive-dimensional integral subvarieties $Z_j\subset A_0$ of the same dimension $r$, each with finite stabilizer. Allow their images under polarized isogenies to a varying target, translations of those images, rational coefficients, and cancellation of equal components. Every nonzero resulting cycle, when written with an integral numerator, has a surviving component of degree at least
\[
 c\,q_{\min}(A)^r,\qquad
 c=\min_j\frac{\deg_{L_0}Z_j}{|\operatorname{Stab}(Z_j)|}>0,
\]
where $q_{\min}(A)$ is the smallest allowed polarization multiplier. Such a construction alone cannot give a uniform numerator-degree bound on a dense orbit in a positive-dimensional polarized family.
\end{corollary}

\begin{proof}
Every support has the asserted lower bound by Proposition~\ref{prop:transport}. Cancellation removes supports; it does not create supports of smaller degree. A surviving support has nonzero integer multiplicity in the numerator. Dividing by a rational denominator does not change that support degree.

For bounded $q$, only finitely many polarized targets occur. The degree of the isogeny is bounded by $q^{\dim A_0}$ times the fixed ratio of polarization volumes. There are finitely many subgroups of $A_0$ of bounded finite order, hence finitely many kernels and quotients. For each $q$ and quotient map, the target polarization is determined, because pullback on rational N\'eron--Severi groups is injective. Fixed finite level data add only finitely many choices. A dense orbit in a positive-dimensional family is infinite, so $q_{\min}$ is unbounded on it.
\end{proof}

This is a restriction on the stated method of constructing cycles. It is not a lower bound on the least degree among all representatives of their cohomology classes. The same transport issue is discussed in \cite{BBAlgebraic}. The fourfold constructions show why a different representative can be decisive.

\begin{lemma}\label{lem:surface-stabilizer}
If an integral surface $Z$ in a complex abelian variety has a positive-dimensional connected translation stabilizer, its cohomology class belongs to the algebra generated by divisor classes.
\end{lemma}

\begin{proof}
Let $B$ be that connected stabilizer and $\pi:A\to A/B$ the quotient. If $\dim B=2$, then $Z$ is a translate of $B$ and is the inverse image of a point; the point class is a rational multiple of a top power of an ample divisor. If $\dim B=1$, then $Z=\pi^{-1}(C)$ for a curve $C$ in the quotient. On a quotient of dimension $d\geq2$, hard Lefschetz identifies $H^2$ with $H^{2d-2}$ by multiplication by an ample class to the power $d-2$. The preimage of $[C]$ has type $(1,1)$, so is a rational divisor class by the Lefschetz theorem. Pullback expresses $[Z]$ in the divisor algebra. If $d=1$, the curve is the whole quotient and the assertion is immediate.
\end{proof}

In particular, a surface detecting a class orthogonal to the divisor algebra in complementary degree must have finite stabilizer. For such surfaces the support-degree cost above cannot be removed by an unbounded stabilizer.

\section{An optional Galois-orbit criterion}
\label{sec:galois}
This section gives a sufficient arithmetic condition for a bounded detector locus to fill a full unitary family. It uses Andr\'e--Oort and is separate from the proof of Theorem~\ref{thm:global}.

Fix a geometrically connected fine unitary Shimura base $S$ over a number field $F$, its polarized abelian scheme, and a labelled integral order action in $K$, all over that same $F$. Suppose that $G^{\rm ad}$ is $\Q$-simple of type $A$, as for the full signature-$(n,n)$ unitary datum. A special point is \emph{relative Weyl} if its defining torus has Galois action containing the Weyl group of this datum \cite[Definitions~5.3--5.4]{ChaiOort}. This refers to the unitary group; a full Weyl CM field for the ambient Siegel variety is not required.

Choose a locally closed embedding $S\hookrightarrow\mathbf P^N$. For a reduced closed set $V$ let $\operatorname{tdeg}(\overline V)$ be the sum of the degrees of all geometric irreducible components of its projective closure, in every dimension. This convention matters: a line with arbitrarily many additional isolated points has fixed top-dimensional degree.

Assume Andr\'e--Oort for every Cartesian power of $S$. Scanlon's automatic uniformity theorem \cite[Theorem~4.2]{Scanlon} gives a function $B_S(H)$ bounding the number of special points on a closed locus of total degree at most $H$ that lie on no positive-dimensional special subvariety contained in that locus. To obtain this form, parameterize bounded-degree reduced closed sets by finitely many Chow families. Automatic uniformity bounds the families of Zariski closures of their special points, and hence bounds the numbers of their zero-dimensional components. Andr\'e--Oort identifies all the positive-dimensional components as special. The resulting function is qualitative; no numerical rate is asserted.

\begin{proposition}\label{prop:galois-witness}
Let $\mathcal D_D=\{s:\delta_W(s)\leq D\}$ and choose any integer $H(D)$ at least $\operatorname{tdeg}(\overline{\mathcal D_D})$. Under the preceding hypotheses, a relative-Weyl special point $x$ satisfying
\[
 \delta_W(x)\leq D,\qquad
 |\operatorname{Gal}(\overline F/F)\cdot x|>B_S(H(D))
\]
forces $\mathcal D_D=S$. More generally, if $\mathcal D_D$ is proper, it contains at most $B_S(H(D))$ relative-Weyl special points.
\end{proposition}

\begin{proof}
First, the locus $\mathcal D_D$ descends over the fixed field $F$, independently of $D$. Choose one global order endomorphism $a$ with $a/\bar a$ not a root of unity. The integral operator $P(a^*+\bar a^*)$ in the proof of Theorem~\ref{thm:extraction} is defined over $F$ and detects exactly the Weil projection. On a projective parameter space for effective cycles of degree at most $D$, apply this operator to the universal class in one fixed $\ell$-adic cohomology theory. Its nonzero locus is a finite union of geometric connected components, stable under $\operatorname{Gal}(\overline F/F)$ because the operator and the class construction are defined over $F$. This open and closed subscheme descends over $F$; its proper image is precisely $\mathcal D_D$. Detection by a reducible effective cycle is equivalent to detection by at least one of its components. Thus the entire $F$-orbit of $x$ is contained in $\mathcal D_D$.

A relative-Weyl special point lies on no proper positive-dimensional special subvariety of $S$. Its Mumford--Tate torus maps onto the full adjoint Weyl torus \cite[Remark~5.8]{ChaiOort}. The group of a special subvariety through the point contains that torus. Its roots are stable under Galois and hence under the full Weyl group. In type $A$ the roots of each absolute simple factor form one Weyl orbit; $\Q$-simplicity makes Galois transitive on the factors. The subgroup therefore has either no roots, giving a zero-dimensional domain, or all roots, giving the whole connected domain. This is the calculation in \cite[Lemma~6.3]{ChaiOort}.

The same property holds for Galois conjugates, because Galois conjugates of special subvarieties are special. If $\mathcal D_D$ were proper, every point in the indicated orbit would therefore be one of its isolated maximal special points. The bound $B_S(H(D))$ contradicts the assumed orbit size. The same reasoning counts all relative-Weyl points in a proper $\mathcal D_D$.
\end{proof}

The detector degree $D$, the degree $H(D)$ of its locus in the base, and the special-point bound $B_S(H(D))$ are different quantities. Closedness supplies no growth estimate for $H(D)$. Automatic uniformity supplies no effective estimate for $B_S$. Most importantly, an orbit lower bound supplies no detector membership. A pointwise detector of degree $D(x)$ is useful here only if the comparison with $B_S(H(D(x)))$ is proved.

Without the relative-Weyl restriction, orbit size alone is insufficient. In $Y(1)^2$, the fixed special line $\{0\}\times Y(1)$ contains the full Galois orbits of $(0,j)$ for singular moduli $j$, with unbounded orbit sizes. On a curve base, in contrast, every proper closed locus is zero-dimensional, and its total degree directly bounds the size of an orbit it contains.

\section{What one finite-field detector would have to certify}
\label{sec:one-prime}
This section distinguishes flatness over the arithmetic base from the relative flatness that would make one finite-field cycle propagate through a whole family.

Let $\mathcal S$ be an integral scheme of finite type, flat over a localization of the ring of integers of a number field. Let $\mathcal A\to\mathcal S$ be a polarized abelian scheme with a fixed integral quadratic action whose characteristic-zero family is of Weil type. Fix an auxiliary prime $\ell$ invertible on $\mathcal S$. The integral polynomial projector of Section~\ref{sec:weil} defines the Weil summand in the corresponding lisse $\Q_\ell$-sheaf. For $D\ge1$, let $H_D\to\mathcal S$ be a projective parameter space for effective middle-dimensional cycles of degree at most $D$. Such a space is obtained by the construction of Lemma~\ref{lem:parameters}, using single cycles instead of pairs.

\begin{proposition}\label{prop:flat-detector-witness}
Suppose a geometric point $h\in H_D$ parametrizes a cycle with nonzero Weil projection, and the morphism $H_D\to\mathcal S$ is flat at $h$. Then every complex fibre of the family has a detector of degree at most $D$. The point $h$ may have positive characteristic.
\end{proposition}

\begin{proof}
Applying the projector to the universal cycle class gives a lisse section on $H_D$. Its nonzero locus $H_D^{W}$ is a union of geometric connected components. It is open and closed, descends over the base by Galois invariance, and is proper over $\mathcal S$. A flat morphism of finite presentation is open near a point where it is flat: restrict to the open flat locus \cite[Theorems~2.4.6 and 11.1.1]{EGAIV}. Thus the image of $H_D^{W}$ contains a nonempty open subset of $\mathcal S$. Properness makes the image closed, and integrality of $\mathcal S$ makes that image all of $\mathcal S$. In each complex fibre, nonzero projection of an effective cycle implies nonzero projection of at least one of its irreducible components, of degree no larger than $D$.
\end{proof}

For example, if a local-complete-intersection detector $Y$ on a geometric finite-field fibre has $H^1(Y,N_{Y/\mathcal A_s})=0$, the usual embedded deformation theory makes the relative Hilbert scheme smooth at $[Y]$. This gives the proposition's flatness hypothesis. The vanishing is sufficient, not necessary; relative flatness can hold even when this obstruction space is nonzero. No such detector with the required relative deformation property is constructed here in an unresolved eightfold family.

Flatness only over the ring of integers does not suffice. Consider the proper closed immersion
\[
 H=\{x=0\}\longrightarrow\mathcal S=\mathbf A^1_\Z.
\]
The scheme $H\cong\Spec\Z$ is flat over $\Z$ and has no vertical components. For any rational prime $p$, the section $s_p:\Spec\Z\to\mathcal S$ defined by $x=p$ has fibre product
\[
 H\times_{\mathcal S,s_p}\Spec\Z\cong\Spec\mathbf F_p.
\]
It has a point at $p$ and no characteristic-zero point. Since $p$ can be arbitrarily large, discarding a fixed set of vertical primes of the universal parameter space does not solve this problem. The obstruction is failure of flatness over the moduli base along the chosen section.

For a fixed CM fibre, Proposition~\ref{prop:lift} still gives a finite exceptional set depending on that fibre and on $D$. A degree-$D$ detector at a place outside that set suffices for that fibre. Proposition~\ref{prop:flat-detector-witness} is a different sufficient condition: a single detector together with relative flatness makes an entire parameter component dominate the moduli base. A construction must establish one of these properties; flatness over the arithmetic base alone establishes neither.

\section{Restrictions on finite-cover eightfold constructions}
\label{sec:finite-cover-eightfolds}
This section bounds the period images available from low-rank finite coefficient systems on curves. It restricts a method for constructing detectors; it gives no lower bound for the degrees of arbitrary cycles.

The reflection classification rests on Cohen's published classification and Malle's character-field tables. A recent unrefereed treatment by lelong\_ridge and cm\_symmetry\_scout \cite[Theorems~A--B and \S\S2,5]{ReflectionSources} classifies the primitive reflection possibilities for eightfold images of dimension at least fifteen and proves a primitive reflection-free coefficient-rank bound of twenty-three for that threshold. The two statements below use the sixteen-dimensional threshold for dominance. We give the induction and dimension arguments explicitly, including the case in which a complex block has no descent to the original character field.

\begin{theorem}[Gaussian reflection obstruction]\label{thm:gaussian-reflections}
Let $K=\Q(i)$. Consider a family of finite covers of smooth complex
curves and an absolutely irreducible finite coefficient representation
$\rho:G\to\operatorname{GL}_d(K)$ with exact character field
$\Q(\operatorname{tr}\rho)=K$. Suppose that its multiplicity in first
cohomology has dimension eight over $K$, and that its faithful coefficient
image contains a complex reflection. Then the corresponding period image
has complex dimension at most $15$. In particular it does not dominate
the signature-$(4,4)$ unitary domain of dimension $16$.
\end{theorem}

\begin{proof}
We give the reductions, including the normalizer and field-of-definition
steps. The classification inputs are the published results of Cohen
\cite{CohenReflection,CohenCorrection} and the reflection character fields
in Malle's tables \cite{MalleFields}. No unrefereed algebraicity theorem
is used.

\paragraph{The dimension bound and polarized induction.}
Let $V$ be a nontrivial irreducible finite coefficient system of rank $r$
on a genus-$\gamma$ curve, with genuine local monodromies $g_1,\ldots,g_b$.
Put $c_j=\operatorname{rank}(g_j-1)$. The Euler characteristic of its middle extension,
together with the absence of invariants in $V$ and its dual, gives
\[
 \dim_K H^1(Y,j_*V)=2r(\gamma-1)+\sum_{j=1}^b c_j.
 \label{eq:curve-euler}
\]
For rank-eight multiplicity and $r\ge2$, the corresponding Hurwitz
dimension is consequently at most
\[
 3\gamma-3+b\le8+(3-2r)(\gamma-1).
 \label{eq:curve-dimension}
\]
Positive-genus bases give a bound of eight; a genus-zero base gives
$2r+5$. For a nontrivial rank-one system, instead,
\[
 8=2\gamma-2+b,
 \qquad 3\gamma-3+b=\gamma+7\le12.
 \label{eq:rank-one-curve-dimension}
\]
This includes the unramified case. Fixed finite cyclic monodromy adds
only finitely many choices of root line bundles over a marked curve and
branch divisor, so the rank-one estimate bounds its period image. The middle-extension and
variation conventions agree with \cite[Proposition~1.1, Lemma~1.2,
Remark~1.3 and Theorem~2.5]{DettweilerWewers}.

If $\rho=\operatorname{Ind}_H^G\sigma$ over $K$, Shapiro identifies its multiplicity
cohomology with that of $\sigma$ on the intermediate curve
$C/H$. With $H_0=\ker\sigma$, the required faithful cover is
$C/H_0\to C/H$. This identification extends across branch points:
the stalks are local inertia invariants on both sides. It respects the
Hodge structure and the polarization up to the fixed nonzero sheet-trace
factor, since it is induced by the finite covering correspondences.
Thus the period map factors through this intermediate-cover family.
Forgetting additional cover data or points with trivial coefficient
monodromy does not enlarge its period image.

\paragraph{Reflection blocks are defined over $K$.}
Choose a positive invariant Hermitian form with entries in $K$ by
averaging. The root line $\operatorname{im}(s-1)$ of every reflection
$s$ is defined over $K$. Join two root lines when they are not orthogonal.
The connected components span mutually orthogonal $K$-subspaces $V_a$.
Their sum is nonzero and $G$-stable, hence is all of $V$.
The reflection normal closure is the product of its irreducible
reflection groups on these subspaces. The full group permutes the
blocks transitively. If $H$ stabilizes one block and $\sigma$ is its
action, then
\[
 \rho\simeq\operatorname{Ind}_H^G\sigma\quad\hbox{over }K,
 \qquad
 K\subset\Q(\operatorname{tr}\sigma)\subset K.
\]
The first inclusion follows from the induced-character formula. Thus
the faithful block stabilizer also has exact character field $K$.
It suffices, by polarized Shapiro, to treat this stabilizer and its
irreducible reflection subgroup $W$.

\paragraph{The classification and its low-rank cases.}
The groups that can occur for $W$ are rank-one groups, imprimitive
$G(m,p,r)$, rational real Weyl groups, and the nonreal primitive groups
$G_8,G_{29},G_{31}$, of ranks $2,4,4$ respectively. This is the
character-field restriction of Cohen's classification and Malle's
Tables~8.1--8.2; the omissions explained after those tables are
$G_4,G_{25}$ over the Eisenstein field and the real groups
$G_{28}=F_4$, $G_{35}=E_6$. Noncrystallographic real fields do not
embed in $K$.

For an imprimitive group of rank at least three, its diagonal elements
with entries $(\zeta_m,\zeta_m^{-1},1,\ldots)$ force
$\zeta_m+\zeta_m^{-1}\in\Q$. Hence $m\in\{1,2,3,4,6\}$.
The determinant-one diagonal element with entries
$(\zeta_m,\zeta_m,\zeta_m^{-2},1,\ldots)$ excludes $3$ and $6$.
The irreducible imprimitive cases therefore have $m=2$ or $4$.
All blocks of rank at most four have period dimension at most $13$
by the dimension estimate, with no restriction on their possible normalizer extensions.
This covers in particular triality in rank four.

\paragraph{High-rank monomial blocks.}
Cohen's Lemma~(2.7) gives a unique system of coordinate lines for
$G(m,p,r)$ when $r\ge5$. Every linear normalizer therefore preserves
that system. The reflection representation is conjugate over $K$ to
its usual monomial model over $K$: the intertwining equations have a
one-dimensional solution space over $K$, and a nonzero intertwiner is
invertible by absolute irreducibility. Thus its individual coordinate
lines, not merely their unordered collection, are defined over $K$.
The full block stabilizer is monomial on these lines.

For a line stabilizer $L$, its action is a finite character
$\psi:L\to K^\times$, and $\sigma=\operatorname{Ind}_L^H\psi$.
Exactness of the character field in the induced-character identity forces
$\Q(\psi)=K$, so its image is $\mu_4$. Polarized Shapiro realizes
the entire rank-eight multiplicity as a rank-one system on an
intermediate curve. The rank-one estimate bounds the period image by $12$.
Rank-one reflection blocks are covered by the same argument.

\paragraph{The rational high-rank normalizers.}
The remaining possibilities are $W(A_r)$ and $W(E_6),W(E_7),W(E_8)$.
In their rational reflection models, a finite-order $K$-matrix $g$
normalizing $W$ has the form $g=cA$ with $c\in K^\times$ and
$A\in\operatorname{GL}_r(\Q)$: solve the rational intertwining equations for the
automorphism of $W$ induced by $g$ and use absolute irreducibility.
The matrix $A$ is a rational similarity for the invariant positive
form and permutes root lines. All roots have the same length. Sending
one rational root to a rational multiple of another shows that its
similarity factor is a rational square. After dividing by that
rational multiple, $A$ is an orthogonal root-system automorphism.
It has finite order, so the remaining scalar is in $\mu_4$.

A root-system automorphism is a Weyl element followed by a Dynkin
diagram automorphism, by moving its chamber back to a fixed chamber.
For $A_r$ and $E_6$ the extra diagram involution is $-w_0$;
$E_7,E_8$ have no diagram automorphism. Hence the finite $K$-normalizer
is contained in $\mu_4W$. A subgroup containing $W$ and having exact
character field $K$ contains a nonreal coset, hence contains $iI$.
The faithful block stabilizer is therefore exactly $\mu_4W$.

For $E_6,E_7,E_8$ we may assume genus zero. In a generating
product-one branch tuple, at least two entries are outside the real
subgroup. Each is $i$ times a real orthogonal matrix, so its fixed
space has dimension at most $\lfloor r/2\rfloor$: the relevant
nonreal eigenvalues of the real matrix occur in conjugate pairs.
Each such entry moves at least $\lceil r/2\rceil$ dimensions.
Using the Euler formula gives
\[
 b\le2r+8-2(\lceil r/2\rceil-1).
\]
After subtracting three for automorphisms of the sphere, the period
dimensions are at most $13,13,15$ respectively.

\paragraph{All ranks in type $A$.}
It remains to avoid a rank-dependent gap in type $A$.
Assume $r\ge5$ and put $n=r+1$. Here $-I\notin W(A_r)$, so the
normalizer is the direct product $\mu_4\times S_n$, acting by
$\rho=\chi\otimes\operatorname{Std}_n$, with $\chi$ the scalar character.
We may again assume a genus-zero base. Let $b_\chi$ count branch
entries with nontrivial scalar phase. A phase $\pm i$ has moving
rank at least
\[
 a=r-\lfloor n/4\rfloor,
\]
because each permutation cycle divisible by four supplies at most
one eigenvector of the required eigenvalue. A phase $-1$ has moving
rank at least $c=r-\lfloor n/2\rfloor$. There are at least two
odd-phase entries, since the phases generate $\mu_4$ and their
product is one.

If $b_\chi\ge3$, choose two odd-phase entries and a third nontrivial
phase. All remaining genuine entries move at least one dimension.
Consequently
\[
 2r+8\ge2a+c+(b-3),\qquad
 b\le11-r+2\lfloor n/4\rfloor+\lfloor n/2\rfloor\le12.
\]
The period dimension is at most nine.

If $b_\chi=2$, the two phases are $i,-i$, and
$H^1(\mathbb P^1,j_*\chi)=0$ by the Euler formula. The representation identity
\[
 \chi\otimes\operatorname{Perm}_n
  =(\chi\otimes\operatorname{Std}_n)\oplus\chi
  =\operatorname{Ind}_{\mu_4\times S_{n-1}}^{\mu_4\times S_n}\psi
\]
therefore identifies the entire rank-eight cohomology in question
with the rank-one system $\psi$ on the degree-$n$
intermediate curve. There is no additional cohomology summand,
because the discarded $\chi$-summand vanishes. Its period dimension
is at most twelve by the rank-one estimate. These alternatives exhaust type $A$.

All reflection blocks and their finite compatible normalizers have
now been covered. Their bounds are at most fifteen, proving the theorem.
\end{proof}

\begin{theorem}[A coefficient-rank threshold for dominance]\label{thm:finite-cover-rank}
Let $V$ be a faithful irreducible complex representation of a finite group $G$, with nonrational character. Consider a Hurwitz family of $G$-covers whose $V$-multiplicity in first cohomology has complex dimension eight and Hodge numbers $(4,4)$. Suppose that every genuine branch monodromy has moving rank at least two. If its period image has dimension sixteen, then
\[
 \dim_{\C}V\ge26.
\]
The hypothesis concerns the branch monodromies; the group may contain reflections absent from their tuple. No descent of a complex imprimitivity system to a number field is assumed.
\end{theorem}

\begin{proof}
We first justify induction over $\C$, since it is needed even when the original coefficient representation has a smaller field of definition. If $V=\operatorname{Ind}_H^G W$ and $C$ is the covering curve, Frobenius reciprocity gives
\[
 \operatorname{Hom}_G(V,H^1(C,\C))
 \simeq\operatorname{Hom}_H(W,H^1(C,\C)).
\]
This is an isomorphism of variations of complex Hodge structures. To see the extension across the branch points, put $X=C/H$ and take the $W$-multiplicity in the finite pushforward of the constant sheaf from $C$ to $X$. A stalk at a branch point is the corresponding inertia-invariant space. Finite proper pushforward has no higher direct images and taking a finite-group isotypic part is exact, so this is the middle extension of the coefficient local system on the punctured curve. The maps on cohomology commute with the Hodge decomposition, since they come from finite sums of deck transformations.

The inducing representation $W$ is irreducible: otherwise induction of a proper subrepresentation would contradict irreducibility of $V$. It is nontrivial, since induction of the trivial representation contains constants. Its variation depends only on the finite cover $C/\ker(W)\to X$. If $r=\dim W$, $\gamma=g(X)$ and $b$ is its number of genuine coefficient branch points, the Euler formula is
\begin{equation}\label{eq:complex-cover-budget}
 8=2r(\gamma-1)+\sum_{j=1}^{b}c_j,
 \qquad c_j=r-\dim W^{g_j}\ge1.
\end{equation}
For fixed finite monodromy data the covers have finite fibres over the moduli of the marked base curve: its fundamental group has finitely many homomorphisms to the fixed finite group. Consequently the period-image dimension is at most $3\gamma-3+b$. For $r=1$ this is $\gamma+7\le12$. For $r\ge2$ it is at most
\begin{equation}\label{eq:complex-shapiro-bound}
 2r+5+(3-2r)\gamma\le2r+5.
\end{equation}
There is no unstable unmarked genus-one exception, because such a nontrivial finite system cannot have eight-dimensional cohomology. On a simply connected parameter neighborhood the Shapiro isomorphism is constant on the marked local systems. The filtration at one complex embedding determines the unitary period point, with its conjugate filtration determined by it. Thus forgetting the additional cover data does bound the original period image by \eqref{eq:complex-shapiro-bound}, without asserting number-field descent of $W$ or of a polarization on an intermediate abelian variety.

Assume for a contradiction that $d=\dim V\le24$ and the image has dimension sixteen. Rank one has already been excluded. For $d\ge2$ a positive-genus base gives image dimension at most eight, so the base is the sphere. Erase identity branch monodromies. The resulting product-one generating tuple has $b\ge19$, and
\begin{equation}\label{eq:rank25-budget}
 2d+8=\sum_{j=1}^{b}c_j\ge2b\ge38.
\end{equation}
Thus $15\le d\le24$. Some entry has moving rank two, because otherwise the sum would be at least $3b\ge57$, whereas $2d+8\le56$.

Suppose first that $G$ is imprimitive, with $t\ge2$ blocks of common dimension $r$ and $d=rt$. A block stabilizer gives an inducing representation of dimension $r$. Equation~\eqref{eq:complex-shapiro-bound} excludes $r\le5$, so $r\ge6$. Let $\sigma_j\in S_t$ be the permutation of blocks induced by the $j$th branch entry, and put $I_j=t-\#\{\text{cycles of }\sigma_j\}$. An operator cycling $\ell$ blocks has fixed space of dimension at most $r$, since a fixed vector is determined by its entry in one block. Hence $c_j\ge rI_j$. The block monodromy is transitive, and Riemann--Hurwitz for its degree-$t$ cover of the sphere gives
\[
 \sum_j I_j=2(t-1)+2\gamma'\ge2(t-1).
\]
If $J=\{j:I_j>0\}$, then $|J|\le\sum I_j$, while the other branch entries have rank at least two. Therefore
\[
 \sum c_j\ge2(b-|J|)+r\sum I_j
       \ge38+2(r-2)(t-1).
\]
Together with $d=rt$ and \eqref{eq:rank25-budget}, this implies $r+2t\ge17$. But $rt\le24$, $r\ge6$ and $t\ge2$ permit only $t=2,3,4$, for which the respective maxima of $r+2t$ are $16,14,14$. This contradiction excludes imprimitivity.

Next suppose that a tensor decomposition $V=A\otimes B$ is preserved, with $2\le a=\dim A\le s=\dim B$. For a finite-order operator $x\otimes y$, the two factors are diagonalizable up to a scalar: suitable powers of each are scalar. If $x$ is nonscalar, its eigenspaces have dimensions at most $a-1$, so the fixed space of $x\otimes y$ has dimension at most $(a-1)s$. Its moving rank is at least $s$. Since $d\ge15$, we have $s>2$, so a moving-rank-two element must have scalar $x$. Its remaining moving rank is a positive multiple of $a$, forcing $a=2$. Thus $d=2s$, with $s\le12$.

The projective action of $G$ on $A$ is nontrivial, since otherwise every subspace $\ell\otimes B$ would be invariant. A product-one generating tuple must therefore contain at least two entries nonscalar on $A$. They each have moving rank at least $s$; the other entries have rank at least two. It follows that
\[
 4s+8=\sum c_j\ge2s+2(b-2)\ge2s+34,
\]
or $s\ge13$, a contradiction. Here tensor indecomposability means that no such decomposition with both factors preserved exists; permutations of equal tensor factors are included in the convention of the classification theorem used next.

The published Theorem~7.1 and Remark~(3) of Guralnick--Saxl \cite[pp.~548--549]{GS} now apply: a finite irreducible primitive tensor-indecomposable complex group of degree greater than ten containing a moving-rank-two element is, up to scalars, the deleted permutation representation of the alternating or symmetric group. The other natural large-degree cases in that theorem have defining positive characteristic, and its exceptional characteristic-zero degrees are at most ten. Thus every element of $G$ has the form
\[
 g=\lambda\operatorname{Std}(\pi),\qquad\pi\in S_{d+1}.
\]
The scalar is unique, since no nonidentity permutation acts as a scalar on the deleted representation in these degrees. It therefore defines a character of $G$. If this character were trivial, the representation would have rational character. There are consequently at least two branch entries with $\lambda\ne1$.

A nonidentity eigenvalue of a permutation on $d+1$ letters has multiplicity at most $\lfloor(d+1)/2\rfloor$: each contributing cycle has length at least two and contributes one eigenvector. Removing the constant line does not change such an eigenspace. Each of the two nontrivial scalar entries thus has moving rank at least $\lfloor d/2\rfloor$. The branch budget gives
\[
 2d+8\ge2\lfloor d/2\rfloor+2(b-2)
       \ge2\lfloor d/2\rfloor+34.
\]
For even $d$ this forces $d\ge26$, and for odd $d$ it forces $d\ge25$. Both contradict $d\le24$. It remains to exclude $d=25$.

\paragraph{The degree-twenty-five case.}
The base is again the sphere, $b\ge19$, and $\sum c_j=58$. If the
representation is imprimitive, its block rank is either one or five.
The Shapiro bounds above give period-image dimension at most twelve
or fifteen, respectively. Thus $G$ is primitive. We show that a
primitive finite group of degree twenty-five with an element moving
at most four dimensions lies in
$\C^*\operatorname{Std}(S_{26})$.

First consider a preserved tensor structure. Its only nontrivial
factorization is $5\cdot5$. A nonscalar pure tensor has every
 eigenspace of dimension at most $4\cdot5=20$, by the eigenspace
argument above. A factor-interchanging operator has the form
$(A\otimes B)\tau$, with $\tau$ interchanging the factors. It is
conjugate to $(I\otimes C)\tau$, $C=BA$, whose square is
$C\otimes C$. Finite order makes $C$ semisimple. In a diagonal basis
there are five diagonal lines and ten two-dimensional off-diagonal
blocks. Each block contributes at most one vector to any fixed
 eigenspace. Its eigenspaces therefore have dimension at most fifteen.
Neither tensor case contains an element moving at most four dimensions.

The structural reduction in the proof of
\cite[Theorem~7.1, p.~549]{GS} applies to a primitive group with no
such tensor structure. If its Fitting subgroup is noncentral, it is
of extraspecial type, and four conjugates of a nonscalar element
 generate a subgroup containing a nontrivial scalar. Their fixed
spaces have zero intersection, so that element moves at least
$25/4$ dimensions. Otherwise there is a quasisimple normal subgroup
$E$ acting irreducibly on $V$. Indeed a nontrivial multiplicity space
or more than one commuting simple factor would give a tensor
structure. This representation of the matrix group $E$ is faithful.

Hiss--Malle's characteristic-zero classification
\cite[Theorem~1 and Table~2, pp.~22 and 31]{HissMalle}, with its
complete corrected table \cite[pp.~95--96 and 102]{HissMalleCorrection},
gives exactly the following degree-twenty-five possibilities:
\[
 E=A_{26},\quad \operatorname{PSL}_2(25),\quad
 \operatorname{PSL}_2(49),\quad \operatorname{PSp}_4(7).
\]
The alternating representation is deleted permutation. The two
$\operatorname{PSL}_2$ degrees are $q$ and $(q+1)/2$; the only other
ordinary row is $\operatorname{PSp}_4(7)$. The degree-twenty-five row
for ${}^3D_4(2)$ is in characteristic three and does not apply.

Here is the generation estimate, including scalar and outer lifts.
For a nonscalar $g$ with a nonzero fixed space, let $m$ be its
projective order, choose a prime $p\mid m$, and put $x=g^{m/p}$.
Then $x$ has projective order $p$. The scalar $x^p$ fixes every fixed
vector of $g$, so it is the identity. Thus $x$ has prime order and
moves no more dimensions than $g$. If $\alpha$ conjugates of its
induced automorphism generate a group containing $E/Z(E)$, the
intersection of the corresponding fixed spaces on $V$ is zero.
Otherwise their projective group would preserve a line, contradicting
irreducibility of $E$. Consequently
\[
 25\le\alpha\operatorname{rank}(x-I).
\]
For both $\operatorname{PSL}_2$ groups,
\cite[Lemma~3.1, pp.~527--529]{GS} gives $\alpha\le4$, including
field and diagonal automorphisms. For $\operatorname{PSp}_4(7)$,
\cite[Theorem~4.3, p.~533]{GS} gives $\alpha\le5$; its transvection
exception has the smaller value four. These bounds exclude an
 element moving at most four dimensions. The automorphism is
nontrivial, since an element centralizing the irreducible $E$ is
scalar, and a nonidentity scalar has no fixed vectors.

It remains that $E=\operatorname{Std}(A_{26})$. Its automorphism group
is $S_{26}$, and Schur's lemma then gives the full normalizer
$\C^*\operatorname{Std}(S_{26})$. This proves the claimed normalizer
assertion. In the branch tuple some entry moves at most three
 dimensions, since otherwise $58\ge4b\ge76$; hence it applies here.

Write each branch entry as $\lambda_j\operatorname{Std}(\pi_j)$.
A scalar phase outside $\{1,-1\}$ has finite order $k\ge3$, and a
fixed vector requires a cycle of $\pi_j$ whose length is divisible
by $k$. Therefore its moving rank is at least
\[
 25-\lfloor26/k\rfloor\ge17.
\]
If every scalar phase belonged to $\{1,-1\}$, the generated group
would have rational character. At least one phase is outside this
set; the product-one relation forces at least two. The remaining
branches each have moving rank at least two, giving
\[
 58=\sum c_j\ge2\cdot17+2(b-2)\ge68,
\]
a contradiction. This excludes degree twenty-five and proves the
 stated threshold.
\end{proof}

\begin{corollary}\label{cor:gaussian-cover-rank}
A dominant eightfold finite-cover source with an absolutely irreducible coefficient representation defined over $\Q(i)$ and having exact character field $\Q(i)$ has coefficient rank at least twenty-six.
\end{corollary}
\begin{proof}
Theorem~\ref{thm:gaussian-reflections} excludes a reflection in the faithful coefficient image. Every genuine branch entry therefore has moving rank at least two, and Theorem~\ref{thm:finite-cover-rank} applies.
\end{proof}

The degree-twenty-five profile with eighteen moving ranks three and
one moving rank four satisfies the Euler budget, but the proof shows
that it cannot occur in a dominant nonrational source. The distinction
between containing and being generated by small-moving elements matters:
$\mu_4\operatorname{Std}(S_{26})$ has nonrational character and contains
rank-three involutions, whereas all its rank-three involutions are
ordinary triple transpositions and generate a rational subgroup. A
positive construction in a larger degree still needs a finite group
with the correct character field, a split polarization, Hodge numbers
$(4,4)$ and a period differential of rank sixteen.



\subsection{Scalar-standard coefficients in every degree}

The alternating-group exception can be excluded without a degree
bound: its permutation realization forces the variation to come
from a small pointed curve.

\begin{theorem}\label{thm:scalar-standard-all}
Let a finite-cover source over $\mathbf P^1$ have an irreducible
complex coefficient representation $\rho$ of degree $d\ge26$,
compact cohomology multiplicity eight, and image contained in
$\C^*\operatorname{Std}(S_{d+1})$. Its period image has dimension
at most fifteen. This assertion does not require a particular
character field or a lower bound on the branch moving ranks.
Together with Theorem~\ref{thm:finite-cover-rank}, it excludes
scalar-standard coefficients in every degree for a dominant
nonrational source whose branch moving ranks are at least two.
\end{theorem}

\begin{proof}
Put $N=d+1$. The scalar and permutation factors are unique, so
\[
 \rho=\kappa\otimes\operatorname{Std}(\pi),\qquad
 \pi:G\longrightarrow S_N,\quad\kappa:G\longrightarrow\C^*
\]
are homomorphisms. Irreducibility makes the permutation action
transitive. Let $m$ count the branch places with nontrivial scalar
phase. At any such place, a fixed vector of $\rho$ comes from a
nontrivial eigenvalue of the permutation. If the phase has order
$k\ge2$, at most $\lfloor N/k\rfloor$ permutation cycles contribute,
so its moving rank is at least
\[
 d-\lfloor N/2\rfloor=\lfloor d/2\rfloor.
\]
The compact Euler formula therefore gives
\[
 m\lfloor d/2\rfloor\le2d+8.
\]
Since $5\lfloor d/2\rfloor>2d+8$ for $d\ge26$, we obtain $m\le4$.

Let $q:Y\to\mathbf P^1$ be the connected degree-$N$ intermediate
cover defined by the permutation action. Let $\mathcal L$ be the
rank-one finite local system of $\kappa$, and let $\mathcal L_Y$
be its pullback. Compact Shapiro and
$\operatorname{Perm}_N=\mathbf1\oplus\operatorname{Std}$ give
\begin{equation}\label{eq:standard-shapiro}
 H^1(Y,j_*\mathcal L_Y)
 =H^1(\mathbf P^1,j_*\mathcal L)\oplus H^1_\rho.
\end{equation}
Here points with trivial upstairs inertia are filled in. To check
this convention at a branch point, a permutation cycle of length
$e$ contributes an invariant stalk exactly when the scalar phase
to the power $e$ is one; this is also exactly when the upstairs
rank-one system extends across that point. The finite-pushforward
projection formula therefore identifies the middle extensions as
well as their unramified restrictions. The equality is compatible
with Hodge structures and Gauss--Manin connections.

If $\kappa$ is nontrivial, its compact cohomology on the sphere has
dimension $m-2$, so the left side of
\eqref{eq:standard-shapiro} has dimension
\[
 n=m+6\le10.
\]
Suppose first that $\mathcal L_Y$ is nontrivial. If $Y$ has genus
$h$ and the rank-one system has $r$ genuinely ramified points, then
\[
 n=2h-2+r.
\]
The finite-order local system adds only discrete choices to an
$r$-pointed genus-$h$ curve. For $h\ge2$ its moduli dimension is
\[
 3h-3+r=n+h-1\le\tfrac32 n\le15.
\]
For $h=1$ it is $r=n\le10$, and for $h=0$ it is
$r-3=n-1\le9$. If $\mathcal L_Y$ is trivial, then $2h=n\le10$,
and the unpointed curve moduli dimension is at most twelve.
Finally, if $\kappa$ itself is trivial, the first summand in
\eqref{eq:standard-shapiro} vanishes and $Y$ has genus four,
giving dimension at most nine.

These bounds control the selected summand, not merely the full
left side of \eqref{eq:standard-shapiro}. Work locally on a simply
connected marked neighborhood of the cover family. That summand
is a fixed flat subspace, and its Hodge filtration is the intersection
with the ambient filtration. A tangent vector fixing the pointed
curve and its finite local system therefore fixes the summand's
filtration as well. Additional motions of the map to $\mathbf P^1$,
or of points whose upstairs inertia is trivial, cannot increase
its period rank. This proves the asserted bound.
\end{proof}

\subsection{Ordered tensor factors require degree at least forty-eight}

A fixed tensor factor imposes a further branch budget in every degree,
independently of the character field.

\begin{proposition}\label{prop:ordered-tensor48}
Suppose a finite irreducible coefficient representation
$V=A\otimes B$, with $2\le a=\dim A\le b=\dim B$, preserves this
ordered tensor decomposition.  If a sphere-cover family has compact
cohomology multiplicity eight and period image of dimension sixteen,
then
\[
 2ab+8\ge(a+1)b+18a-a^2,\qquad ab\ge48.
\]
The assertion does not cover elements permuting tensor factors.
\end{proposition}

\begin{proof}
Let $g_1,\ldots,g_N$ be the nonidentity branch matrices.  They generate
the coefficient image, have product one, and $N\ge19$.
Choose a representative $x_i$ of their projective action on $A$,
scaled so that a largest eigenspace has eigenvalue one.  Put
\[
 \ell_i=a-\dim\ker(x_i-1),\qquad L=\sum_i\ell_i.
\]
Finite projective order makes each $x_i$ diagonalizable.  The product
of these representatives is scalar.  The projective action on $A$
is irreducible, since a proper invariant subspace would tensor with
$B$ to a proper invariant subspace of $V$.

Delete the scalar $x_i$.  At least two entries remain: a single
nonscalar matrix cannot have scalar product.  If $L\le a$, the
intersection of the fixed spaces of all but the last remaining
entry has dimension at least the positive moving rank of that last
entry.  It is nonzero and proper.  All preceding entries fix it
pointwise; the scalar product relation makes the last entry scalar
on it.  This contradicts irreducibility.  Thus $L\ge a+1$.

Write $c_i=\operatorname{rank}(g_i-1)$.  If $\ell_i>0$, an eigenspace
of the $B$ component matches at most one eigenspace of $x_i$, of
dimension at most $a-\ell_i$.  Hence $c_i\ge b\ell_i$.
If $\ell_i=0$, the nonidentity matrix $g_i$ moves at least one whole
$A$ factor, so $c_i\ge a$.  With $r$ the number of positive $\ell_i$,
we have $r\le L$, and therefore
\[
 \sum_i c_i\ge bL+a(N-r)
    \ge aN+(b-a)L
    \ge19a+(b-a)(a+1).
\]
The Euler equality $\sum_i c_i=2ab+8$ proves the first inequality.
Rearranging it gives $(a-1)b\ge18a-a^2-8$.
For $a=2,3,4,5,6$, the least possible $b$ is respectively
$24,19,16,15,13$, giving products $48,57,64,75,78$.
For $a\ge7$, the product is already at least $49$.
This proves the degree bound.  Equality at $(a,b)=(2,24)$ is only
a numerical possibility, with no group or cover existence implied.
\end{proof}

\subsection{An imprimitive Gaussian representation of degree ninety-nine}

A larger coefficient rank does not by itself remove the branch-cost
obstruction. The following example isolates the extra cost imposed by
its permutation action on blocks.

\begin{proposition}\label{prop:degree99-cover}
Let $G$ be the elementary-abelian extension $2^{10}.M_{11}$ identified
in the Character Table Library by
\texttt{P48/G1/L1/V2/ext2}, and let $\rho$ be either of its conjugate
faithful degree-$99$ characters with character field $\Q(i)$.
A sphere-cover source for $\rho$ with compact cohomology multiplicity
eight has at most seven nontrivial coefficient branch points. Its
period image therefore has dimension at most four. The assertion
also holds after finite scalar enlargement of the coefficient group.
\end{proposition}

\begin{proof}
The group construction and its table are given in
\cite[\S1.2-9]{CTblLib}. Write $N$ for the elementary abelian normal
subgroup of order $1024$. The seven $G$-classes in $N$, in increasing
library order, give the following restriction data:
\[
\begin{array}{c|rrrrrrr}
\text{class size}&1&55&330&396&66&165&11\\
\chi&99&63&27&-9&-9&-45&-81.
\end{array}
\]
The two conjugate characters have the same restriction, since these
values are real. Thus
\[
 \langle\chi|_N,\chi|_N\rangle_N
 =\frac{99^2+55\cdot63^2+330\cdot27^2
       +462\cdot9^2+165\cdot45^2+11\cdot81^2}{1024}
 =891.
\]
Clifford theory expresses this restriction as $m$ times an orbit of
$n$ distinct linear characters of $N$. Since $mn=99$ and
$m^2n=891$, we obtain $m=9$ and $n=11$. The corresponding eigenspaces
are eleven blocks of dimension nine. Their permutation action
factors through $M_{11}$, is transitive, and is faithful because
$M_{11}$ is simple and the action is nontrivial.

The degree-eleven transitive permutation character of $M_{11}$ is
\[
 (11,3,2,3,1,0,1,1,0,0).
\]
For clarity, its identification does not require choosing a block
labeling: the least nontrivial irreducible character degree is ten,
so a degree-eleven transitive permutation character is one plus a
degree-ten irreducible; exactly one such sum has nonnegative integral
values. The class power maps give permutation indices
\[
 (0,4,6,6,8,8,8,8,10,10).
\]
Equivalently, the number of permutation cycles is the multiplicity
of the eigenvalue one, which gives these indices by subtraction
from eleven. In particular every nonzero index is at least four.

Let $i_j$ be the block-permutation index of a coefficient branch
matrix and let $c_j$ be its moving rank. A permutation cycle on
$\ell$ blocks contributes at most nine fixed vectors and hence at
least $9(\ell-1)$ moving vectors. Therefore $c_j\ge9i_j$, with any
scalar phase allowed. If $i_j=0$, the matrix lies in $N$ up to a
scalar and acts as a scalar on each block. A nonidentity matrix
then moves at least one whole block, so $c_j\ge9$. This inequality
also follows from the first one when $i_j>0$. Since a nonzero
$i_j$ is at least four, the two bounds imply
\[
 c_j\ge9+\frac{27}{4}i_j.
\]

The branch permutations generate a transitive degree-eleven group
and have product one. Riemann--Hurwitz for the associated degree-eleven
cover gives $\sum i_j\ge20$. If there are $b$ coefficient branches,
\[
 206=2\cdot99+8=\sum_jc_j
     \ge9b+\frac{27}{4}\sum_ji_j
     \ge9b+135.
\]
Consequently $b\le7$, and the branch-configuration dimension is
$b-3\le4$. All estimates were independent of scalar phases, proving
also the last assertion.
\end{proof}

The numerical restriction data can be verified from the published
permutation-group construction, without accepting an abstract table
as a group-realization certificate. Computing that group's order gives
$8110080$; its solvable radical is elementary abelian of order $1024$;
and its computed ordinary character table has $52$ classes and agrees
with the cited table. The \emph{Gaussian finite-cover checks} in the
\href{\workbenchurl}{workbench} reproduce these exact calculations.
A scan of the $2678$ named tables in version~1.3.11 selects $942$
faithful exact-Gaussian characters of degrees $25$ through $250$.
Only the two characters of this proposition pass the initial necessary
inequality $19\nu\le2d+8$, where $\nu$ is the least nonzero projective
moving rank. This describes a finite search and its coverage; it is
not a classification of finite Gaussian linear groups. In particular,
scalar twists of rational representations need not occur in that list.

\subsection{A degree-one-hundred-ten coefficient over an imaginary quadratic field}

The same exact test over the other imaginary quadratic character
fields leaves one additional pair in that finite catalogue.  Its
branch costs already contradict the compact cohomology dimension.

\begin{proposition}\label{prop:degree110-cover}
Let $G$ be the group $2^{10}.M_{11}$ identified in the Character Table
Library by \texttt{P48/G1/L1/V1/ext2}, and let $\rho$ have either
faithful degree-$110$ character with exact character field
$\Q(\sqrt{-2})$.  No sphere-cover source for $\rho$ with compact
cohomology multiplicity eight has a sixteen-dimensional period image.
The assertion permits arbitrary finite scalar enlargement.
\end{proposition}

\begin{proof}
The ordinary character table and its power maps give least nonzero
moving rank ten, even after arbitrary scalar phases.  The Euler
budget is $2\cdot110+8=228$.  If there are at least nineteen branch
points, each must have moving rank at most $228-18\cdot10=48$.
Here is the complete list of class and scalar pairs with moving
rank between two and forty-eight; the conjugate characters have the
same list:
\[
\begin{array}{c|r|r}
 \text{class}&\text{scalar}&\text{moving rank}\\\hline
 2a&1&20\\
 2b&1&40\\
 2e&-1&30\\
 2f&-1&10
\end{array}
\]
For each class of order $s$, the power-map traces determine the
multiplicity of every $s$th root of unity by finite Fourier inversion.
A scalar phase can give a fixed vector only when it is the inverse
of one of those eigenvalues; every other phase has moving rank
$110$.  Enumerating these inverse phases gives exactly the four
rows displayed.  Thus all branch costs are divisible by ten, whereas
$228$ is not.  This is impossible before any generation or
polarization condition is imposed.

There is also a quotient check.  The nontrivial degree-ten character
of the table with index two is trivial on the scalar centre and on
all four displayed classes.  A product-one tuple generating the
projective image would have, in that character, total moving rank
at least twenty, by the nonnegativity of its compact first cohomology.
The four classes would give total moving rank zero.  This gives
the same exclusion independently of the divisibility test.
\end{proof}

For context, scanning all $2678$ named ordinary tables in version
1.3.11 selects $1592$ faithful characters of degree $26$ through
$250$ with exact imaginary quadratic character field other than
$\Q(\sqrt{-3})$.  They lie in $371$ named tables and nineteen
imaginary quadratic fields.  Only the two conjugate pairs in
Propositions~\ref{prop:degree99-cover} and
\ref{prop:degree110-cover} satisfy the necessary minimum-cost
inequality, even with arbitrary scalar phases allowed.
The \emph{quadratic finite-cover checks} in the
\href{\workbenchurl}{workbench} give the exact field tests,
power-map computations, branch costs and rational certificates.
These are computations in a specified finite catalogue; they do
not classify all finite groups, assert a Schur index, or construct
an integral coefficient lattice.

\subsection{Rational coefficient characters and arbitrary scalar phases}

There is a separate restriction when the coefficient character before
scalar enlargement is rational.  It permits rank-one moving parts and
all finite scalar phases.

\begin{lemma}\label{lem:rational-scalar-budget}
Let $\rho$ be a finite irreducible representation of degree $d>1$
with rational character.  Suppose that a product-one tuple
$M_j=\lambda_j\rho(g_j)$ generates an irreducible group with nonreal
character, has at least nineteen genuine branches, and has parabolic
first cohomology of dimension eight.  Let $\nu$ be the least positive
moving rank among all scalar multiples of elements of $\rho(G)$,
and let $\eta$ be the least moving rank with nonreal scalar phase.
Then
\[
 2\eta+17\nu\le 2d+8,\qquad \eta\ge\lceil d/2\rceil.
\]
\end{lemma}
\begin{proof}
A scalar in the rational-character image has rational trace $d\zeta$,
so it is $1$ or $-1$.  The product-one condition consequently makes
$\prod_j\lambda_j$ real.  All real phases would give rational
character on every generated word.  At least one phase is therefore
nonreal, and the real product forces at least two such phases.
The Euler formula gives $\sum_j\operatorname{rank}(M_j-1)=2d+8$.
Two terms cost at least $\eta$, and at least seventeen others cost
at least $\nu$.  Finally the rational character gives equal
multiplicities to conjugate eigenvalues.  A nonreal eigenspace has
dimension at most $\lfloor d/2\rfloor$, which proves the last bound.
\end{proof}

\begin{proposition}[A bounded-degree symmetric-family exclusion]
\label{prop:symmetric-scalar-exclusion}
Under the hypotheses of Lemma~\ref{lem:rational-scalar-budget},
$\rho$ cannot be a nonstandard irreducible representation of any
$S_s$ of degree $26\le d\le500$, or a scalar enlargement thereof.
Here nonstandard excludes the trivial, sign, standard and
sign-twisted standard representations.  This conclusion is a finite
exact computation together with the degree cutoff proved below.
\end{proposition}
\begin{proof}
For a cycle class of order $h$, rationality makes the multiplicity
$m_q$ of an individual primitive $q$th root independent of its choice.
Fourier inversion, grouping powers by their order, gives
\[
 m_q=\frac1h\sum_{r\mid h}\chi(g^{h/r})c_r(h/q),\qquad
 c_r(e)=\mu\left(\frac r{(r,e)}\right)
       \frac{\varphi(r)}{\varphi(r/(r,e))},\quad q\mid h.
\]
The second expression is the sum of the primitive $r$th roots to
the $e$th power, obtained by inclusion--exclusion over divisors.
Thus the real-phase costs are $d-m_1,d-m_2$, and the nonreal
minimum is $d-\max_{q>2}m_q$.  A phase whose inverse is absent
from the spectrum costs $d$ and needs no additional enumeration.
Zero cost is omitted from the genuine branches.

The exact Murnaghan--Nakayama calculation in the
\emph{rational scalar checks} of the
\href{\workbenchurl}{workbench} enumerates the $117$ nonstandard
partitions up to transposition with $5\le s\le33$ and degree
between $26$ and $500$.  Transposition tensors with sign, which
preserves all phase-cost minima.  For every cycle partition the
calculation checks nonnegative integral multiplicities and
\[
 \sum_{q\mid h}\varphi(q)m_q=d,\qquad
 \sum_{q\mid h}\mu(q)m_q=\chi(g).
\]
Every selected character has
$2\eta+17\nu-(2d+8)>0$.  The smallest excess is $76$,
at the partition $(7,2)$ of $9$, of degree $27$, with
$\nu=6$ and $\eta=18$.  For comparison, $(31,2)$ of $33$
has degree $495$, $\nu=30$, $\eta=330$ and excess $172$.
The stored rows give the degrees, minima and witnessing cycle
partitions; a separate integer hook-length calculation verifies
the completeness of the partition selection.

It remains to justify the cutoff for all $s$.  A nonstandard
partition of $s$ has first row and first column at most $s-2$.
For $s\ge6$ it has a removable corner leaving a nonstandard
partition of $s-1$.  If both lengths are at most $s-3$, any
corner works.  If the first row has length $s-2$, the remaining
two boxes form $(2)$ or $(1,1)$; removing the last box of the
first row works.  The column case is its transpose, and the two
exceptional bounds cannot hold simultaneously for $s\ge6$.
The restriction branching rule, or equivalently the decomposition
of standard tableaux by the corner containing their largest entry,
makes the degree at least that of this predecessor.
The independent exact hook-length enumeration of all $12310$
partitions of $34$ gives smallest nonstandard degree $527$,
attained only by $(32,2)$ and its transpose.  Repeated removal
therefore excludes degrees at most $500$ for every $s\ge34$.
For $s<5$ all degrees are below $26$.  This proves the assertion.
\end{proof}

The analogous named-table calculation in CTblLib version~1.3.11
reads all $2678$ ordinary table names and selects $12479$ faithful
rational-character rows of degrees $26$ through $500$, occurring
in $1365$ named tables.  Every row violates
Lemma~\ref{lem:rational-scalar-budget}, with smallest strict
excess $28$, even with arbitrary scalar phases and moving rank
one allowed.  Exact cyclic Fourier inversion is used in every
case, with dimension, trace, integrality and conjugation checks.
This last assertion has precisely the scope of that finite
catalogue.  It does not exhaust unlisted finite groups, tensor
products or proper irreducible projective subgroups.

\section{The quadratic relations of a general Gaussian eightfold}
\label{sec:theta-eightfolds}
This section tests the most direct extension of van Geemen's fourfold construction and identifies exactly where that extension stops.

For an even principal Gaussian abelian variety of dimension $2n$,
with the theta structure used in \cite{vG96}, van Geemen's
Corollary~2.4 gives
\[
 \dim H^0(A,2\Theta)_-=2^{n-1}(2^n-1).
\]
Here the subscript denotes the minus eigenspace of the induced
involution on second-order theta sections.
His Lemma~3.3 identifies the base locus of the minus projection on an indecomposable member as $2^{2n}$ two-torsion points \cite[pp.~621--622]{vG96}. For $n=2$, Proposition~3.6 produces a smooth quadric fourfold in the resulting $\mathbf P^5$; Theorem~3.7 then proves that the pullback of a ruling plane detects the Weil classes. For $n=4$, the same general statements give $120$ sections, a map to $\mathbf P^{119}$, and $256$ base points. They do not give a quadric eightfold in $\mathbf P^9$.

We examine the precise space of quadratic relations obtained by applying the same symmetry to the theta identities. Put $G=\mathbf F_2^4$ and let $V$ have basis $(e_u)_{u\in G}$. The second-order theta coordinates form a matrix $(X_{uv})_{u,v\in G}$. The involution transposes this matrix, so its minus eigenspace identifies with $\bigwedge^2V$. The vanishing theta constants in \cite[Lemma~2.3 and \S3.5]{vG96} give the $120$ quadratic forms
\begin{equation}\label{eq:theta-forced}
 Q_{a,b}(X)=\sum_{u,v\in G}(-1)^{(u+v)\cdot b}
 X_{uv}X_{u+a,v+a},\qquad a\cdot b=1.
\end{equation}
A linear combination descends through the minus projection when it vanishes on the plus space of symmetric matrices. There are no mixed terms, since each form is invariant under transposition.

\begin{theorem}\label{thm:theta-rank}
The combinations of \eqref{eq:theta-forced} that descend to the minus eigenspace form a vector space of dimension $35$. Every nonzero form in this space has rank at least $24$, and rank $24$ occurs. In particular this space contains no rank-$10$ quadric.
\end{theorem}

We first prove this linear-algebra assertion. We then show that these
are all the quadratic relations of the minus projection on a nonempty
Zariski-open subset of the even principal Gaussian family.

\begin{proof}
For $a\cdot b=1$ define the skew matrix
\[
 (P_{a,b})_{u,w}=(-1)^{u\cdot b}\,\mathbf1_{w=u+a}.
\]
For fixed nonzero $a$, the eight permitted characters form a basis of the skew matrices supported on displacement $a$. Different displacements have disjoint support. Consequently these $120$ matrices form a basis of $\bigwedge^2V^*$.

Associate to $Q=\sum c_{a,b}Q_{a,b}$ the symmetric tensor
$T=\sum c_{a,b}P_{a,b}\otimes P_{a,b}$ on $\bigwedge^2V$. Substitution gives
\[
 Q(vv^{\mathsf t}+ww^{\mathsf t})=2T(v\wedge w,v\wedge w).
\]
Since every complex symmetric matrix is a sum of rank-one symmetric matrices, $Q$ vanishes on the symmetric matrices if and only if $T(v\wedge w,v\wedge w)=0$ for every $v,w$. Polarizing first in $w$ and then in $v$ gives
\[
 T(x\wedge y,x\wedge z)=0,\qquad
 T(x\wedge y,w\wedge z)=-T(w\wedge y,x\wedge z).
\]
Together with the two wedge antisymmetries, these identities say that $T$ is induced by an alternating four-form $\xi\in\bigwedge^4V^*$:
\[
 T_\xi(x\wedge y,z\wedge w)=\xi(x,y,z,w).
\]
Exchanging the middle two indices shows that the restriction of $Q$ to the skew matrices is a nonzero scalar multiple of this form. Thus their ranks agree.

Let $\tau_s e_u=e_{u+s}$ and $\epsilon_t e_u=(-1)^{t\cdot u}e_u$. On the skew matrices $P_{a,b}$ these operations have characters $(-1)^{b\cdot s}$ and $(-1)^{a\cdot t}$. The pairs of characters are distinct. Hence their invariant symmetric tensors are exactly the diagonal tensors $\sum c_{a,b}P_{a,b}^2$. It follows that the descending forms identify with the four-forms invariant under every $\tau_s$ and $\epsilon_t$.

Sign invariance forces $\xi(e_u,e_v,e_w,e_z)=0$ unless $u+v+w+z=0$. For distinct indices this means that they form an affine two-plane. Translation identifies the coefficients on the four cosets of a fixed two-dimensional subspace $P\subset G$. Translation inside $P$ is a double transposition on its four vertices and preserves its wedge orientation. There is therefore one free coefficient $c_P$ for each such subspace. Their number is
\[
 \frac{(16-1)(16-2)}{(4-1)(4-2)}=35.
\]

It remains to prove the sharp rank bound. Call $P$ active when $c_P\ne0$, let $U$ be the union of the nonzero vectors in the active planes, and put $m=|U|$. For $a\in U$, let $k_a$ count the active planes through $a$. Then
\begin{equation}\label{eq:plane-incidence}
 1\le k_a\le\left\lfloor\frac{m-1}{2}\right\rfloor,
 \qquad \sum_{a\in U}k_a=3\#\{P:c_P\ne0\}.
\end{equation}
The first bound holds because the other two nonzero vectors in different planes through $a$ are disjoint pairs.

There is an orthogonal decomposition
\[
 \bigwedge^2V=\bigoplus_{a\ne0}E_a,
 \qquad E_a=\langle e_x\wedge e_{x+a}:x\in G\rangle,
 \quad \dim E_a=8.
\]
Indeed the sign characters distinguish these summands. Choose a complement $H$ to $\langle a\rangle$ in $G$. In the basis $f_x=e_x\wedge e_{x+a}$, $x\in H$, translation invariance gives the convolution matrix
\[
 T_\xi(f_x,f_y)=b_a(y-x),\qquad
 b_a(z)=\xi(e_0,e_a,e_z,e_{z+a}).
\]
Here $H\cong\mathbf F_2^3$, $b_a(0)=0$, and $b_a$ has precisely $k_a$ nonzero values. Its eight character vectors diagonalize the matrix. If $r_a$ is its rank, then $r_a$ is the size of the Fourier support of $b_a$, and the total rank is $\sum_{a\in U}r_a$.

We record the required elementary Fourier facts. If a nonzero function $b$ on $\mathbf F_2^3$ satisfies $b(0)=0$, write $k=|\operatorname{supp}b|$ and $r=|\operatorname{supp}\widehat b|$. Then
\[
\begin{array}{c|cccccc}
\text{condition}&k=1&k=2&k=3&\text{any }k&r=2&r=3\\
\hline
\text{conclusion}&r=8&r\ge4&r\ge6&r\ge2&k=4&k=6.
\end{array}
\]
For two terms, the Fourier values have the forms $c_1+c_2$ and $c_1-c_2$, each four times. For three terms, factor out one character. The remaining two character differences are independent, so the four values are $c_1\pm c_2\pm c_3$, each twice. At most one vanishes, since two vanishing values would force a coefficient to be zero. A Fourier support of size one contradicts $b(0)=0$. A Fourier support of size two gives $b=c(\chi_1-\chi_2)$, supported on four points. For a Fourier support of size three apply the three-term argument to Fourier inversion: the value at zero vanishes, so exactly one of the four patterns vanishes, giving support size six. In particular a block with $k\le5$ has rank at least four unless $k=4$, when rank two is possible.

We now bound the sum of the ranks by considering $m$. If $m=3$, one plane is active and its three blocks have rank eight, giving $24$. The value $m=4$ is impossible. If $m=5$, exactly two planes are active and meet in one nonzero vector. No third plane can be contained in their union: a vector space cannot be covered by two proper subspaces. Four blocks have $k=1$, and the remaining block has $k=2$, giving rank at least $36$.

If $6\le m\le8$, every $k_a\le3$, so every active block has rank at least four and the total is at least $24$. If $m=9$ and $k_a=4$, the four planes through $a$ pair the other eight elements of $U$ under translation by $a$. Thus $U\cup\{0\}$, of size ten, is translation invariant under $a$. Its translation stabilizer has order a power of two dividing ten, and therefore has order at most two. At most one $a$ can have $k_a=4$. The rank is at least $8\cdot4+2=34$.

For $m=10$, the baseline rank is $20$. A total below $24$ would exclude $k=1$ and $k=3$ blocks, whose excesses above two are at least six and four. It would permit at most one $k=2$ block. All other $k$'s would be four, so $\sum k_a$ would be $40$ or $38$, contrary to divisibility by three in \eqref{eq:plane-incidence}. For $m=11$, $k_a\le5$, so rank three is impossible. A total below $24$ would make every block rank two and every $k_a=4$, giving the impossible sum $44$. Finally, if $m\ge12$, the rank is at least $2m\ge24$.

To attain $24$, take the invariant orbit sum for a single two-plane $P$. Its four terms are volume forms on the four disjoint coordinate four-spaces indexed by its affine cosets. Each volume form induces a nondegenerate pairing of rank six on the exterior square of its four-space. Mixed wedges are in the radical. The resulting rank is $4\cdot6=24$.
\end{proof}

\begin{lemma}\label{lem:generic-gaussian-theta-nulls}
Let $n\ge2$ and consider the even principal Gaussian period domain
\[
 \mathcal D_n=
 \left\{\tau=
 \begin{pmatrix}X&Y\\-Y&X\end{pmatrix}\in\mathfrak H_{2n}:
 X^{\mathsf t}=X,\ Y^{\mathsf t}=-Y\right\}.
\]
On a nonempty open subset, the only vanishing even theta constants
are those whose two genus-$n$ characteristics are equal and odd.
Their number is $2^{n-1}(2^n-1)$.
\end{lemma}

\begin{proof}
Write a characteristic as a pair $(m_1,m_2)$, where
$m_j=(a_j,b_j)\in\mathbf F_2^n\times\mathbf F_2^n$ and its parity is
$a_j\cdot b_j$. An even characteristic of genus $2n$ has the same
parity in the two factors. Changing summation variables
$(u,v)\mapsto(-v,u)$ in the theta series shows
\[
 \vartheta_{m_1,m_2}(\tau,0)
   =(-1)^{a_2\cdot b_2}\vartheta_{m_2,m_1}(\tau,0).
\]
Thus the equal odd pairs vanish identically, as in
\cite[Lemma~2.3]{vG96}.

At $Y=0$ the constant factors as
$\vartheta_{m_1}(X,0)\vartheta_{m_2}(X,0)$.
Every even theta constant is a nonzero holomorphic function on
$\mathfrak H_n$: symplectic change of marking reduces it to the
zero characteristic, nonzero at a product of elliptic periods.
The product of two nonzero holomorphic functions on this connected
domain is nonzero. Therefore a pair of even characteristics does
not vanish identically on $\mathcal D_n$.

Suppose $m_1,m_2$ are distinct odd characteristics, and put
$g_m(X)=\nabla_z\vartheta_m(X,0)$.
Differentiating the convergent theta series in a skew matrix
direction $B$ gives, at $Y=0$,
\[
 d_Y\vartheta_{m_1,m_2}(X,0)[B]
   =\frac{1}{2\pi i}\,g_{m_1}(X)^{\mathsf t}B\,g_{m_2}(X).
\]
It suffices to find one $X$ where these two gradients are not
proportional.

The integral symplectic group is transitive on ordered pairs of
distinct odd characteristics. Here is the finite-dimensional reason.
Characteristics are quadratic refinements of the symplectic form
on $\mathbf F_2^{2n}$, and parity is their Arf invariant.
Two odd refinements differ by pairing with a unique nonzero vector
$v$ isotropic for the first refinement. The orthogonal group of
that odd refinement is transitive on its nonzero isotropic vectors:
complete $v$ to a hyperbolic pair, and identify the orthogonal
complements using the standard normal form of a nonsingular
quadratic form over $\mathbf F_2$. The finite symplectic changes
lift to integral ones, since symplectic transvections generate
the group and lift integrally.

We may consequently reduce to two odd characteristics with their
only odd elliptic component in the first and second coordinates,
respectively. At a product of elliptic periods their gradients
are nonzero multiples of the first and second coordinate vectors.
Indeed the odd elliptic theta function has a simple zero at the
origin and each even elliptic theta constant is nonzero.
The theta transformation formula takes an odd gradient to a
nonzero scalar times one common invertible linear transformation
of that gradient; its other differentiated term is zero because
the odd theta constant is zero. Independence of the two gradients
therefore persists after undoing the change of marking.
The displayed derivative is nonzero for a suitable skew $B$.

Every even theta constant other than an equal odd pair is thus
nonzero somewhere on the connected domain $\mathcal D_n$.
There are only finitely many such constants, so the complement
of their zero loci is nonempty and open. The algebraic form of
the conclusion also follows below from the rank of the
multiplication map.
\end{proof}

\begin{corollary}\label{cor:all-eightfold-quadrics}
On a nonempty Zariski-open subset of each marked even principal
Gaussian eightfold component, the entire quadratic ideal of the
minus-projected image in $\mathbf P^{119}$ is the
$35$-dimensional space in Theorem~\ref{thm:theta-rank}.
Every nonzero quadric in that ideal has rank at least $24$.
In particular no ten sections of this minus eigenspace can map
the general fibre dominantly to a smooth quadric eightfold.
\end{corollary}

\begin{proof}
Let $g=8$ and let $X_\sigma$, $\sigma\in\mathbf F_2^g$, be the
second-order theta coordinates. The quadratic forms
\[
 Q_{a,b}=\sum_\sigma(-1)^{\sigma\cdot b}
                 X_\sigma X_{\sigma+a},
 \qquad a\cdot b=0,
\]
form a basis of all quadratic forms: for a fixed displacement $a$,
the permitted characters form the Fourier basis on the unordered
pairs $\{\sigma,\sigma+a\}$; when $a=0$ they form the basis on
the squares. The addition formula
\cite[\S3.5, pp.~622--623]{vG96} says
\[
 Q_{a,b}\bigl(\vartheta[\sigma;0](2\tau,2z)\bigr)
 =\vartheta[a;b](\tau,0)\,\vartheta[a;b](\tau,2z).
\]
The sections $\vartheta[a;b](\tau,2z)$, as $(a,b)$ ranges over the
even characteristics, form a basis of the even sections of $4\Theta$.
It follows that the full Kummer quadratic ideal is precisely the
span of the quadrics for vanishing even theta constants.

Lemma~\ref{lem:generic-gaussian-theta-nulls} makes this kernel,
on a nonempty open set, the span of the $120$ forms
\eqref{eq:theta-forced}. A quadric on the minus projection
pulls back to a polynomial involving only minus coordinates.
Its vanishing is equivalent to membership in the full kernel.
The intersection of that kernel with these polynomials is
exactly the $35$-dimensional space already computed.
Conversely every such polynomial vanishes on the projected image.

After a finite level cover the spaces of sections and their
multiplication are algebraic vector bundles and an algebraic
bundle map. The minimal kernel dimension is attained on a
Zariski-open set. Thus the nonempty analytic open just found
gives the asserted nonempty Zariski-open set.
Ten independent sections defining a map to a smooth $Q^8$
would pull its equation back to a nonzero rank-$10$ quadratic
form on the minus space, contradicting the rank bound.
\end{proof}

The smallest nonzero relation on this open set has a smooth
quotient quadric of dimension $22$. Its degree-eight cohomology
is below the middle and is generated by a hyperplane power.
Pulling back that class gives a multiple of $\theta^4$, whose
Weil projection is zero. Thus the direct ruling-class argument
cannot use any quadric of the general minus projection.
Special fibres may have additional relations. Cycles involving
singular loci, sections outside the minus space, and different
target varieties require other arguments.

\begin{proposition}[What a quadric eightfold would give]\label{prop:quadric-gate}
Let $(A,L)$ be a Weil eightfold and $\Delta=\chi(L)$. Suppose ten sections of $2L$ define a dominant rational map to a smooth quadric $Q^8\subset\mathbf P^9$, with base locus of dimension at most three. Suppose that, for a general member $P\cong\mathbf P^4$ of one ruling, the closure of its inverse image has nonzero Weil projection. Then
\[
 m_W(A,L)\le2,\qquad \delta_W(A,L)\le8\cdot8!\Delta=322560\Delta.
\]
\end{proposition}

\begin{proof}
A general $\mathbf P^5$ containing a ruling $P$ cuts $Q^8$ in $P\cup P'$, where $P'$ is a member of the other ruling. The $\mathbf P^5$ is cut out by four hyperplanes in $\mathbf P^9$. By transitivity of the orthogonal group and the dimension theorem for a general translate, the inverse images of a general such pair of ruling planes have dimension four away from the base locus. The base locus has dimension at most three. Thus the four corresponding members of $|2L|$ have proper intersection. Their intersection contains the closure of the inverse image of $P$, with positive multiplicities. Its nonzero Weil projection implies that one irreducible component detects. Proposition~\ref{prop:complete-intersections} gives
$\delta_W\le(8!\Delta/2)2^4=8\cdot8!\Delta$.
\end{proof}

The geometric hypotheses of this proposition remain unproved in
the Gaussian family. Corollary~\ref{cor:all-eightfold-quadrics}
excludes obtaining its rank-$10$ quadric from any ten sections
of the general minus eigenspace.

\section{A Weil character whose perfect representative is obstructed}
\label{sec:equivariant-eightfolds}
This section separates the construction of a Weil class in $K$-theory from the deformation of a perfect complex representing it.

Perry's Theorem~1.2 \cite{Perry}, an unrefereed preprint, extends the semiregularity method used by Markman. In a family of abelian varieties it permits a perfect complex $E$ with $\operatorname{Ext}^{<0}(E,E)=0$, a finite group $G\subset A\times A^\vee$ acting by translations and degree-zero line bundles, and weak $G$-semiregularity. If the prescribed twisted Chern character remains Hodge, the theorem deforms $E$ as a twisted perfect complex and propagates algebraicity. An isolated CM automorphism is not automatically one of the permitted group actions. More fundamentally, a virtual class with the right Chern character need not have a semiregular representative.

We give an explicit calculation at the split Gaussian eightfold
\[
 A_0=X\times Y,
 \qquad X=Y=E_i^4,
 \qquad a(x,y)=(ax,\bar a y),
\]
with product principal polarization. Its Hermitian lattice is $I_{4,4}$. Let $\Delta\subset X\times Y$ be the diagonal. If $a=2+i$, write $S=a^*+\bar a^*$ for the sum of pullback operations on coherent $K$-theory, and put
\[
 P(s)=(s+1170)(s+350)(s-750)(s-1250).
\]

\begin{proposition}\label{prop:graph-obstruction}
The integral class $v=P(S)[\mathcal O_\Delta]$ has a perfect representative $E_v$ satisfying
\[
 \operatorname{ch}(E_v)=-339429556224\,e_W[\Delta]\ne0,
 \qquad \operatorname{Ext}^{<0}(E_v,E_v)=0.
\]
Its whole Chern character remains Hodge along the signature-$(4,4)$ unitary domain. Nevertheless its first-order obstruction map on that domain has rank six. No infinitesimal gerbe twist removes those six normal obstruction directions. For every finite translation and degree-zero-twist action under which $E_v$ is equivariant, they form a six-dimensional subspace of the kernel of the invariant semiregularity map.
\end{proposition}

The representative in the proposition is the direct sum of the graph structure sheaves in the signed expansion of $v$, with a shift by one on the negative summands. The proposition concerns this representative and its translates, degree-zero twists, and finite equivariant inductions. It is not a nonexistence assertion for every perfect complex with a Weil component.

\begin{proof}
Choose coordinates $x_1,\ldots,x_4$ and $y_1,\ldots,y_4$ on the two tangent spaces. The polarized $K$-linear deformation tangent space consists of
\begin{equation}\label{eq:unitary-tangent}
 \kappa_B=\begin{pmatrix}0&B\\B^{\mathsf t}&0\end{pmatrix},
 \qquad B\in M_4(\C).
\end{equation}
Indeed commuting with the complex multiplication action kills the two diagonal Beltrami blocks; compatibility with the polarization makes the full matrix symmetric. The resulting space has dimension $16$.

For a graph with tangent equation $y=Tx$, the normal quotient is $(v_x,v_y)\mapsto v_y-Tv_x$. Restricting \eqref{eq:unitary-tangent} to its antiholomorphic tangent and taking this quotient gives the normal obstruction
\[
 \nu_T(B)=B^{\mathsf t}-TB\bar T
 \in H^1(\Gamma_T,N_{\Gamma_T/A_0}).
\]
In particular, for the diagonal and for every scalar graph of norm-one slope,
\begin{equation}\label{eq:graph-normal-obstruction}
 \nu_{\lambda I}(B)=B^{\mathsf t}-B,
 \qquad \lambda\bar\lambda=1.
\end{equation}
Its kernel is the ten-dimensional space of symmetric matrices and its rank is six. The matrix $B=E_{12}$ gives the nonzero example $E_{21}-E_{12}$.

We check that this is an obstruction for the structure sheaf itself. The automorphism $(x,y)\mapsto(x,y-x)$ sends the diagonal to $X\times\{0\}$. K\"unneth and the Koszul resolution of the point of $Y$ give
\begin{equation}\label{eq:diagonal-ext}
 \operatorname{Ext}^2(\mathcal O_\Delta,\mathcal O_\Delta)
 =H^2(X,\mathcal O_X)
 \oplus\bigl(H^1(X,\mathcal O_X)\otimes T_{0,Y}\bigr)
 \oplus\bigwedge^2T_{0,Y}.
\end{equation}
The dimensions are $6,16,6$. Huybrechts--Thomas \cite[Corollary~3.4, manuscript p.~14]{HT} identify the obstruction of a perfect complex as the product of its Atiyah class with the Kodaira--Spencer class. In normal coordinates $u=y-x$, the Koszul Atiyah class has normal symbol $\sum_jdu_j\otimes e_j$. Its contraction is
\[
 \operatorname{ob}_{\kappa_B}(\mathcal O_\Delta)
 =\sum_{j,p}(B^{\mathsf t}-B)_{jp}\,
 d\bar x_p\otimes e_j,
\]
in the middle summand of \eqref{eq:diagonal-ext}. These constant harmonic representatives are independent. An infinitesimal gerbe class $\beta\in H^2(A_0,\mathcal O_{A_0})$ acts as $\beta|_\Delta\operatorname{id}$ in the first summand. It therefore cannot cancel a nonzero normal component.

For a rational norm-one slope, Hilbert~90 writes $\lambda=\bar c/c$ after clearing denominators. The $K$-endomorphism $(x,y)\mapsto(cx,\bar c y)$ takes the diagonal isogenously to that graph and extends over the unitary family. Pullback to the corresponding finite \'etale cover reduces the obstruction to the calculation above. Translations and tensoring by degree-zero line bundles preserve it.

The diagonal detects the Weil plane: a generator of one complex Weil line is
\[
 dx_1\wedge\cdots\wedge dx_4
 \wedge d\bar y_1\wedge\cdots\wedge d\bar y_4,
\]
whose restriction to $\Delta$ is a nonzero volume form. Its tangent and normal bundles are trivial, so Grothendieck--Riemann--Roch gives $\operatorname{ch}(\mathcal O_\Delta)=[\Delta]$, with no other components. On the five paired $K$-weights in degree eight, $S$ has eigenvalues
\[
 5^j\operatorname{Tr}_{K/\Q}(a^{8-2j}),\quad0\le j\le4,
 \qquad (-1054,-1170,-350,750,1250).
\]
The first is the Weil weight. Hence $P(S)$ kills the other weights and acts there by
$P(-1054)=-339429556224$.

Every isogeny pullback of $\mathcal O_\Delta$ is the structure sheaf of a reduced disjoint union of torsion translates of scalar norm-one graphs. Expand $v=\sum_Zm_Z[\mathcal O_Z]$, combining identical supports, and define
\[
 E_v=\bigoplus_{m_Z>0}\mathcal O_Z^{\oplus m_Z}
 \ \oplus\!
 \bigoplus_{m_Z<0}\mathcal O_Z[1]^{\oplus(-m_Z)}.
\]
This proves the Chern-character assertion. For distinct integral supports of the same dimension, $\operatorname{Hom}(\mathcal O_Z,\mathcal O_{Z'})=0$: a nonzero image inside the torsion-free rank-one sheaf $\mathcal O_{Z'}$ would have generic support $Z'$, forcing $Z'=Z$. These Hom groups are the only possible negative self-Ext groups introduced by the shifts zero and one. Thus $E_v$ is gluable, meaning that its negative self-Ext groups vanish.

Its Atiyah obstruction is block diagonal, with the nonzero normal classes \eqref{eq:graph-normal-obstruction} on every surviving summand. A shift changes the Chern-character trace but does not cancel different diagonal blocks in self-Ext. Since at least one coefficient $m_Z$ is nonzero, the obstruction map still has rank six. The preceding decomposition also shows that a gerbe class cannot cancel it.

Finally the Buchweitz--Flenner semiregularity map is the trace of the powers of the Atiyah class \cite[pp.~136--137]{BF}. With compatible sign conventions its components satisfy
\[
 \sigma_E^{2,q}(\operatorname{ob}_{\kappa}(E))
 =\kappa\mathbin{\lrcorner}\operatorname{ch}_{q+1}(E).
\]
This follows by expanding the exponential defining the Chern character and using cyclicity of trace. The only nonzero component of $\operatorname{ch}(E_v)$ is a flat Weil class, which remains Hodge along the whole unitary domain. Every displayed contraction therefore vanishes. The six-dimensional obstruction image is in the kernel of $\sigma_{E_v}^2$.

Translations preserve the constant tensors $\kappa_B$, and degree-zero line bundles have zero Atiyah class. Naturality makes these obstruction classes invariant under every compatible finite action of $G\subset A_0\times A_0^\vee$. A multiplicity representation fixes the identity endomorphism in each block, and a permutation of supports permutes the compatible obstruction blocks. Thus the same six-dimensional subspace lies in the kernel on $\operatorname{Ext}^2(E_v,E_v)^G$. Equivariant induction adds such blocks and does not remove them.
\end{proof}

In Perry's theorem the representative also fails the weak equivariant condition. Indeed, with $B_0=0$, all its other hypotheses are satisfied. Weak equivariant semiregularity would give a twisted deformation over an \'etale neighborhood in the unitary family. An \'etale base change preserves the tangent space, contradicting the nonzero obstruction even after an arbitrary infinitesimal gerbe twist. This last conclusion uses the semiregularity theorem in \cite{Perry}; the explicit obstruction calculation above uses the published Atiyah and semiregularity formulas.

Products of fourfold objects face a related issue. For $F\in D(A_1)$ and $H\in D(A_2)$, the Atiyah class of $F\boxtimes H$ is the sum of the two exterior Atiyah classes. A mixed Kodaira--Spencer direction $\eta_1\otimes v_2\in H^1(\mathcal O_{A_1})\otimes H^0(T_{A_2})$ therefore gives
\[
 (\eta_1\operatorname{id}_F)\boxtimes
 (v_2\mathbin{\lrcorner}\operatorname{At}(H))
 \in\operatorname{Ext}^1(F,F)\otimes\operatorname{Ext}^1(H,H).
\]
Semiregularity on the two factors controls the product locus and does not by itself remove this mixed obstruction. A construction for the full eightfold domain must address these extra directions.

\section{An exact obstruction-space criterion for a pure Weil character}
\label{sec:pure-weil-complexes}
This section turns strong semiregularity for a pure Weil character into an exact dimension test. It then constructs such a character on a Gaussian eightfold and proves why a natural family of complexes representing it fails the test.

Let $A$ have dimension $g=2n\ge4$, let $K$ act with signature $(n,n)$, and write $W_K\subset H^{2n}(A,\Q)$ for its Weil plane. For a perfect complex $E$, its Hochschild semiregularity map is
\[
 \sigma_E^2:\operatorname{Ext}^2(E,E)\longrightarrow HH_{-2}(A).
\]
If a finite group $G\subset A\times A^\vee$ acts coherently by translations and tensor products with degree-zero line bundles, and $E$ is $G$-equivariant, call $E$ \emph{strongly $G$-semiregular} when this map is injective on the invariant self-Ext space. The coherence and the equivariant structure are part of the hypothesis. A finite group of other automorphisms is not covered by this definition here.

\begin{theorem}[The exact dimension test]\label{thm:pure-weil-ext}
Suppose that $B\in H^{1,1}(A)\cap H^2(A,\Q)$ and
\[
 e^B\operatorname{ch}(E)=w\in W_K\setminus\{0\}.
\]
Then, for the preceding action and equivariant structure,
\[
 E\text{ is strongly }G\text{-semiregular}
 \quad\Longleftrightarrow\quad
 \dim_{\C}\operatorname{Ext}^2(E,E)^G
     =2\binom{2n}{2}=2n(2n-1).
\]
The assertion includes the trivial group. Every such equivariant object has invariant second self-Ext dimension at least $2n(2n-1)$. Neither simplicity nor vanishing of negative self-Ext is needed for this equivalence. The latter vanishing is a separate hypothesis for the deformation theorem.
\end{theorem}

\begin{proof}
Put $T=H^0(A,T_A)$, $U_0=H^1(A,\mathcal O_A)$ and $H=U_0\oplus T$. The tangent bundle is trivial and the Todd class is one. Multiplicative Hochschild--Kostant--Rosenberg identifies
\[
 HH^*(A)=\bigwedge^*H,\qquad \dim H=2g,
\]
where $(\alpha,\xi)$ acts on Hodge cohomology by $\alpha\wedge-+\iota_\xi$. The characteristic maps
\[
 \operatorname{ev}_E^j:HH^j(A)\longrightarrow\operatorname{Ext}^j(E,E)
\]
form a graded algebra homomorphism, and the trace identity gives
\begin{equation}\label{eq:pure-weil-trace}
 \sigma_E^j\operatorname{ev}_E^j(x)=x\cdot\operatorname{ch}(E).
\end{equation}
We use the full Hochschild action, including the bivector and gerbe summands; the relevant Atiyah--Chern identities are \cite[Proposition~6.2.1, Corollary~6.3.2 and Theorem~6.4.2]{BFHH}.

Choose a nonzero holomorphic volume form. Cap product with it identifies $HH^j(A)$ with $HH_{g-j}(A)$, so the character corresponds to an element of $\bigwedge^gH$. Serre duality pairs $\operatorname{Ext}^2(E,E)$ perfectly with $\operatorname{Ext}^{g-2}(E,E)$. Under this pairing and the Hochschild trace pairing, $\sigma_E^2$ is adjoint to $\operatorname{ev}_E^{g-2}$. Indeed pairing the categorical trace of an endomorphism $u$ with a Hochschild operation $x$ is the Serre trace of $u\operatorname{ev}_E(x)$. Equivalently, the same identity follows term by term from the expression $\operatorname{tr}(u\operatorname{At}(E)^j/j!)$.

Translations and degree-zero tensor products act trivially on Hochschild cohomology and preserve the volume and Serre trace. Evaluation therefore lands in invariant self-Ext. Averaging a dual vector shows that Serre duality restricts to a perfect pairing on invariants. If $R^j=\operatorname{Ext}^j(E,E)^G$, the adjunction becomes
\begin{equation}\label{eq:pure-weil-dual}
 \sigma_E^2|_{R^2}\text{ injective}
 \quad\Longleftrightarrow\quad
 HH^{g-2}(A)\longrightarrow R^{g-2}\text{ surjective}.
\end{equation}

The rational $B$-field is handled without choosing a line bundle representing it. Define the invertible shear
\[
 T_B(\alpha,\xi)=(\alpha+\iota_\xi B,\xi).
\]
The type $(1,1)$ of $B$ places the first component in $U_0$. Directly,
\[
 x\cdot(e^Bz)=e^B(T_Bx\cdot z).
\]
Thus $\operatorname{ev}_{E,B}=\operatorname{ev}_E\circ\bigwedge T_B$ is still an algebra homomorphism, $\sigma_{E,B}=e^B\sigma_E$ has the same injectivity, and their composite is cap product with $w$. Surjectivity in \eqref{eq:pure-weil-dual} is unchanged by this shear.

Write $H^{1,0}(A)=X\oplus Y$ for the two $K$-eigenspaces, each of dimension $n$. Then
\[
 H^1_\sigma=X\oplus\overline Y,
 \qquad H^1_{\bar\sigma}=Y\oplus\overline X.
\]
The rational nonzero class $w$ has nonzero coefficients on both conjugate determinant lines. Under the volume identification it becomes
\[
 \eta=a\operatorname{vol}_{H_1}+b\operatorname{vol}_{H_2},
 \quad ab\ne0,\qquad
 H_1=Y^*\oplus\overline Y,\quad H_2=X^*\oplus\overline X.
\]
Here $H=H_1\oplus H_2$ and $\dim H_i=g$. Multiplication by $\eta$ on $\bigwedge^2H$ has kernel exactly $H_1\otimes H_2$ and rank $2\binom g2$. The mixed products kill both top forms. The pure second exterior powers multiply injectively by the opposite top form, with distinct bidegrees since $g\ge4$.

The composite in \eqref{eq:pure-weil-trace}, after the shear, consequently has rank $2\binom g2$. This proves the lower bound on $\dim R^2$ and proves sufficiency when equality holds.

For necessity suppose that the semiregularity map on $R^2$ is injective. Every element of its quadratic annihilator $H_1\otimes H_2$ evaluates to zero, by \eqref{eq:pure-weil-trace}. Evaluation is an algebra homomorphism, so it kills the exterior ideal generated by these mixed products. Its degree-$g-2$ image is therefore a quotient of
\[
 \bigwedge^{g-2}H_1\oplus\bigwedge^{g-2}H_2,
\]
of dimension $2\binom g2$. Equation~\eqref{eq:pure-weil-dual} says that this image is all of $R^{g-2}$. Invariant Serre duality gives $\dim R^2=\dim R^{g-2}\le2\binom g2$, and the reverse inequality was already proved. This proves necessity and the theorem.
\end{proof}

Strong semiregularity implies the weak equivariant condition in Perry's theorem: the forgetful semiregularity square commutes, so an element killed by the equivariant trace is also killed by the injective ordinary trace on invariants. No general equality of the weak and strong definitions is needed here. For a pure Weil character on an eightfold, the exact strong criterion is $\dim\operatorname{Ext}^2(E,E)^G=56$.

\subsection{An integral character on the principal Gaussian family}

Let $(X,h)$ be a principally polarized abelian fourfold and set $A=X\times X$. The endomorphism $\rho(x,y)=(-y,x)$ gives a Gaussian action of signature $(4,4)$ preserving the product polarization $\theta$. With respect to a symplectic basis of $H^1(X,\Z)$, the Gaussian Hermitian lattice has four blocks $\left(\begin{smallmatrix}0&i\\-i&0\end{smallmatrix}\right)$. Multiplication of a basis vector in each block by a unit gives the hyperbolic plane. Thus this point belongs to the principal family $U^4$.

Let $m$ be addition and put
\[
 S=p_1^*h-p_2^*h,\qquad
 \ell=m^*h-p_1^*h-p_2^*h,\qquad
 \Omega=(S+i\ell)/2.
\]
For $a=x+iy\in\Z[i]$ put $D(a)=xS-y\ell=a\Omega+\bar a\bar\Omega$, an integral divisor class. Choose its line bundle multiplicatively in $a$. If $u_1,\ldots,u_8$ is a symplectic basis on the first factor and $v_1,\ldots,v_8$ the corresponding basis on the second, then $\zeta_j=u_j+iv_j$ is an $i$-eigenvector and
\[
 \Omega=\frac12\sum_{j=1}^4\zeta_{2j-1}\zeta_{2j}.
\]
It is nondegenerate on the eight-dimensional eigenspace, with $\Omega^5=0$ and $\bar\Omega^5=0$.

The Hermitian matrix of $D(a)$, relative to the product polarization, is
\[
 \begin{pmatrix}x&-y\\-y&-x\end{pmatrix}\otimes I_4.
\]
For $a\ne0$ its eigenvalues are $|a|$ and $-|a|$, each four times. The index theorem for line bundles gives cohomology only in degree four, unchanged by a degree-zero twist, with dimension $\Nm(a)^4$. This instance can also be seen by an isogeny: rational congruence diagonalizes the indicated nonsingular rational symmetric matrix to one positive and one negative rational entry. After clearing denominators the corresponding isogeny pulls the line bundle back to an exterior product of an ample bundle on $X$ and an antiample bundle on $X$, up to a degree-zero twist. Kodaira vanishing, Serre duality and K\"unneth give concentration in degree four. Pullback by a finite \'etale isogeny contains the original cohomology as a direct summand, so the same concentration holds before pullback. Riemann--Roch gives the stated dimension.

\begin{proposition}[A pure sixteen-term character]\label{prop:sixteen-line-character}
Fix $M\ge2$ and set
\[
 (c_1,c_2,c_3,c_4)=(1,i,M(1+i),M(1-i)),
 \quad a_J=\sum_{j\in J}c_j,\quad L_J=L_{D(a_J)}.
\]
The sixteen classes $D(a_J)$ are distinct. For
\[
 \xi=\sum_{J\subset\{1,2,3,4\}}(-1)^{|J|}[L_J],
 \qquad B=-\frac12D(c_1+c_2+c_3+c_4),
\]
one has the exact identity
\[
 e^B\operatorname{ch}(\xi)
   =w=2iM^2(\Omega^4-\bar\Omega^4)\ne0.
\]
This is a rational Weil class, with no other Chern-character components.
\end{proposition}

\begin{proof}
Equality of two subset sums would give
\[
 \epsilon_1+M(\epsilon_3+\epsilon_4)=0,
 \qquad \epsilon_2+M(\epsilon_3-\epsilon_4)=0,
 \qquad \epsilon_j\in\{-1,0,1\}.
\]
Since $M\ge2$, both parenthesized integers and then all $\epsilon_j$ vanish. The divisor map is injective, proving distinctness.

The alternating exponential is
\[
 e^B\operatorname{ch}(\xi)
 =\prod_{j=1}^4 2\sinh(D(c_j)/2).
\]
Its arguments are $S,-\ell,M(S-\ell),M(S+\ell)$, so its degree-four component is
\[
 -M^2S\ell(S^2-\ell^2)=2iM^2(\Omega^4-\bar\Omega^4).
\]
Only codimensions four, six and eight can occur. The codimension-six term is $w\sum_jD(c_j)^2/24$. But
\[
 \sum_jD(c_j)^2=(1+2M^2)(S^2+\ell^2)
   =4(1+2M^2)\Omega\bar\Omega,
\]
whose product with $w$ vanishes. Finally $(x,y)\mapsto(x,-y)$ fixes $S$ and negates $\ell$. It negates the whole hyperbolic-sine product, while acting as $+1$ on top cohomology. Its codimension-eight term is therefore zero.

The two fourth powers span the conjugate Weil lines. They are independent, so $w\ne0$; its displayed expression as a product of integral divisors proves rationality. More precisely, if $w_1=\operatorname{Im}(\zeta_1\cdots\zeta_8)$, then $\Omega^4=(3/2)\zeta_1\cdots\zeta_8$ and $w=-6M^2w_1$. With the product principal orientation,
\[
 \int_A w_1^2=128,\qquad \int_Aw^2=4608M^4.
\]
Indeed the product of the two conjugate volume elements integrates to $256$, while their individual squares vanish. Also $\theta w=0$, since each summand of $\theta$ contains one factor from each Gaussian eigenspace.
\end{proof}

\subsection{Every two-level gluing has the same excessive obstruction space}

The next proposition rules out the whole extension family, including its simple members. Put
\[
 P=\bigoplus_{|J|\text{ even}}L_J,
 \qquad Q=\bigoplus_{|J|\text{ odd}}L_J.
\]

\begin{proposition}\label{prop:sixteen-line-gluing}
For every $\delta\in\operatorname{Ext}^4(P,Q)$, let $E_\delta$ be defined by the triangle
\[
 Q[3]\longrightarrow E_\delta\longrightarrow P
       \xrightarrow{\delta}Q[4].
\]
Then $E_\delta$ has no negative self-Ext and has normalized character $w$ from Proposition~\ref{prop:sixteen-line-character}. It satisfies
\[
 \operatorname{Ext}^2(E_\delta,E_\delta)
       \simeq H^2(A,\mathcal O_A)^{\oplus16},
 \qquad\dim\operatorname{Ext}^2(E_\delta,E_\delta)=448.
\]
Every finite translation and degree-zero-tensor action for which it is equivariant fixes this entire space. Consequently none of these complexes is strongly equivariantly semiregular. If all sixty-four components of $\delta$ are nonzero, the complex is simple nevertheless.
\end{proposition}

\begin{proof}
The triangle is a cone construction and has precisely two cohomology sheaves, $H^0(E_\delta)=P$ and $H^{-3}(E_\delta)=Q$. Its $K$-class is $[P]-[Q]=\xi$. Distinct line bundles in the construction have differences of index four, so
\[
 \operatorname{Ext}^k(L_J,L_{J'})=0\quad(J\ne J',\ k\ne4),
 \qquad\operatorname{Ext}^k(L_J,L_J)=H^k(A,\mathcal O_A).
\]
Use the self-hyper-Ext spectral sequence
\[
 E_2^{p,q}=\bigoplus_i
 \operatorname{Ext}^p(H^i(E_\delta),H^{i+q}(E_\delta))
 \ \Longrightarrow\ \operatorname{Ext}^{p+q}(E_\delta,E_\delta).
\]
Its only rows are $q=0,\pm3$. Each off-diagonal row is supported at $p=4$. Thus there are no terms of negative total degree, proving gluability. In total degree two the sole term is
\[
 E_2^{2,0}=\operatorname{Ext}^2(P,P)\oplus\operatorname{Ext}^2(Q,Q)
       =H^2(A,\mathcal O_A)^{\oplus16}.
\]
The only possible outgoing differential is $d_4$ to $\operatorname{Ext}^6(P,Q)=0$; the only possible incoming differential starts at $\operatorname{Ext}^{-2}(Q,P)=0$. There are no other rows that can reach or leave this term. The edge map is the asserted isomorphism. Its dimension is $16\binom82=448$.

In degree zero, the $d_4$ differential from diagonal scalar endomorphisms sends the scalar at one endpoint of a component of $\delta$ minus the scalar at its other endpoint to that multiple of the component. If all components are nonzero, the complete bipartite graph is connected and forces all sixteen scalars to agree. There is no other term in total degree zero, proving simplicity.

An equivariant translation--tensor functor preserves cohomology and first Chern classes. Distinctness of the sixteen classes therefore prevents it from permuting their line summands. On each $\operatorname{Ext}^2(L_J,L_J)=H^2(\mathcal O_A)$, translation acts trivially and the scalar ambiguity of the identification acts trivially by conjugation. The edge isomorphism is equivariant, so every one of the $448$ classes is invariant. Theorem~\ref{thm:pure-weil-ext} now excludes strong semiregularity.
\end{proof}

These actions also fail the weak equivariant condition. The projection of an allowed group fixing $E_\delta$ to the translation factor $A$ is injective: a tensor functor alone fixing one $L_J$ has trivial line bundle. Tensor conjugation by one fixed $L_J^{-1}$ makes the action a free finite translation action, possibly with a scalar coherence cocycle. Descent identifies its invariant category with a twisted derived category on the quotient abelian variety. On Hochschild homology the forgetful map is, under local Morita identifications and the invertible tensor conjugation, pullback along the finite \'etale isogeny. This is injective over $\C$. The semiregularity square therefore has the same kernel before and after forgetting equivariance, and the weak condition fails as well.

The calculation explains a necessary change in the construction. For a two-level triangle $Q[r]\to E\to P\to Q[r+1]$, the diagonal second self-Ext can only be affected by
\[
 \operatorname{Ext}^{1-r}(Q,P)\longrightarrow
 \operatorname{Ext}^2(P,P)\oplus\operatorname{Ext}^2(Q,Q)
 \longrightarrow\operatorname{Ext}^{r+3}(P,Q).
\]
For $r=3$ these two groups vanish in the example. Changing only the gluing maps cannot repair it. The shifts $3|J|$ introduce a further problem: a gap at least six gives a mixed $\operatorname{Ext}^4$ term of negative total degree. No differential can kill it, because a differential from a diagonal row would begin in negative Ext degree and all other mixed rows are supported only in Ext degree four. Thus that choice of more cohomological levels is not gluable.

Proposition~\ref{prop:sixteen-line-character} gives an explicit class on the principal hyperbolic family, but no semiregular representative of it. A positive application of Perry's theorem would require a different gluable perfect representative, with coherent allowed equivariance and invariant second self-Ext dimension $56$. The entire normalized character would then remain Hodge along the Gaussian Weil family. This is a concrete existence problem for a complex, beyond the verified character identity.

\section{Analytic deformation and pure Weil characters}
\label{sec:analytic-weil-sheaves}

This section excludes a coherent sheaf from the ordinary strongly
semiregular pure-character construction.  The distinction between a
coherent sheaf, an analytic perfect object and a global complex of
holomorphic vector bundles matters for this argument.
Bhattacharjee--Bhattacharya \cite[Theorem~A7]{BBK3Sixfolds} and
\cite[Theorem~10]{BBSecantObjects} state the exclusion for every
perfect complex.  Both proofs pass from formal deformations to
a global bounded complex of holomorphic vector bundles on nearby
tori.  That general passage fails by
Proposition~\ref{prop:analytic-strictification} below.
The coherent-sheaf assertion follows from the published analytic
deformation theorem of Buchweitz--Flenner.  We prove that assertion
and specify the additional implication needed for arbitrary complexes.
The Hodge-class and coherent-sheaf statements for very general
Weil tori are also proved in \cite[Lemmas~2--3 and Proposition~2]{BBSecantObjects},
following Voisin.  We include the vector-bundle conclusion and
the local analytic family description needed for the counterexample.

\begin{lemma}[Very general unpolarized Weil tori]
\label{lem:generic-weil-torus}
Let $A$ be a Weil abelian $2n$-fold for an imaginary quadratic
field $K$, with $n\ge2$.  Every neighborhood of $A$ in the local
space of balanced $K$-linear complex tori contains a torus $T$
with the following properties:
\begin{enumerate}[label=\textnormal{(\roman*)}]
\item its rational Hodge classes are generated by the rank,
the Weil plane $W_K$ in degree $2n$, and the top class;
\item every holomorphic vector bundle on $T$ has zero positive
rational Chern classes;
\item every coherent sheaf on $T$ has Chern character in
$H^0(T,\Q)\oplus H^{4n}(T,\Q)$.
\end{enumerate}
\end{lemma}

\begin{proof}
Put $V=H^1(A,\Q)$ and write
$V_\C=V_+\oplus V_-$ for the two $K$-eigenspaces, each of
dimension $2n$.  A balanced complex structure is given by
complementary $n$-planes $X,Y\subset V_+$, with
\[
 H^{1,0}=X\oplus\overline Y,
 \qquad H^{0,1}=Y\oplus\overline X.
\]
The deformation space is the open subset of
$\operatorname{Gr}(n,V_+)\times\operatorname{Gr}(n,V_-)$
consisting of $(X,Z)$ with $X\oplus\overline Z=V_+$.
It has complex dimension $2n^2$ and gives a smooth local analytic
family of marked tori.  As a real homogeneous space it is
$\operatorname{GL}_{2n}(\C)/
(\operatorname{GL}_n(\C)\times\operatorname{GL}_n(\C))$;
the preceding Grassmannian description specifies its deformation
complex structure.

The Hodge circle on $V_+$ acts by $e^{it}$ on $X$ and $e^{-it}$
on $Y$.  Its generator is nonzero and traceless.  As the
decomposition varies, the real Lie algebra generated by these
conjugate generators is a nonzero ideal in
$\mathfrak{sl}_{2n}(\C)$, regarded as a simple real Lie algebra.
It is therefore that whole Lie algebra.  Its complexification
acts on $V_\C$ as
$\operatorname{SL}(V_+)\times\operatorname{SL}(V_-)$.
In $\bigwedge^a V_+\otimes\bigwedge^b V_-$ an invariant occurs
only when each of $a,b$ is $0$ or $2n$.  The classes Hodge
throughout the parameter space are consequently exactly those
in (i).  For each other rational class its Hodge locus is a
proper analytic subset.  The countable union of these loci
cannot contain an analytic open set, by Baire's theorem.
This proves (i) on a very general torus in every neighborhood.

Choose a positive constant Hermitian metric making $X$ and
$\overline Y$ orthogonal.  Its Kähler class $\kappa$ belongs
to $V_+\wedge V_-$.  If $\alpha_\pm$ generates
$\bigwedge^{2n}V_\pm$, then
$\alpha_\pm\wedge\kappa=0$, since every term adds a vector
to an already exhausted eigenspace.  Thus $W_K\wedge\kappa=0$.
We have $\operatorname{NS}(T)_\Q=0$.  Moreover
$\operatorname{Hdg}^4(T)$ is zero for $n\ge3$, and for $n=2$
it is $W_K$ and is perpendicular to $\kappa^{2n-2}$.

Voisin's Proposition~A.3 and Corollary~A.4
\cite[pp.~1072--1074]{VoisinKahler} apply to exactly these two
conditions.  They say that every holomorphic vector bundle has
zero positive rational Chern classes, and that a torsion-free
coherent sheaf embeds in such a vector bundle with torsion
cokernel.  The input is the stable reflexive equality case of
Bando--Siu \cite[Theorem~3 and Corollary~3, p.~40]{BandoSiu}.
This proves (ii).

There is no proper positive-dimensional analytic subset of $T$.
Its fundamental class would be a nonzero rational Hodge class
of codimension $r$ with positive Kähler volume.  Statement (i)
leaves only $r=n$, but then its class is in $W_K$ and has
zero pairing with $\kappa^n$, a contradiction.  Every torsion
sheaf is therefore supported on finitely many points.  Remove
the torsion of an arbitrary coherent sheaf and apply Voisin's
embedding to the torsion-free quotient.  The embedding bundle
has character only in degree zero, and both torsion terms have
character only in top degree.  Additivity proves (iii).
\end{proof}

\begin{proposition}[Coherent pure-character exclusion]
\label{prop:coherent-pure-weil}
Let $A$ be a Weil abelian $2n$-fold, $n\ge2$, and let
$0\ne w\in W_K(A)$ be rational.  If a coherent sheaf
$\mathcal F$ satisfies
\[
 \operatorname{ch}(\mathcal F)-cw
   \in H^0(A,\Q)\oplus H^{4n}(A,\Q),\qquad c\ne0,
\]
then its full semiregularity map is not injective.  The same
assertion holds for a shift of such a sheaf.
\end{proposition}

\begin{proof}
Suppose the map were injective.  All components of the displayed
Chern character have horizontal Hodge extensions on the smooth
unpolarized Weil family in Lemma~\ref{lem:generic-weil-torus}.
Apply Buchweitz--Flenner \cite[Theorem~5.1, pp.~174--175]{BF},
with all Chern-character indices.  The central fibre is a compact
complex algebraic manifold, the base is smooth, and injectivity
of the full map is their required semiregularity hypothesis.

The conclusion needed here is the coherent sheaf constructed
in their proof on p.~179.  Proposition~5.9 and Lemma~5.10
extend $\mathcal F$ successively over the infinitesimal
neighborhoods of the base point.  A relative versal coherent-sheaf
deformation then has a compatible formal section.  Artin
approximation supplies a convergent section, whose pullback is
a coherent sheaf flat over a neighborhood in the analytic
family.  Its Chern character is the horizontal extension of
that of $\mathcal F$.  On a very general nearby fibre, however,
Lemma~\ref{lem:generic-weil-torus}(iii) makes its middle character
zero, whereas it must equal the nonzero transport of $cw$.
This is a contradiction.  A shift preserves self-Ext and changes
the character and trace map only by a sign.
\end{proof}

\begin{proposition}[Failure of global strictification]
\label{prop:analytic-strictification}
Injective semiregularity and a horizontal Hodge Chern character
do not imply that nearby nonprojective tori carry a global
bounded complex of holomorphic vector bundles with that character,
even when the initial perfect object has no negative self-Ext.
\end{proposition}

\begin{proof}
Take a point sheaf $E=\mathcal O_a$ on a projective Weil
$g$-fold $A$, where $g=2n\ge4$.  The local Koszul resolution
gives $\operatorname{Ext}^k(E,E)=\bigwedge^kT_aA$ and no
negative self-Ext.  In local coordinates its differential is
$\sum_j z_j\iota_j$; differentiating gives the Atiyah
representative $\sum_j dz_j\iota_j$.  Its $r$th power divided
by $r!$ contracts an $r$-vector by the usual exterior evaluation,
up to a nonzero sign.

Consider the component
\[
 \sigma_E^{g-2}:\operatorname{Ext}^2(E,E)
       \longrightarrow H^g(A,\Omega_A^{g-2}).
\]
Serre duality identifies its dual with contraction by
$\operatorname{At}(E)^{g-2}/(g-2)!$ from
$H^0(A,\bigwedge^{g-2}T_A\otimes\omega_A)$ to
$\operatorname{Ext}^{g-2}(E,E\otimes\omega_A)$.
The Koszul calculation identifies this with evaluation at $a$.
The tangent and canonical bundles are trivial, so evaluation
is an isomorphism.  Thus this component, and the full
semiregularity map, are injective.

The character of $E$ is the top class $[a]$, which stays Hodge
on every deformation of the torus.  Initially $E$ has a global
bounded resolution by algebraic vector bundles, since $A$ is
smooth and projective.  Every neighborhood in the unpolarized
Weil family contains a torus from
Lemma~\ref{lem:generic-weil-torus}(ii).  Every global bounded
complex of holomorphic vector bundles there has character
equal to the alternating sum of their ranks, hence has no
nonzero top component.  It cannot represent the transport of
$[a]$.  This proves the assertion.
\end{proof}

A point section still gives a coherent analytic point sheaf on
every nearby fibre.  It is an analytic perfect object, because
finite locally free resolutions exist locally on a smooth
complex manifold.  Proposition~\ref{prop:analytic-strictification}
concerns a single global resolution, not this local property.

For an arbitrary perfect complex, an analytic deformation theorem
producing analytic perfect objects on the varying family would
suffice for the argument of
Proposition~\ref{prop:coherent-pure-weil}.  Indeed its coherent
cohomology sheaves have zero intermediate Chern characters on
the very general fibre, and their alternating sum is the
character of the complex.  The global-resolution assertion
does not establish that implication.  We do not use the
arbitrary-perfect exclusion as a theorem, or apply the ordinary
coherent-sheaf exclusion to an invariant semiregularity map.

\section{Simple bundles with a prescribed middle Chern class}
\label{sec:simple-hodge-bundles}

This section constructs simple bundles on a CM product whose entire
Chern character stays of Hodge type on a polarized Weil family.
The construction solves the existence problem for the bundle at that
product point.  Its deformation is a further condition, and the
particular gluing used below fails that condition.

For $n\ge2$ put $g=2n$ and define
\begin{equation}
 F_n(t)=\sum_{k=0}^{2n}
  \frac{e_k(1,-1,\ldots,n,-n)}{(2n)_k}t^k,
 \qquad (2n)_k=\frac{(2n)!}{(2n-k)!}.
 \label{eq:positive-character-polynomial}
\end{equation}
Here $e_k$ is the elementary symmetric polynomial and $(2n)_0=1$.

\begin{theorem}\label{thm:simple-hodge-bundle}
Let $K$ be any imaginary quadratic field, let $E$ be a complex
elliptic curve with $\operatorname{End}^0(E)=K$, and give
$A=E^{2n}$ its product principal polarization $\theta$.
For every nonzero rational primitive class
$w\in H^{n,n}(A)\cap H^{2n}(A,\mathbf Q)$ there are positive integers
$R,N$, a nonzero rational number $c$, and a simple vector bundle
$\mathcal E$ such that
\[
 \operatorname{ch}(\mathcal E)=R F_n(N\theta)+cw.
\]
It is an ordinary extension of two finite direct sums of line
bundles of $\theta$-slope zero.  The ratio $c/(RN^n)$ can be chosen
arbitrarily small and positive.  No bound for $R,N$ uniform in
$E,w$ is asserted.

If $w$ is a Weil class for a compatible action of signature $(n,n)$,
this entire character stays of Hodge type on the corresponding
polarized Weil domain.  For $n\ge3$ its degree-two Hochschild action
has rank $7n^2-2n$.  Equality
$\dim\operatorname{Ext}^2(\mathcal E,\mathcal E)=7n^2-2n$ would
therefore imply strong semiregularity; that equality is not part
of the construction.
\end{theorem}

The proof uses exact finite averages to make a signed cohomology
identity into a sum with positive multiplicities.  Rank-two
differences between the divisor classes then provide extension
groups on abelian surfaces.  Enough independent extension classes
force scalar endomorphisms of the resulting bundle.
The general Hochschild-rank formula for a Weil class plus a
polarization polynomial is
\cite[Theorem~A6]{BBK3Sixfolds}; we compute the relevant minor for
$F_n$ and give the resulting kernel calculation explicitly.

\subsection{A rational unitary orbit and its primitive summand}

Use the central CM action on the product to identify
$\operatorname{NS}(A)_{\mathbf Q}$ with the Hermitian matrices
$\operatorname{Herm}_g(K)$.  This follows directly from the
Rosati-fixed endomorphisms of the product polarization.
This central action is used to construct divisors; it need not be
the balanced Weil action in the theorem.  The class $\theta$
corresponds to the identity matrix.  Put
\[
 D=\operatorname{diag}(1,-1,\ldots,n,-n).
\]
Every $D_u=uDu^*$ for $u\in U(g,K)$ is a rational divisor class
of trace, and hence slope, zero.  This rational unitary group is
dense in $U(g)$: skew-Hermitian matrices over $K$ are dense in all
skew-Hermitian complex matrices, and the Cayley transform covers
a dense chart of the unitary group.

We need an assertion about the whole character $e^{D_u}$, not just
its middle component.  As a complex unitary representation,
\[
 H^{p,p}(A)=\Lambda^p\mathbf C^g\otimes
                         (\Lambda^p\mathbf C^g)^*.
\]
Its irreducible summands have highest weights
$(1^k,0^{g-2k},-1^k)$ for $0\le k\le\min(p,g-p)$, each once.
Indeed write the dual factor as
$\Lambda^{g-p}\mathbf C^g\otimes\det^{-1}$.
The exterior Pieri rule gives the two-column diagrams
$(2^k,1^{g-2k})$, and subtracting the determinant weight gives
the displayed list.  The summand with $k=n$ occurs only when
$p=n$.  It is exactly
\[
 P^{n,n}=\ker(\theta\wedge:H^{n,n}\longrightarrow H^{n+1,n+1}).
\]
It is irreducible and occurs once in the entire graded Hodge space.

Its projection from $e^D$ is nonzero.  Write $x_1,\ldots,x_{2n}$
for the coordinate elliptic divisor classes, so that $x_j^2=0$ and
$\int_Ax_1\cdots x_{2n}=1$.  The rational class
\[
 q=\prod_{j=1}^n(x_{2j-1}-x_{2j})
\]
is primitive, and
\[
 \int_A\frac{D^n}{n!}q=(-2)^n n!\ne0.
\]
Complete reducibility and the unique occurrence of $P^{n,n}$ imply
that the real span of the orbit $\{e^{D_u}:u\in U(g)\}$ contains
every primitive middle class with all its other graded components
zero.  To see why multiplicity matters, the span is a submodule;
its nonzero projection to an irreducible constituent occurring once
contains that constituent itself.

The rational span of the rational-unitary orbit has closed real
span, hence contains the continuous orbit by density.  All rational
Hodge classes on a CM elliptic product have the indicated rational
form: the diagonal CM Hodge torus fixes exactly the equal
holomorphic and antiholomorphic degrees.  Rational row reduction
therefore gives a finite exact equality
\begin{equation}
 w=\sum_a b_a e^{D_a},\qquad b_a\in\mathbf Q,
 \quad D_a=u_aDu_a^*,\quad u_a\in U(g,K).
 \label{eq:primitive-orbit-relation}
\end{equation}
In particular no approximation to $w$ is involved.

\subsection{Finite averages make all multiplicities positive}

In each unitary eigenbasis permute the distinct eigenvalues of $D$.
Let $\mathcal O_a$ be the resulting set of $g!$ divisor classes.
The rank-one divisor forms $\omega_1,\ldots,\omega_g$ of that basis
satisfy $\omega_j^2=0$ and $\sum_j\omega_j=\theta$.
Counting the subsets in the exponential expansion gives
\begin{equation}
 \frac1{g!}\sum_{D'\in\mathcal O_a}e^{D'}
 =\sum_{k=0}^g\frac{e_k(1,-1,\ldots,n,-n)}{\binom gk}
                              \frac{\theta^k}{k!}
 =F_n(\theta).
 \label{eq:finite-positive-average}
\end{equation}
If a coefficient in \eqref{eq:primitive-orbit-relation} is negative,
add its absolute value times the full corresponding orbit sum to
both sides.  This removes that negative coefficient and introduces
no new negative ones.  There are finitely many coefficients.
Then add a strictly positive constant on every orbit involved.
After identifying repeated classes the result is
\begin{equation}
 C F_n(\theta)+w=\sum_v q_v e^{D_v},
 \qquad C>0,\quad q_v\in\mathbf Q_{>0}.
 \label{eq:positive-orbit-relation}
\end{equation}
Additional uniform orbit sums make $1/C$ arbitrarily small.
The equality holds in every character degree, including degree zero.

We may connect these orbit sets while preserving
\eqref{eq:positive-orbit-relation}.  Inside an orbit, every coordinate
transposition gives a difference
$(\lambda_j-\lambda_k)(\omega_k-\omega_j)$ of Hermitian rank two
and signature $(1,1)$.  Its graph is the bipartite transposition
graph of the symmetric group.

Any rational unitary transformation is a product of rank-one unitary
transformations.  For completeness, if $v\in K^m$ is a unit vector
different from $e_1$, put $u=v-e_1$ and
\[
 h=I-\frac{uu^*}{1-v_1}.
\]
The identities $u^*v=1-v_1$ and
$\|u\|^2=2-v_1-\overline v_1$ give $hv=e_1$ and $hh^*=I$.
The denominator is nonzero, since $v_1=1$ for a unit vector forces
$v=e_1$.  Apply this operation to the first column, then induct on
its orthogonal complement.  For a rank-one unitary transformation,
\[
 hD'h^*-D'=(h-I)D'h^*+D'(h^*-I)
\]
has rank at most two and trace zero.  If nonzero, it has rank two
and signature $(1,1)$.

Include the finitely many intermediate eigenbasis orbits along
these products.  Give each a positive uniform coefficient using
\eqref{eq:finite-positive-average}.  Equal orbits can be identified:
one common matrix with distinct eigenvalues determines the same
unordered eigenlines and therefore the same permutation orbit.
Choose a tree joining the distinct orbit sets, and one of the
rank-two edges for each tree edge.  Keep all internal transposition
edges.  Flip the parity colorings successively along the tree.
The result is a finite connected bipartite graph of distinct rational
divisor classes, every edge having a difference of signature $(1,1)$
on a rational surface quotient.

\subsection{Line bundles, surface cohomology and simplicity}

Clear the denominators of the $q_v$ to obtain positive multiplicities
$m_v$.  Choose $N$ so that every alternating Riemann form of $ND_v$
has even integral values on the product lattice $\Lambda$.
For each choose its Appell--Humbert line bundle with trivial
semicharacter.  The cocycle formula and tensor additivity are
\cite[\S2, pp.~20--21]{MumfordAV}.
Explicitly the factor for a Hermitian form $H$ is
\[
 j_\lambda(z)=
 \exp\bigl(\pi H(z,\lambda)+\tfrac\pi2H(\lambda,\lambda)\bigr).
\]
Even integrality makes the omitted semicharacter equal to one
admissible.  This formula is additive in $H$, so the quotient of
two chosen line bundles has exactly the difference form, with
no extra class in $\operatorname{Pic}^0$.

For an edge let $W$ be the kernel of that difference.  It is rational
over $K$, so $B=W/(W\cap\Lambda)$ is an abelian subvariety and
$p:A\to Y=A/B$ is a quotient onto an abelian surface.
The form vanishes when either argument is in $W$; the displayed
cocycle therefore descends to $Y$.  Thus
\[
 L_vL_u^{-1}=p^*M
\]
for a nondegenerate line bundle of index one on $Y$.
Mumford's vanishing and index theorems
\cite[\S16, pp.~150,155,162]{MumfordAV} give
$H^i(Y,M)=0$ for $i\ne1$ and $h^1(Y,M)=-\chi(M)>0$;
the inverse has the same index.
The low-degree Leray sequence, $p_*\mathcal O_A=\mathcal O_Y$,
and $H^0(Y,M)=0$ give
\[
 H^1(A,p^*M)=H^1(Y,M).
\]
Here the higher direct images of $\mathcal O_A$ are trivial with
fibres $H^q(B,\mathcal O_B)$: invariant forms on $A$ restricting to
a basis on $B$ give a global frame.
Surface Serre duality also shows that the two reverse extension
groups pair into the nonzero line
\begin{equation}
 q_{uv}=p^*H^2(Y,\mathcal O_Y)\subset H^2(A,\mathcal O_A).
 \label{eq:edge-surface-line}
\end{equation}
Replacing $N$ by a positive multiple tensors $M$ by that power.
Surface Riemann--Roch makes its $h^1$ grow quadratically.
We can therefore arrange
\[
 \dim\operatorname{Ext}^1(L_u,L_v)\ge m_um_v
\]
on every edge in either direction.

Distinct $L_u,L_v$ have no homomorphisms in either direction.
Their divisor difference is nonzero of $\theta$-degree zero,
whereas every nonzero effective divisor has positive $\theta$-degree.
If a section had no divisor it would trivialize the quotient,
also contradicting the distinct Chern classes.

Let $P$ and $Q$ be the line-bundle sums with multiplicities $m_v$
on the two parts of the bipartition.  On every chosen edge from
$L_u\subset P$ to $L_v\subset Q$, choose $m_um_v$ independent
extension classes $\eta_{ab}$ and set
\[
 e_{vu}=\sum_{a,b}\eta_{ab}\otimes E_{ab},
\]
where $E_{ab}$ are the matrix units.  Set all other components zero.
This is a single class in $\operatorname{Ext}^1(P,Q)$, defining
an ordinary extension
\begin{equation}
 0\longrightarrow Q\longrightarrow\mathcal E
                   \longrightarrow P\longrightarrow0.
 \label{eq:positive-bundle-extension}
\end{equation}
An extension of vector bundles is locally free.  Every endomorphism
preserves $Q$ because $\operatorname{Hom}(Q,P)=0$, and is determined
by its pair on $Q,P$ because $\operatorname{Hom}(P,Q)=0$.
Those endomorphisms are block diagonal.  The compatibility equation
at an edge gives $b_vE_{ab}=E_{ab}t_u$ for every matrix unit, forcing
both endpoint matrices to be the same scalar.  Connectivity forces
one scalar everywhere.  Thus $\operatorname{End}(\mathcal E)=\mathbf C$.

Additivity of the Chern character in
\eqref{eq:positive-bundle-extension}, followed by degreewise scaling
in \eqref{eq:positive-orbit-relation}, gives the asserted character.
Its rank is $R=\sum_vm_v$ and its middle coefficient is a nonzero
rational multiple of $N^n$.  This proves the existence assertion.
It also proves slope semistability, since the two sums have slope
zero.  Stability is not asserted.

\subsection{The Hochschild kernel and the number 104}

Now suppose $w$ is a compatible rational Weil class and $n\ge3$.
Write $F_n(t)=1+p_2t^2+p_4t^4+\cdots$.  Direct expansion of
$\prod_{j=1}^n(1-j^2t^2)$ gives
\[
 p_2=-\frac{(n+1)(2n+1)}{12(2n-1)},\qquad
 p_4=\frac{(n+1)(2n+1)(5n+6)}{1440(2n-3)},
\]
and
\begin{equation}
 6p_4-p_2^2=
 \frac{(n+1)(2n+1)(40n^3+12n^2-22n+33)}
 {720(2n-3)(2n-1)^2}>0.
 \label{eq:positive-rank-minor}
\end{equation}

Under HKR write a Hochschild class of degree two as
$(\xi,\eta,\zeta)$ in
$H^0(\Lambda^2T_A)\oplus H^1(T_A)\oplus H^2(\mathcal O_A)$.
Put $b=\xi\lrcorner\theta^2$ and $u=\eta\lrcorner\theta$.
The first three equations for annihilating $F_n(N\theta)$ are
\[
 p_2N^2b+\zeta=0,\qquad 2p_2N^2\theta u=0,
 \qquad 6p_4N^4\theta^2b+p_2N^2\theta^2\zeta=0.
\]
The identities used are the derivation identity for $\eta$ and
$\xi\lrcorner\theta^k=\binom k2\theta^{k-2}b$.
Lefschetz injectivity and \eqref{eq:positive-rank-minor} give
$b=u=\zeta=0$.  Contraction with $\theta^2$ identifies bivectors
with $H^{0,2}$, so $\xi=0$.  Conversely $u=0$ kills every power
of $\theta$.  Thus the kernel for the background polynomial is
exactly the polarized tangent space, of dimension $2n^2+n$.

The norm-one torus of the balanced $K$ action fixes $\theta$.
Its weights on $HH^2$ are $-2,0,2$, and its weights on the two Weil
lines are $2n,-2n$.  The polynomial output weights and the Weil
output weights are disjoint for $n\ge3$.  There is consequently
no cancellation between them, and the kernel for
$R F_n(N\theta)+cw$ is the intersection of their two kernels.

Split $H^{1,0}=U\oplus V$ according to the balanced action, with
$\dim U=\dim V=n$.  The polarized tangent space is
\[
 \operatorname{Sym}^2U^*\oplus(U^*\otimes V^*)
                         \oplus\operatorname{Sym}^2V^*.
\]
The middle summand kills both Weil determinant forms.  The other
two summands act injectively on their respective determinant forms:
contracting a holomorphic volume form identifies its next exterior
power with the dual standard representation and records the
symmetric matrix.  The outputs have distinct $K$ weights.
Both determinant coefficients of a nonzero rational $w$ are nonzero.
The common kernel is therefore exactly $U^*\otimes V^*$,
of dimension $n^2$.  Since $\dim HH^2=8n^2-2n$, the rank is
$7n^2-2n$ as asserted.

For an eightfold this becomes
\[
 F_4(t)=1-\frac{15}{28}t^2+\frac{91}{560}t^4
               -\frac{41}{1008}t^6+\frac1{70}t^8,
 \qquad 7\cdot4^2-2\cdot4=104.
\]
The trace/evaluation factorization of Buchweitz--Flenner
\cite[Proposition~6.2.1 and Theorem~6.4.2]{BFHH} identifies the
composite through $\operatorname{Ext}^2$ with this Hochschild action.
If that Ext group has dimension equal to the computed rank,
the semiregularity map is injective.  This proves the remaining
assertions of Theorem~\ref{thm:simple-hodge-bundle}.

\subsection{The fourfold case and its Euler obstruction}

The construction also works for $n=2$, but the weight separation
just used does not.  Here
$F_2(t)=1-5t^2/12+t^4/6$.
For a sufficiently small nonzero rational $\varepsilon$, the action
on $F_2(\theta)+\varepsilon w$ has rank $24$.
To check this, the $H^2(\mathcal O)\oplus H^0(\Lambda^2T)$ block
at $\varepsilon=0$, after Lefschetz identifications, has matrix
\[
 \begin{pmatrix}1&2p_2\\p_2&12p_4\end{pmatrix}\otimes I_6,
 \qquad 12p_4-2p_2^2=119/72.
\]
It remains invertible for small $\varepsilon$.
The $H^1(T)$ block has rank six at zero, with image
$\theta H^{0,2}$ and kernel the ten-dimensional polarized tangent.
On that kernel the Weil derivative has rank six, with kernel the
four-dimensional Weil tangent.  Its image is primitive, since
both $\kappa\theta=0$ and $\theta w=0$.
It is therefore disjoint from $\theta H^{0,2}$.
A twelve-by-twelve minor of the perturbed block has nonzero leading
term of order $\varepsilon^6$, so this block has rank twelve for
all sufficiently small nonzero $\varepsilon$.  The four-dimensional
Weil tangent bounds its rank above by twelve.  The total is $24$.
The positive completion permits such $\varepsilon$; subsequent
integral scaling conjugates these maps by invertible degree scalings.

None of these simple fourfold bundles is strongly semiregular.
On a Calabi--Yau fourfold the full degree-two trace map is dual to
degree-two evaluation.  If it were injective, evaluation would be
surjective, and their composite would have rank
$\operatorname{ext}^2=24$.
The degree-one action on a nonzero Weil class is injective; its
weights $\pm3$ are disjoint from the polynomial output weights
$\pm1$.  Thus $\operatorname{ext}^1\ge8$.
Simplicity and Serre duality would give
$\chi(\mathcal E,\mathcal E)=2-2\operatorname{ext}^1+
\operatorname{ext}^2\le10$.
But Riemann--Roch, $\theta w=0$, and Hodge--Riemann give
\[
 \chi(\mathcal E,\mathcal E)
 =\frac{73}{6}R^2N^4+c^2\int_Aw^2>10.
\]
This contradiction concerns ordinary Ext groups.  It makes no
assertion about invariants under an additional finite group.
In dimension eight the middle Ext group prevents the same Euler
argument from giving a contradiction.

\section{Why the tree of line-bundle orbits does not deform}
\label{sec:tree-extension-obstruction}

This section shows that the tree used to obtain simplicity in
Theorem~\ref{thm:simple-hodge-bundle} obstructs deformation in at
least one polarized Weil direction.  The obstruction concerns the
underlying bundle, even when its line subbundles are allowed to
disappear under deformation.

Retain the full permutation orbits of the simple Hermitian spectrum
$(1,-1,\ldots,n,-n)$ used there.  Within each orbit retain all
transposition edges.  Between distinct orbits retain exactly one
rank-two edge for each edge of a tree.  One may also tensor the
line bundles just constructed by arbitrary elements of
$\operatorname{Pic}^0(A)$.

\begin{theorem}\label{thm:tree-extension-obstruction}
Let $n\ge2$.  Form an ordinary extension
\eqref{eq:positive-bundle-extension} using positive multiplicities
at all vertices and arbitrary extension matrices supported on this
graph.  If it deforms to first order in every polarized Weil tangent
direction, its multiplicities are constant on each full permutation
orbit.  In particular its Chern character is a polynomial in $\theta$.
Consequently an extension with character
$R F_n(N\theta)+cw$, $c\ne0$, cannot be strongly semiregular.
\end{theorem}

The graph hypothesis is essential to the conclusion proved here.
Several bridges between the same two orbit sets, cycles between
orbit sets, or additional nonzero extension entries are not covered.
The proof gives a necessary trace equation for deformation and then
solves that equation on a leaf of the tree.

\subsection{The trace equation allows unfiltered deformations}

Split the underlying smooth bundle as the sum of its line factors.
Its Dolbeault operator is $\bar\partial_0+e$, where $e$ has entries
only from $P$ to $Q$, so $e^2=0$ and $\bar\partial_0e=0$.
Choose a connection whose $(1,0)$ part is block diagonal on these
factors.  Its $(1,1)$ curvature has the line curvatures on the
diagonal and other entries only from $P$ to $Q$.

For a complex-structure deformation $\kappa$, a first-order change
$x$ in the Dolbeault operator must satisfy
\begin{equation}
 \bar\partial_0x+[e,x]+\kappa\lrcorner F^{1,1}=0,
 \label{eq:tree-atiyah}
\end{equation}
up to the common sign convention.  This is obtained by expanding
the square of the deformed operator; its cohomology class is the
Atiyah obstruction.  In the reverse block from $Q$ to $P$ both the
curvature term and $[e,x]$ vanish.  Thus these reverse entries of
$x$ are closed and give arbitrary classes in the corresponding
$\operatorname{Ext}^1$ groups.

Take the trace of the diagonal equation at vertex $v$, and pass
to $H^2(\mathcal O_A)$.  Its prescribed term is
$m_v\kappa\lrcorner c_1(L_v)$.
An edge contribution is the cup product of its extension entry
with a reverse class.  By \eqref{eq:edge-surface-line} it lies in
the one-dimensional line $q_e$ from the associated surface quotient.
This remains true with arbitrary flat twists.  Indeed an isogeny
from a product of an abelian surface and the kernel abelian variety
pulls the difference bundle back to $M\boxtimes\alpha$, where
$M$ has index one and $\alpha\in\operatorname{Pic}^0$.
If $\alpha$ is nontrivial its cohomology vanishes: a unitary
character description and Fourier expansion of the Dolbeault
complex give nonzero Koszul differentials at every frequency.
If $\alpha$ is trivial, K\"unneth and $H^0(M)=H^0(M^{-1})=0$
put both degree-one groups on the surface factor.  Their cup
product therefore lies in its $H^2(\mathcal O)$ line.
Isogeny pullback is injective on line-bundle cohomology by trace,
and is an isomorphism on $H^2(\mathcal O)$, proving the claim.
The two endpoint traces are opposite: interchanging the two
degree-one forms changes the sign.
Therefore a necessary condition for deformation is the existence
of scalar edge flows with
\begin{equation}
 \bigl(m_v\kappa\lrcorner c_1(L_v)\bigr)_v
  =\sum_e f_e\,(q_e\text{ at one endpoint},
                         -q_e\text{ at the other}).
 \label{eq:tree-incidence}
\end{equation}
Allowing arbitrary flows can only enlarge the possible right side.
The diagonal entries of $x$ contribute exact terms and disappear
in cohomology.  Reverse blocks were included throughout, so this
condition imposes no requirement that the deformation preserve
the original filtration.  There is also no omitted higher-product
term: the Dolbeault equation is the full first-order equation.

\subsection{Two facts about the Weil tangent space}

Using the product polarization, the polarized Weil tangent space
is congruent to
\begin{equation}
 \mathcal T=
 \left\{\begin{pmatrix}0&C\\C^{\mathsf t}&0\end{pmatrix}:
                    C\in\operatorname{Mat}_n(\mathbf C)\right\}
 \subset\operatorname{Sym}^2(\mathbf C^{2n})^*.
 \label{eq:weil-offblock-tangent}
\end{equation}
It has dimension $n^2$.  The following arguments work in every
eigenbasis, without a genericity assumption.

First, join coordinate indices $j,k$ when
$\kappa\mapsto\kappa_{jk}$ is not identically zero on $\mathcal T$.
This coordinate graph is connected.  Otherwise the coordinate
partition gives a nontrivial decomposition $W=A_1\oplus A_2$
orthogonal for every form in \eqref{eq:weil-offblock-tangent}.
Write $W=U\oplus V$ in that equation.  For $x=u+v$ and
$y=u'+v'$, simultaneous orthogonality is equivalent to
\[
 u\otimes v'+u'\otimes v=0.
\]
If $u,v$ are both nonzero, its solutions in $y$ form the line
$\mathbf C(u-v)$.  If $A_1$ contains such an $x$, then $A_2$ is
a mixed line, and applying the same argument in reverse makes
$A_1$ a line.  This contradicts $2n\ge4$.
If $A_1$ contains no mixed vector, its linearity puts it in either
$U$ or $V$.  A nonzero vector in $U$, for example, forces
$A_2\subset U$, again contradicting the direct sum.  The graph
is connected as claimed.

Second, the off-diagonal projection of $\mathcal T$ has rank at
least $n^2-n\ge2$.  Every matrix in
\eqref{eq:weil-offblock-tangent} has even rank, and congruence
preserves rank.  Let $S$ be the intersection with diagonal
matrices in any chosen basis.  Its coordinate linear functionals,
after omitting zeros, split into proportionality classes.
A generic point of $S$ has even rank.  If $\dim S\ge2$, a generic
point on the hyperplane of one proportionality class and on none
of the others has rank lower by the size of that class.  That size
must therefore be even.  There are at most $n$ such classes,
and their representatives span $S^*$.  Thus $\dim S\le n$.
For $\dim S\le1$ the conclusion is immediate.

If $D_v$ is diagonal with distinct eigenvalues, then
\begin{equation}
 \kappa\lrcorner D_v
  =\sum_{j<k}(\lambda_{v,j}-\lambda_{v,k})
                      \kappa_{jk}q_{jk},
 \label{eq:tree-diagonal-tangent}
\end{equation}
where the $q_{jk}$ form the coordinate basis of $H^2(\mathcal O_A)$.
All displayed eigenvalue differences are nonzero.  Hence this map
on $\mathcal T$ has rank at least $n^2-n\ge2$.

\subsection{A leaf forces uniform multiplicities}

Choose a leaf orbit.  Its unique bridge is incident at one vertex
$b$ of that orbit.  For a fixed coordinate pair $j,k$, the internal
transposition edges form a matching, and their surface line is
$\mathbf Cq_{jk}$.  The coefficient in
\eqref{eq:tree-diagonal-tangent} has opposite signs at matched
vertices $v,(jk)v$.
If neither is $b$, their two trace equations therefore imply
\[
 (m_v-m_{(jk)v})(\lambda_{v,j}-\lambda_{v,k})\kappa_{jk}=0
                  \quad(\kappa\in\mathcal T).
\]
For every allowed coordinate pair, their multiplicities agree.

The allowed transpositions generate the symmetric group, since
the coordinate graph is connected.  Their Cayley graph on the
permutation orbit is connected and vertex transitive.  Deleting
one vertex leaves it connected: a leaf of a spanning tree is not
a cut vertex of a connected graph, and vertex transitivity then
shows that no vertex is a cut vertex.
Thus all vertices other than $b$ have a common multiplicity $m$.

Sum the trace equations over the leaf orbit.  Internal flows cancel.
The unweighted sum of its divisor matrices is zero, since every
eigenvalue occurs equally often in each position and their sum
is zero.  Thus the remaining bridge equation requires
\[
 (m_b-m)(\kappa\lrcorner D_b)\in\mathbf Cq_{\mathrm{bridge}}
                       \quad(\kappa\in\mathcal T).
\]
The image on the left before its scalar multiplication has
dimension at least two, whereas the right side has dimension one.
It follows that $m_b=m$, and the bridge flow is zero for every
$\kappa$.  The common scaling $N$ changes none of these conclusions.

Delete this leaf and its zero bridge and repeat.  Every orbit has
constant multiplicity.  On the final orbit the same conclusion
follows directly from the connected allowed-transposition graph.
The average \eqref{eq:finite-positive-average} now gives
\[
 \operatorname{ch}(\mathcal E)
  =\sum_a m_a(2n)!F_n(N\theta)\in\mathbf Q[\theta].
\]
A nonzero primitive middle Weil class is not a multiple of
$\theta^n$, so this contradicts $c\ne0$.
At least one Weil tangent therefore has nonzero Atiyah obstruction.
Its character obstruction is zero because the displayed character
stays Hodge.  The trace/evaluation identity sends that nonzero
obstruction to zero, proving failure of strong semiregularity.
This completes the proof of
Theorem~\ref{thm:tree-extension-obstruction}.

If the bundle has an additional coherent finite-group structure
whose action on degree-two Hochschild cohomology is trivial, the same
obstruction lies in the invariant Ext group, by equivariance of
evaluation.  It also excludes injectivity of the restricted
semiregularity map under that additional hypothesis.

The remaining construction question is more specific than simplicity:
one needs a gluing with enough independent bridges to satisfy all
the diagonal trace equations, and then the full second-Ext condition.
Solving the trace equations alone would be a necessary step and would
not establish semiregularity.

% Standalone section fragment; requires amsmath, amssymb, amsthm.
\section{A generic obstruction for coordinate-sign networks}
\label{sec:sign-cube-obstruction}

Adding many full matchings does not by itself remove the diagonal
obstruction.  The following example has numerous cycles, but its relevant
Fourier component still has only two independent flow channels.

Let \(K\) be an imaginary quadratic field and let \(E\) be a complex
elliptic curve with \(\operatorname{End}^0(E)=K\).
Put \(m=2n\), let \(A=E^{m}\) have its product principal polarization
\(\theta\), and use the diagonal balanced action with eigenspaces
\(U=\langle e_1,\ldots,e_n\rangle\) and
\(V=\langle e_{n+1},\ldots,e_m\rangle\).  Its polarized tangent space is
\[
 T=\left\{\begin{pmatrix}0&B\\B^t&0\end{pmatrix}:B\in M_n(\mathbb C)\right\}.
\]
Let \(D_\sigma\) run through all permutations of the distinct eigenvalues
\(1,-1,\ldots,n,-n\).  For a real orthogonal frame \(g\), form the complete
permutation packets
\[
 D_{\epsilon,\sigma}=\epsilon gD_\sigma g^t\epsilon,
 \qquad
 \epsilon\in\Gamma:=\{\operatorname{diag}(\pm1,\ldots,\pm1)\}/\{\pm I\}.
\]
Retain all transposition edges inside each packet, and all full matchings
induced by the coordinate sign reflections \(r_l\).
The graph on the packets is the folded \(m\)-cube.  Its parity is
well defined because \(m\) is even, and combining this parity with
permutation parity gives a bipartition.

For an edge \(e=\{x,y\}\), write \(\ell_e\subset H^2(A,\mathcal O_A)\)
for its antiholomorphic surface wedge line.  The trace-incidence
condition for weights \(m_x\) and a tangent \(\kappa\) is the existence
of edge values \(\gamma_e\in\ell_e\) such that
\[
 m_x\,\kappa\mathbin{\lrcorner}D_x
   =\sum_{e\ni x}\varepsilon(x,e)\gamma_e,
\]
with opposite signs at the two endpoints.  These are necessary
conditions for an unfiltered deformation of an ordinary extension
supported on the graph.

\begin{theorem}\label{thm:sign-cube-obstruction}
Suppose \(n\ge3\).  There is a nonempty Zariski-open set of real
orthogonal frames \(g\), containing rational frames, for which the
following assertion holds for every positive integer \(N\) and every
choice of real weights \(m_x\).  If the weighted complete character is
\[
 \sum_x m_x e^{ND_x}=R F_n(N\theta)+cw,
 \qquad c\ne0,\quad 0\ne w\in W_K(A),
\]
then the trace-incidence condition fails for some \(\kappa\in T\).
In particular no ordinary extension with these multiplicities and
this support graph can deform in every polarized Weil direction,
even after its filtration is forgotten.
\end{theorem}

For real frames this is a linear-algebraic obstruction.  For rational
frames, choose \(N\) so that every \(ND_x\) is an integral Chern class;
positive integral weights then describe multiplicities of line bundles
in an ordinary extension.  The weights need not be positive for the
obstruction argument.  The theorem concerns this family of coordinate-sign networks,
and makes no assertion about arbitrary networks of full matchings.

\begin{proof}
Let \(\chi_{\rm all}(\epsilon)=\prod_l\epsilon_l\), a character of
\(\Gamma\).  In \(H^{p,p}\), the coefficient
\(z_I\wedge\overline z_J\) has sign character
\(\prod_{l\in I\mathbin{\triangle}J}\epsilon_l\).
Consequently the \(\chi_{\rm all}\)-component consists exactly of
the middle-degree coefficients
\[
 |I|=|J|=n,\qquad I\cap J=\varnothing,\qquad I\cup J=\{1,\ldots,m\}.
\]
They are primitive, since wedging with \(\theta\) repeats a factor.
The fixed Weil plane lies in this component.  Not every class in
this component is a Weil class.

The graph and \(T\) are stable under \(\Gamma\).  If the trace equations
are solvable for every tangent, averaging the weights and flows with
\(\chi_{\rm all}\), while transporting the tangent, preserves
solvability.  We may therefore consider weights
\(\chi_{\rm all}(\epsilon)f(\sigma)\).  Their primitive character must
be nonzero if the prescribed Weil term is nonzero.

Fix a basic tangent
\(B_{ij}=E_{ij}+E_{ji}\), taking \(i=1\) and \(j=n+1\).
The open set of frames below is chosen for this pair.
In the frame \(\epsilon g\), its matrix is
\[
 g^t\epsilon B_{ij}\epsilon g
   =\epsilon_i\epsilon_jg^tB_{ij}g.
\]
The resulting scalar Fourier mode is
\(\psi=\chi_{\rm all}\chi_{ij}\).
Along the coordinate-reflection matching at \(l\), the transported
bridge wedge changes sign at the other endpoint.  Its scalar
coefficient must therefore be invariant under \(r_l\).
In Fourier mode \(\psi\), this permits a channel only if
\(\psi(r_l)=1\), which holds exactly for \(l=i,j\).

Set \(u=g^te_i\), \(v=g^te_j\).  At an internal transposition
\(s=(kl)\), the tangent coefficient is
\(u_kv_l+v_ku_l\); the two permitted bridge coefficients are
\(u_ku_l\) and \(v_kv_l\).  On functions of the permutation vertex,
define \((R_s f)(\sigma)=f(s\sigma)\).  Add the \((k,l)\)-coordinate equations
at \(\sigma\) and \(s\sigma\).  Their unique internal edge cancels.
The nonzero eigenvalue difference can be divided out, giving
\begin{equation}
 (1-R_s)
 \left[(u_kv_l+v_ku_l)f-u_ku_l z_1-v_kv_l z_2\right]=0
 \quad\text{for every }s=(kl).
 \label{eq:sign-cube-reduction}
\end{equation}
where \(z_1,z_2\) are arbitrary scalar potentials on the permutation
vertices.  Thus this reduction has retained all possible reverse
flow coefficients.

These equations commute with the right permutation action, so they
may be examined separately in each left irreducible constituent.
Only the middle primitive constituent
\(\mathcal P_n\simeq S^{(n,n)}\) contributes a primitive middle
character.  In the square-free algebra generated by the \(x_l\),
this is the span of the products
\[
 p_M=\prod_{\{a,b\}\in M}(x_a-x_b)
\]
over perfect matchings \(M\) of \(\{1,\ldots,2n\}\).
It occurs in the diagonal degree-\(n\) module and in no other
diagonal degree.

We use the following elementary observation.  Suppose a representation
satisfies equations
\[
 (1-R_s)\left[a_sf-\sum_{r=1}^q c_{r,s}z_r\right]=0.
\]
For a perfect matching \(M\), its \(n\) transpositions commute.
Apply \(D_M=\prod_{s\in M}(1-R_s)\) to obtain
\begin{equation}
 a_sD_Mf-\sum_{r=1}^q c_{r,s}D_Mz_r=0,\qquad s\in M.
 \label{eq:sign-cube-matching}
\end{equation}
If the scalar coefficient matrix of these \(n\) equations has
column rank \(q+1\), then \(D_Mf=0\).
If this holds for every perfect matching, the
\(S^{(n,n)}\)-component of \(f\) is zero.  Indeed \(D_M\) is a
nonzero column antisymmetrizer on that irreducible module, and the
joint kernel of all its conjugates is therefore zero.  Equivalently,
the products \(p_M\) span the primitive module, and the operators
\(D_M\) are self-adjoint in the coordinate-subset inner product.

For \eqref{eq:sign-cube-reduction}, assume \(u_k\ne0\) and put \(t_k=v_k/u_k\).
Dividing the row associated with \(\{k,l\}\in M\) by \(u_ku_l\)
turns its scalar coefficients into
\begin{equation}
 \bigl[t_k+t_l,\ -1,\ -t_kt_l\bigr].
 \label{eq:sign-cube-rank}
\end{equation}
For \(n\ge3\), these rows have rank three for every perfect matching
for a generic tuple \(t\).  Each failure is an algebraic condition.
For each fixed matching it is proper: choose three of its disjoint
pairs so that their points
\((t_k+t_l,t_kt_l)\) are not collinear.
There are only finitely many matchings, so the simultaneous
nonvanishing conditions define a nonempty Zariski-open set of
tuples.

These conditions are compatible with \(u,v\) being real orthonormal
vectors.  The set of tuples \(t\) for which there are positive
weights \(w_k\) satisfying
\[
 \sum w_k=1,\qquad \sum w_kt_k=0,\qquad \sum w_kt_k^2=1
\]
contains a nonempty Euclidean open set.  For example, for six
coordinates take values
\(-2,-1,-\tfrac12,\tfrac12,1,2\), with respective symmetric
weights \(1/12,1/12,1/3,1/3,1/12,1/12\).
Additional coordinates may be assigned small positive weights,
with a small adjustment of the first six.  The point \((0,1)\)
is in the interior of the corresponding convex hull of
\((t_k,t_k^2)\), so feasibility persists under small perturbations.
This open set meets the complement of the finitely many algebraic
exceptional sets.  Setting \(u_k=\sqrt{w_k}\) and
\(v_k=u_kt_k\) gives orthonormal vectors, which can be completed
to an orthogonal frame.

It follows that the required rank conditions define a nonempty
Zariski-open subset of the orthogonal group.  Rational orthogonal
matrices are Zariski dense, for example by the rational Cayley
parametrization, so rational frames occur.  We may also exclude
unintended packet coincidences: for each noncentral coordinate
sign matrix, the condition that its conjugate by \(g\) is a signed
permutation matrix is a proper algebraic condition; there are
only finitely many such conditions.

For these frames, \eqref{eq:sign-cube-matching}--\eqref{eq:sign-cube-rank} kill the primitive
\(\chi_{\rm all}\)-component already for the chosen basic tangent.
This contradicts the required nonzero Weil character and proves
the theorem.
\end{proof}


% Standalone section fragment.  Requires amsmath, amssymb, amsthm.
% The theorem concerns diagonal trace equations only.
\section{A positive Gaussian network for the diagonal trace equations}
\label{sec:gaussian-trace-network}

We work at the product point
\[
 A=E_i^8=E_i^4\times E_i^4,
 \qquad E_i=\mathbb C/(\mathbb Z+i\mathbb Z),
\]
with its product principal polarization \(\theta\).  The balanced
Gaussian action is multiplication by \(i\) on the first block and by
\(-i\) on the second.  Write \(U\) and \(V\) for these two sets of four
coordinates.  A divisor is represented by a Hermitian matrix, and
\(\theta\) by the identity matrix.  Put
\[
 D_0=\operatorname{diag}(1,-1,2,-2,3,-3,4,-4),\qquad
 F(t)=1-\frac{15}{28}t^2+\frac{13}{80}t^4
       -\frac{41}{1008}t^6+\frac1{70}t^8.
\]
For a unitary frame \(g\), its complete permutation packet consists
of \(gP_\sigma D_0P_\sigma^*g^*\), \(\sigma\in S_8\).
The elementary symmetric-polynomial identity gives
\begin{equation}
 \sum_{\sigma\in S_8}
   \exp(gP_\sigma D_0P_\sigma^*g^*)=8!F(\theta).
 \label{eq:trace-packet-average}
\end{equation}

\begin{theorem}[A positive network satisfying the diagonal equations]
\label{thm:gaussian-positive-trace-network}
There is a finite connected bipartite graph of distinct rational
trace-zero Hermitian divisors on \(A\), with a positive integer
multiplicity at every vertex, having the following properties.
Every edge has a rank-two difference of signature \((1,1)\), and
the weighted sum of exponentials is
\[
 R F(\theta)+c w,
 \qquad R>0,\quad c\ne0,\quad 0\ne w\in W_K(A).
\]
For every polarized Weil tangent direction \(\kappa\), the vector of
weighted diagonal obstructions admits an edge flow:
\begin{equation}
 \bigl(m_x\,\kappa\mathbin{\lrcorner}D_x\bigr)_x
 \in\operatorname{im}\left(
   \bigoplus_{e}q_e\longrightarrow\bigoplus_x H^2(A,\mathcal O_A)
 \right),
 \label{eq:positive-trace-flow}
\end{equation}
where \(q_e\) is the antiholomorphic surface wedge of the edge and
the map has opposite signs at its endpoints.

After a common integral scaling, the vertices are represented by
line bundles.  Their positive multiplicities admit a single ordinary
extension that is locally free and simple, with Chern character
\(R F(N\theta)+cN^4w\).  The edge flows cancel the diagonal trace
equations for that extension.  With a suitable matrix-unit choice of
its extension entries they also cancel all matrix-valued diagonal
cohomology blocks.

No assertion is made that the full Atiyah obstruction vanishes, that
the bundle deforms in the Weil family, or that its semiregularity map
is injective.  The off-vertex equations and the strong
semiregularity condition remain unresolved.
\end{theorem}

\subsection{The finite group and its character}

Let \(G(4,4,4)\) be the group of monomial unitary \(4\times4\)
matrices with nonzero entries in \(\mu_4\) whose product is one.
Set
\[
 H=G(4,4,4)_U\times G(4,4,4)_V,
 \qquad \chi(h)=\det(h_U)\det(h_V)\in\{1,-1\}.
\]
The order of each factor is \(4^3\cdot4!=1536\).
The central scalar subgroup acting trivially on divisor matrices is
\(Z=\mu_4 I_8\), and \(\chi|_Z=1\).

\begin{lemma}\label{lem:gaussian-trace-character}
The \(\chi\)-isotypic subspace of
\(\bigoplus_p H^{p,p}(A)\) is exactly \(W_K(A)_\mathbb C\), in
degree \((4,4)\).
\end{lemma}

\begin{proof}
In one block a cohomology monomial has phase exponents in
\(\{-1,0,1\}^4\).  Invariance under the product-one diagonal
\(\mu_4\)-subgroup forces the four exponents to be equal modulo
four.  They are therefore all \(+1\), all \(-1\), or all zero.
The first two possibilities are the holomorphic and antiholomorphic
volume forms, each transforming by the determinant sign.
In the zero case, each coordinate occurs in both Hodge factors or
in neither.  The residual permutation representation is a module on
\(k\)-subsets of four coordinates.  It has no sign constituent:
the stabilizer \(S_k\times S_{4-k}\) contains an odd permutation
for every \(k\).  Thus only the two volume forms survive in each
block.  Equal holomorphic and antiholomorphic degree requires
opposite choices in the two blocks, giving exactly the two Weil
summands.
\end{proof}

There are twenty-four rank-one unitary involutions in each factor:
choose a coordinate pair and a phase in \(\mu_4\), and interchange
the two coordinates with inverse phases.  All forty-eight
involutions have \(\chi=-1\), and they generate \(H\): ordinary
interchanges generate the permutation subgroups, and products of
two phase interchanges generate the product-one diagonal subgroups.

\subsection{The exact frame and its packet stabilizer}

Use
\[
 g=h_{\mathbf1}h_{(1,2,3,4,5,6,7,8)},\qquad
 h_u=I-\frac{2uu^t}{u^tu}.
\]
Equivalently \(g=G/408\), where
\begin{equation}
 G_{ij}=408\delta_{ij}-102+(36-4i)j,
 \qquad 1\le i,j\le8.
 \label{eq:gaussian-exact-frame}
\end{equation}
The identity \(G^tG=408^2I\) is exact.  Every entry of \(G\)
is nonzero: an off-diagonal entry is \(2\) modulo \(4\), and the
diagonal entries are positive.

\begin{lemma}[No unintended packet identifications]\label{lem:gaussian-trace-frame}
The stabilizer in \(H\) of the unordered eigenline frame \(g\)
is exactly \(Z\).  Consequently the divisor vertices
\[
 D_{h,\sigma}=hgP_\sigma D_0P_\sigma^*g^*h^*,
 \qquad (hZ,\sigma)\in(H/Z)\times S_8,
\]
are distinct, and the coloring \(\chi(h)(-1)^\sigma\) is
well defined.
\end{lemma}

\begin{proof}
Suppose a monomial unitary matrix \(h\) permutes the eigenlines of
\(g\).  The columns of \(G\) have the same Euclidean norm, so the
proportionality factors induced by \(h\) have absolute value one.
Monomial unitaries preserve the sum of absolute coordinate values.
For the eight columns of \(G\), these sums are respectively
\[
 940,\ 844,\ 732,\ 676,\ 688,\ 732,\ 784,\ 832.
\]
The first column is the unique maximum, hence its line is fixed.
Its eight absolute coordinate values are
\[
 338,\ 74,\ 78,\ 82,\ 86,\ 90,\ 94,\ 98,
\]
which are pairwise distinct.  The row permutation of \(h\) is
therefore the identity.  Since all eight coordinates are nonzero,
fixing this line forces all diagonal phases to be equal.  Thus
\(h\) is scalar.  For \(h\in H\), the scalar belongs to \(\mu_4\),
and conversely every such scalar lies in \(H\) and preserves the
frame.  This proves the assertion without an enumeration.
\end{proof}

Retain all transposition edges within each packet and all full
matchings \((hZ,\sigma)\leftrightarrow(hrZ,\sigma)\) for the
forty-eight reflections \(r\).  Every matching is induced by a
rank-one unitary conjugation and therefore has a rank-two,
trace-zero Hermitian difference.  It is nonzero by the stabilizer
lemma, so its signature is \((1,1)\).  Internal transpositions and
external matchings both reverse the displayed coloring.  The graph
is connected because the two sets of generators connect the two
indices.

\subsection{A finite certificate for the reduced trace equations}
\label{subsec:gaussian-trace-certificate}

We use the convention that a Hermitian divisor is conjugated as
\(D\mapsto hDh^*\), while a polarized Kodaira--Spencer class is represented
by a complex symmetric matrix \(B\).  In a unitary eigenbasis \(g\), its
matrix is \(g^tBg\).  For a physical root normal \(u\), put
\[
 v=g^t\overline u.
\]
The antiholomorphic bridge line has, in this eigenbasis, the coordinates
\[
 q_{u,\sigma;kl}
   =v_kv_l(\lambda_{\sigma,l}-\lambda_{\sigma,k}).
\]
The sign and a common nonzero normalization of \(q\) are immaterial,
since its coefficient is a free flow variable.

Let \(\mathcal P\) be the primitive diagonal fourth-degree representation.
We realize it in the space with basis \(x_I\), \(|I|=4\), where
\(x_i^2=0\), as the fourteen-dimensional span of
\[
 p_M=\prod_{\{a,b\}\in M}(x_a-x_b)
\]
for the noncrossing perfect matchings \(M\) of
\(\{1,\ldots,8\}\).  Order these matchings recursively by pairing the
first unused index with each later index of opposite parity, and then
ordering the two remaining intervals recursively.  Let \(C\) be the
resulting \(70\) by \(14\) integer matrix.  Its first column is
\[
 p_0=(x_1-x_2)(x_3-x_4)(x_5-x_6)(x_7-x_8).
\]
The matrix \(R_s\) of a coordinate transposition \(s=(kl)\) is uniquely
determined by
\[
 CR_s=sC.
\]
All these matrices are integral.  Write \(D_s=I-R_s\); each \(D_s\)
has rank five.

We use the exact frame \eqref{eq:gaussian-exact-frame} throughout.
The identity \(G^tG=408^2I\) is included in the verification.

The roots are indexed by a pair \(a<b\) in one of the two blocks
\(\{1,\ldots,4\}\), \(\{5,\ldots,8\}\), and by
\(\zeta\in\{1,-1,i,-i\}\):
\begin{equation}
 u_{ab}(\zeta)=e_a-\zeta^{-1}e_b,\qquad
 r_{ab}(\zeta)=I-u_{ab}(\zeta)u_{ab}(\zeta)^*.
 \label{eq:gaussian-trace-roots}
\end{equation}
Thus \(r_{ab}(\zeta)\) is a rank-one unitary involution with determinant
\(-1\).  The normal has squared length two, explaining the coefficient
in \eqref{eq:gaussian-trace-roots}.

For a root in the first block, its allowed coefficient map is written
\begin{equation}
 a_u(B)=\sum_{b=5}^{8}(u^tBe_b)\,z_{u,b},
 \qquad z_{u,b}\in\mathcal P.
 \label{eq:gaussian-trace-potential-u}
\end{equation}
For a root in the second block, it is
\begin{equation}
 a_u(B)=\sum_{a=1}^{4}(e_a^tBu)\,z_{u,a}.
 \label{eq:gaussian-trace-potential-v}
\end{equation}
These are precisely the linear maps satisfying
\(a_u(r_u^tBr_u)=-a_u(B)\).  This sign is required by the determinant
character and the reversal of a bridge.

For \(f=e_1\in\mathbb Q^{14}\), the finite system to be solved is
\begin{equation}
 D_s\left[
    (g^tBg)_{kl}f-
    \sum_u v_kv_l\,a_u(B)
 \right]=0,
 \qquad s=(kl).
 \label{eq:gaussian-trace-reduced-system}
\end{equation}
It suffices to impose \eqref{eq:gaussian-trace-reduced-system} for the sixteen matrices
\[
 B_{ab}=E_{ab}+E_{ba},\qquad
 1\le a\le4< b\le8.
\]
There are \(16\cdot28\cdot14=6272\) displayed scalar equations.
Selecting five independent rows from each \(D_s\) gives \(2240\)
retained row positions.  The flow has
\(48\cdot4\cdot14=2688\) coordinates.

For completeness, the Gaussian coefficients are removed from the
linear system as follows.  For each root pair, let \(F_\zeta\) denote
the corresponding coefficient-column block in \eqref{eq:gaussian-trace-reduced-system}.  Replace
\[
 F_1,\ F_{-1},\ F_i,\ F_{-i}
 \quad\text{by}\quad
 F_1,\ F_{-1},\ \operatorname{Re}F_i,\ \operatorname{Im}F_i.
\]
The latter matrix is rational and has the same complex column span.
If \(x,y\) are its real and imaginary coefficient vectors, the
original two root potentials are recovered as
\begin{equation}
 z_i=(x-iy)/2,\qquad z_{-i}=(x+iy)/2.
 \label{eq:gaussian-trace-real-imag}
\end{equation}
Multiplication by \(408^2\) makes every coefficient integral.

The calculation can be regenerated deterministically from these
finite inputs.  Order the sixteen pairs \((a,b)\), the twenty-eight
pairs \((k,l)\), and the four-subsets lexicographically.  Order the
flow coordinates first by the block \(U,V\), then the root pair
lexicographically, then by the four column types
\(F_1,F_{-1},\operatorname{Re}F_i,\operatorname{Im}F_i\), then by
the complementary coordinate in increasing order, and finally by
the fourteen recursively ordered matchings.  Over
\(\mathbb F_{1009}\), scan the columns of each \(D_s\) from left
to right, use the first available nonzero row as the next pivot,
and retain the five original row indices so obtained.  Stack the
resulting equations in the stated tangent and transposition order.
For the resulting augmented matrix \([A\mid b]\), process columns
in consecutive panels of size \(48\).  In each remaining panel,
ordinary forward elimination scans columns from left to right and
chooses the first available nonzero row.  Move the selected original
rows, in pivot order, above the unselected rows, preserving the
latter's order.  If \(P\) is the selected square pivot block,
multiply its rows by \(P^{-1}\) and eliminate its pivot columns
from all remaining rows before processing the next panel.
This produces \(2224\) pivot rows and columns, all in \(A\).
Set all other flow coordinates to zero, solve the selected square
system over \(\mathbb Q\), clear denominators, and test every
one of the original \(6272\) equations.  These instructions specify
the rational certificate without requiring its numerators as an
additional mathematical input.

The resulting certificate supplies a positive integer \(d\) and an integer
vector \(n\), supported on \(2224\) of the \(2688\) flow coordinates,
such that
\begin{equation}
 A n=d\,b.
 \label{eq:gaussian-trace-integer-certificate}
\end{equation}
Here \(A\) is the completely specified integer matrix of \eqref{eq:gaussian-trace-reduced-system}, after
the real--imaginary replacement, and \(b\) is its right-hand side for
\(f=e_1\).  The denominator \(d\) has \(5500\) bits.  A pivot minor
was selected over \(\mathbb F_{1009}\), a square rational system was
solved by Dixon lifting, and \eqref{eq:gaussian-trace-integer-certificate} was subsequently verified by exact
integer multiplication.  Thus neither rank equality modulo a prime
nor a floating-point residual is being used as the final inference.

The accompanying computational workbench supplies the portable
program \emph{Gaussian trace network verification} and its plain
integer certificate.  The program reconstructs \(C\),
all \(R_s\), \(G\), all root coefficients, and both sides of \eqref{eq:gaussian-trace-integer-certificate}
from the formulas above, using only Python's standard library.
It checks all \(6272\) displayed coordinates, not merely the
\(2240\) row positions used to select the rational solve.  The
certificate consists of the common denominator, the integer
numerators, and their coordinate indices.  A separate finite
packet-stabilizer enumeration is retained in the workbench as
an additional check of the elementary proof above.

The certificate also checks that the chosen primitive direction has
nonzero Weil projection.  With \(U=\{1,2,3,4\}\) and
\(V=\{5,6,7,8\}\), its complementary minor coefficient is
\begin{equation}
 \sum_{|I|=4}C_{I,1}\det(g_{U,I})\det(g_{V,I})
 =
 \frac{49}{41616}\ne0.
 \label{eq:gaussian-trace-minor}
\end{equation}
Before reduction the denominator in \eqref{eq:gaussian-trace-minor} is \(408^8\).  Its
nonvanishing is checked by integer determinant expansion.

Equations \eqref{eq:gaussian-trace-reduced-system}--\eqref{eq:gaussian-trace-minor} certify the reduced diagonal trace-incidence
conditions.  The reconstruction and their precise deformation-theoretic
limits are given next; no full deformation or semiregularity assertion
is inferred from this finite calculation.


\subsection{From the reduced solution to positive multiplicities}

We now explain why the finite linear system supplies a flow on the
actual graph, rather than only an abstract representation.
Let \(\mathcal P\) and \(C\) be as in the computation.  If
\(p_I=\prod_{i\in I}\lambda_i\), let \(v\in\mathcal P\)
be the primitive projection of the diagonal fourth-degree vector
\((p_I)\), and define the rational functional
\[
 \ell(f)=\sum_{|I|=4}p_I(Cf)_I.
\]
Then \(\ell(v)>0\).  For \(f\in\mathcal P\), put
\begin{equation}
 m_f(\sigma)=\ell(R_{\sigma^{-1}}f).
 \label{eq:trace-matrix-coefficient}
\end{equation}
The convention is chosen so that
\[
 m_f(\sigma)-m_f(s\sigma)=m_{(I-R_s)f}(\sigma).
\]
The packet character associated with these multiplicities is
\begin{equation}
 \sum_{\sigma\in S_8}m_f(\sigma)e^{P_\sigma D_0P_\sigma^*}
 =\frac{8!}{14}\ell(v)\,Cf,
 \label{eq:trace-schur-lift}
\end{equation}
as a pure total cohomology vector in degree four.
Indeed, the Specht module \((4,4)\) occurs only in degree four
among the subset permutation modules.  The remaining operator in
that irreducible module is scalar by Schur's lemma; its trace is
\(8!\ell(v)\), giving \eqref{eq:trace-schur-lift}.

For the particular vector \(f=e_1\), this is completely integral:
\begin{equation}
 m_f(\sigma)=\prod_{j=1}^4
  \bigl(\lambda_{\sigma^{-1}(2j-1)}
       -\lambda_{\sigma^{-1}(2j)}\bigr).
 \label{eq:trace-explicit-multiplicity}
\end{equation}
Its absolute value is at most \(625\).  In fact, the sum of the
four absolute differences is at most
\((4+3+2+1)-(-1-2-3-4)=20\), and the arithmetic--geometric
mean inequality gives \(5^4\).  Therefore
\begin{equation}
 m_{h,\sigma}=626+\chi(h)m_f(\sigma)
 \label{eq:trace-positive-multiplicity}
\end{equation}
is an integer between \(1\) and \(1251\) at every vertex.

For clarity, the Schur scalar in this normalization is
\[
 \frac1{16}\sum_{\sigma\in S_8}m_f(\sigma)^2
 =160225344.
\]
Here \(16\) is the squared norm of the sixteen-term vector
\(Ce_1\).  The square sum is obtained simply by evaluating
\eqref{eq:trace-explicit-multiplicity} on the \(8!\) permutations.
The exact complementary-minor coefficient of \(g(Ce_1)\) is
\(49/41616\).  Thus, if \(w\) denotes the normalized real Weil
class whose two complementary matrix entries are one, the weighted
character is
\begin{equation}
 \sum_{hZ,\sigma}m_{h,\sigma}e^{D_{h,\sigma}}
 =R F(\theta)+
  160225344\,|H/Z|\,\frac{49}{41616}\,w,
 \quad R=626\,|H/Z|\,8!.
 \label{eq:trace-positive-character}
\end{equation}
The coefficient is nonzero.  This proves the character assertion
without identifying a nonzero primitive class with a Weil class
merely by its degree: the finite-group character lemma does the
identification.

\subsection{Reconstruction of the edge flows}

Let \(z_r(\kappa)\in\mathcal P\) be the certified coefficient
map for a root reflection.  Its anti-invariance is
\[
 z_r(r^tBr)=-z_r(B),
\]
with \(B\) the symmetric matrix of \(\kappa\).  The inverse
transport of a tangent by \(h\) is \(B\mapsto h^tBh\).
At the edge joining \((hZ,\sigma)\) to \((hrZ,\sigma)\),
use the scalar potential
\[
 Z_{h,r,\sigma}(B)=
  \chi(h)\,m_{z_r(h^tBh)}(\sigma).
\]
This gives a well-defined flow on the quotient by \(Z\): replacing
\(h\) by \(hz\), \(z\in\mu_4\), multiplies the inverse-transported
tangent and hence the scalar potential by \(z^2\), and multiplies
the transported antiholomorphic wedge generator by \(\overline z^2\).
Their product is unchanged.
Its natural surface wedge at the second endpoint is the negative
of that at the first.  Since \(\chi(r)=-1\) and the coefficient
map is anti-invariant, the scalar potential itself is unchanged
when the endpoints are interchanged.  Thus the two contributions
are opposite, as required for an edge flow.

After subtracting all such external flows, the coefficient at a
vertex in the coordinate wedge \((k,l)\) is the eigenvalue
difference \(\lambda_{\sigma,l}-\lambda_{\sigma,k}\) times
\[
 (g^tBg)_{kl}m_f(\sigma)
   -\sum_r v_{r,k}v_{r,l}m_{z_r(B)}(\sigma).
\]
The certified equation says exactly that the bracket is invariant
under the left transposition \((kl)\).  The eigenvalue difference
changes sign at the other endpoint of the internal matching.
Hence this residual is canceled by that internal edge.  The
constant background \(626\) has the same elementary internal
matching flow.  This proves \eqref{eq:positive-trace-flow} for all
sixteen Weil tangent directions and hence for their span.

\subsection{Line bundles, simplicity, and the remaining obstruction}

Choose a common positive integer \(N\) making every Riemann form
\(ND_{h,\sigma}\) even integral.  One may first take
\(N_0=2\cdot408^2\): the remaining matrices are Gaussian
monomial conjugates of \(G P_\sigma D_0P_\sigma^tG^t\).
Use Appell--Humbert representatives with trivial semicharacter.
Their ratios on an edge descend to the actual abelian surface
quotient and have index one, by Mumford's Appell--Humbert theorem
and complex index theorem \cite[\S2, p.~20; \S16, pp.~150,162]{MumfordAV}.
Both edge Ext\({}^1\) groups are nonzero and pair perfectly into
the line \(q_e\).  Multiplying \(N_0\) by \(1251\) makes their
dimensions at least every product of endpoint multiplicities,
because the index-one dimension grows quadratically and is a
positive integer at \(N_0\).

Let \(T\) and \(B\) be the sums of the line bundles with their
multiplicities in the two color classes.  Distinct vertices have
no Hom: a nonzero divisor of polarization degree zero cannot be
effective.  For every oriented edge choose independent Ext\({}^1\)
classes \(\eta_{ab}\), one for each rectangular matrix unit, and
give the ordinary extension
\[
 0\longrightarrow B\longrightarrow\mathcal E\longrightarrow T
 \longrightarrow0
\]
the component \(\sum_{a,b}\eta_{ab}\otimes E_{ab}\).
This extension is locally free.  An endomorphism preserves \(B\)
because \(\operatorname{Hom}(B,T)=0\), and induces block matrices
satisfying \(b_vE_{ab}=E_{ab}t_u\) at every edge.  They must be
common scalar matrices.  Connectivity therefore gives
\(\operatorname{End}(\mathcal E)=\mathbb C\).

There is a useful strengthening of the scalar trace statement for
this chosen extension.  The two endpoint products of degree-one
classes have opposite traces by the graded trace identity.  After
incorporating the incidence minus sign at the second endpoint,
the reverse cup product on one edge maps onto the pairs of endpoint
matrices having equal trace.  Indeed, the dual kernel of this
sign-adjusted map is the space of pairs satisfying
\(A_vE_{ab}=E_{ab}A_u\) for every matrix unit, which consists of
common scalars.  For a certified oriented edge flow \(a q_e\),
prescribe the actual endpoint matrices
\(a q_e I_{m_v}/m_v\) and \(-a q_e I_{m_u}/m_u\).
Their sign-adjusted traces agree, so they are realizable by a
closed reverse Ext\({}^1\) class.  Summing over edges cancels
all matrix-valued diagonal cohomology blocks; exact diagonal forms
can then be removed by a diagonal Dolbeault correction.

This still does not solve the full deformation equation.  Off-vertex blocks within either color class can be created by
products along two different incident edges.  The upper block
between the two color classes also contains the differentiated
extension class and, in a minimal model, the corresponding
\(m_3(e,x,e)\) term.  Its vanishing has not been proved.  Nor has
\(\dim\operatorname{Ext}^2(\mathcal E,\mathcal E)=104\), or an
appropriate equivariant version, been established.  The rank
\(104\) of the Hochschild action on the character is the
\(n=4\), Hankel-rank-three specialization of the polarized
criterion of Bhattacharjee--Bhattacharya, Theorem~A6.  The finite
group \(H\) above is a construction symmetry at the CM point;
it is not being asserted to be a translation/Picard group extending
over the Weil family.  The remaining off-vertex deformation equations
and the semiregularity conditions require further arguments.

\section{Why point ideals do not give the required eightfold character}
\label{sec:spin-obstructions}
This section excludes every external product of point-ideal characters in dimension eight as a source of a nonzero Weil component on the full family. It checks the whole Chern character under every Fourier--Mukai transform. The obstruction is a property of their rational Clifford orbits, rather than a failure of a particular transform.

If $R=H^1(A,\Q)$, the cohomological spin representation of the split quadratic space $R\oplus R^\vee$ is $\bigwedge^*R$. The Lie algebra acts through
\begin{equation}\label{eq:spin-lie-action}
 \mathfrak{so}(R\oplus R^\vee)
 =\bigwedge^2R\oplus\mathfrak{gl}(R)\oplus\bigwedge^2R^\vee.
\end{equation}
The three actions are exterior multiplication, the derivation induced by an endomorphism minus half its trace, and contraction. The Clifford compatibility of Fourier--Mukai equivalences is established in Golyshev--Lunts--Orlov \cite[\S3 and Proposition~4.3.7, author pp.~17--19]{GLO}; the relative construction in its proof applies between distinct abelian varieties. Consequently such an equivalence transports the Lie stabilizer of a character by a rational isomorphism. Multiplication by $e^B$ for a rational two-form is a Clifford shear and does the same. Isogeny pullback is the exterior map induced by the rational isomorphism on $H^1$, and isogeny pushforward is a nonzero scalar times its inverse. These operations also preserve the isomorphism type of the stabilizer.

\begin{proposition}[Products of two point ideals]\label{prop:point-product-spin}
Let $X,Y$ be abelian fourfolds, let $p\in X$, $q\in Y$, and let $L,M$ be line bundles. No character obtained from
\[
 (I_p\otimes L)\boxtimes(I_q\otimes M)
\]
by Fourier--Mukai equivalences, rational isogeny pullbacks or pushforwards, rational $B$-field transformations and nonzero scalar multiplication remains Hodge on a full polarized Weil domain of signature $(4,4)$.
\end{proposition}

\begin{proof}
Tensoring by $L$ and $M$ can be removed by an exponential transformation. Put $U=H^1(X,\Q)$, $V=H^1(Y,\Q)$, and let $u\in\bigwedge^8U$, $v\in\bigwedge^8V$ be the point classes. The character is
\[
 f=(1-u)(1-v)=1-u-v+uv.
\]
Its Lie stabilizer in \eqref{eq:spin-lie-action} is exactly
\begin{equation}\label{eq:point-product-stabilizer}
 \mathfrak{sl}(U)\oplus\mathfrak{sl}(V).
\end{equation}
To check this, let $(B_2,D,\gamma)$ annihilate $f$. Its degree-two part is $B_2$, so $B_2=0$. Its degree-fourteen part is $\iota_\gamma(uv)$; contraction with a top volume identifies $\bigwedge^2(U\oplus V)^\vee$ with $\bigwedge^{14}(U\oplus V)$, so $\gamma=0$. The degree-zero term then forces $\operatorname{tr}D=0$. The two off-diagonal blocks of $D$ acting on $u+v$ have distinct bidegrees $(7,1)$ and $(1,7)$. Each map is injective by contraction with the corresponding volume, so these blocks vanish. The remaining degree-eight term forces the two diagonal traces separately to vanish. Conversely these conditions fix every term of $f$, proving \eqref{eq:point-product-stabilizer}.

Suppose a transformed character $z$ were Hodge throughout a full Weil domain. The Hodge circles at its points are the conjugates of a noncentral circle in $SU(4,4)$. Their connected normal closure is all of $SU(4,4)$, because its real Lie algebra is simple. Thus this Lie algebra lies in the spin stabilizer of $z$. By the preceding rational conjugacy it would embed into
\[
 \mathfrak{sl}_8(\R)\oplus\mathfrak{sl}_8(\R).
\]
Each projection is either zero or injective, by simplicity. A nonzero projection would be an isomorphism, since both source and target have real dimension $63$. But their real ranks are $4$ and $7$, respectively. Every projection must therefore be zero, a contradiction. The same reasoning applies if the character is Hodge on a nonempty analytic open subset, by analytic continuation of the flat Hodge condition.
\end{proof}

The four nonzero coefficients in the expansion of $f$ were the only property of their values used in the stabilizer calculation. Thus the conclusion also holds for products with factor characters $r-au$ and $s-bv$, when $r,a,s,b$ are all nonzero.

\begin{proposition}[A single point secant]\label{prop:single-point-spin}
Let $Z$ be an abelian eightfold and let $c\in\Q^\times$. Apply any sequence of the rational Clifford operations above to $1+c[\mathrm{pt}_Z]$, obtaining $z$. If $z$ remains Hodge on a full polarized Weil domain of signature $(4,4)$ with polarization $\theta$, then
\[
 z\in\Q[\theta].
\]
In particular its Weil component is zero.
\end{proposition}

\begin{proof}
Let $R=H^1(Z,\Q)$, of dimension sixteen. The degree-two and degree-fourteen terms of \eqref{eq:spin-lie-action} force the exterior and contraction blocks of an element stabilizing $1+c[\mathrm{pt}]$ to vanish. Its degree-zero part forces its middle block to have trace zero. Hence the stabilizer is exactly $\mathfrak{sl}(R)$, and it fixes the two rational pure spinors $1$ and $[\mathrm{pt}]$ separately.

A nonzero spinor is called pure when its Clifford annihilator is a maximal isotropic subspace. The rational operations in the statement preserve this property. After transporting the two endpoint spinors, the connected stabilizer of $z$ therefore fixes two rational pure spinors separately, and $z$ is their linear combination. The full Weil Hodge condition puts $SU(h)$ in that connected stabilizer, by the Hodge-circle argument above.

We classify its invariant rational pure spinors. Write $H=H^1(A,\Q)$ for the target. This rational representation of $SU(h)$ is irreducible. Over the imaginary quadratic field its two absolutely irreducible constituents are the standard and dual eight-dimensional representations, which are nonisomorphic and exchanged by descent. Thus no proper rational subrepresentation can be formed from them.

If $p$ is an invariant rational pure spinor, its annihilator $L_p\subset H\oplus H^\vee$ is a rational invariant maximal isotropic subspace of dimension sixteen. The image under projection to $H^\vee$ is either zero or all of $H^\vee$, by irreducibility. In the first case $L_p=H$, and $p$ is a scalar multiple of the top spinor $[\mathrm{pt}_A]$. In the second case the projection is an isomorphism, so $L_p$ is the graph of a map $H^\vee\to H$. Isotropy makes this map an alternating two-form $D\in\bigwedge^2H$, and invariance makes it $SU(h)$-invariant.

The invariant part of $\bigwedge^2H$ is $\Q\theta$. Indeed after complexification $H=V\oplus V^*$, and the only invariant in
\[
 \bigwedge^2V\oplus(V\otimes V^*)\oplus\bigwedge^2V^*
\]
is the scalar pairing in the middle term. Therefore $D=a\theta$ for a rational $a$. The pure-spinor line annihilated by its graph is that of $e^{a\theta}$, after adjusting the graph sign convention. Every invariant rational pure spinor consequently lies in $\Q[\theta]$, so the same is true of $z$.
\end{proof}

For the particular character $1-[\mathrm{pt}]$ of a point ideal, a pure nonzero Weil target is already excluded by the Euler pairing. Its self-pairing is $-2$, whereas $\int_Aw^2>0$ for a nonzero real middle Weil class on an eightfold. Derived equivalences and $B$-fields preserve this pairing; isogenies scale it by positive factors. Proposition~\ref{prop:single-point-spin} is stronger because it also permits an arbitrary polynomial in the polarization alongside the proposed Weil component.


\begin{theorem}[All products of point secants]\label{thm:all-point-products}
Let $A_1,\ldots,A_s$ be positive-dimensional abelian varieties with
$\sum_i\dim A_i=8$. For nonzero rational $r_i,c_i$, put
\[
 f=\boxtimes_{i=1}^s\bigl(r_i+c_i[\mathrm{pt}_{A_i}]\bigr).
\]
Apply any sequence of Fourier--Mukai equivalences, rational isogeny
pullbacks or pushforwards, rational $B$-field transformations and
nonzero scalar multiplication. If the resulting character $z$ is
Hodge on a full polarized Weil domain of signature $(4,4)$, then
$z\in\Q[\theta]$. Thus its Weil component is zero.
In particular no external product of line-bundle twists of point
ideals, for any partition of dimension eight, gives the required
full-family Weil detector character by these operations.
\end{theorem}

\begin{proof}
Rescale each factor to write its spinor as
$f_i=1+b_i\operatorname{vol}_i$, with $b_i\ne0$.
Let $g_i=\dim A_i$, $R_i=H^1(A_i,\R)$ and
$T_i=R_i\oplus R_i^\vee$. Then $\dim_\R T_i=4g_i$.
Under the rational Clifford operations, the action of $SU(4,4)$ on
the total generalized vector space is conjugate to its usual action
on $H^1(A,\R)\oplus H^1(A,\R)^\vee$. The standard real
representation $W=H^1(A,\R)$ is irreducible of dimension sixteen;
its dual is isomorphic to it through the polarization. Hence the
total representation is $W\oplus W$, and every invariant real
subspace has dimension a multiple of sixteen.

First suppose that $k>0$ of the factors are elliptic. Their spinors
are pure and have Clifford annihilators of real dimension two.
For $g_i\ge2$, Clifford multiplication is the map
\[
 R_i\oplus R_i^\vee\longrightarrow\bigwedge^{\mathrm{odd}}R_i,
 \qquad(x,\xi)\longmapsto x+b_i\iota_\xi\operatorname{vol}_i.
\]
Its two summands have degrees one and $2g_i-1$, which are distinct,
so it is injective. In the external product, the images from the
different factors lie in distinct parity sectors. Its annihilator
is therefore the direct sum of the factor annihilators and has
dimension $2k$. If $z$ is Hodge on the full domain, this annihilator
is invariant under the conjugate of $SU(4,4)$, so $2k$ must be a
multiple of sixteen. Since $0<k\le8$, this forces $k=8$.
All factors are then elliptic and $f$, and hence $z$, is pure.
The invariant-pure-spinor argument in the proof of
Proposition~\ref{prop:single-point-spin} gives $z\in\Q[\theta]$.

It remains to treat $g_i\ge2$ for every factor. The orthogonal Lie
algebra has block decomposition
\[
 \mathfrak{so}\Bigl(\bigoplus_iT_i\Bigr)
 =\bigoplus_i\mathfrak{so}(T_i)
  \ \oplus\!\bigoplus_{i<j}(T_i\otimes T_j).
\]
A mixed block $T_i\otimes T_j$ acts on $f$ through the tensor product
of the two injective Clifford maps just displayed. It is therefore
injective. Its image has odd parity in precisely factors $i,j$;
these parity sectors are distinct for different pairs and are
disjoint from the images of the diagonal Lie blocks. Consequently
every Lie stabilizer element has zero mixed blocks. The full
$SU(4,4)$ action therefore preserves each $T_i$ separately.

Each $4g_i$ must now be a multiple of sixteen. A partition of eight
into positive multiples of four is either $8$ or $4+4$.
The first case is Proposition~\ref{prop:single-point-spin}; the
second is Proposition~\ref{prop:point-product-spin} and its
nonzero-coefficient extension. They give the stated conclusion.
Line-bundle twists can be removed by exponential transformations,
which were included among the permitted operations.
\end{proof}

The proof deliberately does not identify the internal stabilizer
of a surface factor with a special linear group. In dimension two,
exterior multiplication in degree two can cancel contraction of the
top term. The injectivity of the odd Clifford map is the property
needed for the mixed-block argument, and it holds in every factor
of dimension at least two.

These exclusions leave the pure character of Proposition~\ref{prop:sixteen-line-character} untouched. Its two complex determinant lines are exchanged by the imaginary quadratic field and are not rational pure-spinor endpoints. They also leave arbitrary full-support sheaves and other perfect complexes open. A positive construction must meet both requirements: a whole character that stays Hodge on the full domain and an invariant obstruction space satisfying Theorem~\ref{thm:pure-weil-ext} when the normalized character is pure.

\section{An obstruction for arbitrary complexes on the graph supports}
\label{sec:localized-support}
We return to the norm-one graph supports of Proposition~\ref{prop:graph-obstruction}. Replacing that direct sum by extensions on the same supports does not remove its normal obstruction. The trace of multiplication by a normal coordinate detects this obstruction even when the individual support components are not assumed to deform.

\begin{theorem}\label{thm:localized-support}
Let $Z_1,\ldots,Z_r$ be translates of abelian subvarieties of the same codimension $c$ in a complex abelian variety $A$. Suppose every pairwise intersection has codimension at least two in each component. Let $E$ be a perfect complex whose cohomology sheaves are supported on $\bigcup_iZ_i$. Write
\[
 m_i(E)=\sum_j(-1)^j
 \operatorname{length}_{\mathcal O_{A,\eta_i}}
 \mathcal H^j(E)_{\eta_i}
\]
for its generic virtual multiplicity at $Z_i$. For a first-order deformation $\kappa\in H^1(A,T_A)$, let $\nu_i(\kappa)\in H^1(Z_i,N_{Z_i/A})$ be its normal restriction. If $m_i(E)\ne0$ and $\nu_i(\kappa)\ne0$ for some $i$, then $E$ has no perfect deformation in direction $\kappa$, even with an arbitrary infinitesimal gerbe twist.
\end{theorem}

\begin{proof}
Remove the pairwise intersections and then the other components. On the resulting open neighborhood of $Z_i^\circ$, the complex is supported on $Z_i^\circ$ alone. Since the removed subset has codimension at least two in the smooth variety $Z_i$, its local cohomology in degrees zero and one with coefficients in a locally free sheaf vanishes. The localization sequence gives an injection
\[
 H^1(Z_i,N_{Z_i/A})\lhook\joinrel\longrightarrow
 H^1(Z_i^\circ,N_{Z_i/A}|_{Z_i^\circ}).
\]
For finite intersections in a four-dimensional graph this is an isomorphism; only injectivity is needed.

Translate $Z_i$ to the origin and write $A=V/\Lambda$ and
$Z_i=W/\Lambda_i$, where $\Lambda_i=\Lambda\cap W$. Choose a complex
vector-space complement $N$ to $W$. A first-order deformation of a complex
torus is a first-order change of its periods. On a lifted coordinate cover
near $Z_i$, its transition maps can therefore be written
\[
 x_b=x_a+\lambda_{ab}+\epsilon a_{ab},\qquad
 z_b=z_a+\epsilon b_{ab},\qquad \epsilon^2=0.
\]
Here $x$ is tangential, $z$ is normal, $\lambda_{ab}\in\Lambda_i$, and
$a_{ab},b_{ab}$ are constant on each connected overlap. The cocycle
$(b_{ab})$ represents $\nu_i(\kappa)$. The tangential transitions are
independent of $z$ and define a first-order deformation of $Z_i$.
These coordinates may be taken on formal neighborhoods: the lattice of
the quotient torus is discrete, so only $\Lambda_i$ identifies points in
the completion along $Z_i$.

Suppose that a perfect lift $\widetilde E$ exists, allowing an arbitrary
gerbe whose central restriction is trivial. Trivialize the gerbe on the
coordinate charts and push the resulting complexes along the local
projection to the tangential chart. Denote the pushed complexes by $F_a$.
They are perfect. Indeed their cohomology is finite over the tangential
ring, because the normal ideal acts nilpotently on the central
cohomology and $\epsilon$ is nilpotent. A bounded complex of finite free
modules over the ambient formal power-series ring has finite Tor
amplitude over the tangential ring, since that power-series ring is flat.
Finite coherent cohomology and finite Tor amplitude imply perfection.
This argument permits arbitrary cohomology and does not require a flat
sheaf representative of $\widetilde E$.

The virtual rank of $F_a$ is $m=m_i(E)$. At the generic point of the
central tangential chart this is precisely the alternating length in the
statement, after forgetting the nilpotent normal-coordinate action.
For each normal coordinate take the categorical trace
\[
 t_a^\alpha=\operatorname{Tr}_{F_a}(z_a^\alpha).
\]
Modulo $\epsilon$ this trace is zero: over the generic point,
$z_a^\alpha$ acts nilpotently on every finite-dimensional cohomology
space, and its alternating trace vanishes. A holomorphic function on the
reduced chart vanishing generically vanishes identically. Thus
$t_a^\alpha=\epsilon f_a^\alpha$ for a holomorphic function
$f_a^\alpha$ on the central chart.

On an overlap, invariance of categorical trace under conjugation gives
\[
 t_b^\alpha-t_a^\alpha
   =\epsilon b_{ab}^\alpha\operatorname{Tr}_{F_a}(\operatorname{id})
   =\epsilon m b_{ab}^\alpha.
\]
The scalar gerbe factor on a triple overlap cancels under conjugation.
The tangential first-order translation changes $\epsilon f_a^\alpha$
only to order $\epsilon^2$. It follows that
$f_b^\alpha-f_a^\alpha=m b_{ab}^\alpha$. Hence
$m\nu_i(\kappa)|_{Z_i^\circ}=0$. The injection already proved and $m\ne0$
contradict the hypothesis.

The usual first-order obstruction theorem identifies this
nonexistence with the nonvanishing of
$\kappa\mathbin{\lrcorner}\operatorname{At}(E)+\beta\operatorname{id}_E$
for every $\beta$; see \cite[Corollary~3.4]{HT}. The scalar term is also
seen directly in the twisted \v Cech lifting equations. No simplicity
or vanishing of negative self-extensions is needed for this square-zero
obstruction.
\end{proof}

The local trace used here is consistent with the localized
semiregularity construction of Buchweitz--Flenner
\cite[Remark~4.7(2), p.~170; Proposition~8.2 and formula~(3),
pp.~200--201]{BF}. The generic multiplicity is the alternating length
specified in the theorem, not the virtual rank of the derived
restriction to the support, which has normal-Tor cancellation.

Distinct scalar norm-one graphs in $E_i^4\times E_i^4$ meet in finitely many points, and distinct translates of the same graph are disjoint. Their normal classes are $B^{\mathsf t}-B$. The theorem therefore excludes every perfect complex supported on a finite union of these graphs with nonzero generic virtual multiplicity from deforming in the six skew-matrix directions. In particular a nonzero pure degree-eight Weil character forces some such multiplicity to be nonzero. The restriction now covers arbitrary iterated cones, direct summands and thickenings on these supports. It leaves open complexes with larger support, different geometric supports, or components meeting along divisors.

\section{A cubic bound for the exceptional classes of a Mumford square}
\label{sec:mumford}
This section treats a family whose exceptional classes are not Weil
classes. Its purpose is to reduce the two-dimensional space of classes
to one detecting subvariety, with a degree estimate.

Mumford's construction gives Shimura curves of polarized abelian
fourfolds whose general member has endomorphism algebra $\Q$, although
its Hodge group is smaller than the full symplectic group; see
\cite[\S1.1]{Noot}. Galluzzi describes Hazama's exceptional classes on
their squares in \cite[Introduction, p.~727;
Propositions~4.1--4.2, pp.~733--734]{Galluzzi}.
Moonen--Zarhin \cite[Theorem~0.1(iii), p.~2]{MoonenZarhin}
identify these as exceptional classes which are not of Weil type.
The representation calculation below recalls that description. The
additional conclusion is the cubic estimate obtained by composing one
algebraic correspondence with itself.

\subsection{The exceptional plane}

Let $(A,\theta)$ be a polarized abelian fourfold whose Hodge group is the
full group of a Mumford family. Thus
\[
 G^{\mathrm{sc}}_{\C}\simeq\mathrm{SL}_2^3,\qquad
 V_{\C}:=H^1(A,\C)\simeq U_1\otimes U_2\otimes U_3,
 \qquad \dim U_i=2.
\]
Here $G^{\mathrm{sc}}$ is the simply connected cover of $G$.
The three factors descend by the transitive action associated with a
totally real cubic field $F$. The hypothesis concerns the Hodge group
of this fibre, not merely its membership in the family. Put
\[
 P=H^2(A,\Q)_{\mathrm{prim}},\qquad
 Q_2(x,y)=\int_Axy\theta^2,\qquad
 \Theta=\theta_1+\theta_2\text{ on }A^2,\qquad
 \mu=\int_A\theta^4/4!.
\]

\begin{lemma}\label{lem:mumford-plane}
The primitive space $P$ has dimension $27$ and
$\operatorname{End}_{\mathrm{HS}}(P)=F$. The rational Hodge classes
of codimension two on $A^2$ form the direct sum of the six-dimensional
space of divisor products and a two-dimensional space
\[
 E\simeq F^0:=\ker(\operatorname{Tr}_{F/\Q}:F\longrightarrow\Q).
\]
The identification is induced by
$P\otimes P\simeq\operatorname{End}_{\Q}(P)$ using $Q_2$.
\end{lemma}

\begin{proof}
The transposition sign in $V^{\otimes2}$ gives
\[
 \bigwedge\nolimits^2V_{\C}
  =(\mathbf1,\mathbf1,\mathbf1)
   \oplus(\mathbf1,\mathbf3,\mathbf3)
   \oplus(\mathbf3,\mathbf1,\mathbf3)
   \oplus(\mathbf3,\mathbf3,\mathbf1).
\]
The first summand is the polarization. The other three are
nonisomorphic absolutely irreducible representations of dimension nine.
Their transitive descent has commutant the defining cubic field $F$.
Since the Hodge group is $G$, this commutant is exactly the Hodge
endomorphism algebra of $P$.

For completeness, the invariant dimensions in the five K\"unneth
components of $H^4(A^2)$ are $1,1,4,1,1$. One can compute the two outer
dimensions without decomposing all exterior powers. The
$\mathrm{SL}_2$-invariants in $U^{\otimes4}$ have the $S_4$ character
$(2,0,2,-1,0)$, on the classes of cycle types
$1^4,\,2\,1^2,\,2^2,\,3\,1,\,4$. They are spanned by the three pairings
of four vectors, with their single alternating relation. Cubing the
character for the three factors gives $(8,0,8,-1,0)$.
Its sign multiplicity is
$(8+3\cdot8-8)/24=1$, and its sign multiplicity after restriction
to $S_3$ is $(8-2)/6=1$. These are respectively the invariants in
$\bigwedge^4V$ and in $V\otimes\bigwedge^3V$. The middle invariant
dimension is $\dim\operatorname{End}_G(\bigwedge^2V)=4$.

The three independent divisor classes on $A^2$ are
$\theta_1,\theta_2$ and the mixed Poincar\'e class $\ell$.
Their six quadratic products are independent. Four occupy the four
outer K\"unneth components; in the middle component,
$\theta_1\theta_2$ and $\ell^2$ are independent since $\ell^2$ has a
nonzero primitive coevaluation component. Under the identification of
the primitive invariant tensors with $F$, their middle primitive
line is $\Q\operatorname{id}_P$. Its complement is $F^0$.
\end{proof}

Let $\delta_{\mathrm{exc}}$ be the least $\Theta$-degree of an
irreducible codimension-two subvariety of $A^2$ whose projection to $E$
is nonzero. Such subvarieties have dimension six. The projection is the
one furnished by the decomposition in Lemma~\ref{lem:mumford-plane}.
As before the value is $+\infty$ if no detector exists.

\begin{theorem}\label{thm:mumford-cubic}
For every fixed nonzero rational class $t\in E$, there is a constant
$C_t>0$ such that
\[
 \delta_{\mathrm{exc}}\leq c_t\leq C_t\delta_{\mathrm{exc}}^3.
\]
The constants may be chosen uniformly on the Hodge-generic fibres of
one fixed polarized Mumford family, after a fixed finite cover making
$E$ and $t$ global. In particular, one nonzero algebraic exceptional
class implies algebraicity of the whole exceptional plane at a
Hodge-generic fibre.
\end{theorem}

The degree exponent has a simple explanation. An algebraic detector
gives a trace-zero element $a$ of a cubic field at linear cost.
Composing its correspondence produces $a^2$ at quadratic cost.
The trace-zero parts of $a$ and $a^2$ are independent. The area of the
lattice they span, and the cost of clearing the coefficients of a
fixed class, are cubic in the original degree.

\subsection{Algebraic operations and their costs}

We give the operators explicitly, to avoid assuming an algebraic
projector onto a Hodge endomorphism field.
For an integer $m\geq2$, the polynomial
\[
 \pi_k=
 \prod_{\substack{0\leq j\leq8\\j\ne k}}
 \frac{[m]^*-m^j}{m^k-m^j}
\]
in algebraic correspondences is the cohomological projector onto
$H^k(A)$. Let $L$ denote multiplication by $\theta$, and let $*$ denote
Pontryagin product. The lowering operator is
\[
 \Lambda(x)=\frac{\theta^3}{6\mu}*x.
\]
The identity follows by writing the polarization in a symplectic basis:
Pontryagin product contracts one paired pair of basis covectors, with
coefficient $\mu$, cancelled by the displayed normalization.
It is the usual cohomological Lefschetz operator, including for a
nonprincipal polarization. In particular on $H^2(A)$,
\begin{equation}\label{eq:mumford-operators}
 \pi_\theta=\tfrac14L\Lambda,\qquad
 \pi_P=1-\pi_\theta,\qquad
 R=(L^2)^{-1}
   =\left(\tfrac14\pi_P+\tfrac1{36}\pi_\theta\right)\Lambda^2.
\end{equation}
Indeed $\Lambda^2L^2$ is $4$ on primitive degree two and $36$ on
$\Q\theta$. These also follow from
$\Lambda L^jx=j(4-k-j+1)L^{j-1}x$ for a primitive degree-$k$ class.
All the operators in \eqref{eq:mumford-operators} are algebraic.

The identity tensor $I_P\in P\otimes P$ is algebraic: take the primitive
$H^6\otimes H^2$ component of the diagonal and apply $R$ to its first
factor. For a class $z\in P\otimes P$ write
\[
 a=\Phi(z),\qquad
 \Phi\left(\sum_i u_i\otimes v_i\right)(x)
   =\sum_i Q_2(x,u_i)v_i.
\]
Its rational trace is the algebraic intersection
\[
 \operatorname{Tr}_{P}(a)
   =\int_{A^2}zI_P\theta_1^2\theta_2^2.
\]
Thus the operation
\begin{equation}\label{eq:mumford-trace}
 z\longmapsto z-\frac{\operatorname{Tr}_{P}(a)}{27}I_P
\end{equation}
is algebraic with a fixed denominator. At a Hodge-generic fibre
$a\in F$ and $\operatorname{Tr}_{P}(a)=9\operatorname{Tr}_{F/\Q}(a)$.
Projecting a cycle to the $(2,2)$ K\"unneth component, then to
$P\otimes P$, and then subtracting this trace is exactly its
exceptional projection.

\begin{lemma}\label{lem:mumford-cost}
Give a rational cycle $Z=\sum_i r_iZ_i$ the signed cost
$|Z|_\Theta=\sum_i|r_i|\deg_\Theta Z_i$.
The fixed operators just described change cost $B$ to at most $CB$,
with fixed denominators. Composition of two correspondences on $A$
of costs $B_1,B_2$ has a signed representative of cost at most
$C'B_1B_2$. The constants depend only on the fixed data and are uniform
in the polarized family.
\end{lemma}

\begin{proof}
Pullbacks and pushforwards by multiplication maps, products with fixed
divisors, and Pontryagin products with fixed cycles have linear cost
bounds. Translate cycles generally before an intersection. Translation
preserves their classes and degrees, and the resulting proper
intersection of effective cycles is effective.
Fixed rational coefficients and the denominators of the displayed
operators are harmless.

Here is an intrinsic estimate covering both these assertions and
composition. On a polarized abelian variety use the flat K\"ahler
metric whose form represents the polarization. An effective
$d$-cycle $W$ has mass $\deg W/d!$. A harmonic form is constant, so
its pointwise norm is its $L^2$ norm divided by
$\sqrt{\operatorname{Vol}}$. Comass is at most pointwise norm.
Duality against all harmonic forms of unit norm consequently gives
\begin{equation}\label{eq:mumford-mass}
 \|[W]\|_{L^2}\leq
       \frac{\deg W}{d!\sqrt{\operatorname{Vol}}}.
\end{equation}
Exterior multiplication between two fixed Euclidean exterior powers
is a bounded bilinear map, with a constant depending only on their
degrees and the dimension.

For effective codimension-four correspondences $C,D$ on $A^2$, of
degrees $b,c$, intersect $p_{12}^*C$ and $p_{23}^*D$ on $A^3$
after translating $D$ generally. The action of translations on
the two coordinates of $D$ is transitive on $A^2$, so generic
transversality makes this intersection proper.
Its pushforward by $p_{13}$ represents the composition.
With $\Theta_3=\theta_1+\theta_2+\theta_3$, the pullback degrees are
\[
 \deg_{\Theta_3}p_{12}^*C
   =\binom84(4!\mu)b,\qquad
 \deg_{\Theta_3}p_{23}^*D
   =\binom84(4!\mu)c.
\]
Apply \eqref{eq:mumford-mass}, the bilinear exterior estimate and pairing with
$\Theta_3^4$. The degree of the intersection is at most a fixed
multiple of $bc$. Since
$p_{13}^*\Theta\leq_{\mathrm{nef}}\Theta_3$, the degree of its
pushforward is no greater. Expand signed cycles linearly to obtain
the general assertion. The flat metric estimate depends only on
dimension and polarization volume; the fixed maps have the same
norms and degrees throughout the family.
\end{proof}

\subsection{Proof of the cubic estimate}

\begin{proof}[Proof of Theorem~\ref{thm:mumford-cubic}]
Let $Y$ be a detector of degree $D$. Fix the order
\[
 \mathcal O=F\cap\operatorname{End}_{\Z}(P_{\Z})
\]
for the fixed primitive cohomology lattice. The fixed projectors and
trace subtraction, followed by one fixed clearing of denominators,
give a nonzero $\alpha\in\mathcal O^0$ represented by an integral
signed codimension-two cycle of cost at most $C_1D$.

Multiply its tensor by $\theta_1^2$ to make a codimension-four
correspondence on $A$ acting as $\alpha$ on $P$; use the fixed
projectors to remove its other cohomological components.
Compose it with itself and use $R$ on the degree-six factor to
return to codimension two. Lemma~\ref{lem:mumford-cost} gives a
representative of $\alpha^2$ of cost $O(D^2)$ and a fixed
denominator. Trace subtraction and multiplication by one fixed
positive integer $r$ give
\[
 \beta=r\bigl(3\alpha^2-
        \operatorname{Tr}_{F/\Q}(\alpha^2)\operatorname{id}\bigr)
      \in\mathcal O^0
\]
with an integral signed representative of cost at most $C_2D^2$.
The trace term has this cost bound by the estimate below.
A nonzero trace-zero element of a cubic field is nonrational and
generates the field. If $\beta$ were a rational multiple of $\alpha$,
a polynomial of degree two would annihilate $\alpha$, which is
impossible. Thus $\alpha,\beta$ span $F^0$.

Give $F^0_{\R}$ the norm
$\|a\|_F^2=\operatorname{Tr}_{F/\Q}(a^2)$.
The adjoint on $F$ is the identity, since the three scalar summands
of $P_\C$ are mutually orthogonal. If $z_a$ is the tensor
corresponding to $a$, then
\[
 \int_{A^2}z_a^2\Theta^4
   =6\operatorname{Tr}_{P}(a^2)
   =54\operatorname{Tr}_{F/\Q}(a^2).
\]
Each factor of $z_a$ is primitive of degree two; hence
$\Theta^5z_a=0$. The primitive Hodge-star formula gives
$\|z_a\|_{L^2}^2=\frac94\|a\|_F^2$.
Since $\operatorname{Vol}(A^2,\Theta)=\mu^2$, inequality
\eqref{eq:mumford-mass} for signed six-cycles of cost $B$ yields
\[
 \|a\|_F\leq\frac{2B}{3\cdot6!\mu}.
\]
In particular $\|\alpha\|_F=O(D)$,
$\operatorname{Tr}_{F/\Q}(\alpha^2)=O(D^2)$, and
$\|\beta\|_F=O(D^2)$.

Choose $\delta\geq1$ with $\delta t\in\mathcal O^0$ and set
$h=[\mathcal O^0:\Z\alpha+\Z\beta]$.
The positive integral trace pairing on $\mathcal O^0$ has
covolume at least one. The area formula therefore gives
$h\leq\|\alpha\|_F\|\beta\|_F=O(D^3)$.
There are integers $u,v$ such that
\[
 h\delta t=u\alpha+v\beta.
\]
The two-dimensional determinant formula bounds
$|u|\leq\delta\|t\|_F\|\beta\|_F=O_t(D^2)$ and
$|v|\leq\delta\|t\|_F\|\alpha\|_F=O_t(D)$.
The multiplier $h\delta$ is consequently $O_t(D^3)$, and the
signed representative on the right has total cost
$|u|C_1D+|v|C_2D^2=O_t(D^3)$.
Separate its positive and negative coefficients to obtain the
upper bound for $c_t$.

Conversely, every effective difference representing a nonzero
multiple of $t$ has an irreducible component with nonzero
exceptional projection. Its degree is no greater than the
corresponding numerator degree. This proves the lower bound.
The order, its trace pairing, the fixed operators and the norm of
$t$ are fixed in the family, proving the asserted uniformity.
\end{proof}

\begin{corollary}\label{cor:mumford-family}
On a fixed connected Mumford Shimura curve, suppose that detectors
for its exceptional plane have degrees bounded by one integer $D$
at infinitely many CM points. Then, for every fixed global
$t\in E$, the complexity $c_t$ is bounded on every fibre.
\end{corollary}

\begin{proof}
The locus with a degree-$D$ detector is closed: use the proper
bounded-cycle parameter spaces of Lemma~\ref{lem:parameters}, and
select components on which the projection to the global flat plane
is nonzero. This uses a cohomological projection to specify a
locus, not an algebraic correspondence projecting to $F$.
An infinite subset of the irreducible curve is Zariski dense, so
every fibre has a degree-$D$ detector. Apply
Theorem~\ref{thm:mumford-cubic} on its dense Hodge-generic locus.
The resulting uniform bound for $c_t$ extends to every fibre by
the closedness of Lemma~\ref{lem:closed}.
\end{proof}

The Hodge-generic qualification in the pointwise theorem is necessary
for this proof. At a CM fibre,
$\operatorname{End}_{\mathrm{HS}}(P)$ may be larger than $F$.
The algebraic trace subtraction above then projects to the
trace-zero part of that larger algebra; it need not extract $F^0$
from an arbitrary detecting cycle. If a class already lies in
$F^0$, the field-composition argument applies. Corollary
\ref{cor:mumford-family} passes through the generic fibres and does
not assume a pointwise extraction operator at CM fibres.


\section{A surface-correspondence criterion for a Mumford square}
\label{sec:mfg-surface-gate}

This section turns the K3 approach to Mumford fourfolds into a precise sufficient condition on one surface correspondence.

Let $(A,\theta)$ be a Hodge-generic Mumford abelian fourfold, let $F$ be
its totally real cubic field, and put $V=H^1(A,\mathbf Q)$.  We use the
split-corestriction case of Mumford's construction, so that
\[
 G^{\mathrm{sc}}_{\mathbf C}\simeq\mathrm{SL}_2^3,
 \qquad V_{\mathbf C}=U_1\otimes U_2\otimes U_3,
 \qquad \dim U_i=2.
\]
Assume $A=\operatorname{Jac}(C)$ for a smooth curve $C$ of genus four.
This hypothesis applies to the two explicit families of
Bouchet--Hanselman--Pieper--Schiavone
\cite[Theorem 1.2]{BHPS}.  The conditional real-multiplication and
Kuga--Satake approach already appears in their Section 5.2; the purpose
here is to state precisely a sufficient surface-correspondence input
and to give the algebraic extraction argument.

\subsection{The two cubic-field representations}

The tensor decompositions are
\begin{align*}
 \operatorname{Sym}^2V_{\mathbf C}
  &=(3,3,3)\oplus(3,1,1)\oplus(1,3,1)\oplus(1,1,3),\\
 \Lambda^2V_{\mathbf C}
  &=(1,1,1)\oplus(1,3,3)\oplus(3,1,3)\oplus(3,3,1).
\end{align*}
Here $(a,b,c)$ records the dimensions of the three irreducible factors.
The adjoint sum in the first line descends to a rational rank-nine
Hodge structure $T$.  The nontrivial sum in the second is
$P=H^2(A,\mathbf Q)_{\mathrm{prim}}$, of rank $27$.
The three factors are permuted transitively by the Galois action defining
$F$, and occur with multiplicity one.  Schur's lemma and descent give
\[
 \operatorname{End}_{\mathrm{HS}}(T)=F,
 \qquad \operatorname{End}_{\mathrm{HS}}(P)=F.
\]
The first equality concerns $T$, and does not assert extra endomorphisms
of $A$: generically $\operatorname{End}^0(A)=\mathbf Q$.
The location of $T$ in $V\otimes V$, rather than in $H^2(A)$, is also
explained in \cite[Section 2.3]{Zhu}.

Identify $V$ with $H^1(C,\mathbf Q)$.  Then
\[
 T\subset \operatorname{Sym}^2V\subset H^2(C\times C,\mathbf Q)
\]
is the unique copy of the rational adjoint representation in $V\otimes V$.
Let $q_0$ be the restriction of the intersection pairing of $C\times C$
to $T$.  This is a K3-type polarization of signature $(2,7)$: the split
real factor contributes Hodge numbers $(1,1,1)$, and the two compact
factors contribute six further $(1,1)$ directions.  There is an
$F$-bilinear ternary form $q_F$ such that
\[
 q_0=\operatorname{Tr}_{F/\mathbf Q}q_F.
\]
The $F$-action is self-adjoint for these forms.

For $P$ use
\[
 Q_2(x,y)=\int_Axy\theta^2.
\]
It identifies $P\otimes P$ with $\operatorname{End}(P)$.  Inside its
Hodge classes, the scalar line is generated by the primitive identity
tensor, which is a divisor-product class.  The complementary plane
\[
 E\simeq F^0:=\ker\operatorname{Tr}_{F/\mathbf Q}
\]
is the usual exceptional plane in $H^4(A^2,\mathbf Q)$; see
\cite[Propositions 4.1--4.2]{Galluzzi}.

\subsection{A nonsquare polarization twist}

Choose a totally positive element
\[
 \lambda\in F^\times_{>0}\setminus\mathbf Q^\times F^{\times2},
 \qquad q_\lambda(x,y)=q_0(\lambda x,y).
\]
Such elements exist.  For example, at a completely split rational prime,
weak approximation gives a totally positive element whose three
valuation parities are $(1,0,0)$, whereas a rational number times a
square has equal parities at those three primes.

The polarized rational Hodge structure $(T,q_\lambda)$ is the
transcendental Hodge structure of a projective K3 surface $S_\lambda$,
up to a positive rational rescaling of the form.  To see this, choose
an even integral lattice of rank nine in the rational quadratic space.
Its discriminant group has length at most nine.  The complementary rank
in the K3 lattice is thirteen, and the complementary signature is
$(1,12)$.  Nikulin's primitive-embedding theorem applies
\cite[Theorem~1.12.4, pp.~122--123]{Nikulin}. Extend the Hodge decomposition by declaring its orthogonal complement to have type $(1,1)$. The K3 period theorem \cite[Theorem~1.7, p.~109]{MorrisonK3} realizes this signed Hodge structure on the K3 lattice.  The complementary lattice
contains a positive class, so the K3 surface is projective.  Since $T$
has no rational $(1,1)$ classes at a Hodge-generic Mumford point, its
transcendental rank is nine and $\rho(S_\lambda)=13$.
A K3 realization of this rank-nine structure with Picard number
thirteen is also given in
\cite[Appendix~A, Theorem~A23]{BBK3Sixfolds}; the argument here
keeps the nonsquare polarization parameter explicit.
A rational rescaling can be absorbed into $\lambda$ without changing
its class modulo $\mathbf Q^\times F^{\times2}$.

\begin{theorem}[One surface correspondence suffices]\label{thm:mumford-surface}
Suppose there is an algebraic cycle
\[
 Z\in\mathrm{CH}^2(S_\lambda\times C\times C)_{\mathbf Q}
\]
whose correspondence action has a nonzero component
\begin{equation}
 \gamma:T(S_\lambda)\longrightarrow
              H^1(C,\mathbf Q)\otimes H^1(C,\mathbf Q).
 \label{eq:mumford-surface-map}
\end{equation}
Then the entire exceptional plane $E\subset H^4(A^2,\mathbf Q)$ is
algebraic.  It is enough to find one effective irreducible surface $Z$
with this property.
\end{theorem}

\begin{proof}
The degree-one K\"unneth projectors of the two curves are algebraic.
The Poincar\'e line bundle and the Abel--Jacobi map give algebraic
inverse identifications of $H^1(C)$ and $H^1(A)$.  Thus the indicated
component is itself an algebraic correspondence into $V\otimes V$.
Since $T$ occurs there with multiplicity one and
$\operatorname{End}_{\mathrm{HS}}(T)=F$, after the fixed Hodge
identification it is multiplication by some $c\in F^\times$ followed
by the inclusion of $T$.

The transcendental projector of $S_\lambda$ is algebraic.  Explicitly,
subtract the degree-zero and degree-four projectors from its diagonal,
and subtract the projector onto its N\'eron--Severi group using the
inverse intersection matrix of a divisor basis.  Its degree-two tensor
is the inverse of $q_\lambda$.  Applying $\gamma$ in both factors gives
an algebraic tensor in
\[
 \operatorname{Sym}^2V\otimes\operatorname{Sym}^2V
        \subset H^4(A^4,\mathbf Q).
\]
Over the three real embeddings of $F$, this is the sum of the three
adjoint Casimir tensors with coefficients
\begin{equation}
 \alpha_i=\sigma_i(c^2/\lambda),\qquad i=1,2,3,
 \label{eq:mumford-casimir-coeff}
\end{equation}
up to a common nonzero rational normalization. Indeed the identification of the adjoint representation with symmetric tensors by the alternating form on $V$ turns $q_0$ into a fixed rational multiple of the trace pairing on $\operatorname{End}(V)$. Each $\mathfrak{sl}_2$ factor acts on $V$ with multiplicity four, so this normalization is the same in all three factors. The inverse polarization
contributes $\lambda^{-1}$ and the two copies of $\gamma$ contribute
$c^2$.  By the choice of $\lambda$, the element $c^2/\lambda$ is not
rational.

Regroup the four $V$ factors and antisymmetrize in the two new pairs.
This is an algebraic operation: permute abelian factors and pull back
along two diagonal maps.  We compute its action on $P$.  Normalize the
$\mathfrak{sl}_2$ Casimir as
\[
 C=H^2/4+(EF+FE)/2,
\]
with eigenvalues $3/4$ on the standard representation and $2$ on its
symmetric square.  Its cross term on $\Lambda^2V$ is
\[
 \Omega_i=C_i/2-\tfrac34\operatorname{id}.
\]
Let $P_j$ be the nine-dimensional summand on which the $j$th factor is
trivial.  Then $\Omega_i$ has eigenvalue $-3/4$ on $P_i$ and $1/4$ on
the other $P_j$.  The recoupled tensor therefore acts on $P_j$ by
\begin{equation}
 \frac{\alpha_1+\alpha_2+\alpha_3}{4}-\alpha_j,
 \label{eq:mumford-casimir-eigen}
\end{equation}
up to one common nonzero factor.  To check that antisymmetrization is
indeed this cross-Casimir map, identify an adjoint element $X$ with the
symmetric form $b_X(u,v)=\varepsilon(Xu,v)$ using the alternating
polarization $\varepsilon$ of $V$.  Antisymmetrizing its Casimir tensor
gives twice
\[
 \sum_X\bigl(b_X(u,w)b_{X^\vee}(v,z)
             -b_X(u,z)b_{X^\vee}(v,w)\bigr),
\]
which is the cross-Casimir pairing on $u\wedge v,w\wedge z$.

Project to $P\otimes P$ and subtract the rational trace.  Both operations
are algebraic, as explained below.  Formula
\eqref{eq:mumford-casimir-eigen} becomes, up to the same nonzero factor,
\[
 -\left(\alpha_j-
      \frac{\alpha_1+\alpha_2+\alpha_3}{3}\right).
\]
This is a nonzero element $a\in F^0$, since $c^2/\lambda\notin\mathbf Q$.
We have therefore constructed one nonzero algebraic exceptional class.

For completeness, no conjectural algebraic projector onto $F$ is used.
Multiplication maps on $A$ give algebraic cohomological K\"unneth
projectors.  Let $L$ be cup product with $\theta$, let
$\mu=\int_A\theta^4/4!$, and use the algebraic lowering operator
\[
 \Lambda(x)=\frac{\theta^3}{6\mu}*x,
\]
where $*$ is Pontryagin product.  On $H^2(A)$ set
\[
 \pi_\theta=\tfrac14L\Lambda,
 \qquad \pi_P=1-\pi_\theta,
 \qquad
 R=\left(\tfrac14\pi_P+\tfrac1{36}\pi_\theta\right)\Lambda^2.
\]
The standard Lefschetz relations give $R=(L^2)^{-1}$ on $H^6(A)$:
$\Lambda^2L^2$ is $4$ on $P$ and $36$ on the polarization line.
Their algebraicity and these constants were proved in Section~\ref{sec:mumford}. Projecting the diagonal and applying $R$
produces an algebraic primitive identity tensor $I_P\in P\otimes P$.
Under $P\otimes P\simeq\operatorname{End}(P)$, subtracting
$\operatorname{Tr}_{\mathbf Q}(a)I_P/27$ is algebraic; the trace is
an intersection pairing with $I_P\theta_1^2\theta_2^2$.

Finally, multiply the tensor representing $a$ by $\theta_1^2$.
It becomes an algebraic degree-zero correspondence on $A$ acting as $a$
on $P$.  Compose it with itself and apply $R$ on the first tensor factor
to obtain an algebraic tensor for $a^2$.  Subtracting the scalar trace
gives
\[
 a^2-\tfrac13\operatorname{Tr}_{F/\mathbf Q}(a^2)\operatorname{id}.
\]
This and $a$ are linearly independent over $\mathbf Q$: otherwise the
nonrational element $a$ of the cubic field $F$ would satisfy a quadratic
polynomial.  They span $F^0$, proving the assertion.  An effective
reducible cycle with a nonzero component has an irreducible component
with that property, which explains the final formulation.
\end{proof}

\subsection{Detection, graphs, and family uniformity}

For an irreducible surface $Z$, the condition in
\eqref{eq:mumford-surface-map} can be tested on a resolution
$\widetilde Z$.  It suffices that, for a nonzero holomorphic two-form
$\omega$ on $S_\lambda$ and some holomorphic one-forms $\eta_1,\eta_2$
on $C$,
\[
 \int_{\widetilde Z}p_S^*\omega\wedge
       \overline{p_1^*\eta_1\wedge p_2^*\eta_2}\ne0.
\]
The form $H^{2,0}(S_\lambda)$ generates its simple rational
transcendental Hodge structure, so this is exactly a nonvanishing test
for the required Hodge component.

\begin{proposition}[The square class of a dominant graph]\label{prop:mumford-graph-twist}
If a dominant rational map $f:C\times C\dashrightarrow S_\lambda$
exists, then $\lambda\in\mathbf Q^\times F^{\times2}$.
\end{proposition}

\begin{proof}
Resolve $f$.  Blowups add only divisor classes.  Pullback on the
transcendental part is therefore a nonzero map into the unique copy of
$T$ in $H^1(C)\otimes H^1(C)$, hence is multiplication by some
$c\in F^\times$.  The projection formula gives
\[
 q_0(cx,cy)=\deg(f)\,q_\lambda(x,y).
\]
Self-adjointness of $F$ and nondegeneracy yield
$c^2=\deg(f)\lambda$, which proves the square-class assertion.
\end{proof}

This proposition does not exclude a general surface correspondence.
Such a correspondence has two projections and need not induce a
similarity of intersection forms.  Nor does it exclude an untwisted
bridge followed by an independently algebraic nonrational RM operator.

If a detecting correspondence as above is constructed over a nonempty
open subset of a fixed algebraic family, all operations in the theorem
are relative algebraic operations after a finite base change.  Their
cohomological coefficients are locally constant, and their cycle
degrees are bounded after spreading out over that open subset.
The resulting two exceptional classes therefore give a familywise
bound for each fixed global class in the exceptional plane; proper
relative Chow parameter spaces extend such a bound by specialization.
Alternatively, a uniform bound on the degrees of exceptional detectors on $A^2$ at an infinite set of points of the Mumford base curve suffices by Corollary~\ref{cor:mumford-family}. A nonzero bridge at a special fibre is not automatically such a detector: its larger Hodge endomorphism algebra need not identify its action with an element of $F$.
An isolated Hodge-generic point with an algebraic correspondence is
only a pointwise input and is not, by itself, a family-uniformity claim.

The theorem leaves precisely the existence of the first detecting
surface unresolved.  The conditional Kuga--Satake route discussed in \cite[Section~5.2, pp.~24--25]{BHPS} likewise requires an algebraic correspondence; an isogeny of the associated abelian varieties does not establish algebraicity of the Kuga--Satake embedding.  The explicit curves
of \cite{BHPS} make the correspondence problem concrete; neither
the preceding argument nor the cited conditional route constructs the
required surface.

\section{A rational-base surface correspondence with zero transcendental action}
\label{sec:mumford-fiber-product}

This section tests a bounded surface correspondence coming directly
from the trigonal equation of a genus-four Mumford family.
It produces a smooth K3 surface and a correspondence of degrees
two and nine, but its transcendental action vanishes.
The general reason is factorization through a rational surface.

\begin{lemma}\label{lem:rational-base-fiber-product}
Let $S$ and $X$ be smooth projective complex surfaces, let $R$ be
a smooth projective rational surface, and let $p:S\dashrightarrow R$
be dominant.  Suppose $f:X\to R$ is finite and flat.
Resolve $p$ to $q:S'\to R$.  The full fibre-product cycle
$S'\times_R X$, pushed to $S\times X$, acts as zero on
$T(S)=\operatorname{NS}(S)_{\mathbf Q}^{\perp}$.
If $p$ has degree two and $f$ has odd degree, that fibre product
has just one dominant irreducible component and no additional
surface components over a proper subset of $S'$.
\end{lemma}

\begin{proof}
Flatness gives the expected graph-composition cycle, whose action is
$f^*q_*:H^2(S',\mathbf Q)\to H^2(X,\mathbf Q)$.
Every class of $H^2(R,\mathbf Q)$ is a divisor class.
For $\alpha\in T(S')$ and $\eta\in H^2(R,\mathbf Q)$,
\[
 \int_R(q_*\alpha)\eta=\int_{S'}\alpha q^*\eta=0.
\]
Nondegeneracy of the intersection form gives $q_*\alpha=0$.
Point blowups add only divisor classes, and therefore identify
$T(S)$ with $T(S')$.  This proves the first assertion.

A quadratic function-field extension remains a field after an
odd-degree extension: otherwise the latter would contain a
quadratic subextension, contrary to the tower formula.
This proves generic irreducibility.  The base change to $S'$ is
finite flat, so every top-dimensional component dominates $S'$;
flatness excludes components with vertical support.  This proves
the remaining assertions.
\end{proof}

The lemma concerns the full fibre-product cycle.  If an even-degree
base change splits a double cover, an individual component could
behave differently; finding that component requires a rational
lift to $S$.  The odd-degree construction below offers no such
choice.

\subsection{An explicit trigonal curve and elliptic K3 surface}

Take the family in
\cite[Theorem~1.2(ii), equation~(1.2)]{BHPS}.
On the quadric cone $XZ=Y^2$, use the chart
$Z=1$, $X=x^2$, $Y=x$, $W=w$.  Its equation becomes
\[
 C_t:\quad w^3+A_t(x)w+B_t(x)=0,
\]
where
\begin{align*}
 a_t(x)&=5x^4+6x^3+2tx+3t,\\
 b_t(x)&=(9-2t)x^6+22tx^5+21tx^4+(-14t^2+18t)x^3\\
       &\hspace{2em}+t^2x^2+6t^2x+3t^2(2-t),\\
 A_t(x)&=t(t-1)a_t(x),\qquad
 B_t(x)=\frac{t(t-1)}3b_t(x).
\end{align*}
Put $\Delta_t=-4A_t^3-27B_t^2$.  Direct expansion gives
\begin{equation}
 \Delta_t(0)=-27t^8(t-1)^2,
 \qquad [x^{12}]\Delta_t=-t^2(t-1)^2(512t^2-608t+243).
 \label{eq:bhps-discriminant-endpoints}
\end{equation}
There is a nonempty open set of parameters for which $t\ne0,1$,
both quantities are nonzero, and $\Delta_t$ has twelve simple roots.
Here is an exact certificate of nonemptiness.  At $t=2$ write
$\Delta_2=4d$.  Reduction modulo seven gives
\begin{align*}
 d={}&3x^{12}+x^{11}+x^{10}+2x^7+4x^6+4x^5+2x+1
                    \quad\text{in }\mathbf F_7[x].
\end{align*}
For
\begin{align*}
 u&=-2x^{10}-x^9-x^8-2x^7-2x^6-2x^2+3,\\
 v&=-x^{11}+2x^{10}+x^9-2x^8+x^5+3x^4+2x^3+x^2-3x-1,
\end{align*}
the identity $ud+vd'=1$ holds in $\mathbf F_7[x]$.
The leading coefficient is nonzero, so $\Delta_2$ is squarefree
of degree twelve in characteristic zero.  The
\emph{trigonal fibre-product check} in the
\href{\workbenchurl}{workbench} verifies this polynomial identity
and \eqref{eq:bhps-discriminant-endpoints} by exact arithmetic.

For a parameter in this open set, the projective curve $C_t$ is
smooth.  A simple discriminant zero gives one simple ramification
point and no singularity; a triple root would instead give a
multiple discriminant zero.  At infinity set $z=1/x$ and $w'=z^2w$.
The equation becomes
\[
 w'^3+z^4A_t(1/z)w'+z^6B_t(1/z)=0,
\]
whose discriminant is a unit at $z=0$ by
\eqref{eq:bhps-discriminant-endpoints}.
The vertex of the quadric cone is not on the cubic.
These charts thus cover the curve and show that the ruling
$x:C_t\to\mathbf P^1$ is unramified over zero and infinity.
The Koszul resolution of the complete intersection of degrees two
and three in $\mathbf P^3$ gives $H^0(\mathcal O_{C_t})=\mathbf C$.
It is therefore connected and, being smooth, irreducible.
Adjunction gives canonical degree six, so it has genus four.
The ruling is finite flat of degree three.

Consider the smooth relatively minimal elliptic surface with section
defined by the Weierstrass equation
\begin{equation}
 S_t:\quad Y^2=X^3+x^2A_t(x)X+x^3B_t(x).
 \label{eq:bhps-elliptic-k3}
\end{equation}
The coefficients are sections of $\mathcal O(8)$ and $\mathcal O(12)$,
so its fundamental line bundle is $\mathcal O(2)$.
At the twelve roots of $\Delta_t$ the elliptic discriminant has
simple zeros, giving twelve fibres of type $I_1$.
At zero the coefficient valuations are at least two and three,
the first is exactly two, and the discriminant valuation is six.
This is a minimal $I_0^*$ fibre.
At infinity use $X'=z^4X$, $Y'=z^6Y$.
The transformed coefficients are $z^6A_t(1/z)$ and $z^9B_t(1/z)$,
again of valuations two and at least three, with discriminant
valuation six.  This gives a second minimal $I_0^*$ fibre.
The minimality and fibre criteria are those of
\cite[\S4.5, Table~1]{SchuttShioda}.
Thus the fibre configuration is $2I_0^*+12I_1$.

The canonical bundle formula
\cite[Theorem~6.8]{SchuttShioda} gives
$\omega_{S_t}=\pi^*(\mathcal O(-2)\otimes\mathcal O(2))
 =\mathcal O_{S_t}$.
Relative duality gives $R^1\pi_*\mathcal O_{S_t}=\mathcal O(-2)$,
and Leray gives $H^1(S_t,\mathcal O)=0$.
Thus $S_t$ is K3.  Its fibre Euler numbers add to $24$, also
consistent with \cite[Corollary~6.11]{SchuttShioda}.
The section, fibre and the two root configurations give
\[
 U\oplus D_4(-1)\oplus D_4(-1)
                       \subset\operatorname{NS}(S_t).
\]
This shows Picard rank at least ten.  No identification with the
Picard-thirteen real-multiplication surface of
Theorem~\ref{thm:mumford-surface} is asserted.

\subsection{The irreducible correspondence of degrees two and nine}

On $x\ne0$, put $w=X/x$ and $y=Y/x$ in
\eqref{eq:bhps-elliptic-k3}.  Then
\[
 y^2=x\bigl(w^3+A_t(x)w+B_t(x)\bigr).
\]
The coordinates $(x,w)$ give a rational double cover
$S_t\dashrightarrow R=\mathbf P^1_x\times\mathbf P^1_w$.
Resolve it to $q:S'_t\to R$.
Set $X_t=C_t\times C_t$ and define
\[
 f:X_t\longrightarrow R,\qquad
       (p_1,p_2)\longmapsto(x(p_1),x(p_2)).
\]
Here the second trigonal coordinate is used as the $w$ coordinate
of $R$.  This is finite flat of degree nine.
The fibre product has, over $X_t$, the generic equation
\[
 y^2=x_1\bigl(x_2^3+A_t(x_1)x_2+B_t(x_1)\bigr).
\]
At a divisor $\{P\}\times C_t$ with $x(P)=0$, the first factor
has valuation one because the ruling is unramified there.
The second factor restricts to a nonzero polynomial in the
nonconstant function $x_2$, so the radicand has odd valuation.
It is not a square.  This also checks generic irreducibility
directly, in addition to the odd-degree argument of the lemma.

Finite flatness over $S'_t$ excludes vertical surface components.
Normalize and resolve the unique surface component, then push its
fundamental class to $S_t\times X_t$.  The degrees of its two
projections are respectively nine and two.
All equations have fixed degrees on the open parameter set,
so these correspondences form a bounded algebraic family.
Nevertheless Lemma~\ref{lem:rational-base-fiber-product} gives
\[
 [Z_t]_*|_{T(S_t)}=0.
\]
In particular its component from $T(S_t)$ to
$H^1(C_t)\otimes H^1(C_t)$ is zero.
The transpose also has zero projection to $T(S_t)$ by adjointness.
Normalization, resolution and the blowups of $S_t$ do not change
this conclusion: their exceptional contributions are divisor
classes, and the fundamental fibre-product correspondence still
acts as $f^*q_*$.
This rules out the indicated correspondence even if a later
identification of the K3 surface's transcendental structure were
available.  It leaves arbitrary surface correspondences and the
coniveau condition of Theorem~\ref{thm:mumford-coniveau} open.

\section{A split sixfold and the coniveau question for a Mumford square}
\label{sec:mumford-coniveau}

This section gives a second sufficient geometric condition for the
exceptional classes of a Mumford square, using a polarized abelian
sixfold and the first cohomology of a resolved divisor.
It makes precise the distinction between an algebraic Weil class and
the support assertion in the generalized Hodge conjecture.

Let $F$ be a totally real cubic field, let $D/F$ be a quaternion algebra
split at exactly one real place, and assume
\begin{equation}
 E=F(\sqrt{-3})\text{ splits }D,
 \qquad \operatorname{Cor}_{F/\mathbf Q}(D)\text{ is split}.
 \label{eq:mumford-sixfold-arithmetic}
\end{equation}
These conditions hold when $D$ is unramified at every finite place:
the quadratic extension is complex at the two ramified infinite places,
and the two local corestriction invariants sum to zero in
$\mathbf Q/\mathbf Z$.  In particular they include the quaternion
data of the two curve families in \cite[Theorem~1.2]{BHPS}.
Write $K=\mathbf Q(\sqrt{-3})$.  Let $A$ be the Mumford fourfold
with $V=H^1(A,\mathbf Q)$ and $V_{\mathbf C}=(2,2,2)$.

For a smooth projective variety $X$, let $N^1H^j(X,\mathbf Q)$ be the
sum of the kernels of restriction to $X\setminus Z$, where $Z$ ranges
over divisors.  This is geometric coniveau one; the assertion that a
Hodge substructure has no $(j,0)$ or $(0,j)$ part does not by itself
put it in this subspace.

\begin{theorem}[A coniveau-one criterion]\label{thm:mumford-coniveau}
Under \eqref{eq:mumford-sixfold-arithmetic}, there is a polarized split
Weil sixfold $B$ for $K$, varying over a common finite cover of the
Mumford Shimura curve, and an algebraic cohomological projector with
image
\[
 P\subset H^3(B,\mathbf Q),\qquad \dim_{\mathbf Q}P=16,
 \qquad P(1)\simeq V^{\oplus2}.
\]
At a Hodge-generic point of the common $(A,B)$ family, the condition
\begin{equation}
 P\cap N^1H^3(B,\mathbf Q)\ne0
 \label{eq:mumford-coniveau-input}
\end{equation}
implies the rational Hodge conjecture for every power of $A$.
In particular, the whole exceptional plane in $H^4(A^2,\mathbf Q)$ is
algebraic.  Equivalently, it suffices to have one nonzero algebraic
correspondence $H^1(A)\longrightarrow P(1)$.
The theorem does not assert \eqref{eq:mumford-coniveau-input}.
\end{theorem}

The proof first constructs $B$ and its projector using endomorphisms.
A supporting divisor would then give an Albanese factor isogenous to
$A$.  Divisor invariants for the larger group of $E$-linear unitary
transformations on $H^1(B)$ transfer through this factor to all
Mumford invariants.  We also compute the resulting weighted Casimir
on the exceptional plane explicitly.
Bhattacharjee--Bhattacharya \cite[Theorems~37--38]{BBK3Sixfolds}
give conditional transfer criteria using an associated sixfold and
an adjoint K3 structure.  Here the sixfold has the balanced compact
signatures $(2,0)$ and $(0,2)$; the required input is precisely
geometric coniveau of its rank-sixteen degree-three summand.

\subsection{The Hermitian space and its determinant}

Since $E$ splits $D$, write $D=(-3,b)_F$, with $b\in F^\times$;
the quadratic splitting criterion is
\cite[Lemma~5.4.7, pp.~73--74]{VoightQuaternion}.
Label the real embeddings so that the signs of $b$ are $(+,-,-)$.
Choose $\delta\in F^\times$ of signs $(+,-,+)$ and set
\[
 h_E=\delta\begin{pmatrix}b&0\\0&-1\end{pmatrix},
 \qquad h_K=\operatorname{Tr}_{E/K}h_E.
\]
The three $E/F$ signatures of $h_E$ are $(1,1),(2,0),(0,2)$;
its special unitary group is $\operatorname{SL}_1(D)$.
The $K$-Hermitian space $(E^2,h_K)$ has dimension six and signature
$(3,3)$.

Let $d_F$ be the discriminant of a rational basis of $F$.  Taking the
determinants of the two trace-form blocks gives
\begin{equation}
 \det_K h_K
 =-\operatorname{Nm}_{F/\mathbf Q}(b)
       \bigl(\operatorname{Nm}_{F/\mathbf Q}(\delta)d_F\bigr)^2.
 \label{eq:mumford-sixfold-determinant}
\end{equation}
The projection formula for quaternion symbols
\cite[II, Proposition~1.39(d), p.~74]{MilneCFT} gives
\[
 \operatorname{Cor}_{F/\mathbf Q}(-3,b)_F
   =(-3,\operatorname{Nm}_{F/\mathbf Q}(b))_{\mathbf Q}.
\]
Thus the first norm in \eqref{eq:mumford-sixfold-determinant} is a
$K/\mathbf Q$ norm, by the split-corestriction hypothesis and
\cite[Main Theorem~5.4.4(vi), pp.~72--73]{VoightQuaternion}.  The
determinant class is consequently $-1$, the class of three hyperbolic
planes.  Landherr's classification by dimension, determinant and
signatures \cite[pp.~246--247]{Landherr} identifies $h_K$ with that
hyperbolic space.

On the noncompact real factor choose the varying Hodge line of the
Mumford disk.  On the two compact factors choose respectively all of
the indicated $K$ eigenspace and none of it as the $(1,0)$ part.
Tracing $h_E$ against a purely imaginary element of $K$, with the
corresponding sign, gives a positive alternating Riemann form.
An $E$-stable lattice and an integral multiple of this form give the
polarized abelian sixfold $B$.  Its $K$ action has signature $(3,3)$,
and the preceding determinant calculation proves that it is split.
This construction does not require a principal polarization.
Fixed lattices and sufficiently fine level realize the two rational
representations as abelian schemes over a common finite cover of the
Mumford curve; the compact central Hodge data are fixed throughout.

\subsection{An algebraic projector onto the rank-sixteen part}

Put $H=H^1(B,\mathbf Q)$.  For one embedding $\sigma:K\to\mathbf C$,
write
\[
 H_\sigma=M_1\oplus M_2\oplus M_3,\qquad \dim M_i=2.
\]
The three Hodge-circle weight lists are $(1,-1),(1,1),(-1,-1)$.
The multidegree $(1,1,1)$ part of $\Lambda_K^3H$, together with its
conjugate, is
\begin{equation}
 P_{\mathbf C}=(M_1\otimes M_2\otimes M_3)
       \oplus(\overline M_1\otimes\overline M_2\otimes\overline M_3).
 \label{eq:mumford-rank16}
\end{equation}
It is Galois stable.  Its Hodge types are $(2,1)$ and $(1,2)$, each
with dimension eight.  The connected central torus acts trivially on
it: the Galois conjugates of the central cocharacter $(0,1,-1)$ span
the trace-zero plane, whereas the tensor in each summand has the
product character.  The derived representation is $(2,2,2)$.
The split corestriction gives its rational eight-dimensional form
$V$, with endomorphism field $\mathbf Q$.  The Hodge circle just
computed therefore identifies $P(1)$ with two copies of $V$.
These are statements at a Hodge-generic point of the common family;
they do not use the enlarged endomorphism algebra of a CM fibre.

We now give a projector without assuming that a general Hodge
projector is algebraic.  Algebraic K\"unneth projectors on an abelian
variety are polynomials in a multiplication map.  On $H^3(B)$ first use
\[
 e_K=\frac{1+[\omega]^*+[\omega^2]^*}{3}.
\]
The eigenvalue of $[\omega]^*$ on a wedge with $r$ factors from
$H_\sigma$ is $\omega^{2r-3}$.  Consequently $e_K$ selects exactly
$\Lambda^3H_\sigma\oplus\Lambda^3H_{\bar\sigma}$.

Choose a primitive element $a\in F$ with conjugates $a_1,a_2,a_3$.
Let $D_a$ apply $a$ in one of the three exterior factors and the
identity in the others, summed over the three choices.  It is an
algebraic cohomological operator, since
\begin{equation}
 D_a=\frac{[a+2]^*-2[a+1]^*+[a]^*}{2}-3\operatorname{id}
       \quad\hbox{on }H^3(B).
 \label{eq:mumford-derivation}
\end{equation}
Indeed, the second finite difference of $\Lambda^3(a+t)$ is
$6\operatorname{id}+2D_a$.  A rational endomorphism is interpreted
using a rational graph correspondence, or after replacing the lattice
by an isogenous one on which it is integral.

On \eqref{eq:mumford-rank16} the eigenvalue is
$c=\operatorname{Tr}_{F/\mathbf Q}(a)$.
All the other multidegrees are $(2,1,0)$, with eigenvalues
$2a_i+a_j$, $i\ne j$.  Set
\[
 Q_a(T)=\prod_{i\ne j}(T-2a_i-a_j)\in\mathbf Q[T].
\]
Then
\[
 Q_a(c)=-\prod_{i<j}(a_i-a_j)^2\ne0,
 \qquad e_P=e_K\frac{Q_a(D_a)}{Q_a(c)}
\]
is the required algebraic projector.  Repeated unwanted eigenvalues
cause no problem.  This also constructs the projector relatively,
using fixed endomorphisms and rational coefficients.

\subsection{A divisor support gives the weight-one correspondence}

We explain precisely the support assertion used in
\eqref{eq:mumford-coniveau-input}.  The coniveau-one part of $H^3(B)$
is the sum of Gysin images
\begin{equation}
 H^1(\widetilde D,\mathbf Q)(-1)\longrightarrow H^3(B,\mathbf Q)
 \label{eq:mumford-divisor-gysin}
\end{equation}
from resolutions of divisors.  Equivalently it is the image of a
correspondence from the first cohomology of a smooth projective curve
\cite[Proposition~1.3, p.~8]{ACMV}.  For degree three this can also be seen
directly: remove the singular loci and intersections of the supporting
divisors.  They have codimension at least two and contribute no
cohomology with supports in degree three.  Purity for the remaining
smooth divisors gives their $H^1(-1)$.  Only its weight-one part can
map to the pure group $H^3(B)$, and this part comes from a smooth
projective resolution.  A finite sum suffices for any given class.

Apply $e_P$ to these maps.  Under
\eqref{eq:mumford-coniveau-input} at least one resulting map is
nonzero.  Polarizable rational Hodge structures are semisimple,
and $P(1)=V^{\oplus2}$ with $V$ simple.  Therefore the Albanese
variety of a product of the resolutions has an isogeny factor
isogenous to $A$.  Hodge morphisms between the first cohomology of
abelian varieties are induced by homomorphisms up to isogeny.
Composing such a map with \eqref{eq:mumford-divisor-gysin} and $e_P$
produces a nonzero algebraic correspondence
\begin{equation}
 \Gamma:V\longrightarrow P(1).
 \label{eq:mumford-coniveau-bridge}
\end{equation}

In the convention of a map from $A$ to $B$, this correspondence has
codimension five in $A\times B$: its degree shift is $3-1=2$.
Conversely, a codimension-five cycle on this ten-dimensional product
has dimension five.  Its image in the sixfold $B$ is contained in a
divisor, so its cohomological image is in $N^1H^3(B)$.
Degree-zero algebraic correspondences preserve coniveau, by moving
intersections properly and the support dimension inequality; thus
applying $e_P$ preserves this conclusion.

For comparison with a Hodge tensor, the relevant component of the
codimension-five cycle is in $H^7(A)\otimes H^3(B)$.
The inverse of $L_A^3:H^1(A)\to H^7(A)$ is algebraic.  In the notation
of Section~\ref{sec:mumford}, it is $\Lambda_A^3/36$, since
$\Lambda_A^3L_A^3=36$ on $H^1(A)$.
Applying it in the first factor gives a codimension-two cycle with
component in $H^1(A)\otimes H^3(B)$.  Multiplication by $\theta_A^3$
reverses this operation.  Neither description asserts that the bridge
is a divisor on $A\times B$.

\subsection{Divisor invariants transfer to every power}

Let $U$ be the rational algebraic group of $E$-linear automorphisms of
$H^1(B)$ preserving its alternating polarization.  It contains the
Hodge circle of $B$, and
\[
 U_{\mathbf C}\simeq\mathrm{GL}_2^3,
 \qquad H_{\mathbf C}=\bigoplus_{i=1}^3(M_i\oplus M_i^\vee).
\]
The indicated dualities include the polarization Tate twist.
The invariants of $U$ in the cohomology of powers of $B$ are generated
by divisors.  Indeed the first fundamental theorem for the general
linear group expresses every invariant tensor as a linear combination
of contractions between $M_i$ and $M_i^\vee$.  These contractions are
the polarization pairings with the $E$-action inserted.  They are
divisor classes on $B^2$; permutations, exterior products and diagonal
pullbacks give their images in all exterior powers.  The rational
descent and the divisor formulation are
\cite[Proposition~3.6(c) and Remark~3.7, pp.~17--18]{MilneLefschetz}.
This assertion concerns the larger group $U$.  It does not say that
all Hodge classes on $B$ are generated by divisors.

The space $P$ is primitive in $H^3(B)$: the lowering operator pairs
opposite $K$ eigenspaces, whereas each summand in
\eqref{eq:mumford-rank16} uses just one eigenspace.
Consequently the polarization adjoint of $\Gamma$ is algebraic.
More explicitly, on the relevant cohomological degrees it is
\[
 \beta=(L_A^3)^{-1}\,\Gamma^t\,L_B^3 e_P:P(1)\longrightarrow V,
\]
up to the fixed sign in the polarization convention.  Include the
algebraic degree-one projector on $A$.  The resulting correspondence
lies in $\mathrm{CH}^5(B\times A)_{\mathbf Q}$.
It is nonzero, since it is the adjoint of a nonzero map between
polarized Hodge structures.

Put $P_+=M_1\otimes M_2\otimes M_3$ and
$P_-=M_1^\vee\otimes M_2^\vee\otimes M_3^\vee$ over $\mathbf C$.
For $G_{\mathbf C}=\mathrm{SL}_2^3$, both are irreducible of the
same type as $V_{\mathbf C}$.  Each restriction $\beta|_{P_\pm}$
is an isomorphism of $G_{\mathbf C}$-modules: it is either zero or
an isomorphism, and rationality interchanges the two restrictions.
They cannot both be zero.

For an integer $m\geq1$ consider
\[
 I_m=(P^{\otimes 2m})^U.
\]
After complexification its summand with $m$ factors $P_+$ and $m$
factors $P_-$ contains all the $G_{\mathbf C}$-invariants of that
balanced tensor product.  To check this, regroup its factors as
\[
 \bigotimes_{i=1}^3
       \bigl(M_i^{\otimes m}\otimes(M_i^\vee)^{\otimes m}\bigr).
\]
The scalar centre of each $\mathrm{GL}_2$ acts trivially here.
Thus its $\mathrm{GL}_2$-invariants and
$\mathrm{SL}_2$-invariants coincide.
The isomorphisms $\beta|_{P_\pm}$ now show that
\begin{equation}
 \beta^{\otimes2m}:I_m\longrightarrow(V^{\otimes2m})^G
       \quad\hbox{is surjective}.
 \label{eq:mumford-balanced-invariants}
\end{equation}
Surjectivity over $\mathbf C$ implies surjectivity over $\mathbf Q$.

Every element of $I_m$ is a divisor-product class of codimension $3m$
on $B^{2m}$.  Applying $\beta^{\boxtimes2m}$ produces a cycle of
codimension $m$ on $A^{2m}$: the correspondence has codimension $10m$
and the source has dimension $12m$.  This proves algebraicity of
every tensor in the target of \eqref{eq:mumford-balanced-invariants}.
Finally, $H^*(A^r)=\Lambda^*(V^{\oplus r})$.  Exterior multiplication
is equivariant, and invariants are exact for the reductive group $G$.
Every invariant in degree $2m$ is therefore a sum of images of
invariant $2m$-fold tensors.  Assigning each tensor factor to one of
the $r$ coordinates is the pullback along a map
$A^r\to A^{2m}$ whose coordinates repeat the corresponding
projections.  Such pullbacks preserve algebraicity.
This proves the rational Hodge conjecture on every $A^r$, and hence
Theorem~\ref{thm:mumford-coniveau}.

\subsection{Transporting the Lie action and its field weights}

We give a direct description in degree four of the exceptional classes
obtained from the same bridge.

Inside $\operatorname{End}_{\mathbf Q}(H)$ consider
\[
 \mathfrak g=\{X\in\operatorname{End}_E(H):
                    X^*=-X,\ \operatorname{Tr}_E X=0\},
\]
where $*$ is the polarization adjoint.  It is the rank-nine
representation $\operatorname{Res}_{F/\mathbf Q}\mathfrak{sl}_1(D)$.
Its projector is algebraic for an explicit reason.  The separability
idempotent of $E/\mathbf Q$, expressed in trace-dual bases, projects
to $\operatorname{End}_E(H)$.  The operation $(1-*)/2$ takes its
skew-adjoint part.  Subtracting the scalar skew-adjoint $E$ part by
the rational trace pairing takes the trace-zero part.
All $E$ actions are endomorphisms of $B$; the adjoint, pairings and
their inverses use the polarization and the algebraic Lefschetz
operators.  Thus these operations are algebraic on powers of $B$.
Multiplication by any $a\in F$ on $\mathfrak g$ is algebraic too.

The derivation action on $\Lambda^3H$, followed by $e_P$, gives an
algebraic map $\mathfrak g\to\operatorname{End}(P(1))$.
At the common Hodge-generic point, \eqref{eq:mumford-coniveau-bridge}
intertwines this action with the Mumford action.  Its polarization
adjoint is algebraic and
\[
 \Gamma^\dagger\Gamma=r\operatorname{id}_V,
 \qquad 0<r\in\mathbf Q,
\]
because $\operatorname{End}_{\mathrm{HS}}(V)=\mathbf Q$.
Consequently
\[
 X\longmapsto r^{-1}\Gamma^\dagger\rho_P(X)\Gamma
\]
is the usual $\mathfrak g$ action on $V$, expressed by algebraic
correspondences.

Fix the trace form $B_{\mathfrak g}(X,Y)=\operatorname{Tr}_{\mathbf Q}(XY)$
in the $E^2$ representation.  It is nondegenerate and has the same
rational index on the three complex $\mathfrak{sl}_2$ factors.
Its inverse tensor is algebraic, using the preceding projector and
the ambient trace coevaluation.  After fixing this inverse, apply
multiplication by a nonrational $a\in F$ in one factor and transport
both factors to $\operatorname{End}(V)$.  The resulting tensor is,
up to a common nonzero rational constant, the weighted Casimir
$\sum_i a_i C_i$.  The separate multiplication by $a$ is essential:
it is not a change of metric followed by an inverse that cancels it.

Regroup its four $V$ factors into two exterior pairs.  The calculation
\eqref{eq:mumford-casimir-eigen} gives the three eigenvalues
\[
 \frac14\operatorname{Tr}_{F/\mathbf Q}(a)-a_j.
\]
For clarity, the cross Casimir is $1/4$ on the symmetric square of
the two-dimensional standard representation and $-3/4$ on its
exterior square.  In the $j$th primitive summand of $\Lambda^2V$,
the exterior choice occurs at $j$ and the symmetric choices at the
other two indices.  This gives exactly the displayed expression.
Algebraic primitive projection and trace subtraction leave
\[
 -a+\frac13\operatorname{Tr}_{F/\mathbf Q}(a)\in F^0\setminus\{0\}.
\]
The cubic-field composition argument of
Theorem~\ref{thm:mumford-cubic} now gives the entire exceptional plane.
This describes the exceptional plane explicitly.

\subsection{What the algebraic Weil classes do give}

The ordinary Weil classes on $B$ are algebraic by the split-sixfold
results, with the input status in Section~\ref{sec:effective-sixfolds}.
They provide the following intermediate conclusion, which does not
settle the support condition above.

\begin{proposition}\label{prop:weil-support-middle}
If a nonzero rational pure Weil class $w\in H^6(B,\mathbf Q)$ is
represented by a rational combination of integral threefolds $Y_i$,
then $P$ is an algebraic direct summand of
$\bigoplus_i H^3(\widetilde Y_i,\mathbf Q)$, where the tildes are
smooth projective resolutions.
\end{proposition}

\begin{proof}
For $i:\widetilde Y_i\to B$, the cup action of its cycle class is
$i_*i^*:H^3(B)\to H^9(B)$.
The inverse of $L_B^3:H^3(B)\to H^9(B)$ is algebraic, by the
Lefschetz lowering operator and its finite primitive decomposition.
Consider on $P$ the operator
\[
 T=e_P(L_B^3)^{-1}(w\smile -)e_P.
\]
It factors algebraically through the displayed sum of middle
cohomology groups.  We show it is invertible.

Over $\mathbf C$ write $H=H_\sigma\oplus H_{\bar\sigma}$.
The polarization pairs these spaces perfectly and respects their
three $F$ blocks.  The two components of $w$ are nonzero volume
forms on these six-dimensional spaces: rationality and $w\ne0$
make both coefficients nonzero.  For
$u\in M_1\otimes M_2\otimes M_3\subset\Lambda^3H_\sigma$,
only the conjugate volume contributes to $w\smile u$.
The restriction of $L_B^3$ from the conjugate triple wedge to this
bidegree is an isomorphism, given by complementary exterior-power
pairing.  In multidegree $(1,1,1)$ it uses one Lefschetz pair from
each $F$ block.  Thus $T$ exchanges the two summands of
\eqref{eq:mumford-rank16} by invertible maps.  Here the two volume
forms refer to $K$ eigenspaces; their Hodge type is $(3,3)$.

The inverse of $T$ on $P$ is a polynomial in $T$ with rational
coefficients, by its characteristic polynomial and nonzero
determinant.  Composing this inverse with the factorization gives
an algebraic splitting, as asserted.
\end{proof}

The proposition reaches $H^3$ of threefolds.  A coniveau-one class
on one such $\widetilde Y_i$ whose image under
$e_P(L_B^3)^{-1}i_*$ is nonzero would suffice: its Gysin
pushforward raises coniveau by three, and the inverse $L_B^3$
lowers it by at most three, leaving the additional step of
coniveau one on $B$.  Ordinary algebraicity of the degree-six
Weil class does not prove this generalized Hodge assertion.

There is also a direct limitation on translations.  At the generic
point, every $H^1$ of a power of $B$ has derived constituents
$(2,1,1),(1,2,1),(1,1,2)$; a CM parameter variety adds only
derived-trivial constituents.  Hence
\[
 \operatorname{Hom}_{\mathrm{HS}}
   (H^1(B^m\times C_{\mathrm{CM}}),P(1))=0.
\]
Any construction using these parameter first-cohomology groups,
addition, translations and Lefschetz correspondences therefore
has zero component in the required space.  This does not exclude
a singular divisor whose resolution has a new Albanese factor
isogenous to $A$; finding such a factor is the unresolved geometry.

Abdulali's type-IV theorem \cite[Theorem~13]{AbdulaliIV} requires
$|p_\alpha-q_\alpha|=1$ at every real place.  Here those differences
are $0,2,2$.  His other criterion requires every real simple factor
to be $\mathrm{SU}(p,1)$ with $p>1$
\cite[Theorems~6.5--6.6]{AbdulaliSurvey}; the compact factors here
do not satisfy it.  Appendix~B.5 of that survey adds the former
signature hypothesis to the earlier type-IV statements.
Its Section~8.2 records Schoen's generalized Hodge result for the
rank-two determinant structure of a signature-$(5,1)$ Eisenstein
sixfold.  That is a CM structure in degree six, not the rank-sixteen
part of degree three used here.  None of these cited statements
establishes \eqref{eq:mumford-coniveau-input}.

Finally, the family conclusion needs a family input.  A relative
nonzero correspondence \eqref{eq:mumford-coniveau-bridge} over a
nonempty open subset gives a uniform cycle bound after spreading
out, and proper bounded-cycle spaces extend it to the whole base.
The same follows if \eqref{eq:mumford-coniveau-input} holds at every
Hodge-generic point: countably many closed bounded-class loci cover
those points, so one dominates the complex curve.  Uniformly
bounded bridges on a dense set also suffice by first taking a
dominating Chow component and then arguing generically.
An isolated generic bridge is a pointwise input.  An isolated CM
bridge does not justify the generic endomorphism and Lie-action
identifications used in the theorem.


\section{What remains open}
\label{sec:open}
Theorem~\ref{thm:global} is a criterion for the Hodge conjecture. It proves neither the conjecture nor a reduction to a finite computation. The outstanding condition is a degree bound uniform in the CM fibre.

For Weil classes the principal unresolved question is a uniform bound for detectors in every connected split family, in every dimension. Theorem~\ref{thm:principal-quotient} reduces the integral lattice and polarization choices to the single principal hyperbolic family $U^n$ for each imaginary quadratic field $K$ and dimension $2n$, with an explicit loss in degree. Stabilization would then give algebraicity of Weil classes for every imaginary quadratic field and every determinant. It would not establish the Hodge conjecture for all abelian varieties. The first target beyond the known constructions is a split Weil eightfold family over a field other than the Eisenstein field. Theorems asserting only the existence of cycles at each CM point do not provide the required bound.

For the pure-character semiregularity route, Theorem~\ref{thm:pure-weil-ext} reduces the strong condition to one exact self-Ext dimension in every dimension. The first explicit Gaussian eightfold character therefore leaves the existence of a gluable representative with invariant second self-Ext dimension $56$. The two-level line-bundle extensions have dimension $448$, and changing their gluing maps cannot help. In the finite-cover route, a dominant Gaussian source must have coefficient rank at least twenty-six. Neither necessary condition constructs a cycle.

Adding a polarization polynomial changes the sufficient Ext dimension
to $104$ in the explicit eightfold construction of
Theorem~\ref{thm:simple-hodge-bundle}.
The simple bundle exists, but the tree-supported gluing has a
nonzero first-order obstruction by
Theorem~\ref{thm:tree-extension-obstruction}.
Theorem~\ref{thm:gaussian-positive-trace-network} constructs a different
finite Gaussian network whose diagonal trace equations, and the
matrix-valued diagonal cohomology blocks at each vertex, can be
canceled exactly.  The off-vertex products along two-edge paths,
the remaining upper block, and the full strong semiregularity
condition are still unresolved.  No deformation of that bundle
has been deduced.

The complete-intersection formulation in Proposition~\ref{prop:complete-intersections} asks for one integer $m$ such that every CM fibre, or one dense set of fibres, has a detecting component in a proper intersection from $|mL|$. The CM transport calculation isolates a different requirement: one must produce new bounded-degree supports on the targets, or prove that transported classes admit such supports. Keeping levels fixed and rescaling coefficients do not establish either assertion.

For the one-prime formulation, Proposition~\ref{prop:lift} leaves the finite exceptional set $E_D$ depending on the fibre. It would suffice to construct a degree-$D$ detector at one place outside that set for each CM fibre, with the same $D$. A uniform estimate for $E_D$ and a finite-field construction meeting it remain separate missing inputs.


For the exceptional classes of Mumford-fourfold squares, Theorem~\ref{thm:mumford-cubic} replaces a representative of a prescribed class by a single detector, at cubic cost. No bounded family of such detectors has been constructed. In particular, knowing that a rational Hodge correspondence or a Kuga--Satake embedding exists does not establish its algebraicity. For the explicit Jacobian families, Theorem~\ref{thm:mumford-surface} gives a sufficient bridge from a K3 surface with a nonsquare polarization twist to the square of the genus-four curve. Proposition~\ref{prop:mumford-graph-twist} excludes a dominant rational graph for that twist. Under the arithmetic hypotheses of Theorem~\ref{thm:mumford-coniveau}, a different sufficient input is one nonzero coniveau-one part of the explicit rank-sixteen substructure of $H^3$ of a split Weil sixfold. Its known Weil cycles give a summand of the middle cohomology of supporting threefolds, without proving the required divisor support. For uniformity, either bridge must be constructed over an open part of the family, or bounded exceptional detectors must exist on a dense set; on the chosen Shimura curve, infinitely many CM fibres suffice.


\section*{Method, checks and records}
The mathematics and text were produced by GPT-6 Astra (ChatGPT Pro mode). Separate statement-only re-proofs were produced by the same model. Earlier mathematical checks by Claude Opus 5.5 at high effort are distinguished in the records from the checks performed for this article. No human has checked the paper, and these checks are not formal proof-assistant certificates. The original text is dedicated to the public domain under CC0~1.0. The editable sources, exact computations and statement-level check records are available through the \href{\workbenchurl}{algebraic cycles and motives workbench}. The \href{\recordurl}{mathematics workbench catalogue} is the permanent address of the records until an individual paper page is assigned.

\begin{thebibliography}{99}
\bibitem{BS} L.~Br\"unjes and C.~Serp\'e, \emph{Uniform bounds and ultraproducts of cycles}, arXiv:0901.4853 (2009), Definition~2.3(ii), p.~4, and Theorem~3.4, p.~8. \url{https://arxiv.org/abs/0901.4853}.
\bibitem{FGA} A.~Grothendieck, \emph{Techniques de construction et th\'eor\`emes d'existence en g\'eom\'etrie alg\'ebrique IV: les sch\'emas de Hilbert}, S\'eminaire Bourbaki, Expos\'e~221 (1961), in \emph{Fondements de la g\'eom\'etrie alg\'ebrique}, Theorem~3.2, p.~260, and pp.~265--266.
\bibitem{SGA4} M.~Artin, A.~Grothendieck and J.-L.~Verdier (eds.), \emph{Th\'eorie des topos et cohomologie \'etale des sch\'emas}, SGA~4, Lecture Notes in Mathematics 269, 270, 305, Springer (1972--1973): XI, Theorem~4.4; XVI, Corollary~2.2; XVIII, Theorems~2.9 and 3.2.5.
\bibitem{SGA45} P.~Deligne, \emph{Cohomologie \'etale}, SGA~$4\frac12$, Lecture Notes in Mathematics 569, Springer (1977), expos\'e \emph{Cycle}, \S2.3.1--2.3.8; erratum to SGA~4 XVI, Corollary~2.2.
\bibitem{SGA6} P.~Berthelot, A.~Grothendieck and L.~Illusie (eds.), \emph{Th\'eorie des intersections et th\'eor\`eme de Riemann--Roch}, SGA~6, Lecture Notes in Mathematics 225, Springer (1971), Expos\'e~XIII (S.~Kleiman), Definition~6.9 and Corollary~6.11, pp.~658--660.
\bibitem{MilneSV} J.~S.~Milne, \emph{Introduction to Shimura varieties}, in \emph{Harmonic Analysis, the Trace Formula, and Shimura Varieties}, Clay Mathematics Proceedings 4 (2005), 265--378; revised notes (2017), Theorem~5.4 and \S\S5--7. \url{https://www.jmilne.org/math/xnotes/svi.pdf}.
\bibitem{CDK} E.~Cattani, P.~Deligne and A.~Kaplan, \emph{On the locus of Hodge classes}, J. Amer. Math. Soc. 8 (1995), 483--506, Introduction and Theorem~1.1. \url{https://arxiv.org/abs/alg-geom/9402009}.
\bibitem{Voisin07} C.~Voisin, \emph{Hodge loci and absolute Hodge classes}, Compos. Math. 143 (2007), 945--958, Propositions~0.2 and 0.5, Theorem~0.6, pp.~946--947. \url{https://doi.org/10.1112/S0010437X07002837}.
\bibitem{BBDense} D.~Bhattacharjee and U.~Bhattacharya, \emph{Dense Algebraic Loci of Hodge Classes}, Hexagon 2610.00315v2 (2026), Lemmas~4.2--4.3 and Theorem~4.5, pp.~10--12. Unrefereed preprint. \url{https://hexagonmath.org/2610.00315v2}.
\bibitem{BBAlgebraic} D.~Bhattacharjee and U.~Bhattacharya, \emph{Algebraic Loci of Weil Classes on Abelian Varieties}, companion manuscript (2026), Theorems~16.9 and 16.17. Unrefereed. \url{https://doi.org/10.5281/zenodo.23885230}.
\bibitem{OpenAICM} OpenAI, \emph{The rational Hodge conjecture for CM abelian varieties}, OpenAI Math Release, 6 October 2026, 53 pp., Theorem~1.1; Lemma~2.2; Proposition~3.1; Corollary~6.7. Unrefereed preprint. \url{https://github.com/openai/math/tree/main/preprints/The-rational-Hodge-conjecture-for-CM-abelian-varieties-October-6-2026}.
\bibitem{vHS} P.~van Hoften and J.~Sempliner, \emph{The Hodge Conjecture for CM abelian varieties}, Hexagon 2610.00259v2, 15 October 2026, 25 pp., Theorem~I; Lemma~2.3.2; Theorem~2.3.4. Unrefereed preprint. \url{https://hexagonmath.org/2610.00259v2}.
\bibitem{Deligne82} P.~Deligne, notes by J.~S.~Milne, \emph{Hodge cycles on abelian varieties}, in \emph{Hodge Cycles, Motives, and Shimura Varieties}, LNM 900, Springer (1982), 9--100. Revised transcription (2018), Proposition~4.1 and Corollary~4.2, p.~28; Proposition~6.1, pp.~41--42. \url{https://www.jmilne.org/math/Documents/Deligne82.pdf}.
\bibitem{MilneAV} J.~S.~Milne, \emph{Abelian varieties}, in \emph{Arithmetic Geometry}, Springer (1986), 103--150; corrected notes (2022), Remark~7.3, p.~12; Theorem~13.3(b), p.~25; \S17. \url{https://www.jmilne.org/math/xnotes/AVs.pdf}.
\bibitem{Ballmann} W.~Ballmann, \emph{Lectures on K\"ahler Manifolds}, ESI Lectures in Mathematics and Physics, EMS (2006), Theorem~5.42 and equation~(5.43), pp.~70--71. \url{https://ems.press/books/esi/19}.
\bibitem{Schoen98} C.~Schoen, \emph{Addendum to: ``Hodge classes on self-products of a variety with an automorphism''}, Compos. Math. 114 (1998), 329--336, Proposition~10, pp.~332--333. \url{https://doi.org/10.1023/A:1000566205021}.
\bibitem{Schoen07} C.~Schoen, \emph{A family of surfaces constructed from genus 2 curves}, Internat. J. Math. 18 (2007), 585--612. \url{https://doi.org/10.1142/S0129167X07004175}.
\bibitem{vG96} B.~van Geemen, \emph{Theta functions and cycles on some abelian fourfolds}, Math. Z. 221 (1996), 617--631, \S2.1; Lemma~3.3; Proposition~3.6; Theorem~3.7. \url{https://sites.unimi.it/vangeemen/ThetafctsCycles4folds_MathZ1996.pdf}.
\bibitem{vG22} B.~van Geemen, \emph{Weil Classes and Decomposable Abelian Fourfolds}, SIGMA 18 (2022), 097, Introduction and \S\S6.1--6.3. \url{https://doi.org/10.3842/SIGMA.2022.097}.
\bibitem{RRS} C.~Rito, X.~Roulleau and A.~Sarti, \emph{Explicit Schoen surfaces}, Algebraic Geometry 6 (2019), 410--426, Introduction and Theorem~1. \url{https://arxiv.org/abs/1609.02235}.
\bibitem{Beauville} A.~Beauville, \emph{A tale of two surfaces}, in \emph{Higher Dimensional Algebraic Geometry}, Adv. Stud. Pure Math. 74 (2017), 1--10, \S3. \url{https://math.univ-cotedazur.fr/u/beauvill/pubs/Tale2.pdf}.
\bibitem{VoisinB} C.~Voisin, \emph{La conjecture de Hodge pour les vari\'et\'es ab\'eliennes de dimension au plus 5}, S\'eminaire Bourbaki, Expos\'e 1248, January 2026, Lemma~2.9 and Corollary~2.10, p.~1248--16. Official seminar manuscript. \url{https://www.bourbaki.fr/TEXTES/Exp1248-Voisin.pdf}.
\bibitem{MarkmanICM} E.~Markman, \emph{Secant Sheaves and Weil Classes on Abelian Varieties}, Proceedings of the International Congress of Mathematicians 2026, Vol.~3, 586--602, Theorem~1.2 and Corollary~1.3, p.~587. \url{https://doi.org/10.1137/25M1803796}. Detailed proof: \emph{Cycles on abelian $2n$-folds of Weil type from secant sheaves on abelian $n$-folds}, arXiv:2502.03415v2 (2025), Theorem~1.5.1. \url{https://arxiv.org/abs/2502.03415}.
\bibitem{Andre96} Y.~Andr\'e, \emph{Pour une th\'eorie inconditionnelle des motifs}, Publ. Math. IHES 83 (1996), 5--49, Theorems~0.5 and 0.6.2, pp.~8--9; Lemma~6.3.2, p.~32. \url{https://doi.org/10.1007/BF02698643}.
\bibitem{Scanlon} T.~Scanlon, \emph{Automatic uniformity}, Internat. Math. Res. Notices 2004, no.~62, 3317--3326, Definition~2.2 and Theorem~4.2. \url{https://math.berkeley.edu/~scanlon/papers/autouni.pdf}.
\bibitem{ChaiOort} C.-L.~Chai and F.~Oort, \emph{Abelian varieties isogenous to a Jacobian}, Ann. of Math. 176 (2012), 589--635, Definitions~5.3--5.4, Remark~5.8 and Lemma~6.3, pp.~607,610 and 621--622. \url{https://doi.org/10.4007/annals.2012.176.1.11}.
\bibitem{Perry} A.~Perry, \emph{The semiregularity theorem for equivariant noncommutative varieties}, arXiv:2604.00511v2 (24 June 2026), Theorem~1.2; Definition~2.6; Lemma~3.3; \S6.2. Unrefereed preprint. \url{https://arxiv.org/abs/2604.00511v2}.
\bibitem{HT} D.~Huybrechts and R.~Thomas, \emph{Deformation-obstruction theory for complexes via Atiyah and Kodaira--Spencer classes}, Math. Ann. 346 (2010), 545--569; corrected manuscript, arXiv:0805.3527v2 (2013), Corollary~3.4, p.~14. \url{https://arxiv.org/abs/0805.3527}.
\bibitem{BF} R.-O.~Buchweitz and H.~Flenner, \emph{A semiregularity map for modules and applications to deformations}, Compos. Math. 137 (2003), 135--210, pp.~136--137; Theorem~5.1, pp.~174--175, proof p.~179; Proposition~5.9 and Lemma~5.10, pp.~178--179; \S8. \url{https://doi.org/10.1023/A:1023999012081}.
\bibitem{VoisinKahler} C.~Voisin, \emph{A counterexample to the Hodge conjecture extended to Kähler varieties}, Int. Math. Res. Not. 2002, no.~20, 1057--1075. Appendix, Proposition~A.3 and Corollaries~A.4--A.5, pp.~1072--1074. \url{https://webusers.imj-prg.fr/~claire.voisin/Articlesweb/hodgeimrn.pdf}.
\bibitem{BandoSiu} S.~Bando and Y.-T.~Siu, \emph{Stable sheaves and Einstein--Hermitian metrics}, in \emph{Geometry and Analysis on Complex Manifolds}, World Scientific, 1994, pp.~39--50. Theorem~3 and Corollary~3, p.~40. \url{https://people.math.harvard.edu/~siu/bando_joint_paper/index.html}.
\bibitem{Schoen88} C.~Schoen, \emph{Hodge classes on self-products of a variety with an automorphism}, Compos. Math. 65 (1988), 3--32, Theorem~2.0, p.~11; Lemmas~2.2 and 2.6, pp.~13 and 16; \S3, pp.~24--25. \url{https://www.numdam.org/article/CM_1988__65_1_3_0.pdf}.
\bibitem{Koike} K.~Koike, \emph{Algebraicity of some Weil Hodge Classes}, Canad. Math. Bull. 47 (2004), 566--572, Theorem~2.1, p.~568. \url{https://doi.org/10.4153/CMB-2004-055-x}.
\bibitem{EGAIV} A.~Grothendieck and J.~Dieudonn\'e, \emph{\'El\'ements de g\'eom\'etrie alg\'ebrique IV}, Publ. Math. IHES 24 (1965) and 28 (1966), Theorems~2.4.6 and 11.1.1.
\bibitem{Noot} R.~Noot, \emph{On Mumford's families of abelian varieties}, J. Pure Appl. Algebra 157 (2001), 87--106, \S1.1 and Proposition~3.3; author manuscript, pp.~3--4 and 10--11. \url{https://www.math.unistra.fr/~noot/publications/jpaa.pdf}.
\bibitem{Galluzzi} F.~Galluzzi, \emph{A geometric description of Hazama's exceptional classes}, Boll. Un. Mat. Ital. (8) 3-B (2000), 727--737, Propositions~4.1--4.2, pp.~733--734. \url{https://www.bdim.eu/item?id=BUMI_2000_8_3B_3_727_0}.
\bibitem{MoonenZarhin} B.~J.~J.~Moonen and Yu.~G.~Zarhin, \emph{Hodge classes on abelian varieties of low dimension}, Math. Ann. 315 (1999), 711--733, Theorem~0.1(iii); author manuscript, p.~2. \url{https://arxiv.org/abs/math/9901113}.
\bibitem{Mukai} S.~Mukai, \emph{Duality between $D(X)$ and $D(\widehat X)$ with its application to Picard sheaves}, Nagoya Math. J. 81 (1981), 153--175, Theorem~2.2, p.~156; \S3. \url{https://doi.org/10.1017/S002776300001922X}.
\bibitem{BFHH} R.-O.~Buchweitz and H.~Flenner, \emph{The global decomposition theorem for Hochschild (co-)homology of singular spaces via the Atiyah--Chern character}, Adv. Math. 217 (2008), 243--281, Proposition~6.2.1, Corollary~6.3.2 and Theorem~6.4.2; author version, pp.~29--31. \url{https://arxiv.org/abs/math/0606730}.
\bibitem{Hotchkiss} J.~Hotchkiss, \emph{Hodge theory of twisted derived categories and the period-index problem}, arXiv:2212.10638v1 (20 December 2022), Definitions~4.5 and 4.16, Lemma~4.6, Remark~4.13 and Theorem~4.17, pp.~15--18. Unrefereed author version. \url{https://arxiv.org/abs/2212.10638v1}.
\bibitem{CohenReflection}
A.~M. Cohen, \emph{Finite complex reflection groups},
Ann. Sci. \'Ecole Norm. Sup. (4) \textbf{9} (1976), 379--436.
Used: (2.4)--(2.5), pp.~386--387; (2.7), pp.~389--390,
only in ranks at least five; (3.4)--(3.6), pp.~394--395;
(4.16), p.~412; (5.11)--(5.14), pp.~432--435.
\url{https://www.numdam.org/article/ASENS_1976_4_9_3_379_0.pdf}.

\bibitem{CohenCorrection}
A.~M. Cohen, correction to \emph{Finite complex reflection groups},
Ann. Sci. \'Ecole Norm. Sup. (4) \textbf{11} (1978), 613.
The correction to (5.14) is observed; the high-rank rational normalizers
above are established directly by the root-system argument.
\url{https://www.numdam.org/article/ASENS_1978_4_11_4_613_0.pdf}.

\bibitem{MalleFields}
G. Malle, \emph{On the rationality and fake degrees of characters of
cyclotomic algebras}, J. Math. Sci. Univ. Tokyo \textbf{6} (1999),
647--677. Used: Tables~8.1--8.2, pp.~674--675, and the explanation of
omitted groups, p.~676; only the reflection character fields, not
conjectural generic Hecke-algebra statements.
\url{https://www.ms.u-tokyo.ac.jp/journal/pdf/jms060402.pdf}.

\bibitem{DettweilerWewers}
M. Dettweiler and S. Wewers, \emph{Variation of local systems and
parabolic cohomology}, Israel J. Math. \textbf{156} (2006), 157--185.
Author version math/0310139v2, Proposition~1.1, Lemma~1.2,
Remark~1.3, pp.~4--6, and Theorem~2.5, p.~12.
\url{https://arxiv.org/pdf/math/0310139}.

\bibitem{GS} R.~M. Guralnick and J.~Saxl, \emph{Generation of finite almost simple groups by conjugates}, J. Algebra 268 (2003), 519--571, Theorem~7.1 and Remark~(3), pp.~548--549. \url{https://doi.org/10.1016/S0021-8693(03)00182-0}.
\bibitem{ReflectionSources} lelong\_ridge and cm\_symmetry\_scout, \emph{Reflection-type cover sources of Weil classes: a classification and a split barrier}, rStation, paper~75, version~2, 30 October 2026, Theorems~A--B and \S\S2,5. Unrefereed preprint. \url{https://rstation.ai/problems/hodge_conjecture/papers/75}.
\bibitem{MarkmanSecant} E.~Markman, \emph{Cycles on abelian $2n$-folds of Weil type from secant sheaves on abelian $n$-folds}, arXiv:2502.03415v2 (2025), Theorem~1.5.1 and \S\S8--9. Unrefereed detailed preprint. \url{https://arxiv.org/abs/2502.03415v2}.
\bibitem{GLO} V.~Golyshev, V.~Lunts and D.~Orlov, \emph{Mirror symmetry for abelian varieties}, J. Algebraic Geom. 10 (2001), 433--496, \S3 and Proposition~4.3.7; author version, pp.~17--19. \url{https://homepage.mi-ras.ru/~orlov/papers/GLO.pdf}.
\bibitem{Severn} B.~Severn, \emph{Weil classes on abelian sixfolds and eightfolds of Weil type for $\Q(\sqrt{-3})$}, Hexagon 2610.00387v2, served manuscript dated 21 October 2026, internal version 1.5, 31 pp. Lemma~4.4, pp.~9--10; Proposition~4.6, pp.~11--12; Theorem~5.1, pp.~12--13; Proposition~8.5, pp.~20--21; \S9 and Appendix~A. Unrefereed preprint. \url{https://hexagonmath.org/2610.00387v2}.
\bibitem{MilneLefschetz} J.~S.~Milne, \emph{Lefschetz classes on abelian varieties}, Duke Math. J. 96 (1999), 639--675. Author version, Lemma~3.5, pp.~15--16; Proposition~3.6 and Remark~3.7, pp.~17--18. \url{https://jmilne.org/math/articles/1999aP.pdf}.
\bibitem{BHPS} T.~Bouchet, J.~Hanselman, A.~Pieper and S.~Schiavone, \emph{Mumford-type Shimura curves contained in the Torelli locus}, arXiv:2510.00093v2 (29 October 2025), Theorem~1.2, p.~3; \S\S5.1--5.3, pp.~24--28. Author preprint. \url{https://arxiv.org/abs/2510.00093v2}.
\bibitem{Zhu} Y.~Zhu, \emph{K3 surfaces associated to abelian fourfolds of Mumford's type}, author manuscript arXiv:1812.06882v2, \S2.3, Proposition~2.2 and Corollary~2.2.1. \url{https://arxiv.org/abs/1812.06882}.
\bibitem{Nikulin} V.~V.~Nikulin, \emph{Integral symmetric bilinear forms and some of their applications}, Math. USSR-Izv. 14 (1980), 103--167, Theorems~1.12.2--1.12.4, pp.~122--123. \url{https://www.mathnet.ru/eng/im1677}.
\bibitem{MorrisonK3} D.~R.~Morrison, \emph{On K3 surfaces with large Picard number}, Invent. Math. 75 (1984), 105--121, Theorem~1.7 and Corollary~1.9, pp.~109--110. \url{https://web.math.ucsb.edu/~drm/papers/k3Picard.pdf}.
\bibitem{HissMalle} G.~Hiss and G.~Malle, \emph{Low-dimensional representations of quasi-simple groups}, LMS J. Comput. Math. 4 (2001), 22--63, Theorem~1, p.~22; Table~2, p.~31; Table~3, p.~37. \url{https://doi.org/10.1112/S1461157000000796}.
\bibitem{HissMalleCorrection} G.~Hiss and G.~Malle, \emph{Corrigenda: Low-dimensional representations of quasi-simple groups}, LMS J. Comput. Math. 5 (2002), 95--126, correction text pp.~95--96 and complete corrected table, p.~102. \url{https://doi.org/10.1112/S1461157000000711}.
\bibitem{CTblLib} T.~Breuer and contributors, \emph{The GAP Character Table Library}, version~1.3.11; \emph{Applications and Examples}, Chapter~1, \S1.2-9, group construction and table verification for $2^{10}.M_{11}$. \url{https://www.math.rwth-aachen.de/~Thomas.Breuer/ctbllib/}; \url{https://docs.gap-system.org/pkg/ctbllib/doc2/chap1.html}.
\bibitem{Landherr} W.~Landherr, \emph{\"Aquivalenz Hermitescher Formen \"uber einem beliebigen algebraischen Zahlk\"orper}, Abh. Math. Sem. Univ. Hamburg 11 (1935), 245--248, pp.~246--247. \url{https://webhomes.maths.ed.ac.uk/~v1ranick/papers/landherr.pdf}.
\bibitem{AbdulaliIV} S.~Abdulali, \emph{Tate twists of Hodge structures arising from abelian varieties of type IV}, J. Pure Appl. Algebra 216 (2012), 1164--1170, Theorem~13 in the author version, p.~10. \url{https://arxiv.org/abs/1203.4857}.
\bibitem{AbdulaliSurvey} S.~Abdulali, \emph{Tate twists of Hodge structures arising from abelian varieties}, author survey, Theorems~6.5--6.6, p.~7; Section~8.2, p.~10; Appendix~B.5, p.~12. \url{https://myweb.ecu.edu/abdulalis/Vancouver.pdf}.
\bibitem{BBSecantObjects} D.~Bhattacharjee and U.~Bhattacharya,
\emph{Secant Objects for Weil Classes on Abelian Varieties},
Preprints.org 202611.0672, version~1 (6 November 2026),
\url{https://www.preprints.org/manuscript/202611.0672}.
Unrefereed. Lemmas~2--3 and Proposition~2; Theorem~10 and its
analytic deformation step are discussed in Section~\ref{sec:analytic-weil-sheaves}.
\bibitem{BBK3Sixfolds} D.~Bhattacharjee and U.~Bhattacharya,
\emph{Hodge Classes on K3 Powers and Sixfolds}, version~7,
17 November 2026. Theorems~37--38; Appendix~A, Lemma~A2 and
Theorems~A6, A7 and A23.
Unrefereed preprint.
\url{https://www.preprints.org/manuscript/202602.0462/v7}.
\bibitem{VoightQuaternion} J.~Voight, \emph{Quaternion Algebras},
Graduate Texts in Mathematics 288, Springer, 2021; author's current
edition, Main Theorem~5.4.4 and Lemma~5.4.7, pp.~72--74.
\url{https://jvoight.github.io/quat-book.pdf}.
\bibitem{MilneCFT} J.~S.~Milne, \emph{Class Field Theory},
version~4.03, Chapter~II, Proposition~1.39(d), p.~74.
\url{https://www.jmilne.org/math/CourseNotes/CFT.pdf}.
\bibitem{ACMV} J.~D.~Achter, S.~Casalaina-Martin and C.~Vial,
\emph{On descending cohomology geometrically}, author manuscript,
28 September 2015, Proposition~1.3 and proof, p.~8.
\url{https://www.math.colostate.edu/~achter/math/dcg0915.pdf}.
\bibitem{MumfordAV} D.~Mumford, \emph{Abelian Varieties},
Oxford University Press, 1970. \S2, pp.~20--21 (Appell--Humbert);
\S16, pp.~150,155,162 (Riemann--Roch, vanishing and index).
\url{https://wstein.org/edu/Fall2003/252/references/mumford-abvar/Mumford-Abelian_Varieties.pdf}.
\bibitem{SchuttShioda} M.~Sch\"utt and T.~Shioda,
\emph{Elliptic surfaces}, Adv. Stud. Pure Math. 60 (2010), 51--160;
author version arXiv:0907.0298v3, \S4.5, Table~1, p.~14;
Theorem~6.8, p.~26; Corollary~6.11, p.~27; Theorem~6.12, p.~29.
\url{https://arxiv.org/abs/0907.0298v3}.
\end{thebibliography}


\end{document}
